% !TEX program = lualatex
%% OptiTAB - Première spécialité - Livre complet.
\documentclass[11pt,twoside,openany]{book}
\usepackage{optitab-livre}
\usepackage{figures-geometrie}
\usepackage{relecture-livre}
\usepackage{edition-france}
\usepackage{longtable}
\pgfplotsset{otGraphe/.append style={clip=true}}
% Repères lisibles : les graduations indiquent les échelles de chaque axe.
\pgfplotsset{otLivreGraphe/.style={
  axis lines=middle,
  axis line style={otMuted,line width=0.5pt,-{Stealth[length=2mm]}},
  tick style={otMuted,line width=0.4pt},major tick length=2pt,
  ticklabel style={font=\otLight\fontsize{7}{8}\selectfont,text=otMuted,/pgf/number format/use comma,/pgf/number format/1000 sep={\,}},
  grid=major,grid style={otRule!70,line width=0.3pt},
  xlabel={$x$},ylabel={$y$},
  every axis x label/.style={at={(current axis.right of origin)},anchor=west,font=\small,text=otMuted},
  every axis y label/.style={at={(current axis.above origin)},anchor=south,font=\small,text=otMuted},
  clip=true}}
\newcommand{\otDecoseconddegre}{%
  \otDecoGrille
  \draw[white,opacity=.25] (132mm,47mm)--(210mm,47mm);
  \draw[white,opacity=.25] (171mm,18mm)--(171mm,118mm);
  \draw[otGreen,line width=1.8pt] plot[domain=-1.8:1.8,samples=75]
    ({(171+20*\x)*1mm},{(47+16*\x*\x)*1mm});
  \fill[white] (171mm,47mm) circle (1.5mm);}
\newcommand{\otDecovecteursplan}{%
  \otDecoGrille
  \draw[white,opacity=.7,line width=1.5pt,-{Stealth[length=3mm]}] (130mm,47mm)--(199mm,61mm);
  \draw[otGreen,line width=1.8pt,-{Stealth[length=3mm]}] (130mm,47mm)--(157mm,105mm);
  \draw[white,opacity=.35,dashed] (157mm,105mm)--(168mm,55mm);
  \fill[white] (130mm,47mm) circle (1.5mm);}
\newcommand{\otDecocercleplan}{%
  \otDecoGrille
  \draw[white,opacity=.25] (130mm,73mm)--(210mm,73mm) (170mm,36mm)--(170mm,116mm);
  \draw[otGreen,line width=1.8pt] (170mm,73mm) circle (30mm);
  \draw[white,opacity=.75,line width=1.2pt] (162mm,116.5mm)--(210mm,80.5mm);
  \draw[white,opacity=.45,dashed] (170mm,73mm)--(188mm,97mm);
  \fill[white] (188mm,97mm) circle (1.6mm);}
\newcommand{\otDecoensembles}{%
  \otDecoGrille
  \fill[white,opacity=.10] (155mm,76mm) circle (22mm);
  \fill[otGreen,opacity=.2] (181mm,76mm) circle (22mm);
  \draw[white,opacity=.8,line width=1.3pt] (155mm,76mm) circle (22mm);
  \draw[otGreen,line width=1.7pt] (181mm,76mm) circle (22mm);}
\newcommand{\otBrouillon}[2]{%
  \par\vspace{6mm}\otIntertitre{#1}
  \begin{tikzpicture}
    \foreach \k in {0,...,#2}{\draw[otRule,line width=.35pt] (0.2pt,-\k*7mm)--({\linewidth-0.2pt},-\k*7mm);}
  \end{tikzpicture}\par}
\hypersetup{pdftitle={Mathématiques - Première spécialité - Livre complet},
  pdfsubject={Première France 2026-2027 : cours, 224 exercices corrigés et deux sujets blancs},pdfauthor={OptiTAB}}
\begin{document}
\pagenumbering{alph}
%% =====================================================================
%%  1re de couverture
%% =====================================================================
\thispagestyle{empty}
\begin{tikzpicture}[remember picture,overlay]
\begin{scope}[shift={(current page.south west)}]
  % Fond
  \fill[otNavy] (0,0) rectangle (210mm,297mm);
  \shade[bottom color=otNavyDark,top color=otNavy] (0,0) rectangle (210mm,160mm);
  \draw[white,opacity=0.06,line width=0.3pt,step=7mm] (0,0) grid (210mm,297mm);
  % Repère, tangente et courbe de la fonction du second degré (unité : 16 mm)
  \draw[white,opacity=0.22,line width=0.6pt] (0,42mm) -- (210mm,42mm);
  \draw[white,opacity=0.22,line width=0.6pt] (150mm,0) -- (150mm,297mm);
  \draw[white,opacity=0.45,line width=0.8pt]
    (104mm,0mm) -- (210mm,212mm);
  \draw[otGreen,line width=3pt,line cap=round]
    plot[domain=100:213,samples=180] ({\x*1mm},{(42+0.02*(\x-100)*(\x-100))*1mm});
  \fill[white] (150mm,92mm) circle (2.1mm);
  \fill[otGreen] (150mm,92mm) circle (1.1mm);

  % Logo
  \node[anchor=north west,inner sep=0pt] at (22mm,279mm)
    {\includegraphics[width=60mm]{logo-blanc}};

  % Titre
  \node[anchor=base west,inner sep=0pt,text=otGreen,
    font=\otSemi\fontsize{11}{13}\selectfont] at (24mm,236mm) {\otCaps[22]{\otNiveau}};
  \node[anchor=base west,inner sep=0pt,text=white,
    font=\otBold\fontsize{50}{54}\selectfont] at (21.5mm,212mm) {\otTitreLivre};
  \node[anchor=base west,inner sep=0pt,text=white,
    font=\otLight\fontsize{32}{36}\selectfont] at (23mm,194mm) {\otSousTitre};
  \fill[otGreen] (24mm,181mm) rectangle (54mm,182.4mm);
  \node[anchor=base west,inner sep=0pt,text=white,opacity=0.85,
    font=\otRegular\fontsize{11}{14}\selectfont] at (24mm,168mm)
    {Année complète\enspace·\enspace Exercices progressifs\enspace·\enspace Corrigés détaillés\enspace·\enspace Méthodes};

  % Points forts
  \node[anchor=north west,inner sep=0pt,text width=100mm,text=white,
    font=\otRegular\fontsize{10.5}{14}\selectfont] at (24mm,146mm)
    {\otCheck\hspace{3mm}Exercices classés par difficulté\par\vspace{3.2mm}
     \otCheck\hspace{3mm}La méthode expliquée à chaque étape\par\vspace{3.2mm}
     \otCheck\hspace{3mm}14 défis avec indications progressives\par\vspace{3.2mm}
     \otCheck\hspace{3mm}Deux sujets blancs originaux};

  % Pastille
  \fill[otGreen] (146mm,116mm) circle (18mm);
  \node[text=white,align=center,inner sep=0pt] at (146mm,116mm)
    {{\otMedium\fontsize{7}{9}\selectfont Année complète}\\[0.6mm]
     {\otBold\fontsize{30}{30}\selectfont 224}\\[1.4mm]
     {\otSemi\fontsize{7.5}{9}\selectfont\otCaps[10]{exercices}}\\
     {\otSemi\fontsize{7.5}{9}\selectfont\otCaps[10]{corrigés}}};

  % Pied
  \node[anchor=base west,inner sep=0pt,text=white,opacity=0.75,
    font=\otMedium\fontsize{10}{12}\selectfont] at (24mm,16mm) {\otSite};
\end{scope}
\end{tikzpicture}
\clearpage

\pagenumbering{arabic}
%% =====================================================================
%%  Page de titre
%% =====================================================================
\thispagestyle{empty}
\begin{tikzpicture}[remember picture,overlay]
\begin{scope}[shift={(current page.south west)}]
  \node[anchor=north,inner sep=0pt] at (105mm,262mm)
    {\includegraphics[width=64mm]{logo-couleur}};
  \node[anchor=base,inner sep=0pt,text=otGreenDark,
    font=\otSemi\fontsize{10}{12}\selectfont] at (105mm,182mm) {\otCaps[22]{\otNiveau}};
  \node[anchor=base,inner sep=0pt,text=otNavy,
    font=\otBold\fontsize{42}{46}\selectfont] at (105mm,163mm) {\otTitreLivre};
  \node[anchor=base,inner sep=0pt,text=otInk,
    font=\otLight\fontsize{24}{28}\selectfont] at (105mm,148mm) {\otSousTitre};
  \fill[otGreen] (93mm,138mm) rectangle (117mm,139.2mm);
  \node[anchor=base,inner sep=0pt,text=otMuted,
    font=\itshape\fontsize{12}{15}\selectfont] at (105mm,126mm)
    {224 exercices progressifs et deux sujets blancs corrigés};
  \node[anchor=base,inner sep=0pt,text=otInk,
    font=\otMedium\fontsize{10}{12}\selectfont] at (105mm,42mm)
    {L'équipe pédagogique OptiTAB};
  \node[anchor=base,inner sep=0pt,text=otMuted,
    font=\otLight\fontsize{9}{11}\selectfont] at (105mm,35mm) {\otSite};
\end{scope}
\end{tikzpicture}
\clearpage

\thispagestyle{empty}
\vspace*{\fill}
\begin{minipage}{128mm}
\fontsize{9}{13}\selectfont\color{otMuted}
{\otSemi\color{otInk}Mathématiques · Première spécialité · Livre complet}\par
\vspace{3mm}
Conception et mise en page : collection OptiTAB.\par
224 exercices originaux et deux sujets blancs avec leurs corrigés.\par
Charte graphique reprise du livre OptiTAB de Terminale.\par
\vspace{5mm}
Édition du 23 octobre 2026, revue après audit de conformité au programme\par
et à la note de service de l'épreuve anticipée du 11 septembre 2026.\par
Algèbre, analyse, géométrie, probabilités et méthodes transversales.\par
Ce livre rassemble les thèmes du programme et des entraînements transversaux.\par
\vspace{5mm}
Référence pédagogique : programme de spécialité mathématiques de Première
générale publié au Bulletin officiel du 2 avril 2026, applicable à la rentrée
2026-2027. L'ordre des chapitres est un choix pédagogique de ce livre.\par
\vspace{5mm}
Les figures sont construites à partir des données des énoncés. Une lecture
graphique fournit une conjecture ou une valeur approchée ; les résultats exacts
sont établis par le calcul. Sauf indication contraire, les calculs sont exacts.\par
\vspace{5mm}
© OptiTAB, 2026.\par
\otSite
\end{minipage}
\clearpage

%% =====================================================================
%%  Avant-propos
%% =====================================================================
\thispagestyle{otFront}
\phantomsection\pdfbookmark[0]{Avant-propos}{bm-avant-propos}
\otPageTitre{Avant-propos}{}

\begin{minipage}{128mm}
{\fontsize{14}{20}\selectfont\itshape\color{otNavy}
Ce cahier a été conçu pour t'accompagner tout au long de ton année de
Première, en spécialité mathématiques.\par}
\vspace{7mm}

\fontsize{11}{16}\selectfont
Il réunit 210 exercices dans 14 chapitres, puis 14 défis d'approfondissement :
\textbf{224 exercices corrigés}, des calculs accessibles aux recherches très exigeantes.
Deux sujets blancs originaux complètent l'entraînement.\par\vspace{3.5mm}

Chaque chapitre suit la même progression. Les exercices de
\textbf{Découverte} vérifient que les notions du cours sont en place ; ceux
d'\textbf{Entraînement} installent les méthodes incontournables ; les
\textbf{Problèmes} relient plusieurs notions. Les \textbf{Défis} invitent à chercher
plus loin, avec des indications séparées pour progresser sans lire toute la solution.\par\vspace{3.5mm}

Les corrigés ne se contentent pas de donner la réponse : ils indiquent, à
chaque étape, la méthode utilisée. Prends le temps de chercher avant de les
consulter : c'est en cherchant que l'on progresse.\par\vspace{3.5mm}

Le tableau de bord, au début du livre, te permet de suivre ta progression
chapitre par chapitre et de repérer les exercices à retravailler.\par\vspace{9mm}

{\otMedium\fontsize{10}{13}\selectfont Bon travail, et bonne réussite au bac !}\par
\vspace{3mm}
{\otLight\fontsize{9}{12}\selectfont\color{otMuted}L'équipe pédagogique OptiTAB}\par
\vspace{8mm}
\includegraphics[height=9mm]{icone-couleur}
\end{minipage}
\clearpage

%% =====================================================================
%%  Mode d'emploi
%% =====================================================================
\thispagestyle{otFront}
\phantomsection\pdfbookmark[0]{Mode d'emploi}{bm-mode-emploi}
\otPageTitre{Mode d'emploi}{Comment ce livre est organisé}

%% --- Les trois rubriques -------------------------------------------
\otIntertitre{Chaque chapitre, trois étapes}
\begin{tcbraster}[raster columns=3,raster equal height,raster column skip=4mm,raster width=163mm]
\tcbset{colback=otNavyTint,colframe=otNavyTint,boxrule=0pt,arc=2mm,
  left=4mm,right=4mm,top=4mm,bottom=4mm}
\begin{tcolorbox}
  \otRond[otNavy]{1}\hspace{2mm}{\otSemi\fontsize{8}{10}\selectfont\color{otNavy}\otCaps[10]{Découverte}}\par
  \vspace{2.5mm}{\raggedright\fontsize{10}{13}\selectfont Des applications directes du cours, pour vérifier que les bases sont en place.\par}
\end{tcolorbox}
\begin{tcolorbox}
  \otRond[otNavy]{2}\hspace{2mm}{\otSemi\fontsize{8}{10}\selectfont\color{otNavy}\otCaps[10]{Entraînement}}\par
  \vspace{2.5mm}{\raggedright\fontsize{10}{13}\selectfont Les méthodes incontournables, à travailler jusqu'à les maîtriser.\par}
\end{tcolorbox}
\begin{tcolorbox}
  \otRond[otNavy]{3}\hspace{2mm}{\otSemi\fontsize{8}{10}\selectfont\color{otNavy}\otCaps[10]{Problèmes}}\par
  \vspace{2.5mm}{\raggedright\fontsize{10}{13}\selectfont Des sujets plus complets, dans l'esprit de l'épreuve anticipée de mathématiques.\par}
\end{tcolorbox}
\end{tcbraster}

\vspace{8mm}
%% --- Anatomie d'un exercice ----------------------------------------
\otIntertitre{Anatomie d'un exercice}
\begin{tcolorbox}[colback=white,colframe=otRule,boxrule=0.4pt,arc=2mm,
  left=5mm,right=5mm,top=9mm,bottom=4mm]
\begin{tikzpicture}
  \useasboundingbox (0,-9mm) rectangle (\linewidth,4mm);
  \node[anchor=west,fill=otNavy,rounded corners=1.6pt,inner sep=0pt,minimum width=7mm,
    minimum height=6.2mm,text=white,font=\otBold\fontsize{10}{10}\selectfont] (num) at (0,0) {6};
  \node[anchor=west,inner sep=0pt,text=otInk,font=\otSemi\fontsize{10.5}{13}\selectfont]
    (tit) at (10.5mm,0) {Choisir la bonne forme};
  \node[anchor=east,inner sep=0pt,text=otMuted,font=\otMedium\fontsize{7}{9}\selectfont]
    (ref) at (\linewidth,0) {Corrigé p.~\pageref*{cor-1-6}};
  \node[anchor=east,inner sep=0pt] (dots) at ([xshift=-3mm]ref.west) {\otDots{2}};
  \node[anchor=west,inner sep=0pt,text=otInk,font=\fontsize{10.5}{13}\selectfont]
    at (10.5mm,-7mm) {On considère la fonction $f$ définie sur $\R$ par $f(x)=x^{2}-6x+5$.};
  \foreach \n/\c in {1/num,2/tit,3/dots,4/ref}{%
    \node[circle,fill=otGreen,inner sep=0pt,minimum size=4.4mm,text=white,
      font=\otBold\fontsize{7}{8}\selectfont] at ([yshift=6.5mm]\c.north) {\n};
    \draw[otGreen,line width=0.5pt] ([yshift=4.3mm]\c.north) -- ([yshift=0.6mm]\c.north);}
\end{tikzpicture}
\end{tcolorbox}
\vspace{2mm}
{\fontsize{10}{13.5}\selectfont
\begin{tabular}{@{}p{77mm}@{\hspace{6mm}}p{77mm}@{}}
\otRond{1}\enspace Le numéro, à reporter dans le tableau de bord. &
\otRond{3}\enspace Le niveau de difficulté, de 1 à 3 points.\\[2mm]
\otRond{2}\enspace Le titre, qui annonce la notion travaillée. &
\otRond{4}\enspace La page du corrigé (lien cliquable dans le PDF).
\end{tabular}}

\vspace{8mm}
%% --- Les corrigés ---------------------------------------------------
\otIntertitre{Des corrigés qui expliquent}
\begin{minipage}[c]{96mm}
\begin{tcolorbox}[colback=white,colframe=otRule,boxrule=0.4pt,arc=2mm,
  left=5mm,right=5mm,top=4mm,bottom=4mm]
\otLabel{1.}\enspace\methode{factoriser un trinôme}\\[1mm]
\hspace*{1.4em}$x^{2}-4x+3=\res{(x-1)(x-3)}$
\end{tcolorbox}
\end{minipage}\hfill
\begin{minipage}[c]{62mm}
\fontsize{10}{13.5}\selectfont
\otRond{5}\enspace La méthode à retenir, rappelée avant chaque calcul.\par\vspace{2.5mm}
\otRond{6}\enspace Le résultat final, surligné en vert.
\end{minipage}

\vspace{8mm}
%% --- En ligne -------------------------------------------------------
\begin{tcolorbox}[colback=otNavyTint,colframe=otNavyTint,boxrule=0pt,arc=2mm,
  left=5mm,right=5mm,top=4mm,bottom=4mm]
\begin{minipage}[c]{24mm}
{\color{otNavy}\qrcode[height=21mm,nolink]{\otURL}}
\end{minipage}\hspace{5mm}%
\begin{minipage}[c]{120mm}
{\otSemi\fontsize{11}{14}\selectfont\color{otNavy}Et aussi en ligne}\par\vspace{1.5mm}
{\fontsize{10}{13.5}\selectfont Retrouve les cours, les fiches de synthèse et d'autres
exercices sur \textbf{\otSite}. Scanne le code avec ton téléphone pour y accéder.}
\end{minipage}
\end{tcolorbox}
\clearpage

\thispagestyle{otFront}
\phantomsection\pdfbookmark[0]{Parcours de travail}{bm-parcours}
\otPageTitre{Parcours de travail}{Des bases aux problèmes de recherche}
\fontsize{11}{15}\selectfont
\otIntertitre{Un livre de Première générale, spécialité mathématiques}
Les chapitres 1 à 11 développent l'algèbre, l'analyse, la géométrie, les probabilités
et les statistiques. Les chapitres 12 à 14 donnent les outils transversaux : Python,
raisonnement et automatismes. Ces derniers sont à reprendre au fil de l'année.
L'ordre des thèmes peut être adapté à la progression du professeur.\par\vspace{5mm}
\otIntertitre{Choisir le bon point de départ}
\begin{tabularx}{\linewidth}{@{}p{38mm}Y@{}}\toprule
\textbf{Parcours} & \textbf{Objectif}\\\midrule
Découverte, 1 à 5 & Vérifier les bases avec des applications directes.\\[2mm]
Entraînement, 6 à 11 & Installer les méthodes, puis choisir une démarche.\\[2mm]
Problèmes, 12 à 15 & Relier les notions, modéliser et rédiger.\\[2mm]
Défis D01 à D14 & Prendre des initiatives et traiter des discussions exigeantes.\\[2mm]
Sujets blancs A et B & Travailler en temps limité, sans calculatrice.\\\bottomrule
\end{tabularx}
\vspace{6mm}
\otIntertitre{Une méthode de travail en quatre temps}
\begin{questions}
\item Lire les objectifs et le cours. Reformuler une propriété avec ses hypothèses.
\item Chercher sans corrigé. Écrire les données, un schéma utile et les premières relations.
\item Comparer la démarche à la correction. Repérer précisément l'étape qui manquait.
\item Refaire quelques jours plus tard, puis expliquer la solution à voix haute.
\end{questions}
\vspace{4mm}
\otIntertitre{Se repérer et mesurer ses progrès}
Le tableau de bord suit les 210 exercices des chapitres. Pour les 14 défis,
noter dans son cahier : réussi seul, réussi avec un indice, ou à reprendre.
Les indications sont séparées des énoncés et des corrigés.
Les démonstrations de référence et les repères du programme servent de points d'appui.
\editionRetenir{La difficulté augmente par les choix de méthode et le raisonnement.
Un défi facultatif peut demander du temps, tout en restant fondé sur les outils de Première.}
\clearpage

\thispagestyle{otFront}
\phantomsection\pdfbookmark[0]{Sommaire}{bm-sommaire}
\otPageTitre{Sommaire}{Première générale · Spécialité mathématiques}
\vspace{-4mm}\sommaire\clearpage
\thispagestyle{otFront}
\otPageTitre{Parcours et ressources}{Compléments, entraînements et outils de révision}
\fontsize{11}{16}\selectfont
\newcommand{\editionSomLigne}[2]{\hyperref[#1]{\textcolor{otNavy}{#2}}\hfill p.~\pageref*{#1}\par\vspace{4mm}}
\otIntertitre{Approfondir}
\editionSomLigne{defis-debut}{14 défis d'approfondissement}
\editionSomLigne{indices-defis}{Indications progressives pour les défis}
\otIntertitre{Préparer l'épreuve anticipée}
\editionSomLigne{epreuve-guide}{Format, conseils et barème d'entraînement}
\editionSomLigne{sujet-A}{Sujet blanc A : automatismes, analyse, probabilités}
\editionSomLigne{sujet-B}{Sujet blanc B : automatismes, suites, géométrie}
\otIntertitre{Comprendre et corriger}
\editionSomLigne{demonstrations-debut}{Démonstrations de référence}
\editionSomLigne{corriges-debut}{Corrigés des 210 exercices de chapitre}
\editionSomLigne{defi-cor-1-1}{Corrigés détaillés des 14 défis}
\editionSomLigne{sujet-A-cor-auto}{Corrigés des deux sujets blancs}
\otIntertitre{Vérifier sa couverture du programme}
\editionSomLigne{couverture-programme}{Repères du programme : contenus, méthodes, preuves}
\clearpage

%% =====================================================================
%%  Mon tableau de bord (suivi de progression)
%% =====================================================================
\thispagestyle{otFront}
\phantomsection\pdfbookmark[0]{Mon tableau de bord}{bm-tableau-de-bord}
\otPageTitre{Mon tableau de bord}{}
\vspace{-4mm}
{\fontsize{10}{13.5}\selectfont
Après chaque exercice, colorie sa pastille :\enspace
\tikz[baseline=-0.6ex]\fill[otGreen] (0,0) circle (1.75mm);\ réussi seul\quad
\tikz[baseline=-0.6ex]\fill[otOrange] (0,0) circle (1.75mm);\ réussi avec le corrigé\quad
\tikz[baseline=-0.6ex]\fill[otRed] (0,0) circle (1.75mm);\ à refaire\par}
\vspace{1mm}
\tableauDeBord

\clearpage

% !TEX root = ../main.tex
\begin{ouverture}[seconddegre]{1}{Polynômes du second degré}{Partie 1 · Algèbre}
{Formes d'un trinôme\enspace·\enspace Équations\enspace·\enspace Signe\enspace·\enspace Modélisation}
\objectif{Choisir entre forme développée, factorisée et canonique.}
\objectif{Résoudre une équation et une inéquation du second degré.}
\objectif{Lire une parabole et déterminer son sommet.}
\objectif{Modéliser un problème d'aire, de coût ou de trajectoire.}
\prerequis{Identités remarquables et équations produit nul.}
\prerequis{Intervalles et repérage dans le plan.}
\end{ouverture}
\essentiel
\essSection{1}{Trois écritures, trois usages}
\begin{definition}[Trinôme]
Une fonction du second degré s'écrit $f(x)=ax^2+bx+c$, avec $a\neq0$.
Sa courbe est une parabole. La forme développée donne immédiatement $f(0)=c$.
\end{definition}
\begin{propriete}[Formes utiles]
La forme canonique $a(x-\alpha)^2+\beta$ donne le sommet $S(\alpha\pv\beta)$
et l'axe de symétrie $x=\alpha$. La forme factorisée $a(x-r_1)(x-r_2)$, lorsqu'elle
existe, donne les racines $r_1$ et $r_2$.
Alors $r_1+r_2=-b/a$ et $r_1r_2=c/a$, comme on le vérifie en développant.
Les fonctions du second degré qui s'annulent en deux réels distincts $r_1$ et $r_2$
sont les fonctions $x\mapsto a(x-r_1)(x-r_2)$, avec $a\neq0$.
\end{propriete}
\begin{exemple}
$x^2-4x+3=(x-2)^2-1=(x-1)(x-3)$.
Les racines sont 1 et 3 ; le minimum vaut $-1$, atteint en 2.
Pour une forme développée $ax^2+bx+c$, le sommet a pour abscisse $\alpha=-b/(2a)$
et pour ordonnée $f(\alpha)$. On peut aussi compléter un carré dans les cas simples.
La courbe de $x\mapsto f(x-m)$ se déduit de celle de $f$ par une translation de $m$
unités parallèlement à l'axe des abscisses : la parabole de $(x-2)^2-1$ est celle de $x^2$
décalée de 2 vers la droite et de 1 vers le bas.
\end{exemple}
\begin{center}
\begin{tikzpicture}\begin{axis}[otLivreGraphe,width=104mm,height=63mm,xmin=-0.5,xmax=4.5,ymin=-2,ymax=5,xtick={0,1,2,3,4},ytick={-1,1,3,5}]
\addplot[otNavy,thick,domain=-0.4:4.4,samples=90]{(x-2)^2-1};
\addplot[otGreen,dashed] coordinates {(2,-2)(2,5)};
\addplot[only marks,otNavy] coordinates {(1,0)(3,0)(2,-1)};
\node[below right,font=\small] at (axis cs:2,-1){$S(2;-1)$};
\end{axis}\end{tikzpicture}\par
\legende{La parabole de $f(x)=x^2-4x+3$ et son axe de symétrie.}
\end{center}
\begin{attention}
Une fonction peut ne pas avoir de racine réelle. La forme factorisée sur $\R$
n'existe alors pas. La forme canonique existe toujours.
\end{attention}
\clearpage
\essSection{2}{Résoudre une équation}
\begin{theoreme}[Discriminant]
Pour $ax^2+bx+c=0$, on pose $\Delta=b^2-4ac$.
Si $\Delta>0$, les deux solutions sont $\dfrac{-b-\sqrt\Delta}{2a}$ et
$\dfrac{-b+\sqrt\Delta}{2a}$. Si $\Delta=0$, la seule solution est
$-\dfrac b{2a}$. Si $\Delta<0$, il n'y a aucune solution réelle.
\end{theoreme}
\begin{savoirfaire}[Choisir le calcul le plus court]
Avant le discriminant, chercher une factorisation évidente :
$x^2-9=(x-3)(x+3)$, ou $2x^2-6x=2x(x-3)$.
Si une racine $r_1$ est évidente, l'autre se déduit de $r_1r_2=c/a$.
Pour compléter un carré, utiliser $x^2+2ux=(x+u)^2-u^2$.
\end{savoirfaire}
\essSection{3}{Étudier un signe et résoudre une inéquation}
\begin{propriete}
Lorsque deux racines $r_1<r_2$ existent, $f$ est du signe de $a$ à l'extérieur
de $[r_1;r_2]$, du signe opposé entre les racines, et s'annule aux racines.
Si $\Delta<0$, $f$ garde le signe de $a$. Si $\Delta=0$, $f$ garde le signe
de $a$ sauf à sa racine où elle vaut 0.
\end{propriete}
\begin{center}\begin{tableau}
\otTabL[2.6]{2.7}{$x$/0.8,$(x-1)(x-3)$/0.9}{$-\infty$,$1$,$3$,$+\infty$}
\tkzTabLine{,+,z,-,z,+,}
\end{tableau}\end{center}
\begin{savoirfaire}[Passer du signe à l'ensemble des solutions]
$x^2-4x+3\leq0\iff x\in[1;3]$.
Pour une inégalité stricte, les racines sont exclues :
$x^2-4x+3<0\iff x\in]1;3[$.
\end{savoirfaire}
\begin{attention}
Dans un quotient, interdire les valeurs qui annulent le dénominateur,
même si une simplification paraît possible. Pour un problème concret,
intersecter les solutions avec le domaine du modèle.
\end{attention}
\exercices

\rubrique{Découverte}
\begin{blocLivre}{ex-1-1}
\begin{exercice}[1]{Reconnaître un trinôme}
\begin{questions}
\item Parmi $3x^2-2x+5$, $4x-1$, $(x-2)^2$, $x^3+x$, repérer les polynômes du second degré.
\item Pour chacun d'eux, donner les coefficients $a$, $b$ et $c$.
\end{questions}
\end{exercice}
\end{blocLivre}
\begin{blocLivre}{ex-1-2}
\begin{exercice}[1]{Développer ou factoriser}
\begin{questions}
\item Développer $(x-3)(x+2)$ et $2(x-1)^2-8$.
\item Factoriser $x^2-25$ et $3x^2-12x$.
\end{questions}
\end{exercice}
\end{blocLivre}
\begin{blocLivre}{ex-1-3}
\begin{exercice}[1]{Lire la forme canonique}
Soit $f(x)=2(x-3)^2-8$.
\begin{questions}
\item Donner le sommet, l'axe de symétrie et le minimum de $f$.
\item Calculer $f(0)$ et résoudre $f(x)=0$.
\end{questions}
\end{exercice}
\end{blocLivre}
\begin{blocLivre}{ex-1-4}
\begin{exercice}[1]{Calculer un discriminant}
Résoudre dans $\R$ :
\begin{tasks}(3)
\task $x^2-5x+6=0$.
\task $x^2+4x+4=0$.
\task $2x^2+x+2=0$.
\end{tasks}
\end{exercice}
\end{blocLivre}
\begin{blocLivre}{ex-1-5}
\begin{exercice}[1]{Signe d'une forme factorisée}
Soit $P(x)=-2(x+1)(x-4)$.
\begin{questions}
\item Dresser le tableau de signes de $P$.
\item Résoudre $P(x)\geq0$, puis $P(x)<0$.
\end{questions}
\end{exercice}
\end{blocLivre}
\clearpage
\rubrique{Entraînement}
\begin{blocLivre}{ex-1-6}
\begin{exercice}[2]{Choisir la bonne forme}
On pose $f(x)=x^2-6x+5$.
\begin{questions}
\item Montrer que $f(x)=(x-3)^2-4$ et que $f(x)=(x-1)(x-5)$.
\item Résoudre $f(x)=5$, puis $f(x)\leq0$.
\item Résoudre $f(x)\leq-3$ en choisissant une forme adaptée.
\end{questions}
\end{exercice}
\end{blocLivre}
\begin{blocLivre}{ex-1-7}
\begin{exercice}[2]{Deux paraboles}
On considère $f(x)=x^2-4x+3$ et $g(x)=-x^2+2x+3$.
\begin{center}\begin{tikzpicture}\begin{axis}[otLivreGraphe,width=101mm,height=48mm,xmin=-1,xmax=4,ymin=-2,ymax=6,xtick={-1,0,1,2,3,4},ytick={-2,2,4,6}]
\addplot[otNavy,thick,domain=-0.6:4,samples=75]{x^2-4*x+3};
\addplot[otGreen,thick,domain=-1:3.4,samples=75]{-x^2+2*x+3};
\node[otEtiquette,text=otNavy] at(axis cs:3.7,2.7){$\mathcal C_f$};
\node[otEtiquette,text=otGreenDark] at(axis cs:2,4.7){$\mathcal C_g$};
\end{axis}\end{tikzpicture}\end{center}
\begin{questions}
\item Conjecturer les abscisses des points d'intersection.
\item Les calculer exactement, puis déterminer quand $f(x)\leq g(x)$.
\end{questions}
\end{exercice}
\end{blocLivre}
\begin{blocLivre}{ex-1-8}
\begin{exercice}[2]{Un quotient et une valeur interdite}
Résoudre $\dfrac{x^2-5x+6}{x-1}\geq0$.
Préciser le domaine, factoriser le numérateur, dresser un tableau de signes
et écrire l'ensemble des solutions avec les bons crochets.
\end{exercice}
\end{blocLivre}
\begin{blocLivre}{ex-1-9}
\begin{exercice}[2]{Retrouver une fonction}
Une parabole coupe l'axe des abscisses en $-2$ et $4$, et passe par $A(0;-8)$.
\begin{questions}
\item Justifier que sa fonction s'écrit $f(x)=a(x+2)(x-4)$, puis calculer $a$.
\item Développer $f$ et déterminer son sommet.
\item Résoudre $f(x)=-5$.
\end{questions}
\label{figlivre-r01-09-parabole-ex}
\figureLivre{r01-09-parabole}

\end{exercice}
\end{blocLivre}
\begin{blocLivre}{ex-1-10}
\begin{exercice}[2]{Somme et produit des racines}
\begin{questions}
\item Deux réels ont pour somme 7 et pour produit 10. Construire une équation
dont ils sont les racines et retrouver ces nombres.
\item Même question avec une somme de 4 et un produit de 5. Que peut-on conclure ?
\item Factoriser mentalement $x^2+x-12$ en cherchant les deux racines.
\end{questions}
\end{exercice}
\end{blocLivre}
\begin{blocLivre}{ex-1-11}
\begin{exercice}[2]{Discuter selon un paramètre}
Pour $m\in\R$, on considère $x^2-4x+m=0$.
\begin{questions}
\item Écrire $x^2-4x+m$ sous forme canonique.
\item Selon $m$, déterminer le nombre de solutions réelles.
\item Donner les solutions pour $m=3$, $m=4$ et $m=5$.
\item Pour quelles valeurs de $m$ l'expression est-elle positive ou nulle pour tout $x$ ?
\end{questions}
\end{exercice}
\end{blocLivre}
\clearpage
\rubrique{Problèmes}
\begin{blocLivre}{ex-1-12}
\begin{exercice}[3]{Un jardin contre un mur}
On dispose de 40 m de clôture pour fermer trois côtés d'un rectangle,
le quatrième étant un mur. On note $x$ la largeur perpendiculaire au mur.
\begin{center}\begin{tikzpicture}[x=1cm,y=1cm]
\fill[otNavyTint](0,0)rectangle(6,1.7);
\draw[otNavy,thick](0,1.7)--(0,0)--node[below]{$40-2x$}(6,0)--node[right]{$x$}(6,1.7);
\draw[otMuted,line width=2pt](-0.4,1.7)--(6.4,1.7);
\node[above,font=\small]at(3,1.7){mur};
\end{tikzpicture}\par\legende{Schéma de la situation, non à l'échelle.}\end{center}
\begin{questions}
\item Justifier $0<x<20$ et montrer que l'aire vaut $A(x)=-2x^2+40x$.
\item Écrire $A$ sous forme canonique ; en déduire l'aire maximale et les dimensions du jardin.
\item Pour quelles largeurs l'aire est-elle d'au moins $150$ m$^2$ ?
\end{questions}
\end{exercice}
\end{blocLivre}
\begin{blocLivre}{ex-1-13}
\begin{exercice}[3]{Une vente rentable}
Pour $q\in[0;100]$, on modélise le prix de vente unitaire par $p(q)=50-0{,}2q$
et le coût total par $C(q)=10q+500$, en euros.
\begin{questions}
\item Exprimer la recette $R(q)$ et le bénéfice $B(q)$.
\item Vérifier $B(q)=-0{,}2(q-100)^2+1500$ ; déterminer le bénéfice maximal sur le domaine.
\item Résoudre $B(q)\geq0$ et en déduire les quantités entières rentables. Donner une approximation à $10^{-2}$ des seuils utiles.
\end{questions}
\end{exercice}
\end{blocLivre}
\begin{blocLivre}{ex-1-14}
\begin{exercice}[3]{Un ballon au-dessus d'un filet}
La hauteur d'un ballon, en mètres, est $h(t)=-5t^2+10t+1$,
où $t\geq0$ désigne le temps en secondes après le lancer.
Le modèle est utilisé jusqu'au premier contact avec le sol.
\begin{questions}
\item Déterminer la hauteur initiale et la hauteur maximale.
\item Résoudre $h(t)=0$ ; à quel instant le ballon touche-t-il le sol ?
\item Déterminer les instants où $h(t)\geq4$.
\item Le ballon atteint le plan vertical d'un filet de 4 m à $t=1{,}5$ s.
Le franchit-il ? Justifier par un calcul et interpréter le résultat.
\end{questions}
\label{figlivre-r01-14-ballon-ex}
\figureLivre{r01-14-ballon}

\end{exercice}
\end{blocLivre}
\begin{blocLivre}{ex-1-15}
\begin{exercice}[3]{Une famille de paraboles}
Pour un réel $m$, on pose $P_m(x)=x^2-2mx+m+2$.
\partieExo{A}{Étude des racines}
\begin{questions}
\item Calculer le discriminant et le factoriser.
\item Déterminer les valeurs de $m$ pour lesquelles $P_m$ possède deux racines distinctes,
une racine double, ou aucune racine réelle.
\end{questions}
\partieExo{B}{Un point commun et un cas particulier}
\begin{questions}[resume]
\item Montrer que toutes les paraboles de la famille passent par un même point d'abscisse $\frac12$.
\item Pour $m=3$, résoudre $P_3(x)\leq0$ et déterminer le sommet.
\item Pour quelles valeurs de $m$ a-t-on $P_m(x)>0$ pour tout réel $x$ ?
\end{questions}
\end{exercice}
\end{blocLivre}

% !TEX root = ../main.tex
\begin{ouverture}[suites]{2}{Suites numériques}{Partie 1 · Algèbre}
{Termes\enspace·\enspace Récurrence\enspace·\enspace Suites arithmétiques et géométriques\enspace·\enspace Sommes}
\objectif{Calculer les termes d'une suite explicite ou récurrente.}
\objectif{Reconnaître et utiliser une suite arithmétique ou géométrique.}
\objectif{Étudier un sens de variation et calculer une somme.}
\objectif{Modéliser une évolution et programmer un seuil.}
\prerequis{Puissances, pourcentages et calcul littéral.}
\prerequis{Notions de variable, de boucle et de liste en Python.}
\end{ouverture}
\essentiel
\essSection{1}{Définir une suite et lire ses indices}
\begin{definition}
Une suite associe un nombre $u_n$ à chaque entier $n$ d'un ensemble d'indices.
On la note $(u_n)$ ou $(u(n))$ ; le terme de rang $n$ se note $u_n$ ou $u(n)$.
Une formule explicite donne directement $u_n$ ; une relation de récurrence
donne un terme à partir du précédent et nécessite un terme initial.
\end{definition}
\begin{exemple}
Si $u_n=2n+1$ pour $n\geq0$, alors $u_0=1$, $u_1=3$, $u_2=5$.
Si $v_0=1$ et $v_{n+1}=2v_n+1$, alors $v_1=3$, $v_2=7$, $v_3=15$.
\end{exemple}
\begin{attention}
$u_{n+1}$ est le terme de rang $n+1$, tandis que $u_n+1$ signifie « le terme
de rang $n$ auquel on ajoute 1 ». De $u_0$ à $u_n$, il y a $n+1$ termes.
\end{attention}
\essSection{2}{Étudier les variations}
\begin{propriete}
Si $u_{n+1}-u_n\geq0$ pour tout $n$, la suite est croissante ; si cette
différence est négative ou nulle, elle est décroissante. Pour une suite à
termes strictement positifs, on peut aussi comparer $\dfrac{u_{n+1}}{u_n}$ à 1.
\end{propriete}
\begin{center}\begin{tikzpicture}\begin{axis}[otLivreGraphe,width=108mm,height=54mm,xmin=-0.4,xmax=6.5,ymin=-0.5,ymax=14,xtick={0,1,2,3,4,5,6},ytick={2,6,10,14},xlabel={$n$},ylabel={$u_n$}]
\addplot[only marks,otNavy,mark=*,mark size=2pt]coordinates{(0,1)(1,3)(2,5)(3,7)(4,9)(5,11)(6,13)};
\end{axis}\end{tikzpicture}\par\legende{Les points $(n;u_n)$ de $u_n=2n+1$ : les indices sont entiers.}\end{center}
\clearpage
\essSection{3}{Les deux modèles fondamentaux}
\begin{definition}[Suite arithmétique]
$u_{n+1}=u_n+r$, où $r$ est constant. Alors $u_n=u_0+nr$, ou plus généralement
$u_n=u_p+(n-p)r$. Un accroissement constant correspond à ce modèle.
Les points $(n;u_n)$ sont alignés sur la droite de la fonction affine $x\mapsto u_0+rx$.
\end{definition}
\begin{definition}[Suite géométrique]
$v_{n+1}=qv_n$, où $q$ est constant. Alors $v_n=v_0q^n$, ou
$v_n=v_pq^{n-p}$. Un taux constant $t$ correspond à une raison $q=1+t$.
\end{definition}
\begin{exemple}
Ajouter 30 chaque mois donne une suite arithmétique de raison 30.
Augmenter de 3\,\% chaque année donne une suite géométrique de raison $1{,}03$.
Une diminution de 3\,\% correspond à $0{,}97$.
\end{exemple}
\essSection{4}{Calculer une somme sans oublier les bornes}
\begin{theoreme}
Pour une suite arithmétique, la somme de termes consécutifs est
\[\text{nombre de termes}\times\frac{\text{premier terme}+\text{dernier terme}}2.\]
Pour $q\neq1$, $1+q+\cdots+q^n=\dfrac{1-q^{n+1}}{1-q}$.
Pour $q=1$, la somme vaut $n+1$.
\end{theoreme}
\begin{savoirfaire}[Chercher un seuil]
Écrire une boucle qui actualise le terme et son indice tant que le seuil n'est
pas atteint. Vérifier le terme retenu \emph{et son prédécesseur} pour démontrer
qu'il s'agit du premier rang, après avoir établi le sens de variation.
\end{savoirfaire}
\begin{attention}
Des premiers termes suggèrent un comportement ; ils ne démontrent pas une
formule générale. Une suite définie par récurrence n'est pas nécessairement
arithmétique ou géométrique. L'intuition d'une limite ne remplace pas un calcul de seuil.
\end{attention}
\clearpage
\essSection{5}{Observer un comportement à long terme}
\begin{definition}[Une première intuition de limite]
On peut observer qu'une suite se rapproche d'un nombre, augmente au-delà
de tout seuil fixé, ou ne se rapproche d'aucune valeur unique. Cette approche
reste intuitive en Première : aucune définition formelle de limite ni
théorème de convergence n'est nécessaire ici.
\end{definition}
\begin{exemple}
Pour $u_n=1/(n+1)$, les valeurs deviennent proches de 0 tout en restant positives.
Pour $v_n=n^2$, les valeurs dépassent tout seuil fixé si l'on va assez loin.
Pour $w_n=(-1)^n$, les valeurs alternent entre 1 et $-1$ : elles ne se
rapprochent pas d'un seul nombre.
\end{exemple}
\begin{center}\begin{tikzpicture}\begin{axis}[otLivreGraphe,width=110mm,height=61mm,xmin=-.4,xmax=8.5,ymin=-1.4,ymax=1.6,xtick={0,1,2,3,4,5,6,7,8},ytick={-1,0,1},xlabel={$n$},ylabel={valeur}]
\addplot[only marks,otNavy,mark=*]coordinates{(0,1)(1,.5)(2,.333333)(3,.25)(4,.2)(5,.166667)(6,.142857)(7,.125)(8,.111111)};
\addplot[only marks,otGreenDark,mark=square*]coordinates{(0,1)(1,-1)(2,1)(3,-1)(4,1)(5,-1)(6,1)(7,-1)(8,1)};
\end{axis}\end{tikzpicture}\par
\legende{Points ronds : $1/(n+1)$. Carrés : $(-1)^n$.}
\end{center}
\begin{attention}
Une suite croissante n'est pas forcément proche d'un nombre fini : $n^2$
est croissante et grandit sans borne. Un graphique de quelques termes aide
à conjecturer ; la formule de la suite reste indispensable à l'analyse.
\end{attention}
\exercices

\rubrique{Découverte}
\begin{blocLivre}{ex-2-1}
\begin{exercice}[1]{Substituer un indice}
On pose $u_n=n^2-3n+2$ pour $n\geq0$.
\begin{questions}\item Calculer $u_0$, $u_1$, $u_2$ et $u_5$.
\item Exprimer $u_{n+1}$, puis calculer $u_{n+1}-u_n$.
\end{questions}
\end{exercice}
\end{blocLivre}
\begin{blocLivre}{ex-2-2}
\begin{exercice}[1]{Calculer de proche en proche}
$v_0=2$ et $v_{n+1}=3v_n-1$.
\begin{questions}\item Calculer $v_1$, $v_2$, $v_3$ et $v_4$.
\item Les trois premières différences sont-elles constantes ? Et les rapports ?
\end{questions}
\end{exercice}
\end{blocLivre}
\begin{blocLivre}{ex-2-3}
\begin{exercice}[1]{Une suite arithmétique}
$u_0=7$ et la raison est $-3$.
\begin{questions}\item Donner $u_n$, $u_8$ et le sens de variation.
\item Résoudre $u_n\leq-20$ pour $n$ entier naturel.
\end{questions}
\end{exercice}
\end{blocLivre}
\begin{blocLivre}{ex-2-4}
\begin{exercice}[1]{Une suite géométrique}
$v_0=5$ et la raison est 2.
\begin{questions}\item Donner $v_n$, $v_4$ et $v_6$.
\item Calculer $v_0+v_1+v_2+v_3+v_4$ de deux façons.
\end{questions}
\end{exercice}
\end{blocLivre}
\begin{blocLivre}{ex-2-5}
\begin{exercice}[1]{Du pourcentage à la formule}
Un appareil coûte initialement 800 euros. Sa valeur diminue de 15\,\% par an.
\begin{questions}\item Donner la relation entre $V_{n+1}$ et $V_n$ puis $V_n$ en fonction de $n$.
\item Calculer la valeur après deux ans. La baisse globale est-elle de 30\,\% ?
\end{questions}
\end{exercice}
\end{blocLivre}
\clearpage
\rubrique{Entraînement}
\begin{blocLivre}{ex-2-6}
\begin{exercice}[2]{Retrouver les paramètres}
\begin{questions}
\item Une suite arithmétique vérifie $u_3=11$ et $u_8=31$. Déterminer sa raison, $u_0$ et $u_n$.
\item Une suite géométrique de raison positive vérifie $v_2=12$ et $v_5=96$.
Déterminer sa raison, $v_0$ et $v_n$.
\end{questions}
\end{exercice}
\end{blocLivre}
\begin{blocLivre}{ex-2-7}
\begin{exercice}[2]{Compter les termes d'une somme}
\begin{questions}
\item Calculer $4+7+10+\cdots+61$ en déterminant d'abord le nombre de termes.
\item Calculer $3+6+12+\cdots+384$.
\item Calculer $\sum_{k=3}^{10}(2k+1)$.
\end{questions}
\end{exercice}
\end{blocLivre}
\begin{blocLivre}{ex-2-8}
\begin{exercice}[2]{Comparer des variations}
Pour $n\geq0$, on pose $u_n=n^2-4n$ et $v_n=\dfrac{n+1}{n+2}$.
\begin{questions}
\item Calculer $u_{n+1}-u_n$. La suite $(u_n)$ est-elle monotone sur tous les indices ?
\item Montrer que $v_{n+1}-v_n=\dfrac1{(n+2)(n+3)}$ et conclure.
\item Écrire $v_n=1-\dfrac1{n+2}$. Quel comportement semble se dessiner lorsque $n$ grandit ?
\end{questions}
\end{exercice}
\end{blocLivre}
\begin{blocLivre}{ex-2-9}
\begin{exercice}[2]{Lire puis justifier un nuage de points}
Les premiers termes d'une suite sont représentés ci-dessous.
\begin{center}\begin{tikzpicture}\begin{axis}[otLivreGraphe,width=105mm,height=51mm,xmin=-0.4,xmax=5.5,ymin=-0.5,ymax=18,xtick={0,1,2,3,4,5},ytick={2,4,8,16},xlabel={$n$},ylabel={$u_n$}]
\addplot[only marks,otGreenDark,mark=*,mark size=2.3pt]coordinates{(0,16)(1,8)(2,4)(3,2)(4,1)(5,0.5)};
\end{axis}\end{tikzpicture}\end{center}
\begin{questions}\item Conjecturer la nature de la suite et une formule explicite.
\item On précise $u_{n+1}=\frac12u_n$ et $u_0=16$. Justifier la formule et les variations.
\item Déterminer le premier rang pour lequel $u_n<0{,}1$.
\end{questions}
\end{exercice}
\end{blocLivre}
\begin{blocLivre}{ex-2-10}
\begin{exercice}[2]{Comprendre une boucle Python}
\begin{python}
u = 3
termes = [u]
for n in range(4):
    u = 2*u + 1
    termes.append(u)
\end{python}
\begin{questions}\item Donner la liste finale \code{termes}.
\item Quelle relation de récurrence est programmée ? Quel est l'indice du dernier terme ?
\item Adapter le programme pour obtenir les termes de $u_0$ à $u_{10}$.
\end{questions}
\end{exercice}
\end{blocLivre}
\begin{blocLivre}{ex-2-11}
\begin{exercice}[2]{Un seuil d'épargne}
$C_n=1000\times1{,}05^n$ représente un capital après $n$ années.
\begin{questions}\item Justifier que $(C_n)$ est strictement croissante.
\item Écrire une boucle qui détermine la première année où le capital atteint 1500 euros.
\item Exécuter la boucle ou faire des essais, puis vérifier les deux termes qui encadrent le seuil.
\end{questions}
\end{exercice}
\end{blocLivre}
\clearpage
\rubrique{Problèmes}
\begin{blocLivre}{ex-2-12}
\begin{exercice}[3]{Des rangées de sièges}
Une salle a 20 rangées. La première comporte 18 sièges ; chaque rangée suivante
en comporte deux de plus. On note $u_n$ le nombre de sièges de la rangée $n$, pour $1\leq n\leq20$.
\begin{questions}
\item Exprimer $u_n$ et calculer le nombre de sièges de la dernière rangée.
\item Déterminer la capacité totale de la salle.
\item Les dix dernières rangées sont vendues 12 euros la place, les dix premières 9 euros.
Calculer la recette si toutes les places sont vendues.
\item On ajoute des rangées en poursuivant la règle. Quel nombre total minimal de rangées
faut-il pour atteindre au moins 1000 places ?
\end{questions}
\end{exercice}
\end{blocLivre}
\begin{blocLivre}{ex-2-13}
\begin{exercice}[3]{Un traitement répété}
Après chaque heure, l'organisme conserve 80\,\% de la quantité présente, puis
reçoit une nouvelle dose de 6 mg. Juste après la première injection, $u_0=10$ mg.
On mesure $u_n$ juste après la nouvelle dose reçue à l'heure $n$.
\begin{questions}
\item Justifier $u_{n+1}=0{,}8u_n+6$, puis calculer $u_1$ et $u_2$.
\item Poser $v_n=u_n-30$. Montrer que $v_{n+1}=0{,}8v_n$.
\item En déduire $u_n=30-20\times0{,}8^n$ et justifier que $(u_n)$ est croissante et reste inférieure à 30.
\item Avec des essais ou Python, déterminer le premier rang où $u_n\geq28$.
\end{questions}
\label{figlivre-r02-13-cycle-ex}
\figureLivre{r02-13-cycle}

\end{exercice}
\end{blocLivre}
\begin{blocLivre}{ex-2-14}
\begin{exercice}[3]{Deux contrats de rémunération}
La première année, le salaire annuel vaut 24\,000 euros. Le contrat A prévoit
une augmentation de 1000 euros chaque année ; le contrat B, une hausse de 4\,\%.
On note $A_n$ et $B_n$ les salaires de l'année $n+1$ ; ainsi $A_0=B_0=24000$.
\begin{questions}
\item Donner $A_n$ et $B_n$. Comparer les salaires de la sixième année.
\item Calculer le total perçu pendant les dix premières années avec chaque contrat,
puis comparer. Arrondir les montants au centime.
\item Utiliser un tableau ou un programme pour déterminer le premier indice $n\geq1$
où $B_n>A_n$. Vérifier tous les indices antérieurs concernés.
\end{questions}
\end{exercice}
\end{blocLivre}
\begin{blocLivre}{ex-2-15}
\begin{exercice}[3]{Des carrés construits avec des allumettes}
On aligne $n$ carrés unitaires, deux carrés voisins partageant un côté.
\begin{center}\begin{tikzpicture}[x=1.1cm,y=1.1cm]
\foreach \i in {0,1,2,3}{\draw[otNavy,line width=1.1pt](\i,0)rectangle(\i+1,1);}
\node[below,font=\small]at(2,-0.1){Exemple : quatre carrés alignés.};
\end{tikzpicture}\end{center}
\begin{questions}
\item Donner le nombre $a_n$ d'allumettes nécessaires et expliquer la relation de récurrence.
\item Combien de carrés peut-on aligner avec 100 allumettes, sans les casser ?
\item On construit séparément une ligne de 1 carré, puis de 2, jusqu'à $N$ carrés.
Montrer que le nombre total d'allumettes est $T_N=\dfrac{3N^2+5N}{2}$.
\item Déterminer le plus grand $N$ permettant ces constructions avec 1000 allumettes.
\end{questions}
\end{exercice}
\end{blocLivre}

% !TEX root = ../main.tex
\begin{ouverture}[derivation]{3}{Dérivation}{Partie 2 · Analyse}
{Taux de variation\enspace·\enspace Nombre dérivé\enspace·\enspace Tangentes\enspace·\enspace Calculs}
\objectif{Calculer et interpréter un taux de variation.}
\objectif{Comprendre le nombre dérivé et lire une pente.}
\objectif{Déterminer l'équation d'une tangente.}
\objectif{Dériver des sommes, produits et quotients sur leur domaine.}
\prerequis{Équations de droites et coefficient directeur.}
\prerequis{Calcul littéral, puissances et fractions.}
\end{ouverture}
\essentiel
\essSection{1}{D'une sécante à une tangente}
\begin{definition}
Pour $h\neq0$, le taux de variation entre $a$ et $a+h$ est
$\dfrac{f(a+h)-f(a)}h$. Lorsqu'il se rapproche d'un nombre réel quand $h$
se rapproche de 0, ce nombre est le \textbf{nombre dérivé} $f'(a)$.
\end{definition}
\begin{center}\begin{tikzpicture}\begin{axis}[otLivreGraphe,width=108mm,height=62mm,xmin=-0.5,xmax=3.1,ymin=-1,ymax=8,xtick={0,1,2,3},ytick={1,2,4,6,8}]
\addplot[otNavy,thick,domain=-0.5:2.8,samples=90]{x^2};
\addplot[otGreen,thick,domain=-0.2:3]{2*x-1};
\addplot[otOrange,dashed,domain=0:2.6]{3*x-2};
\addplot[only marks,otNavy]coordinates{(1,1)(2,4)};
\node[above left,font=\small]at(axis cs:1,1){$A$};
\node[above left,font=\small]at(axis cs:2,4){$B$};
\node[otEtiquette,text=otGreenDark]at(axis cs:2.6,3.5){$T$};
\end{axis}\end{tikzpicture}\par
\legende{Pour $f(x)=x^2$, la sécante $(AB)$ a pour pente 3 ; la tangente $T$ en $A$ a pour pente 2.}
\end{center}
\begin{theoreme}[Équation de la tangente]
Si $f$ est dérivable en $a$, la tangente au point $A(a;f(a))$ a pour équation
\[y=f'(a)(x-a)+f(a).\]
Pour un petit $h$, on utilise l'approximation $f(a+h)\approx f(a)+hf'(a)$.
\end{theoreme}
\begin{exemple}
Pour $f(x)=x^2$ et $a=1$, $f(1)=1$ et $f'(1)=2$, donc $T:y=2x-1$.
Ainsi $1{,}02^2\approx1+2\times0{,}02=1{,}04$.
\end{exemple}
\clearpage
\essSection{2}{Les dérivées de référence}
{\renewcommand{\arraystretch}{1.7}\begin{tabularx}{\linewidth}{@{}Y Y Y@{}}
\toprule \textbf{Fonction} & \textbf{Dérivée} & \textbf{Domaine de dérivabilité}\\\midrule
$k$ (constante) & $0$ & $\R$\\
$x^n$, entier $n\geq1$ & $nx^{n-1}$ & $\R$\\
$1/x$ & $-1/x^2$ & $\R\setminus\{0\}$\\
$x^n$, entier $n<0$ & $nx^{n-1}$ & $\R\setminus\{0\}$\\
$\sqrt x$ & $1/(2\sqrt x)$ & $]0;+\infty[$\\\bottomrule
\end{tabularx}}
\essSection{3}{Sommes, produits et quotients}
\begin{propriete}
Pour des fonctions $u$ et $v$ dérivables, et une constante $k$ :
\[(ku+v)'=ku'+v',\qquad(uv)'=u'v+uv'.\]
Sur les intervalles où $v$ ne s'annule pas :
\[\left(\frac uv\right)'=\frac{u'v-uv'}{v^2},\qquad
\left(\frac1v\right)'=-\frac{v'}{v^2}.\]
\end{propriete}
\begin{savoirfaire}[Dériver puis simplifier]
Pour $f(x)=\dfrac{x+1}{x-2}$, on exclut $2$, puis
$f'(x)=\dfrac{(x-2)-(x+1)}{(x-2)^2}=\dfrac{-3}{(x-2)^2}$.
Toujours annoncer le domaine avant de dériver.
\end{savoirfaire}
\begin{attention}
$(uv)'$ n'est pas $u'v'$. Une fonction peut être définie en un point sans y
être dérivable : $\sqrt x$ en 0, ou $|x|$ en 0. Les taux à gauche et à droite
de 0 de $|x|$ valent respectivement $-1$ et $1$.
\end{attention}
\exercices

\rubrique{Découverte}
\begin{blocLivre}{ex-3-1}
\begin{exercice}[1]{Calculer une pente moyenne}
Soit $f(x)=x^2-2x$.
\begin{questions}\item Calculer le taux de variation entre 1 et 3, puis entre 1 et 2.
\item Calculer le taux entre 1 et $1+h$ pour $h\neq0$. En déduire $f'(1)$.
\end{questions}
\end{exercice}
\end{blocLivre}
\begin{blocLivre}{ex-3-2}
\begin{exercice}[1]{Lire la dérivée sur une tangente}
La tangente à une courbe en $A(2;3)$ passe aussi par $B(4;7)$.
\begin{questions}\item Déterminer $f(2)$ et $f'(2)$.
\item Écrire l'équation de cette tangente.
\end{questions}
\label{figlivre-r03-02-tangente-ex}
\figureLivre{r03-02-tangente}

\end{exercice}
\end{blocLivre}
\begin{blocLivre}{ex-3-3}
\begin{exercice}[1]{Dérivées immédiates}
Donner les dérivées de $f(x)=5x^4-3x^2+7x-8$, $g(x)=3/x$,
et $h(x)=\sqrt x+2x$. Préciser les domaines de dérivabilité.
Donner aussi la dérivée de $k(x)=x^{-3}$ sur $\R\setminus\{0\}$.
\end{exercice}
\end{blocLivre}
\begin{blocLivre}{ex-3-4}
\begin{exercice}[1]{Une tangente à un polynôme}
$f(x)=x^2+3x-1$.
\begin{questions}\item Calculer $f'(x)$, $f(1)$ et $f'(1)$.
\item Déterminer la tangente au point d'abscisse 1.
\end{questions}
\end{exercice}
\end{blocLivre}
\begin{blocLivre}{ex-3-5}
\begin{exercice}[1]{Reconnaître l'opération}
Sans développer, dériver $u(x)=(2x+1)(x^2-3)$ et
$v(x)=\dfrac{x+1}{x-2}$. Préciser les domaines.
\end{exercice}
\end{blocLivre}
\clearpage
\rubrique{Entraînement}
\begin{blocLivre}{ex-3-6}
\begin{exercice}[2]{Revenir à la définition}
On considère $f(x)=x^2$ et un réel $a$.
\begin{questions}
\item Montrer que $\dfrac{f(a+h)-f(a)}h=2a+h$ pour $h\neq0$.
\item En déduire $f'(a)$, puis la tangente au point d'abscisse $a$.
\item Pour $a=2$, calculer les pentes des sécantes correspondant à $h=1$, $0{,}1$, $-0{,}1$ et $0{,}01$.
\item Pour un pas $h>0$, écrire une fonction Python \code{secantes(h, n)} qui renvoie la liste
des coefficients directeurs des sécantes passant par les points d'abscisses $kh$ et $(k+1)h$,
pour $k$ allant de 0 à $n-1$. Prévoir \code{secantes(0.5, 4)} et comparer à $f'(kh)$.
\end{questions}
\end{exercice}
\end{blocLivre}
\begin{blocLivre}{ex-3-7}
\begin{exercice}[2]{Dériver trois produits}
Calculer et simplifier les dérivées de :
\begin{questions}\item $f(x)=(x^2+1)(3x-2)$ sur $\R$.
\item $g(x)=x^2\sqrt x$ sur $]0;+\infty[$.
\item $h(x)=(x-1)^2(x+2)$ sur $\R$, en développant d'abord ou en utilisant la règle du produit.
\end{questions}
\end{exercice}
\end{blocLivre}
\begin{blocLivre}{ex-3-8}
\begin{exercice}[2]{Dériver trois quotients}
Préciser le domaine de dérivabilité et calculer la dérivée de :
\begin{tasks}(3)
\task $f(x)=\dfrac{2x-1}{x+3}$.
\task $g(x)=\dfrac{x^2+1}{x}$.
\task $h(x)=\dfrac{x}{x^2+1}$.
\end{tasks}
Expliquer pourquoi une simplification préalable peut être utile pour $g$.
\end{exercice}
\end{blocLivre}
\begin{blocLivre}{ex-3-9}
\begin{exercice}[2]{Des tangentes parallèles}
$f(x)=x^3-3x+1$.
\begin{questions}\item Calculer $f'(x)$.
\item Déterminer les abscisses où la tangente est parallèle à $y=9x-2$.
\item Donner les équations des deux tangentes et vérifier qu'elles sont distinctes.
\end{questions}
\end{exercice}
\end{blocLivre}
\begin{blocLivre}{ex-3-10}
\begin{exercice}[2]{Approcher une racine carrée}
Soit $f(x)=\sqrt x$.
\begin{questions}\item Déterminer la tangente à sa courbe en $x=4$.
\item Utiliser cette tangente pour approcher $\sqrt{4{,}08}$.
\item Comparer à une valeur donnée par la calculatrice et calculer l'erreur absolue à $10^{-6}$ près.
\end{questions}
\end{exercice}
\end{blocLivre}
\begin{blocLivre}{ex-3-11}
\begin{exercice}[2]{Deux exemples de non-dérivabilité}
\begin{questions}
\item Pour $f(x)=|x|$, calculer $\dfrac{f(h)-f(0)}h$ lorsque $h>0$, puis $h<0$.
Conclure sur la dérivabilité en 0.
\item Pour $g(x)=\sqrt x$ et $h>0$, montrer que le taux entre 0 et $h$ vaut $1/\sqrt h$.
Pourquoi cela ne donne-t-il pas un nombre dérivé réel en 0 ?
\item Ces fonctions sont-elles malgré tout définies en 0 ?
\end{questions}
\label{figlivre-r03-11-nonderivable-ex}
\figureLivre{r03-11-nonderivable}

\end{exercice}
\end{blocLivre}
\clearpage
\rubrique{Problèmes}
\begin{blocLivre}{ex-3-12}
\begin{exercice}[3]{Tangentes à une parabole}
On considère $\Cf:y=x^2$ et le point $P(0;-1)$.
\begin{center}\begin{tikzpicture}\begin{axis}[otLivreGraphe,width=97mm,height=56mm,xmin=-2.3,xmax=2.3,ymin=-2,ymax=5,xtick={-2,-1,0,1,2},ytick={-1,1,3,5}]
\addplot[otNavy,thick,domain=-2.2:2.2,samples=80]{x^2};
\addplot[otGreen,domain=-2.3:2.3]{2*x-1};
\addplot[otGreen,domain=-2.3:2.3]{-2*x-1};
\addplot[only marks,otNavy]coordinates{(0,-1)(1,1)(-1,1)};
\node[below right,font=\small]at(axis cs:0,-1){$P$};
\end{axis}\end{tikzpicture}\end{center}
\begin{questions}
\item Écrire l'équation $T_a$ de la tangente au point d'abscisse $a$.
\item Montrer que $P$ appartient à $T_a$ si et seulement si $a^2=1$.
\item Donner les points de contact et les deux équations. Retrouver le dessin.
\end{questions}
\end{exercice}
\end{blocLivre}
\begin{blocLivre}{ex-3-13}
\begin{exercice}[3]{Vitesse moyenne et vitesse instantanée}
Un mobile a pour position $s(t)=t^3-6t^2+12t$, en mètres, pour $t\in[0;4]$ secondes.
\begin{questions}
\item Calculer la vitesse moyenne entre $t=1$ et $t=3$.
\item Calculer la vitesse instantanée $v(t)=s'(t)$ et l'écrire sous forme factorisée.
\item Quand le mobile est-il instantanément à l'arrêt ? Calculer sa position à cet instant.
\item Le mobile peut-il avoir une vitesse négative selon ce modèle ? Justifier.
\end{questions}
\end{exercice}
\end{blocLivre}
\begin{blocLivre}{ex-3-14}
\begin{exercice}[3]{Retrouver un polynôme grâce à une tangente}
On cherche $f(x)=ax^2+bx+c$ telle que $f(0)=1$, $f(2)=3$ et $f'(2)=4$.
\begin{questions}
\item Traduire les trois conditions en équations sur $a$, $b$ et $c$.
\item Déterminer ces coefficients et vérifier les trois conditions.
\item Écrire la tangente $T$ en $x=2$.
\item Factoriser $f(x)-(4x-5)$. La courbe est-elle au-dessus ou au-dessous de $T$ ?
\end{questions}
\end{exercice}
\end{blocLivre}
\begin{blocLivre}{ex-3-15}
\begin{exercice}[3]{Coût marginal et coût supplémentaire}
Le coût de fabrication de $q$ objets est modélisé, pour $q\in[0;100]$, par
$C(q)=0{,}02q^3-0{,}6q^2+10q+100$ euros.
\partieExo{A}{Calculer et interpréter}
\begin{questions}\item Calculer $C'(q)$ et le coût marginal $C'(20)$.
\item Calculer exactement $C(21)-C(20)$ ; comparer à $C'(20)$.
\item À l'aide de l'approximation affine, estimer $C(20{,}1)-C(20)$.
\end{questions}
\partieExo{B}{Contrôler une approximation}
\begin{questions}[resume]
\item Développer $C(20+h)-C(20)$ et montrer que cette différence vaut
$10h+0{,}6h^2+0{,}02h^3$.
\item Pour $h=0{,}1$, calculer l'erreur de l'approximation de la question 3.
Expliquer pourquoi une dérivée ne donne pas toujours exactement le coût d'un objet supplémentaire.
\end{questions}
\end{exercice}
\end{blocLivre}

% !TEX root = ../main.tex
\begin{ouverture}[derivation]{4}{Variations et optimisation}{Partie 2 · Analyse}
{Signe de la dérivée\enspace·\enspace Tableaux\enspace·\enspace Extremums\enspace·\enspace Modèles}
\objectif{Relier le signe de la dérivée aux variations.}
\objectif{Dresser un tableau avec les valeurs aux bornes.}
\objectif{Déterminer les extremums et comparer deux fonctions.}
\objectif{Résoudre un problème d'optimisation sur un domaine précis.}
\prerequis{Calcul des dérivées, tangentes et second degré.}
\prerequis{Aires, volumes et mise en équation d'un problème.}
\end{ouverture}
\essentiel
\essSection{1}{Du signe de la dérivée au sens de variation}
\begin{theoreme}
Soit $f$ dérivable sur un intervalle $I$.
Si $f'(x)\geq0$ sur $I$, $f$ y est croissante ; si $f'(x)\leq0$, elle y est
décroissante ; si $f'(x)=0$ partout, elle est constante.
Si la dérivée est strictement positive sauf en un nombre fini de points,
la croissance est stricte ; de même pour la décroissance.
\end{theoreme}
\begin{savoirfaire}[Les quatre étapes]
Déterminer le domaine. Calculer et factoriser la dérivée. Étudier son signe.
Tracer les flèches de variation et calculer les valeurs remarquables.
\end{savoirfaire}
\begin{exemple}
Pour $f(x)=x^3-3x$ sur $[-2;2]$, $f'(x)=3(x-1)(x+1)$.
\end{exemple}
\begin{center}\begin{tableau}
\otTab[2.4]{$x$/0.8,$f'(x)$/0.8,$f$/1.8}{$-2$,$-1$,$1$,$2$}
\tkzTabLine{,+,z,-,z,+,}
\tkzTabVar{-/$-2$,+/$2$,-/$-2$,+/$2$}
\end{tableau}\end{center}
\begin{attention}
$f'(a)=0$ ne suffit pas à prouver un extremum : $x^3$ a une dérivée nulle en 0
et reste strictement croissante. Il faut examiner le changement de signe.
Un maximum sur un intervalle peut aussi être atteint à une borne.
\end{attention}
\clearpage
\essSection{2}{Extremums et optimisation}
\begin{propriete}
Lorsque $f'$ passe de $+$ à $-$ en $a$, $f(a)$ est un maximum local.
Lorsqu'elle passe de $-$ à $+$, c'est un minimum local. Pour déterminer un
extremum global sur un segment, comparer également les valeurs aux bornes.
Si $f$ admet un extremum local en un point $a$ intérieur à l'intervalle, alors
$f'(a)=0$ : la tangente en ce point est horizontale.
\end{propriete}
\begin{exemple}
Pour $f(x)=ax^2+bx+c$, $f'(x)=2ax+b$ s'annule en $-b/(2a)$ en changeant de signe.
Si $a>0$, $f$ décroît puis croît : le sommet est un minimum ; si $a<0$, c'est un maximum.
\end{exemple}
\begin{savoirfaire}[Modéliser un problème]
Choisir une variable et son domaine. Exprimer la grandeur à optimiser avec
cette seule variable. Étudier ses variations. Conclure avec les dimensions,
la valeur optimale et l'unité demandées. Pour une quantité entière, comparer
les entiers admissibles voisins de l'optimum réel.
\end{savoirfaire}
\essSection{3}{Comparer deux fonctions et exploiter une symétrie}
\begin{propriete}
Pour comparer $f$ et $g$, étudier le signe de $f-g$. Si $f-g\geq0$, la courbe
de $f$ est au-dessus de celle de $g$.
Sur un domaine symétrique par rapport à 0, $f(-x)=f(x)$ caractérise une fonction
paire ; $f(-x)=-f(x)$, une fonction impaire.
\end{propriete}
\begin{center}
\begin{tikzpicture}\begin{axis}[otLivreGraphe,width=74mm,height=52mm,xmin=-2.2,xmax=2.2,ymin=-0.5,ymax=4,xtick={-2,-1,1,2},ytick={1,2,3}]
\addplot[otNavy,thick,domain=-1.9:1.9]{x^2};
\addplot[only marks,otGreenDark]coordinates{(-1.5,2.25)(1.5,2.25)};
\end{axis}\end{tikzpicture}\hspace{3mm}
\begin{tikzpicture}\begin{axis}[otLivreGraphe,width=74mm,height=52mm,xmin=-2.2,xmax=2.2,ymin=-3,ymax=3,xtick={-2,-1,1,2},ytick={-2,2}]
\addplot[otGreenDark,thick,domain=-2:2]{0.3*x^3};
\addplot[only marks,otNavy]coordinates{(-1.5,-1.0125)(1.5,1.0125)};
\end{axis}\end{tikzpicture}\par
\legende{À gauche : une fonction paire. À droite : une fonction impaire.}
\end{center}
\exercices

\rubrique{Découverte}
\begin{blocLivre}{ex-4-1}
\begin{exercice}[1]{Du signe aux variations}
Une fonction $f$ est dérivable sur $[-3;4]$, avec $f'(x)>0$ sur $[-3;1[$,
$f'(1)=0$ et $f'(x)<0$ sur $]1;4]$.
\begin{questions}\item Donner les variations de $f$ et la position de son maximum.
\item Avec $f(-3)=2$, $f(1)=7$ et $f(4)=-1$, donner aussi son minimum global.
\end{questions}
\end{exercice}
\end{blocLivre}
\begin{blocLivre}{ex-4-2}
\begin{exercice}[1]{Un trinôme sur un segment}
Étudier les variations de $f(x)=x^2-4x+1$ sur $[-1;5]$.
Déterminer son minimum et son maximum sur cet intervalle.
\end{exercice}
\end{blocLivre}
\begin{blocLivre}{ex-4-3}
\begin{exercice}[1]{Une dérivée ne change pas de signe}
Calculer la dérivée de $g(x)=x^3+3x$, justifier ses variations sur $\R$
et déterminer ses extremums sur $[-1;2]$.
\end{exercice}
\end{blocLivre}
\begin{blocLivre}{ex-4-4}
\begin{exercice}[1]{Un quotient simple}
Soit $h(x)=\dfrac{x+1}{x+2}$ sur $[0;4]$.
Calculer $h'$, dresser le tableau de variations et donner les deux extremums.
\end{exercice}
\end{blocLivre}
\begin{blocLivre}{ex-4-5}
\begin{exercice}[1]{Dérivée nulle et extremum}
Pour $f(x)=x^3$, calculer $f'(0)$ et étudier le signe de $f'$.
Le nombre 0 est-il l'abscisse d'un extremum ? Justifier par les variations.
\end{exercice}
\end{blocLivre}
\clearpage
\rubrique{Entraînement}
\begin{blocLivre}{ex-4-6}
\begin{exercice}[2]{Un polynôme du troisième degré}
$f(x)=x^3-3x^2-9x+5$ sur $[-3;5]$.
\begin{questions}\item Montrer que $f'(x)=3(x+1)(x-3)$.
\item Dresser le tableau complet de variations avec les valeurs aux bornes.
\item Distinguer les extremums locaux et les extremums globaux sur $[-3;5]$.
\end{questions}
\end{exercice}
\end{blocLivre}
\begin{blocLivre}{ex-4-7}
\begin{exercice}[2]{Un minimum avec un quotient}
$f(x)=x+4/x$ sur $]0;+\infty[$.
\begin{questions}\item Calculer $f'(x)$ et en étudier le signe.
\item En déduire le minimum de $f$ et l'inégalité $x+4/x\geq4$ pour tout $x>0$.
\item Retrouver cette inégalité en écrivant $f(x)-4$ sous la forme d'un quotient dont le numérateur est un carré.
\end{questions}
\end{exercice}
\end{blocLivre}
\begin{blocLivre}{ex-4-8}
\begin{exercice}[2]{Étudier une fonction paire}
$f(x)=x^4-4x^2$ sur $[-3;3]$.
\begin{questions}\item Vérifier la parité de $f$.
\item Factoriser $f'(x)$ ; étudier les variations sur $[0;3]$.
\item Compléter le tableau sur $[-3;3]$ par symétrie et déterminer les extremums globaux.
\item Vérifier également que $g(x)=x^3-3x$ est impaire sur $\R$ et préciser sa symétrie.
\end{questions}
\end{exercice}
\end{blocLivre}
\begin{blocLivre}{ex-4-9}
\begin{exercice}[2]{Position par rapport à une tangente}
Pour $f(x)=x^3$ sur $[0;3]$, on considère la tangente $T$ en $x=1$.
\begin{questions}\item Écrire l'équation de $T$.
\item Montrer que $x^3-3x+2=(x-1)^2(x+2)$.
\item Comparer la courbe et $T$ sur $[0;3]$. Où se rencontrent-elles ?
\end{questions}
\end{exercice}
\end{blocLivre}
\begin{blocLivre}{ex-4-10}
\begin{exercice}[2]{Une inégalité grâce aux variations}
Soit $f(x)=\sqrt x-\frac x2$ sur $[1;9]$.
\begin{questions}\item Calculer $f'(x)$ pour $x\in]1;9[$ et justifier que $f$ est décroissante.
\item Déterminer ses extremums sur $[1;9]$.
\item En déduire l'encadrement $-\frac32\leq\sqrt x-\frac x2\leq\frac12$ sur cet intervalle.
\end{questions}
\end{exercice}
\end{blocLivre}
\begin{blocLivre}{ex-4-11}
\begin{exercice}[2]{Optimiser une quantité entière}
Pour un nombre entier $n$ d'articles, $0\leq n\leq20$, le bénéfice est
$B(n)=-2n^2+27n-40$ euros.
\begin{questions}\item Étudier d'abord $B(x)$ sur $[0;20]$ pour une variable réelle.
\item Quels sont les deux entiers à comparer pour maximiser $B(n)$ ?
\item Donner la quantité entière optimale et le bénéfice correspondant.
\end{questions}
\end{exercice}
\end{blocLivre}
\clearpage
\rubrique{Problèmes}
\begin{blocLivre}{ex-4-12}
\begin{exercice}[3]{Une boîte sans couvercle}
Dans une feuille carrée de 24 cm de côté, on découpe quatre carrés de côté $x$,
puis on replie les bords. Le volume de la boîte est $V(x)$.
\begin{center}\begin{tikzpicture}[x=1cm,y=1cm]
\fill[otNavyTint](0,0)rectangle(3.7,3.7);
\foreach \a/\b in {0/0,3/0,0/3,3/3}{\fill[white](\a,\b)rectangle++(.7,.7);\draw[otMuted](\a,\b)rectangle++(.7,.7);}
\draw[otNavy,dashed](.7,.7)rectangle(3,3);
\node at(1.85,1.85){fond};\node at(.35,.35){$x$};
\draw[otMuted,<->](0,-.3)--node[below]{$24$ cm}(3.7,-.3);
\end{tikzpicture}\par\legende{Patron schématique, non à l'échelle.}\end{center}
\begin{questions}\item Justifier le domaine $]0;12[$ et établir $V(x)=x(24-2x)^2$.
\item Montrer que $V'(x)=12(x-4)(x-12)$ et étudier les variations.
\item Donner les dimensions de la boîte de volume maximal et ce volume en litres.
\end{questions}
\end{exercice}
\end{blocLivre}
\begin{blocLivre}{ex-4-13}
\begin{exercice}[3]{Un rectangle sous une parabole}
Les sommets supérieurs d'un rectangle sont $(-x;9-x^2)$ et $(x;9-x^2)$,
les deux autres étant sur l'axe des abscisses, avec $0<x<3$.
\begin{questions}\item Faire un schéma et exprimer l'aire $A(x)$.
\item Calculer $A'$ puis déterminer le rectangle d'aire maximale.
\item Donner la largeur, la hauteur et l'aire exactes, puis une valeur approchée de l'aire à $10^{-2}$ près.
\end{questions}
\label{figlivre-r04-13-rectangle-ex}
\figureLivre{r04-13-rectangle}

\end{exercice}
\end{blocLivre}
\begin{blocLivre}{ex-4-14}
\begin{exercice}[3]{Un emballage de périmètre imposé}
Un rectangle a un périmètre de 40 cm. On le fait tourner autour de l'un de
ses côtés pour obtenir un cylindre : le rayon vaut $r$ et la hauteur vaut $h$.
\begin{questions}\item Justifier $h=20-r$ et $0<r<20$.
\item Exprimer le volume $V(r)$ du cylindre et sa dérivée.
\item Déterminer le rayon et la hauteur donnant le volume maximal.
\item Calculer ce volume et expliquer pourquoi les deux limites géométriques
$r=0$ et $r=20$ ne donnent pas un cylindre utilisable.
\end{questions}
\label{figlivre-r04-14-cylindre-ex}
\figureLivre{r04-14-cylindre}

\end{exercice}
\end{blocLivre}
\begin{blocLivre}{ex-4-15}
\begin{exercice}[3]{Distance minimale à un point}
Un point mobile $M(x;x^2)$ se déplace sur la parabole $y=x^2$, avec $x\in[-2;2]$.
On cherche sa distance minimale au point $A(0;2)$.
\partieExo{A}{Éviter de dériver une racine}
\begin{questions}\item Montrer que $AM^2=D(x)=x^4-3x^2+4$.
\item Expliquer pourquoi minimiser $AM$ revient à minimiser $D(x)$.
\end{questions}
\partieExo{B}{Trouver tous les points optimaux}
\begin{questions}[resume]\item Factoriser $D'(x)$ et étudier ses variations sur $[-2;2]$.
\item Déterminer les deux positions de $M$ qui minimisent $AM$.
\item Donner la distance minimale exacte et vérifier qu'elle est inférieure à $AO$.
\end{questions}
\label{figlivre-r04-15-distance-ex}
\figureLivre{r04-15-distance}

\end{exercice}
\end{blocLivre}

% !TEX root = ../main.tex
\begin{ouverture}[exp]{5}{Fonction exponentielle}{Partie 2 · Analyse}
{Propriétés algébriques\enspace·\enspace Dérivées\enspace·\enspace Équations\enspace·\enspace Modélisation}
\objectif{Transformer une expression contenant des exponentielles.}
\objectif{Résoudre des équations et inéquations par comparaison.}
\objectif{Dériver et étudier des fonctions avec une exponentielle.}
\objectif{Relier les modèles continus et les suites géométriques.}
\prerequis{Calcul des puissances, dérivation et variations.}
\prerequis{Second degré et suites géométriques.}
\end{ouverture}
\essentiel
\essSection{1}{Définition et règles de calcul}
\begin{definition}
La fonction exponentielle, notée $\exp$ ou $x\mapsto e^x$, est l'unique
fonction dérivable sur $\R$ vérifiant $\exp'=\exp$ et $\exp(0)=1$.
Son existence et son unicité sont admises. Le nombre $e=\exp(1)$ vaut environ $2{,}71828$.
\end{definition}
\begin{propriete}
Pour tous réels $a$ et $b$,
\[e^{a+b}=e^ae^b,\quad e^{-a}=\frac1{e^a},\quad
e^{a-b}=\frac{e^a}{e^b},\quad (e^a)^n=e^{na}\ (n\in\mathbb Z).\]
On a $e^x>0$ pour tout réel $x$ : une exponentielle ne s'annule jamais.
\end{propriete}
\begin{attention}
$e^{a+b}$ est un produit, pas une somme : en général $e^{a+b}\neq e^a+e^b$.
Ne pas confondre $e^{2x}$ et $2e^x$.
\end{attention}
\essSection{2}{Une fonction strictement croissante}
\begin{theoreme}
La dérivée de $e^x$ est $e^x>0$, donc l'exponentielle est strictement croissante
sur $\R$. Ainsi $e^a=e^b\iff a=b$ et $e^a<e^b\iff a<b$.
En particulier, $e^x=1\iff x=0$.
\end{theoreme}
\begin{center}\begin{tikzpicture}\begin{axis}[otLivreGraphe,width=102mm,height=59mm,xmin=-2.5,xmax=2,ymin=-0.5,ymax=6,xtick={-2,-1,1},ytick={1,2,3,4,5,6}]
\addplot[otNavy,thick,domain=-2.5:1.75,samples=90]{exp(x)};
\addplot[otGreen,domain=-1.4:2]{x+1};
\addplot[only marks,otNavy]coordinates{(0,1)(1,2.718281828)};
\node[above left,font=\small]at(axis cs:0,1){$(0;1)$};
\end{axis}\end{tikzpicture}\par\legende{La courbe de l'exponentielle et sa tangente $y=x+1$ en 0.}\end{center}
\clearpage
\essSection{3}{Dériver une exponentielle dans un produit}
\begin{propriete}
Pour une constante réelle $a$, la dérivée de $t\mapsto e^{at}$ est $t\mapsto ae^{at}$.
Comme $e^{ax+b}=e^be^{ax}$, on en déduit $\big(e^{ax+b}\big)'=a e^{ax+b}$.
Dans un produit, on conserve les deux termes de la dérivée :
\[\big((x+1)e^{-x}\big)'=e^{-x}-(x+1)e^{-x}=-xe^{-x}.\]
\end{propriete}
\begin{savoirfaire}[Étudier un signe]
Factoriser la dérivée par l'exponentielle. Comme celle-ci est strictement
positive, le signe dépend du facteur restant. Pour l'exemple précédent,
le signe est celui de $-x$.
\end{savoirfaire}
\essSection{4}{Modéliser une évolution continue}
\begin{definition}
Un modèle $Q(t)=Q_0e^{kt}$, avec $Q_0>0$, décrit une croissance lorsque
$k>0$ et une décroissance lorsque $k<0$. Pour des instants entiers,
$Q(n)=Q_0(e^k)^n$ : les valeurs forment une suite géométrique de raison $e^k$.
\end{definition}
\begin{center}\begin{tikzpicture}\begin{axis}[otLivreGraphe,width=110mm,height=51mm,xmin=-0.2,xmax=4.4,ymin=-0.1,ymax=4.1,xtick={0,1,2,3,4},ytick={1,2,3,4},xlabel={$t$},ylabel={$Q(t)$}]
\addplot[otNavy,thick,domain=0:4,samples=75]{exp(0.3*x)};
\addplot[only marks,otGreenDark,mark=*,mark size=2.1pt]coordinates{(0,1)(1,1.349859)(2,1.822119)(3,2.459603)(4,3.320117)};
\end{axis}\end{tikzpicture}\par\legende{Modèle continu $Q(t)=e^{0{,}3t}$ et valeurs aux instants entiers.}\end{center}
\begin{attention}
Le taux entre deux instants séparés d'une unité vaut $e^k-1$, et non $k$.
Pour chercher un seuil qui ne se résout pas par comparaison d'exposants,
on utilise des essais, un tableau ou Python. Aucun logarithme n'est nécessaire ici.
\end{attention}
\exercices

\rubrique{Découverte}
\begin{blocLivre}{ex-5-1}
\begin{exercice}[1]{Simplifier une expression}
Écrire sous la forme d'une seule exponentielle :
\begin{tasks}(3)
\task $e^{2x}e^{3-x}$.
\task $e^{5x}/e^{2x-1}$.
\task $(e^{-x})^3e^{2x}$.
\end{tasks}
\end{exercice}
\end{blocLivre}
\begin{blocLivre}{ex-5-2}
\begin{exercice}[1]{Des équations immédiates}
Résoudre dans $\R$ : $e^{2x-1}=e^3$, puis $e^{x+2}=1$ et $e^x=-2$.
Justifier chaque réponse à l'aide d'une propriété de l'exponentielle.
\end{exercice}
\end{blocLivre}
\begin{blocLivre}{ex-5-3}
\begin{exercice}[1]{Comparer des exponentielles}
Résoudre $e^{3x+1}\leq e^{x+5}$ et $e^{-2x}>1$.
Ranger aussi $e^{-1}$, $1$ et $e^{1/2}$ dans l'ordre croissant.
\end{exercice}
\end{blocLivre}
\begin{blocLivre}{ex-5-4}
\begin{exercice}[1]{Dériver et donner le signe}
Calculer les dérivées de $f(x)=e^{2x}$, $g(x)=3e^{-x}$ et $h(x)=e^{x+1}-4$.
Donner le sens de variation de chacune sur $\R$.
\end{exercice}
\end{blocLivre}
\begin{blocLivre}{ex-5-5}
\begin{exercice}[1]{La tangente en zéro}
Pour $f(x)=e^x$, calculer $f(0)$ et $f'(0)$, puis écrire la tangente en 0.
Utiliser cette droite pour approcher $e^{0{,}03}$.
\end{exercice}
\end{blocLivre}
\clearpage
\rubrique{Entraînement}
\begin{blocLivre}{ex-5-6}
\begin{exercice}[2]{Factoriser avant de résoudre}
Résoudre dans $\R$ :
\begin{questions}\item $e^{2x}-e^x=0$.
\item $(x-2)e^x=0$.
\item $(x^2-3x+2)e^{-x}\leq0$.
\end{questions}
\end{exercice}
\end{blocLivre}
\begin{blocLivre}{ex-5-7}
\begin{exercice}[2]{Étudier un produit}
On pose $f(x)=(x+1)e^{-x}$ sur $[-2;4]$.
\begin{questions}\item Montrer que $f'(x)=-xe^{-x}$.
\item Dresser le tableau de variations et déterminer les extremums sur ce segment.
\item Écrire la tangente en 1 et déterminer où $f$ est positive ou nulle.
\end{questions}
\end{exercice}
\end{blocLivre}
\begin{blocLivre}{ex-5-8}
\begin{exercice}[2]{Une différence d'exponentielles}
On considère $g(x)=e^{2x}-2e^x$ sur $\R$.
\begin{questions}\item Factoriser $g'(x)$ et étudier son signe.
\item En déduire le minimum de $g$.
\item Retrouver ce résultat avec $g(x)=(e^x-1)^2-1$.
\end{questions}
\end{exercice}
\end{blocLivre}
\begin{blocLivre}{ex-5-9}
\begin{exercice}[2]{Un quotient toujours défini}
Soit $f(x)=\dfrac{e^x}{1+e^x}$.
\begin{questions}\item Justifier que $f$ est définie et dérivable sur $\R$.
\item Montrer que $f'(x)=\dfrac{e^x}{(1+e^x)^2}$ et en déduire les variations.
\item Montrer $0<f(x)<1$, puis résoudre $f(x)=\frac12$.
\end{questions}
\end{exercice}
\end{blocLivre}
\begin{blocLivre}{ex-5-10}
\begin{exercice}[2]{Passer du continu au discret}
Une population est modélisée par $N(t)=500e^{0{,}2t}$, $t$ en années.
\begin{questions}\item Calculer $N(0)$ et $N(5)$ (valeur exacte, puis arrondie à l'unité).
\item Pour $u_n=N(n)$, donner la nature, le premier terme et la raison de $(u_n)$.
\item Déterminer le pourcentage d'augmentation annuel à $0{,}1$ point de pourcentage près.
\end{questions}
\end{exercice}
\end{blocLivre}
\begin{blocLivre}{ex-5-11}
\begin{exercice}[2]{Construire une approximation d'Euler}
Pour approcher $e^1$, on part de $y_0=1$ et on effectue $N$ étapes de longueur
$h=1/N$ en posant $y_{k+1}=y_k+hy_k$.
\begin{questions}\item Donner la nature de $(y_k)$ et exprimer $y_N$ en fonction de $N$.
\item Calculer les approximations pour $N=2$ et $N=10$, puis comparer à $e$.
\item Écrire une fonction Python qui calcule l'approximation par une boucle de $N$ étapes.
\end{questions}
\end{exercice}
\end{blocLivre}
\clearpage
\rubrique{Problèmes}
\begin{blocLivre}{ex-5-12}
\begin{exercice}[3]{Le refroidissement d'une boisson}
La température d'une boisson est modélisée par $T(t)=20+70e^{-0{,}1t}$,
où $t\geq0$ est exprimé en minutes et $T$ en degrés Celsius.
\begin{center}\begin{tikzpicture}\begin{axis}[otLivreGraphe,width=107mm,height=51mm,xmin=-1,xmax=31,ymin=-3,ymax=100,xtick={0,5,10,15,20,25,30},ytick={20,40,60,80,100},xlabel={$t$},ylabel={$T$}]
\addplot[otNavy,thick,domain=0:30,samples=100]{20+70*exp(-0.1*x)};
\addplot[otOrange,dashed,domain=0:30]{40};
\end{axis}\end{tikzpicture}\end{center}
\begin{questions}\item Calculer la température initiale et justifier la décroissance.
\item Montrer que la température reste strictement supérieure à $20\,{}^\circ$C.
\item Trouver la première minute entière où $T(t)\leq40$, en vérifiant la minute précédente.
\item Programmer ce seuil avec une boucle \code{while}.
\end{questions}
\end{exercice}
\end{blocLivre}
\begin{blocLivre}{ex-5-13}
\begin{exercice}[3]{Une concentration après injection}
Pour $t\in[0;8]$ heures, la concentration est $C(t)=8te^{-0{,}5t}$ mg/L.
\begin{questions}\item Calculer $C'(t)$ et dresser le tableau de variations.
\item Déterminer l'instant et la valeur de la concentration maximale.
\item Écrire la tangente en $t=0$ et interpréter sa pente.
\item Entre $t=2$ et $t=3$, quel est le taux de variation moyen ? Donner une valeur approchée à $10^{-2}$ et son unité.
\end{questions}
\end{exercice}
\end{blocLivre}
\begin{blocLivre}{ex-5-14}
\begin{exercice}[3]{Prouver une inégalité sans convexité}
On pose $F(x)=e^x-1-x$ sur $\R$.
\partieExo{A}{Une inégalité fondamentale}
\begin{questions}\item Calculer $F'(x)$ et étudier son signe.
\item Déterminer le minimum de $F$. En déduire $e^x\geq1+x$ pour tout réel $x$,
et préciser le cas d'égalité.
\end{questions}
\partieExo{B}{Une application numérique}
\begin{questions}[resume]\item Appliquer cette inégalité à $x=0{,}1$ puis $x=-0{,}1$.
\item En utilisant $e^{-0{,}1}=1/e^{0{,}1}$, en déduire
$1{,}1\leq e^{0{,}1}\leq10/9$.
\item Comparer cet encadrement à l'approximation affine de $e^{0{,}1}$ en 0.
\end{questions}
\end{exercice}
\end{blocLivre}
\begin{blocLivre}{ex-5-15}
\begin{exercice}[3]{Choisir une durée de campagne}
Le gain d'une campagne, pour $t\in[0;6]$ semaines, est modélisé par
$G(t)=100(t+1)e^{-0{,}5t}$, en centaines d'euros.
On pourra utiliser l'encadrement $2<e<3$.
\partieExo{A}{Optimiser le gain}
\begin{questions}\item Montrer que $G'(t)=50(1-t)e^{-0{,}5t}$.
\item Dresser le tableau de variations et déterminer la durée optimale et le gain maximal en euros.
\end{questions}
\partieExo{B}{Comparer deux modèles}
\begin{questions}[resume]\item Un autre prestataire propose un gain constant $H(t)=100$ dans la même unité.
Comparer exactement $G(0)$ et $G(1)$ à 100.
\item Calculer $G(2)$ et $G(3)$ à $10^{-2}$ près. En utilisant les variations,
expliquer pourquoi la campagne rapporte davantage que $H$ pour toute durée $0<t\leq2$.
\item Pour des durées entières de 1 à 6 semaines, lesquelles donnent un gain strictement supérieur à $H$ ?
\end{questions}
\end{exercice}
\end{blocLivre}

% !TEX root = ../main.tex
\begin{ouverture}[trigo]{6}{Trigonométrie}{Partie 2 · Analyse}
{Radian\enspace·\enspace Cercle unité\enspace·\enspace Sinus et cosinus\enspace·\enspace Angles associés}
\objectif{Passer des degrés aux radians et calculer une longueur d’arc.}
\objectif{Placer un point sur le cercle trigonométrique.}
\objectif{Déterminer des valeurs exactes par symétrie.}
\objectif{Utiliser les coordonnées du cercle dans une situation géométrique.}
\prerequis{Triangles rectangles et théorème de Pythagore.}
\prerequis{Coordonnées et symétries dans le plan.}
\end{ouverture}
\essentiel
\essSection{1}{Enrouler une droite sur le cercle}
\begin{definition}
Le cercle trigonométrique est le cercle de centre $O$ et de rayon 1, orienté
dans le sens inverse des aiguilles d'une montre. Un réel $t$ correspond à un
déplacement de longueur $|t|$ sur ce cercle : sens positif si $t>0$, sens
négatif si $t<0$. Un tour correspond à $2\pi$ radians, soit $360^\circ$.
Ainsi $180^\circ=\pi$ rad et $90^\circ=\pi/2$ rad.
\end{definition}
\begin{propriete}
Sur un cercle de rayon $R$, un angle de mesure $\theta\geq0$ radians intercepte
un arc de longueur $R\theta$. Deux réels ont le même point-image exactement
lorsque leur différence est un multiple entier de $2\pi$.
\end{propriete}
\essSection{2}{Les coordonnées du point-image}
\begin{definition}
Le point-image de $t$ a pour coordonnées $(\cos t;\sin t)$ dans le repère
orthonormé $(O;I;J)$. On a $-1\leq\cos t\leq1$, $-1\leq\sin t\leq1$ et
$\cos^2t+\sin^2t=1$.
\end{definition}
\begin{center}
\begin{tikzpicture}[scale=1.45]
\draw[otMuted,->](-1.45,0)--(1.6,0)node[right]{$\cos t$};
\draw[otMuted,->](0,-1.25)--(0,1.45)node[above]{$\sin t$};
\draw[otNavy,thick](0,0)circle(1);
\draw[otGreen,thick](0,0)--(.5,.8660254);
\draw[otGreenDark,->](.3,0)arc(0:60:.3);
\draw[otMuted,dashed](.5,0)--(.5,.8660254)--(0,.8660254);
\fill[otGreenDark](.5,.8660254)circle(1.4pt)node[above right]{$M$};
\node[below left]at(0,0){$O$};\node[below right]at(1,0){$I$};
\node[font=\small]at(.45,.22){$t$};
\end{tikzpicture}
\par\legende{Pour $t=\pi/3$, le point $M$ a pour coordonnées $(1/2;\sqrt3/2)$.}
\end{center}

\clearpage
\essSection{3}{Les valeurs remarquables}
{\renewcommand{\arraystretch}{1.9}\begin{tabularx}{\linewidth}{@{}l*{5}{>{\centering\arraybackslash}X}@{}}
\toprule $t$ & $0$ & $\pi/6$ & $\pi/4$ & $\pi/3$ & $\pi/2$\\\midrule
$\cos t$ & $1$ & $\sqrt3/2$ & $\sqrt2/2$ & $1/2$ & $0$\\
$\sin t$ & $0$ & $1/2$ & $\sqrt2/2$ & $\sqrt3/2$ & $1$\\\bottomrule
\end{tabularx}}
\begin{propriete}[Les symétries du cercle]
Pour tout réel $t$,
\[\cos(-t)=\cos t,\quad\sin(-t)=-\sin t,\]
\[\cos(\pi-t)=-\cos t,\quad\sin(\pi-t)=\sin t,\]
\[\cos(t+\pi)=-\cos t,\quad\sin(t+\pi)=-\sin t.\]
Ajouter $2k\pi$, avec $k$ entier relatif, ne change ni sinus ni cosinus.
\end{propriete}
\begin{savoirfaire}[Déterminer une valeur exacte]
Réduire l'angle d'un nombre entier de tours. Repérer le quadrant. Utiliser
une symétrie pour revenir à une valeur remarquable, puis vérifier le signe.
Par exemple, $\cos(5\pi/6)=-\cos(\pi/6)=-\sqrt3/2$.
\end{savoirfaire}
\begin{attention}
Les mesures utilisées dans les calculs sont en radians. Un angle peut avoir
plusieurs mesures, mais un point-image est unique. Une valeur décimale à la
calculatrice ne remplace pas une valeur exacte demandée.
\end{attention}
\begin{exemple}[Un argument géométrique]
Dans un triangle rectangle isocèle de côtés de l'angle droit égaux à 1,
l'hypoténuse vaut $\sqrt2$ ; donc $\cos(\pi/4)=\sin(\pi/4)=1/\sqrt2$.
La démonstration des autres valeurs est détaillée dans les démonstrations de référence.
\end{exemple}

\exercices

\rubrique{Découverte}
\begin{blocLivre}{ex-6-1}
\begin{exercice}[1]{Convertir des angles}
\begin{questions}
\item Convertir $30^\circ$, $45^\circ$, $120^\circ$ et $270^\circ$ en radians.
\item Convertir $\pi/3$, $5\pi/6$ et $7\pi/4$ en degrés.
\end{questions}
\end{exercice}
\end{blocLivre}
\begin{blocLivre}{ex-6-2}
\begin{exercice}[1]{Même point-image}
\begin{questions}
\item Les réels $\pi/3$ et $13\pi/3$ ont-ils le même point-image ?
\item Même question pour $-\pi/2$ et $3\pi/2$, puis pour $\pi/4$ et $5\pi/4$.
\end{questions}
\label{figlivre-r06-02-images-ex}
\figureLivre{r06-02-images}

\end{exercice}
\end{blocLivre}
\begin{blocLivre}{ex-6-3}
\begin{exercice}[1]{Lire les valeurs remarquables}
\begin{questions}
\item Donner les coordonnées des points-images de $0$, $\pi/6$, $\pi/4$ et $\pi/2$.
\item Donner $\cos\pi$, $\sin\pi$, $\cos(3\pi/2)$ et $\sin(3\pi/2)$.
\end{questions}
\label{figlivre-r06-03-valeurs-ex}
\figureLivre{r06-03-valeurs}

\end{exercice}
\end{blocLivre}
\begin{blocLivre}{ex-6-4}
\begin{exercice}[1]{Utiliser une symétrie}
\begin{questions}
\item Calculer exactement $\cos(-\pi/3)$ et $\sin(-\pi/3)$.
\item Calculer $\cos(2\pi/3)$ et $\sin(2\pi/3)$.
\end{questions}
\label{figlivre-r06-04-symetries-ex}
\figureLivre{r06-04-symetries}

\end{exercice}
\end{blocLivre}
\begin{blocLivre}{ex-6-5}
\begin{exercice}[1]{Une longueur d'arc}
Un cercle a pour rayon 6 cm. Calculer la longueur des arcs associés aux angles $\pi/3$ et $3\pi/4$ radians, puis la longueur du cercle entier.
\label{figlivre-r06-05-arcs-ex}
\figureLivre{r06-05-arcs}

\end{exercice}
\end{blocLivre}
\clearpage
\rubrique{Entraînement}
\begin{blocLivre}{ex-6-6}
\begin{exercice}[2]{Réduire des grandes mesures}
\begin{questions}
\item Ramener $17\pi/6$, $-11\pi/4$ et $29\pi/3$ dans $\mathopen{]}-\pi;\pi]$.
\item En déduire le sinus et le cosinus de chacune.
\end{questions}
\end{exercice}
\end{blocLivre}
\begin{blocLivre}{ex-6-7}
\begin{exercice}[2]{Retrouver une coordonnée}
On sait que $\cos t=-3/5$ et que $t\in[\pi/2;\pi]$.\begin{questions}
\item Calculer $\sin^2t$ puis $\sin t$.
\item Déterminer les coordonnées des points-images de $-t$ et $t+\pi$.
\end{questions}
\label{figlivre-r06-07-coordonnee-ex}
\figureLivre{r06-07-coordonnee}

\end{exercice}
\end{blocLivre}
\begin{blocLivre}{ex-6-8}
\begin{exercice}[2]{Deux angles associés}
\begin{questions}
\item Calculer $\cos(5\pi/4)$, $\sin(5\pi/4)$, $\cos(7\pi/6)$ et $\sin(7\pi/6)$.
\item Calculer $A=2\cos(5\pi/4)+\sqrt2\sin(7\pi/6)$.
\end{questions}
\end{exercice}
\end{blocLivre}
\begin{blocLivre}{ex-6-9}
\begin{exercice}[2]{Retrouver des angles sur le cercle}
À l’aide du cercle, déterminer tous les réels $t\in[0;2\pi[$ tels que :\begin{questions}
\item $\cos t=1/2$.
\item $\sin t=-\sqrt2/2$.
\item $\cos t=0$.
\end{questions}

\label{figgeo-06-09-ex}
\figureGeo{06-09}
\end{exercice}
\end{blocLivre}
\begin{blocLivre}{ex-6-10}
\begin{exercice}[2]{Une portion de cercle}
Un disque de rayon 10 cm est partagé en un secteur d’angle $2\pi/5$ et son complément.\begin{questions}
\item Calculer la fraction du disque occupée par le premier secteur.
\item Calculer son aire et son périmètre, en incluant les deux rayons.
\end{questions}

\label{figgeo-06-10-ex}
\figureGeo{06-10}
\end{exercice}
\end{blocLivre}
\begin{blocLivre}{ex-6-11}
\begin{exercice}[2]{Une identité à démontrer}
\begin{questions}
\item Développer $(\cos t+\sin t)^2+(\cos t-\sin t)^2$.
\item En déduire sa valeur pour tout réel $t$.
\item Vérifier le résultat pour $t=\pi/4$ sans calculatrice.
\end{questions}
\end{exercice}
\end{blocLivre}
\clearpage
\rubrique{Problèmes}
\begin{blocLivre}{ex-6-12}
\begin{exercice}[3]{Une grande roue}
Une roue tourne à vitesse de rotation uniforme. Son rayon est de 12 m et son centre est à 14 m du sol. Le point $M$ a pour angle $t$ à partir de l’horizontale positive. Ses coordonnées dans un repère au sol sont $M(12\cos t;14+12\sin t)$.\begin{questions}
\item Donner les positions de $M$ pour $t=0$, $\pi/2$, $\pi$ et $3\pi/2$.
\item Calculer la hauteur pour $t=\pi/6$ et $7\pi/6$.
\item Si un tour dure 48 s, quelle durée correspond à un angle parcouru de $5\pi/6$ ?
\end{questions}

\label{figgeo-06-12-ex}
\figureGeo{06-12}
\end{exercice}
\end{blocLivre}
\begin{blocLivre}{ex-6-13}
\begin{exercice}[3]{Un hexagone régulier}
On inscrit un hexagone régulier dans le cercle unité, un sommet étant $A(1;0)$.\begin{questions}
\item Donner les angles puis les coordonnées de ses six sommets.
\item Calculer la longueur d’un côté avec la distance entre les deux premiers sommets.
\item En déduire le périmètre et l’aire exacte de l’hexagone.
\end{questions}

\label{figgeo-06-13-ex}
\figureGeo{06-13}
\end{exercice}
\end{blocLivre}
\begin{blocLivre}{ex-6-14}
\begin{exercice}[3]{Une zone autorisée}
Un point $M(\cos t;\sin t)$ appartient au cercle unité, avec $t\in[0;2\pi[$.\begin{questions}
\item Représenter la zone où $\sin t\geq0$.
\item Déterminer les angles où $\cos t\leq-1/2$.
\item Déterminer les angles qui satisfont simultanément les deux conditions et calculer la longueur de l’arc correspondant.
\end{questions}

\label{figgeo-06-14-ex}
\figureGeo{06-14}
\end{exercice}
\end{blocLivre}
\begin{blocLivre}{ex-6-15}
\begin{exercice}[3]{Approcher un cercle par des polygones}
Un polygone régulier à $n\geq3$ côtés est inscrit dans le cercle unité. En partageant un triangle central en deux, on obtient un demi-côté de longueur $\sin(\pi/n)$.\begin{questions}
\item Exprimer le périmètre $P_n$ de ce polygone.
\item Calculer exactement $P_6$ et $P_{12}$ en utilisant $\sin(\pi/12)=(\sqrt6-\sqrt2)/4$, fourni ici.
\item Comparer numériquement ces périmètres à $2\pi$.
\item Méthode d’Archimède. On admet que le côté $c_{2n}$ du polygone à $2n$ côtés vérifie $c_{2n}=\sqrt{2-\sqrt{4-c_n^2}}$. En partant de l’hexagone, de côté $c_6=1$, écrire une fonction Python, sans utiliser $\pi$, qui renvoie le demi-périmètre du polygone à $6\times2^k$ côtés. Donner une valeur approchée pour 96 côtés.
\end{questions}

\label{figgeo-06-15-ex}
\figureGeo{06-15}
\end{exercice}
\end{blocLivre}

% !TEX root = ../main.tex
\begin{ouverture}[vecteursplan]{7}{Vecteurs et produit scalaire}{Partie 3 · Géométrie}
{Coordonnées\enspace·\enspace Projection\enspace·\enspace Orthogonalité\enspace·\enspace Longueurs et angles}
\objectif{Choisir une méthode de calcul du produit scalaire.}
\objectif{Démontrer une orthogonalité et calculer un angle.}
\objectif{Développer des expressions vectorielles.}
\objectif{Résoudre des problèmes de longueurs et de lieux géométriques.}
\prerequis{Vecteurs, coordonnées et théorème de Pythagore.}
\prerequis{Cosinus d’un angle et valeurs remarquables.}
\end{ouverture}
\essentiel
\essSection{1}{Calculer avec des vecteurs}
\begin{propriete}
Si $A(x_A;y_A)$ et $B(x_B;y_B)$, alors $\vect{AB}=(x_B-x_A;y_B-y_A)$.
La relation de Chasles donne $\vect{AB}+\vect{BC}=\vect{AC}$.
Pour $u=(x;y)$ et $v=(x';y')$, $u+v=(x+x';y+y')$ et $ku=(kx;ky)$.
Dans un repère orthonormé, $\|u\|=\sqrt{x^2+y^2}$.
Deux vecteurs du plan sont colinéaires si et seulement si $xy'-yx'=0$.
\end{propriete}
\essSection{2}{Définir le produit scalaire}
\begin{definition}
Pour deux vecteurs non nuls $u$ et $v$, formant un angle $\theta$,
$u\cdot v=\|u\|\,\|v\|\cos\theta$. Le produit est nul si l'un des vecteurs est nul.
Le résultat est un nombre réel, qui peut être positif, négatif ou nul.
\end{definition}
\begin{propriete}[Coordonnées et projection]
En base orthonormée, $u\cdot v=xx'+yy'$. Les vecteurs sont orthogonaux
exactement lorsque $u\cdot v=0$. Si $H$ est le projeté orthogonal de $C$
sur $(AB)$, $\vect{AB}\cdot\vect{AC}=AB\times\overline{AH}$,
où la mesure algébrique est orientée de $A$ vers $B$.
\end{propriete}
\begin{center}
\begin{tikzpicture}[x=9mm,y=9mm]
\draw[otNavy,thick,->](0,0)--(6,0)node[right]{$B$};
\draw[otGreenDark,thick,->](0,0)--(2,3)node[above]{$C$};
\draw[otMuted,dashed](2,3)--(2,0)node[below]{$H$};
\draw[otMuted](2,.3)--(2.3,.3)--(2.3,0);
\node[left]at(0,0){$A$};\node[below]at(3,-.3){$\vect{AB}\cdot\vect{AC}=AB\times\overline{AH}$};
\end{tikzpicture}
\par\legende{$\overline{AH}$ est la projection algébrique : elle peut être négative.}
\end{center}

\clearpage
\essSection{3}{Développer et transformer}
\begin{propriete}[Symétrie et bilinéarité]
$u\cdot v=v\cdot u$, $u\cdot(v+w)=u\cdot v+u\cdot w$ et
$(ku)\cdot v=k(u\cdot v)$. En particulier $u\cdot u=\|u\|^2$ et
\[\|u+v\|^2=\|u\|^2+2u\cdot v+\|v\|^2,\]
\[\|u-v\|^2=\|u\|^2-2u\cdot v+\|v\|^2.\]
Ainsi $u\cdot v=\frac12\left(\|u+v\|^2-\|u\|^2-\|v\|^2\right)$ : on peut calculer
un produit scalaire à l'aide de normes seulement.
Dans une base orthonormée $(i;j)$, si $u=xi+yj$, alors $u\cdot i=x$ et $u\cdot j=y$.
\end{propriete}
\begin{theoreme}[Al-Kashi]
Dans un triangle $ABC$, $BC^2=AB^2+AC^2-2AB\,AC\cos\widehat{BAC}$.
Cette formule prolonge Pythagore et permet de calculer une longueur ou un angle.
\end{theoreme}
\begin{savoirfaire}[Choisir une méthode]
Utiliser les coordonnées lorsqu'elles sont connues ; les longueurs et l'angle
lorsqu'ils sont donnés ; la projection dans une configuration rectangle ;
les normes seules lorsque trois longueurs sont connues ;
la bilinéarité dans un parallélogramme. Pour un angle, isoler son cosinus,
puis utiliser une valeur remarquable ou la calculatrice.
\end{savoirfaire}
\begin{propriete}[Un lieu géométrique]
Si $A\neq B$, l'ensemble des points $M$ tels que
$\vect{MA}\cdot\vect{MB}=0$ est le cercle de diamètre $[AB]$,
en incluant $A$ et $B$. La démonstration figure dans les démonstrations de référence.
\end{propriete}
\begin{attention}
Une base doit être orthonormée pour utiliser $xx'+yy'$.
Le produit scalaire ne se confond ni avec une longueur ni avec un vecteur.
Deux vecteurs colinéaires non nuls ne sont pas orthogonaux.
\end{attention}


\clearpage
\essSection{4}{Longueurs, angles et cercle}
Les deux configurations à reconnaître avant de commencer un calcul.
\figureGeo[corrige]{07-cours-angle}
\vspace{5mm}
\figureGeo[corrige]{07-cours-cercle}
\vspace{5mm}
\exercices

\rubrique{Découverte}
\begin{blocLivre}{ex-7-1}
\begin{exercice}[1]{Coordonnées et norme}
On donne $A(1;2)$ et $B(4;6)$.\begin{questions}
\item Calculer $\vect{AB}$ et $\vect{BA}$.
\item Calculer $AB$ et les coordonnées du milieu de $[AB]$.
\end{questions}

\label{figgeo-07-01-ex}
\figureGeo{07-01}
\end{exercice}
\end{blocLivre}
\begin{blocLivre}{ex-7-2}
\begin{exercice}[1]{Un calcul avec les coordonnées}
Dans une base orthonormée, $u=(2;-1)$ et $v=(3;4)$.\begin{questions}
\item Calculer $u\cdot v$, $u\cdot u$ et $v\cdot v$.
\item Donner $\|u\|$ et $\|v\|$.
\end{questions}

\label{figgeo-07-02-ex}
\figureGeo{07-02}
\end{exercice}
\end{blocLivre}
\begin{blocLivre}{ex-7-3}
\begin{exercice}[1]{Longueurs et angle}
Deux vecteurs ont pour normes 5 et 6. Calculer leur produit scalaire lorsque l’angle vaut $\pi/3$, $\pi/2$ puis $2\pi/3$.
\label{figgeo-07-03-ex}
\figureGeo{07-03}
\end{exercice}
\end{blocLivre}
\begin{blocLivre}{ex-7-4}
\begin{exercice}[1]{Orthogonaux ou colinéaires}
On considère $u=(3;2)$, $v=(2;-3)$ et $w=(-6;-4)$.\begin{questions}
\item Quels vecteurs sont orthogonaux parmi les paires $(u,v)$ et $(u,w)$ ?
\item Quelle relation de colinéarité peut-on écrire ?
\end{questions}

\label{figgeo-07-04-ex}
\figureGeo{07-04}
\end{exercice}
\end{blocLivre}
\begin{blocLivre}{ex-7-5}
\begin{exercice}[1]{Lire une projection}
On donne $A(0;0)$, $B(6;0)$ et $C(2;3)$.\begin{questions}
\item Donner le projeté orthogonal $H$ de $C$ sur $(AB)$.
\item Calculer $\vect{AB}\cdot\vect{AC}$ par projection, puis par coordonnées.
\end{questions}

\label{figgeo-07-05-ex}
\figureGeo{07-05}
\end{exercice}
\end{blocLivre}
\clearpage
\rubrique{Entraînement}
\begin{blocLivre}{ex-7-6}
\begin{exercice}[2]{Les normes d'une somme}
On sait que $\|u\|=3$, $\|v\|=4$ et $u\cdot v=6$.\begin{questions}
\item Calculer $\|u+v\|^2$ puis $\|u-v\|^2$.
\item En déduire les deux normes et vérifier l’identité du parallélogramme : la somme des deux carrés vaut $2\|u\|^2+2\|v\|^2$.
\end{questions}

\label{figgeo-07-06-ex}
\figureGeo{07-06}
\end{exercice}
\end{blocLivre}
\begin{blocLivre}{ex-7-7}
\begin{exercice}[2]{Calculer un angle de triangle}
On donne $A(0;0)$, $B(4;0)$ et $C(1;\sqrt3)$.\begin{questions}
\item Calculer $\vect{AB}\cdot\vect{AC}$, $AB$ et $AC$.
\item En déduire l’angle $\widehat{BAC}$.
\item Calculer $BC$ et déterminer la nature du triangle.
\end{questions}

\label{figgeo-07-07-ex}
\figureGeo{07-07}
\end{exercice}
\end{blocLivre}
\begin{blocLivre}{ex-7-8}
\begin{exercice}[2]{Déterminer un paramètre}
Les vecteurs $u_m=(m;2)$ et $v=(3;-1)$ sont exprimés dans une base orthonormée.\begin{questions}
\item Pour quelle valeur de $m$ sont-ils orthogonaux ?
\item Pour quelle valeur sont-ils colinéaires ?
\item Pour $m=0$, calculer le cosinus de leur angle.
\end{questions}

\label{figgeo-07-08-ex}
\figureGeo{07-08}
\end{exercice}
\end{blocLivre}
\begin{blocLivre}{ex-7-9}
\begin{exercice}[2]{Les diagonales d'un rectangle}
$ABCD$ est un rectangle avec $AB=6$ et $AD=4$. On pose $u=\vect{AB}$ et $v=\vect{AD}$.\begin{questions}
\item Exprimer $\vect{AC}$ et $\vect{BD}$ avec $u$ et $v$.
\item Calculer $\vect{AC}\cdot\vect{BD}$ et $\vect{AB}\cdot\vect{AC}$.
\item Les diagonales sont-elles perpendiculaires ?
\end{questions}

\label{figgeo-07-09-ex}
\figureGeo{07-09}
\end{exercice}
\end{blocLivre}
\begin{blocLivre}{ex-7-10}
\begin{exercice}[2]{Un parallélogramme non rectangle}
$ABCD$ est un parallélogramme avec $AB=3$, $AD=4$ et $\widehat{BAD}=\pi/3$.\begin{questions}
\item Calculer $\vect{AB}\cdot\vect{AD}$.
\item Calculer les longueurs des deux diagonales.
\item Comparer leur somme de carrés à $2AB^2+2AD^2$.
\end{questions}

\label{figgeo-07-10-ex}
\figureGeo{07-10}
\end{exercice}
\end{blocLivre}
\begin{blocLivre}{ex-7-11}
\begin{exercice}[2]{Projeter sans coordonnées}
Un triangle vérifie $AB=8$, $AC=5$ et $\widehat{BAC}=\pi/3$. On projette $C$ en $H$ sur $(AB)$.\begin{questions}
\item Calculer le produit $\vect{AB}\cdot\vect{AC}$.
\item En déduire $AH$, puis $CH$.
\item Calculer $BC$ par Al-Kashi.
\end{questions}

\label{figgeo-07-11-ex}
\figureGeo{07-11}
\end{exercice}
\end{blocLivre}
\clearpage
\rubrique{Problèmes}
\begin{blocLivre}{ex-7-12}
\begin{exercice}[3]{Un triangle et sa hauteur}
On considère $A(1;1)$, $B(5;3)$ et $C(2;4)$.\begin{questions}
\item Montrer que le triangle est rectangle en $C$.
\item Calculer l’angle $\widehat{BAC}$.
\item Déterminer le milieu $H$ de $[AB]$ et vérifier que $CH\perp AB$.
\item Calculer la hauteur $CH$ et l’aire du triangle.
\end{questions}

\label{figgeo-07-12-ex}
\figureGeo{07-12}
\end{exercice}
\end{blocLivre}
\begin{blocLivre}{ex-7-13}
\begin{exercice}[3]{Le cercle de diamètre donné}
On donne $A(-2;0)$, $B(2;0)$ et $M(x;y)$.\begin{questions}
\item Exprimer $\vect{MA}\cdot\vect{MB}$ en fonction de $x$ et $y$.
\item Décrire l’ensemble où ce produit est nul.
\item Décrire les zones où il est négatif ou positif.
\item Situer $P(0;2)$ et $Q(1;1)$ dans ces zones.
\end{questions}

\label{figgeo-07-13-ex}
\figureGeo{07-13}
\end{exercice}
\end{blocLivre}
\begin{blocLivre}{ex-7-14}
\begin{exercice}[3]{Travail d'une force}
Une force constante est représentée par $F=(30;40)$ en newtons. Un déplacement est $d=(6;0)$ en mètres. Le travail est $W=F\cdot d$, en joules.\begin{questions}
\item Calculer $W$ et la norme de la force.
\item Pour un déplacement de 6 m d’angle $\pi/6$ avec l’horizontale, calculer le travail exact.
\item Quelle direction d’un déplacement de norme 6 maximise le travail ? Donner ce maximum.
\end{questions}

\label{figgeo-07-14-ex}
\figureGeo{07-14}
\end{exercice}
\end{blocLivre}
\begin{blocLivre}{ex-7-15}
\begin{exercice}[3]{Une identité avec le milieu}
$A$ et $B$ sont distincts et $I$ est le milieu de $[AB]$.\begin{questions}
\item En écrivant $\vect{MA}=\vect{MI}+\vect{IA}$ et $\vect{MB}=\vect{MI}-\vect{IA}$, démontrer $\vect{MA}\cdot\vect{MB}=MI^2-AB^2/4$.
\item Pour $A(0;0)$ et $B(6;0)$, déterminer le lieu des points tels que $\vect{MA}\cdot\vect{MB}=5$.
\item Quel est le minimum possible du produit ? Où est-il atteint ?
\end{questions}

\label{figgeo-07-15-ex}
\figureGeo{07-15}
\end{exercice}
\end{blocLivre}

% !TEX root = ../main.tex
\begin{ouverture}[cercleplan]{8}{Géométrie repérée}{Partie 3 · Géométrie}
{Droites\enspace·\enspace Vecteurs normaux\enspace·\enspace Projections\enspace·\enspace Cercles}
\objectif{Déterminer une équation cartésienne de droite.}
\objectif{Calculer les coordonnées d’un projeté orthogonal.}
\objectif{Reconnaître et utiliser une équation de cercle.}
\objectif{Résoudre des problèmes d’intersection et de configuration.}
\prerequis{Repère orthonormé, distances et milieux.}
\prerequis{Produit scalaire et résolution du second degré.}
\end{ouverture}
\essentiel
\essSection{1}{Équation d'une droite et vecteur normal}
\begin{definition}
Un vecteur non nul $n=(a;b)$ est normal à une droite $d$ s'il est orthogonal
à ses vecteurs directeurs. Si $A(x_A;y_A)$ appartient à $d$, alors
\[M(x;y)\in d\iff a(x-x_A)+b(y-y_A)=0.\]
Une équation $ax+by+c=0$, avec $(a;b)\neq(0;0)$, définit une droite de
vecteur normal $(a;b)$ et de vecteur directeur $(-b;a)$.
\end{definition}
\begin{propriete}
Deux droites sont parallèles si leurs vecteurs normaux sont colinéaires.
Elles sont perpendiculaires si leurs vecteurs normaux sont orthogonaux.
Deux équations proportionnelles peuvent représenter la même droite.
\end{propriete}
\essSection{2}{Projeter un point sur une droite}
\begin{savoirfaire}
Pour projeter $A$ sur $d$, on cherche $H\in d$ tel que $\vect{AH}$ soit
colinéaire à un vecteur normal de $d$. Cela donne deux équations.
Pour $d:ax+by+c=0$, on peut écrire $H=A-\lambda(a;b)$, avec
\[\lambda=\frac{ax_A+by_A+c}{a^2+b^2}.\]
On vérifie ensuite $H\in d$ et $AH\perp d$.
\end{savoirfaire}
\begin{exemple}
Pour $A(3;1)$ et $d:y=x$, le normal $(1;-1)$ donne $\lambda=1$,
donc $H(2;2)$. On a bien $2=2$ et $\vect{AH}=(-1;1)$, normal à $d$.
\end{exemple}
\begin{attention}
Une équation cartésienne est valable pour une droite verticale, contrairement
à une écriture $y=mx+p$. Dans ce chapitre, tous les repères sont orthonormés.
\end{attention}

\clearpage
\essSection{3}{Équation d'un cercle}
\begin{theoreme}
Le cercle de centre $\Omega(\alpha;\beta)$ et de rayon $R>0$ a pour équation
\[(x-\alpha)^2+(y-\beta)^2=R^2.\]
En effet cette égalité exprime $\Omega M^2=R^2$.
\end{theoreme}
\begin{savoirfaire}[Reconnaître un cercle]
Compléter les carrés. Par exemple,
$x^2+y^2-2x-2y-2=0$ devient $(x-1)^2+(y-1)^2=4$.
Le centre est $(1;1)$ et le rayon vaut 2. Un second membre négatif donne
un ensemble vide ; un second membre nul donne un seul point.
\end{savoirfaire}
\begin{center}
\begin{tikzpicture}[scale=.7]
\draw[otMuted,->](-1.6,0)--(4,0)node[right]{$x$};
\draw[otMuted,->](0,-1.6)--(0,3.7)node[above]{$y$};
\draw[otNavy,thick](1,1)circle(2);
\fill[otGreenDark](1,1)circle(1.7pt)node[above]{$\Omega(1;1)$};
\draw[otGreenDark](1,1)--node[above]{$2$}(3,1);
\fill[otNavy](3,1)circle(1.7pt)node[right]{$M(3;1)$};
\end{tikzpicture}
\par\legende{Le cercle $(x-1)^2+(y-1)^2=4$.}
\end{center}

\begin{propriete}[Tangente et intersections]
Au point $A$ d'un cercle de centre $\Omega$, la tangente est perpendiculaire
à $\Omega A$. Pour chercher une intersection avec une droite, remplacer
une coordonnée par son expression dans l'équation du cercle et résoudre.
\end{propriete}


\clearpage
\essSection{4}{Projections, tangentes et intersections}
Vérifier les appartenances et les angles droits dans chaque construction.
\figureGeo[corrige]{08-cours-normal}
\vspace{5mm}
\figureGeo[corrige]{08-cours-positions}
\vspace{5mm}
\exercices

\rubrique{Découverte}
\begin{blocLivre}{ex-8-1}
\begin{exercice}[1]{Lire une équation de droite}
On considère $d:3x-2y+5=0$.\begin{questions}
\item Donner un vecteur normal et un vecteur directeur de $d$.
\item Vérifier si $A(1;4)$ et $B(0;2)$ appartiennent à $d$.
\end{questions}

\label{figgeo-08-01-ex}
\figureGeo{08-01}
\end{exercice}
\end{blocLivre}
\begin{blocLivre}{ex-8-2}
\begin{exercice}[1]{Construire une équation}
Déterminer une équation cartésienne de la droite passant par $A(2;-1)$ et de vecteur normal $(3;4)$. Donner ensuite un vecteur directeur et l’ordonnée à l’origine.
\label{figgeo-08-02-ex}
\figureGeo{08-02}
\end{exercice}
\end{blocLivre}
\begin{blocLivre}{ex-8-3}
\begin{exercice}[1]{Écrire un cercle}
Le cercle $\mathcal C$ a pour centre $(2;-1)$ et pour rayon 3.\begin{questions}
\item Donner son équation.
\item Les points $A(5;-1)$, $B(2;2)$ et $C(0;0)$ sont-ils sur le cercle, à l’intérieur ou à l’extérieur ?
\end{questions}

\label{figgeo-08-03-ex}
\figureGeo{08-03}
\end{exercice}
\end{blocLivre}
\begin{blocLivre}{ex-8-4}
\begin{exercice}[1]{Reconnaître une équation de cercle}
Identifier le centre et le rayon du cercle $x^2+y^2-4x+6y-12=0$. Donner ses points d’intersection avec la droite $x=2$.
\label{figgeo-08-04-ex}
\figureGeo{08-04}
\end{exercice}
\end{blocLivre}
\begin{blocLivre}{ex-8-5}
\begin{exercice}[1]{Deux projections simples}
\begin{questions}
\item Projeter $P(3;4)$ sur l’axe des abscisses et calculer la distance correspondante.
\item Projeter $A(3;1)$ sur $y=x$ et calculer $AH$.
\end{questions}

\label{figgeo-08-05-ex}
\figureGeo{08-05}
\end{exercice}
\end{blocLivre}
\clearpage
\rubrique{Entraînement}
\begin{blocLivre}{ex-8-6}
\begin{exercice}[2]{Projeter et mesurer une distance}
On donne $A(4;3)$ et $d:2x+y-5=0$.\begin{questions}
\item Écrire $H=A-\lambda(2;1)$ et calculer $\lambda$ pour que $H\in d$.
\item Donner $H$ et vérifier l’orthogonalité.
\item Calculer la distance de $A$ à $d$.
\end{questions}

\label{figgeo-08-06-ex}
\figureGeo{08-06}
\end{exercice}
\end{blocLivre}
\begin{blocLivre}{ex-8-7}
\begin{exercice}[2]{Droites parallèles et perpendiculaires}
On donne $d:3x-2y+1=0$, $e:6x-4y-7=0$ et $f:2x+3y-5=0$.\begin{questions}
\item Comparer les directions de $d$ et $e$. Sont-elles confondues ?
\item Montrer que $d\perp f$.
\item Donner une équation de la perpendiculaire à $d$ passant par $(1;2)$.
\end{questions}

\label{figgeo-08-07-ex}
\figureGeo{08-07}
\end{exercice}
\end{blocLivre}
\begin{blocLivre}{ex-8-8}
\begin{exercice}[2]{Intersections avec les axes décalés}
On considère $\mathcal C:(x-1)^2+(y+2)^2=25$.\begin{questions}
\item Calculer les intersections avec $y=1$.
\item Calculer les intersections avec $x=1$.
\item Faire un schéma avec le centre et les quatre points obtenus.
\end{questions}

\label{figgeo-08-08-ex}
\figureGeo{08-08}
\end{exercice}
\end{blocLivre}
\begin{blocLivre}{ex-8-9}
\begin{exercice}[2]{Une famille d'ensembles}
Pour $m\in\R$, étudier l’ensemble $E_m$ d’équation $x^2+y^2-2x+4y+m=0$.\begin{questions}
\item Mettre l’équation sous forme de carrés.
\item Selon $m$, dire si $E_m$ est un cercle, un point ou l’ensemble vide.
\item Donner le rayon pour $m=1$.
\end{questions}

\label{figgeo-08-09-ex}
\figureGeo{08-09}
\end{exercice}
\end{blocLivre}
\begin{blocLivre}{ex-8-10}
\begin{exercice}[2]{Déterminer une tangente}
Le cercle a pour centre $\Omega(1;2)$ et pour rayon 5. On donne $A(4;6)$.\begin{questions}
\item Vérifier que $A$ appartient au cercle.
\item Déterminer la tangente au cercle en $A$.
\item Vérifier que $A$ est le projeté orthogonal de $\Omega$ sur cette tangente.
\end{questions}

\label{figgeo-08-10-ex}
\figureGeo{08-10}
\end{exercice}
\end{blocLivre}
\begin{blocLivre}{ex-8-11}
\begin{exercice}[2]{Une médiatrice}
On donne $A(-2;1)$ et $B(4;3)$.\begin{questions}
\item Calculer le milieu $I$ et le vecteur $\vect{AB}$.
\item Déterminer la médiatrice de $[AB]$.
\item Vérifier par les distances que $C(0;5)$ appartient à cette médiatrice.
\end{questions}

\label{figgeo-08-11-ex}
\figureGeo{08-11}
\end{exercice}
\end{blocLivre}
\clearpage
\rubrique{Problèmes}
\begin{blocLivre}{ex-8-12}
\begin{exercice}[3]{Hauteurs et cercle circonscrit}
On donne $A(0;0)$, $B(6;0)$ et $C(2;4)$.\begin{questions}
\item Déterminer les hauteurs issues de $C$ et de $A$, puis leur intersection $H$.
\item Calculer le pied $K$ de la hauteur issue de $A$.
\item Déterminer le centre $\Omega$ du cercle circonscrit avec deux médiatrices.
\item Écrire l’équation du cercle circonscrit.
\end{questions}

\label{figgeo-08-12-ex}
\figureGeo{08-12}
\end{exercice}
\end{blocLivre}
\begin{blocLivre}{ex-8-13}
\begin{exercice}[3]{Une droite coupe un cercle}
On étudie le cercle $\mathcal C:(x-2)^2+(y-1)^2=5$ et la droite $d:x+y=4$.\begin{questions}
\item Substituer $y=4-x$ et obtenir une équation du second degré.
\item Calculer les deux points d’intersection.
\item Projeter le centre sur $d$ et vérifier que la distance obtenue est inférieure au rayon.
\end{questions}

\label{figgeo-08-13-ex}
\figureGeo{08-13}
\end{exercice}
\end{blocLivre}
\begin{blocLivre}{ex-8-14}
\begin{exercice}[3]{Un cercle et des droites horizontales}
Le cercle $\mathcal C$ a pour centre $(1;3)$ et pour rayon 3. Pour un réel $k$, on considère la droite $d_k:y=k$.\begin{questions}
\item Écrire l’équation qui détermine les abscisses des intersections.
\item Selon $k$, déterminer le nombre de points communs.
\item Donner les points de tangence dans les cas où il existe un seul point commun.
\item Calculer les intersections pour $k=2$.
\end{questions}

\label{figgeo-08-14-ex}
\figureGeo{08-14}
\end{exercice}
\end{blocLivre}
\begin{blocLivre}{ex-8-15}
\begin{exercice}[3]{Une parabole comme lieu géométrique}
On cherche les points $M(x;y)$ équidistants du point $F(0;2)$ et de l’axe des abscisses.\begin{questions}
\item Justifier que la distance à l’axe vaut $|y|$ et écrire l’égalité des carrés des distances.
\item Montrer que le lieu est la parabole $y=x^2/4+1$.
\item Déterminer les points du lieu d’ordonnée 5.
\item Au point d’abscisse 4, déterminer la tangente à cette parabole et vérifier l’égalité des deux distances.
\end{questions}

\label{figgeo-08-15-ex}
\figureGeo{08-15}
\end{exercice}
\end{blocLivre}

% !TEX root = ../main.tex
\begin{ouverture}[conditionnelles]{9}{Probabilités conditionnelles}{Partie 4 · Probabilités et statistiques}
{Arbres\enspace·\enspace Conditionnement\enspace·\enspace Indépendance\enspace·\enspace Épreuves répétées}
\objectif{Calculer une probabilité conditionnelle dans un tableau ou un arbre.}
\objectif{Utiliser une partition et les probabilités totales.}
\objectif{Distinguer indépendance et incompatibilité.}
\objectif{Calculer des probabilités sur deux, trois ou quatre épreuves.}
\prerequis{Probabilités sur un univers fini et évènement contraire.}
\prerequis{Fractions, proportions et lecture de tableaux.}
\end{ouverture}
\essentiel
\essSection{1}{Conditionner une probabilité}
\begin{definition}
Si $P(A)>0$, la probabilité de $B$ sachant $A$ est
$P_A(B)=P(A\cap B)/P(A)$. On se restreint à la sous-population $A$.
Ainsi $P(A\cap B)=P(A)P_A(B)$ et $P_A(\overline B)=1-P_A(B)$.
\end{definition}
\begin{attention}
$P_A(B)$ et $P_B(A)$ ne sont pas interchangeables : le groupe de référence
change. Une probabilité conditionnelle n'est pas une probabilité d'intersection.
\end{attention}
\essSection{2}{Arbres et probabilités totales}
\begin{theoreme}
Si $A_1,\ldots,A_k$ forment une partition de l'univers et ont des probabilités
strictement positives, alors
$P(B)=\sum_iP(A_i)P_{A_i}(B)$.
Pour $A$ et son contraire, on obtient
$P(B)=P(A)P_A(B)+P(\overline A)P_{\overline A}(B)$.
\end{theoreme}
\begin{center}
\begin{tikzpicture}[otArbre]
\node {} child {node {$A$} child {node {$B$} edge from parent node[otPoids] {$0{,}2$}}
child {node {$\overline B$} edge from parent node[otPoids] {$0{,}8$}}
edge from parent node[otPoids] {$0{,}6$}}
child {node {$\overline A$} child {node {$B$} edge from parent node[otPoids] {$0{,}5$}}
child {node {$\overline B$} edge from parent node[otPoids] {$0{,}5$}}
edge from parent node[otPoids] {$0{,}4$}};
\end{tikzpicture}
\par\legende{Les poids issus d’un même nœud ont pour somme 1.}
\end{center}

\begin{savoirfaire}
Multiplier les poids le long d'un chemin pour calculer une intersection.
Additionner les probabilités des chemins incompatibles qui réalisent l'évènement.
Pour inverser le conditionnement, revenir à la définition avec l'intersection.
\end{savoirfaire}

\clearpage
\essSection{3}{Indépendance et incompatibilité}
\begin{definition}
Deux évènements $A$ et $B$ sont indépendants si $P(A\cap B)=P(A)P(B)$.
Lorsque $P(A)>0$, cela équivaut à $P_A(B)=P(B)$.
Ils sont incompatibles si $A\cap B=\varnothing$.
\end{definition}
\begin{propriete}
Si $A$ et $B$ sont indépendants, $A$ et $\overline B$ le sont aussi.
En effet $P(A\cap\overline B)=P(A)-P(A\cap B)=P(A)(1-P(B))$.
Deux évènements incompatibles de probabilités non nulles ne sont pas indépendants.
\end{propriete}
\essSection{4}{Répéter des épreuves indépendantes}
\begin{definition}
Une épreuve de Bernoulli a deux issues : succès de probabilité $p$ et échec
de probabilité $1-p$. Pour des répétitions indépendantes, le poids d'un chemin
est le produit des poids de ses branches. Dans ce livre, les arbres de
répétition sont étudiés jusqu'à quatre épreuves, comme dans le programme.
\end{definition}
\begin{exemple}
Sur trois épreuves, le chemin succès, échec, succès a pour probabilité
$p(1-p)p=p^2(1-p)$. Il y a trois chemins avec exactement deux succès :
$SSE$, $SES$, $ESS$. Leur probabilité totale est $3p^2(1-p)$.
\end{exemple}
\begin{attention}
L'indépendance doit être donnée ou justifiée dans le modèle. Un tirage sans
remise modifie en général les probabilités de l'épreuve suivante.
Les fréquences observées dans un tableau ne démontrent pas une relation de cause à effet.
\end{attention}

\exercices

\rubrique{Découverte}
\begin{blocLivre}{ex-9-1}
\begin{exercice}[1]{Appliquer la définition}
On sait $P(A)=0{,}4$ et $P(A\cap B)=0{,}1$.\begin{questions}
\item Calculer $P_A(B)$ et $P_A(\overline B)$.
\item Calculer $P(A\cap\overline B)$.
\end{questions}
\end{exercice}
\end{blocLivre}
\begin{blocLivre}{ex-9-2}
\begin{exercice}[1]{Lire un tableau d'effectifs}
Parmi 200 élèves, on compte 120 filles. Une option est suivie par 50 filles et 30 garçons. On note $F$ : « être une fille » et $A$ : « suivre l’option ».\begin{questions}
\item Compléter un tableau croisant le sexe et l’option.
\item Calculer $P(F)$, $P_F(A)$ et $P_A(F)$.
\end{questions}
\label{figlivre-r09-02-ex}
\figureLivre{r09-02}

\end{exercice}
\end{blocLivre}
\begin{blocLivre}{ex-9-3}
\begin{exercice}[1]{Deux évènements indépendants}
Les évènements $A$ et $B$ sont indépendants, avec $P(A)=0{,}3$ et $P(B)=0{,}5$.\begin{questions}
\item Calculer $P(A\cap B)$ et $P(A\cup B)$.
\item Calculer $P_A(B)$ et $P(A\cap\overline B)$.
\end{questions}
\end{exercice}
\end{blocLivre}
\begin{blocLivre}{ex-9-4}
\begin{exercice}[1]{Incompatibles ne signifie pas indépendants}
Les évènements $A$ et $B$ sont incompatibles, avec $P(A)=0{,}2$ et $P(B)=0{,}3$.\begin{questions}
\item Calculer $P(A\cap B)$ et $P_A(B)$.
\item Sont-ils indépendants ? Justifier avec le critère du produit.
\end{questions}
\end{exercice}
\end{blocLivre}
\begin{blocLivre}{ex-9-5}
\begin{exercice}[1]{Exploiter un arbre}
On donne $P(A)=0{,}6$, $P_A(B)=0{,}2$ et $P_{\overline A}(B)=0{,}5$.\begin{questions}
\item Reproduire l’arbre pondéré et compléter les poids manquants.
\item Calculer $P(A\cap B)$, $P(B)$ et $P_B(A)$.
\end{questions}
\label{figlivre-r09-05-ex}
\figureLivre{r09-05}

\end{exercice}
\end{blocLivre}
\clearpage
\rubrique{Entraînement}
\begin{blocLivre}{ex-9-6}
\begin{exercice}[2]{Deux ateliers de fabrication}
Un atelier A produit 70\,\% des pièces, un atelier B les autres. Les taux de défaut sont respectivement 2\,\% et 5\,\%. Une pièce est choisie au hasard.\begin{questions}
\item Calculer la probabilité $P(D)$ qu’elle soit défectueuse.
\item Sachant qu’elle est défectueuse, calculer la probabilité qu’elle provienne de B.
\end{questions}
\label{figlivre-r09-06-ex}
\figureLivre{r09-06}

\end{exercice}
\end{blocLivre}
\begin{blocLivre}{ex-9-7}
\begin{exercice}[2]{Une partition en trois groupes}
Trois machines A, B et C produisent respectivement 50\,\%, 30\,\% et 20\,\% des pièces. Les taux de défaut sont 1\,\%, 2\,\% et 4\,\%.\begin{questions}
\item Calculer la probabilité de défaut.
\item Parmi les pièces défectueuses, quelle proportion provient de C ?
\item Quelle machine contribue le plus au nombre total de défauts ?
\end{questions}
\label{figlivre-r09-07-ex}
\figureLivre{r09-07}

\end{exercice}
\end{blocLivre}
\begin{blocLivre}{ex-9-8}
\begin{exercice}[2]{Inverser un conditionnement}
Dans un modèle fictif, un caractère concerne 1\,\% d’une population. Un test est positif pour 95\,\% des personnes ayant ce caractère et pour 4\,\% des autres.On note $C$ : « avoir le caractère », et $T$ : « avoir un test positif ».\begin{questions}
\item Pour une population théorique de 10000 personnes, calculer les effectifs positifs dans les deux groupes.
\item Calculer la probabilité qu’un individu positif ait le caractère.
\item Expliquer pourquoi cette probabilité ne vaut pas 95\,\%.
\end{questions}
\label{figlivre-r09-08-ex}
\figureLivre{r09-08}

\end{exercice}
\end{blocLivre}
\begin{blocLivre}{ex-9-9}
\begin{exercice}[2]{Deux épreuves de Bernoulli}
On répète deux fois indépendamment une épreuve de succès de probabilité $0{,}4$.\begin{questions}
\item Dessiner l’arbre et donner les probabilités des quatre chemins.
\item Calculer la probabilité d’exactement un succès, puis d’au moins un succès.
\end{questions}
\label{figlivre-r09-09-ex}
\figureLivre{r09-09}

\end{exercice}
\end{blocLivre}
\begin{blocLivre}{ex-9-10}
\begin{exercice}[2]{Tirer sans remise}
Une urne contient 3 boules rouges et 2 bleues. On tire successivement deux boules sans remise.\begin{questions}
\item Donner les probabilités des chemins $RR$, $RB$, $BR$ et $BB$.
\item Calculer la probabilité de deux couleurs différentes.
\item Les évènements « première rouge » et « deuxième rouge » sont-ils indépendants ?
\end{questions}
\label{figlivre-r09-10-ex}
\figureLivre{r09-10}

\end{exercice}
\end{blocLivre}
\begin{blocLivre}{ex-9-11}
\begin{exercice}[2]{Tirer avec remise}
On reprend l’urne de 3 rouges et 2 bleues, mais la première boule est remise avant le deuxième tirage.\begin{questions}
\item Justifier l’indépendance dans ce modèle.
\item Calculer la probabilité de deux rouges, puis de deux couleurs différentes.
\item Comparer ces résultats au tirage sans remise.
\end{questions}
\label{figlivre-r09-11-ex}
\figureLivre{r09-11}

\end{exercice}
\end{blocLivre}
\clearpage
\rubrique{Problèmes}
\begin{blocLivre}{ex-9-12}
\begin{exercice}[3]{Trois tentatives indépendantes}
Une tentative réussit avec probabilité $0{,}7$. Les trois tentatives sont indépendantes.\begin{questions}
\item Construire l’arbre à trois niveaux et lister les trois chemins avec exactement deux succès.
\item Calculer la probabilité d’exactement deux succès.
\item Calculer la probabilité d’au moins un succès.
\item Calculer la probabilité que le premier succès ait lieu à la troisième tentative.
\end{questions}
\label{figlivre-r09-12-ex}
\figureLivre{r09-12}

\end{exercice}
\end{blocLivre}
\begin{blocLivre}{ex-9-13}
\begin{exercice}[3]{Quatre essais et six chemins}
On effectue quatre épreuves de Bernoulli indépendantes, de probabilité de succès $1/4$.\begin{questions}
\item Lister les six chemins ayant exactement deux succès.
\item Calculer leur probabilité totale.
\item Calculer la probabilité d’au moins un succès.
\item Calculer la probabilité d’au plus un succès en comptant les chemins.
\end{questions}
\label{figlivre-r09-13-ex}
\figureLivre{r09-13}

\end{exercice}
\end{blocLivre}
\begin{blocLivre}{ex-9-14}
\begin{exercice}[3]{Retrouver une proportion inconnue}
Un groupe A représente une proportion $x\in]0;1[$ d’une population. Un évènement B concerne 10\,\% du groupe A et 40\,\% du reste de la population. On sait $P(B)=0{,}22$.\begin{questions}
\item Écrire une équation sur $x$ et la résoudre.
\item Calculer $P_B(A)$.
\item Les évènements A et B sont-ils indépendants ? Justifier de deux façons.
\end{questions}
\label{figlivre-r09-14-ex}
\figureLivre{r09-14}

\end{exercice}
\end{blocLivre}
\begin{blocLivre}{ex-9-15}
\begin{exercice}[3]{Comparer des taux sans se tromper}
Deux centres proposent une formation. Dans A, 90 personnes sur 100 réussissent le parcours facile et 20 sur 100 le parcours difficile. Dans B, 19 sur 20 réussissent le parcours facile et 18 sur 60 le parcours difficile.\begin{questions}
\item Comparer les taux de réussite dans chaque parcours.
\item Calculer le taux global dans chaque centre.
\item Calculer la proportion de parcours faciles dans chaque centre.
\item Expliquer le renversement de la comparaison globale. Peut-on en déduire à lui seul un effet causal du centre ?
\end{questions}
\end{exercice}
\end{blocLivre}

% !TEX root = ../main.tex
\begin{ouverture}[sommes]{10}{Variables aléatoires}{Partie 4 · Probabilités et statistiques}
{Lois\enspace·\enspace Espérance\enspace·\enspace Variance\enspace·\enspace Linéarité}
\objectif{Déterminer la loi d’une variable aléatoire réelle.}
\objectif{Calculer et interpréter une espérance et une dispersion.}
\objectif{Utiliser la linéarité et la formule de König-Huygens.}
\objectif{Étudier un gain net et un choix de prévision.}
\prerequis{Univers fini, probabilités et évènements.}
\prerequis{Moyennes pondérées et calculs de carrés.}
\end{ouverture}
\essentiel
\essSection{1}{Associer un nombre à chaque issue}
\begin{definition}
Sur un univers fini $\Omega$, une variable aléatoire réelle est une fonction
$X:\Omega\to\R$. L'évènement $\{X=a\}$ rassemble les issues qui donnent la valeur $a$.
La loi de $X$ liste ses valeurs distinctes $x_i$ et les probabilités
$p_i=P(X=x_i)$, toutes positives ou nulles et de somme 1.
\end{definition}
\begin{exemple}
Pour deux lancers indépendants d'une pièce équilibrée, le nombre de faces
vaut 0, 1 ou 2, avec les probabilités $1/4$, $1/2$ et $1/4$.
$P(X\leq1)=P(X=0)+P(X=1)=3/4$.
\end{exemple}
\begin{center}
\begin{tikzpicture}\begin{axis}[otLivreGraphe,width=101mm,height=51mm,xmin=-.5,xmax=2.7,ymin=0,ymax=.65,xtick={0,1,2},ytick={.25,.5},xlabel={$x_i$},ylabel={$P(X=x_i)$}]
\addplot[ycomb,otNavy,thick,mark=*,mark size=2pt]coordinates{(0,.25)(1,.5)(2,.25)};
\end{axis}\end{tikzpicture}
\par\legende{Une loi discrète se représente par des masses aux valeurs possibles.}
\end{center}

\essSection{2}{Calculer une espérance}
\begin{definition}
L'espérance est la moyenne théorique pondérée : $E(X)=\sum_i x_ip_i$.
Elle peut être une valeur que $X$ ne prend jamais. Pour un gain net, un jeu
est équitable si son espérance est nulle.
\end{definition}
\begin{attention}
L'espérance n'est pas une promesse sur une réalisation. Pour établir la loi,
regrouper toutes les issues qui produisent la même valeur.
\end{attention}

\clearpage
\essSection{3}{Mesurer la dispersion}
\begin{definition}
Avec $\mu=E(X)$, la variance est $V(X)=E((X-\mu)^2)=\sum_i(x_i-\mu)^2p_i$.
L'écart type est $\sigma(X)=\sqrt{V(X)}$ ; il s'exprime dans la même unité que $X$.
La variance est positive ou nulle et vaut 0 pour une variable constante.
\end{definition}
\begin{theoreme}[König-Huygens]
$V(X)=E(X^2)-E(X)^2$, avec $E(X^2)=\sum_i x_i^2p_i$.
La preuve repose sur le développement du carré et figure dans les démonstrations de référence.
\end{theoreme}
\essSection{4}{Transformer et additionner des variables}
\begin{propriete}[Linéarité de l'espérance]
Pour deux variables définies sur le même univers et des réels $a,b$,
$E(aX+bY)=aE(X)+bE(Y)$. L'indépendance n'est pas nécessaire.
Pour une constante $c$, $E(c)=c$ et donc $E(aX+c)=aE(X)+c$.
\end{propriete}
\begin{propriete}
$V(aX+c)=a^2V(X)$ et $\sigma(aX+c)=|a|\sigma(X)$.
Ajouter une constante décale les valeurs sans changer la dispersion.
Pour une variable de Bernoulli de paramètre $p$, $E(X)=p$ et $V(X)=p(1-p)$.
\end{propriete}
\begin{savoirfaire}[Vérifier un calcul]
Contrôler que les probabilités totalisent 1. Calculer d'abord $E(X)$ et $E(X^2)$,
puis la variance. Vérifier que celle-ci est positive ou nulle.
Conserver les fractions exactes jusqu'à l'arrondi final.
\end{savoirfaire}

\exercices

\rubrique{Découverte}
\begin{blocLivre}{ex-10-1}
\begin{exercice}[1]{Une variable indicatrice}
On lance un dé équilibré. La variable $X$ vaut 1 si le résultat est pair et 0 sinon.\begin{questions}
\item Donner la loi de $X$.
\item Calculer $E(X)$ et $P(X\leq0)$.
\end{questions}
\end{exercice}
\end{blocLivre}
\begin{blocLivre}{ex-10-2}
\begin{exercice}[1]{Lire et exploiter une loi}
$X$ prend les valeurs $-1$, 0 et 2 avec les probabilités $0{,}2$, $0{,}5$ et $0{,}3$.\begin{questions}
\item Calculer $P(X\leq0)$ et $P(X>1)$.
\item Calculer $E(X)$ et $E(X^2)$.
\end{questions}
\end{exercice}
\end{blocLivre}
\begin{blocLivre}{ex-10-3}
\begin{exercice}[1]{Compléter une loi}
Une variable prend les valeurs 0, 1, 2 et 3 avec les probabilités $0{,}1$, $0{,}2$, $p$ et $0{,}4$.\begin{questions}
\item Déterminer $p$.
\item Calculer l’espérance puis la variance.
\end{questions}
\end{exercice}
\end{blocLivre}
\begin{blocLivre}{ex-10-4}
\begin{exercice}[1]{Une loi de Bernoulli}
$X$ vaut 1 avec probabilité $0{,}3$ et 0 avec probabilité $0{,}7$.\begin{questions}
\item Calculer $E(X)$ et $V(X)$ à partir des définitions.
\item En déduire l’écart type, à $10^{-3}$ près.
\end{questions}
\end{exercice}
\end{blocLivre}
\begin{blocLivre}{ex-10-5}
\begin{exercice}[1]{Un jeu équitable}
Un dé équilibré donne un gain net de 5 euros si le 6 sort et une perte de 1 euro sinon. On note $G$ le gain net.\begin{questions}
\item Donner la loi et l’espérance de $G$.
\item Calculer sa variance. Que signifie l’équité du jeu ?
\end{questions}
\end{exercice}
\end{blocLivre}
\clearpage
\rubrique{Entraînement}
\begin{blocLivre}{ex-10-6}
\begin{exercice}[2]{Calculer la dispersion de deux façons}
On reprend $X$ de valeurs $-1$, 0 et 2, de probabilités $1/5$, $1/2$ et $3/10$.\begin{questions}
\item Calculer la variance par König-Huygens.
\item Refaire le calcul avec la moyenne des carrés des écarts.
\item Donner l’écart type exact.
\end{questions}
\end{exercice}
\end{blocLivre}
\begin{blocLivre}{ex-10-7}
\begin{exercice}[2]{Un changement d'unité}
On sait $E(X)=2/5$ et $V(X)=31/25$. On définit $Y=3X-2$.\begin{questions}
\item Calculer $E(Y)$, $V(Y)$ et $\sigma(Y)$.
\item Si les valeurs de $X$ sont $-1$, 0 et 2, donner les valeurs de $Y$ et leurs probabilités, celles de $X$ étant $1/5$, $1/2$ et $3/10$.
\end{questions}
\end{exercice}
\end{blocLivre}
\begin{blocLivre}{ex-10-8}
\begin{exercice}[2]{Même moyenne, risques différents}
Le gain A vaut $-2$ ou 2 avec probabilités égales ; le gain B vaut $-1$ ou 1 avec probabilités égales.\begin{questions}
\item Comparer leurs espérances, variances et écarts types.
\item Comparer leurs probabilités d’un gain strictement positif.
\item Expliquer ce que permet de comparer la variance dans cet exemple.
\end{questions}
\end{exercice}
\end{blocLivre}
\begin{blocLivre}{ex-10-9}
\begin{exercice}[2]{Linéarité sans indépendance}
On lance un dé équilibré et on pose $X$ égal au résultat, puis $Y=7-X$.\begin{questions}
\item Donner $E(X)$ et $E(Y)$.
\item Calculer $E(X+Y)$ et $V(X+Y)$.
\item Montrer que $X$ et $Y$ ne sont pas indépendantes en comparant les évènements $X=1$ et $Y=6$.
\end{questions}
\end{exercice}
\end{blocLivre}
\begin{blocLivre}{ex-10-10}
\begin{exercice}[2]{Deux variables sur le même univers}
On lance deux pièces équilibrées indépendantes. $X$ est le nombre de faces et $Y$ vaut 1 si la première pièce donne face, 0 sinon. On pose $Z=X-2Y$.\begin{questions}
\item Construire le tableau des valeurs de $X$, $Y$ et $Z$ sur les quatre issues.
\item Calculer $E(Z)$ par linéarité.
\item Donner la loi de $Z$ et vérifier son espérance et sa variance.
\end{questions}
\end{exercice}
\end{blocLivre}
\begin{blocLivre}{ex-10-11}
\begin{exercice}[2]{Choisir une valeur de prévision}
$X$ prend les valeurs $-1$, 1 et 3 avec probabilités $1/4$, $1/2$ et $1/4$.\begin{questions}
\item Calculer $E(X)$ et $E(X^2)$.
\item Pour un réel $a$, montrer que $E((X-a)^2)=(a-1)^2+2$.
\item Quelle valeur de $a$ minimise cette erreur quadratique moyenne ?
\end{questions}
\end{exercice}
\end{blocLivre}
\clearpage
\rubrique{Problèmes}
\begin{blocLivre}{ex-10-12}
\begin{exercice}[3]{Une loterie et son prix d'entrée}
Parmi 100 tickets équiprobables, 3 rapportent 20 euros, 10 rapportent 5 euros et les autres ne rapportent rien. Le prix d’un ticket est $m$ euros.\begin{questions}
\item Donner la loi du lot brut $L$ et son espérance.
\item Exprimer l’espérance du gain net $G=L-m$ et déterminer le prix équitable.
\item Pour $m=2$, calculer $E(G)$ et $V(G)$.
\item Calculer la probabilité d’un gain net strictement positif pour $m=2$.
\end{questions}
\end{exercice}
\end{blocLivre}
\begin{blocLivre}{ex-10-13}
\begin{exercice}[3]{Un jeu sur trois pièces}
On lance trois pièces équilibrées indépendantes. Le lot brut est de 4 euros par face obtenue. On paie 5 euros pour jouer et on note $X$ le nombre de faces.\begin{questions}
\item Établir la loi de $X$ en comptant les huit issues.
\item Calculer $E(X)$ et $V(X)$.
\item Calculer l’espérance et la variance du gain net $G=4X-5$.
\item Quelle mise rendrait le jeu équitable ? Quelle est la probabilité de gain positif avec la mise de 5 euros ?
\end{questions}
\end{exercice}
\end{blocLivre}
\begin{blocLivre}{ex-10-14}
\begin{exercice}[3]{Un modèle de garantie}
Une entreprise reçoit 15 euros pour une garantie. Une réparation de 500 euros est nécessaire avec probabilité $0{,}02$, et aucune dépense sinon. $G$ désigne le résultat net d’une garantie.\begin{questions}
\item Donner la loi de $G$ et calculer son espérance.
\item Calculer sa variance et son écart type.
\item Pour 100 garanties, quel est le résultat total espéré ? L’indépendance est-elle nécessaire pour cette dernière réponse ?
\end{questions}
\end{exercice}
\end{blocLivre}
\begin{blocLivre}{ex-10-15}
\begin{exercice}[3]{Une aire aléatoire}
Le triangle a pour sommets $A(0;0)$, $B(1;0)$ et $C(0;H)$, où $H$ vaut 1, 2 ou 4 avec probabilités $1/2$, $1/4$ et $1/4$. On note $S$ son aire.\begin{questions}
\item Exprimer $S$ en fonction de $H$ et établir sa loi.
\item Calculer $E(S)$, $V(S)$ et $P(S\geq1)$.
\item Pour un réel $a$, montrer que $E((S-a)^2)=(a-1)^2+3/8$.
\item Quelle aire fixe constitue la meilleure prévision au sens de cette erreur quadratique ?
\end{questions}
\label{figlivre-r10-15-triangle-ex}
\figureLivre{r10-15-triangle}

\end{exercice}
\end{blocLivre}

% !TEX root = ../main.tex
\begin{ouverture}[grandsnombres]{11}{Statistiques et simulations}{Partie 4 · Probabilités et statistiques}
{Indicateurs\enspace·\enspace Échantillons\enspace·\enspace Moyennes observées\enspace·\enspace Monte-Carlo}
\objectif{Lire des graphiques et calculer des indicateurs statistiques.}
\objectif{Simuler une variable et la moyenne d’un échantillon.}
\objectif{Comparer des moyennes observées à une espérance.}
\objectif{Expérimenter un critère de dispersion et une estimation d’aire.}
\prerequis{Moyennes pondérées, fréquences et variables aléatoires.}
\prerequis{Listes et boucles Python ; le chapitre 12 sert de référence.}
\end{ouverture}
\essentiel
\essSection{1}{Décrire une série statistique}
\begin{definition}
Pour des valeurs $x_i$ d'effectifs $n_i$, l'effectif total est $N=\sum_i n_i$
et la fréquence de $x_i$ vaut $f_i=n_i/N$.
La moyenne est $\overline x=\sum_i f_ix_i$ et la variance descriptive est
$v=\sum_i f_i(x_i-\overline x)^2=\sum_i f_ix_i^2-\overline x^2$.
L'écart type descriptif vaut $\sqrt v$ ; on divise ici par $N$, pas par $N-1$.
\end{definition}
\begin{propriete}[Médiane et quartiles]
Dans une série ordonnée, la médiane est la valeur centrale si $N$ est impair,
et la moyenne des deux valeurs centrales si $N$ est pair.
On prend $Q_1$ au rang $\lceil N/4\rceil$ et $Q_3$ au rang $\lceil3N/4\rceil$.
L'étendue est le maximum moins le minimum ; l'écart interquartile vaut $Q_3-Q_1$.
\end{propriete}
\essSection{2}{Lire les bons graphiques}
\begin{savoirfaire}
Un diagramme en barres convient aux valeurs discrètes. Dans un histogramme,
l'aire de chaque rectangle représente l'effectif ou la fréquence ; pour des
classes de largeurs inégales, les hauteurs sont des densités.
Un diagramme en boîte représente minimum, $Q_1$, médiane, $Q_3$, maximum.
\end{savoirfaire}
\begin{center}
\begin{tikzpicture}[x=1cm,y=1cm]
\draw[otMuted,->](0,0)--(7.5,0);
\foreach \x in {1,2,3,4,5,6,7}{\draw[otMuted](\x,.08)--(\x,-.08)node[below,font=\small]{\x};}
\draw[otNavy,thick](1,.8)--(2,.8) (4,.8)--(7,.8);
\draw[otNavy,fill=otNavyTint](2,.4)rectangle(4,1.2);
\draw[otGreenDark,thick](3.5,.4)--(3.5,1.2);
\draw[otNavy](1,.5)--(1,1.1) (7,.5)--(7,1.1);
\end{tikzpicture}
\par\legende{Boîte de minimum 1, quartiles 2 et 4, médiane 3,5, maximum 7.}
\end{center}

\clearpage
\essSection{3}{Simuler une variable et une moyenne}
\begin{propriete}
En Python, \code{random()} renvoie un réel pseudo-aléatoire de $[0;1[$.
L'expression \code{random() < p} simule un succès de probabilité $p$.
Une fonction de simulation peut être répétée pour créer une liste de $n$ valeurs,
puis une moyenne observée \code{sum(L)/len(L)}.
\end{propriete}
\begin{savoirfaire}[Distinguer théorie et observation]
Calculer d'abord l'espérance $\mu$ et l'écart type $\sigma$ du modèle.
Simuler plusieurs échantillons de même taille et comparer leurs moyennes à $\mu$.
Une moyenne observée fluctue ; elle n'est pas obligatoirement égale à $\mu$.
Fixer une graine permet de reproduire une expérience informatique.
\end{savoirfaire}
\essSection{4}{Expérimenter avec de nombreux échantillons}
\begin{definition}[Un critère d'observation]
Pour $N$ échantillons indépendants de taille $n$, on peut compter la proportion
des moyennes $m$ vérifiant $|m-\mu|\leq2\sigma/\sqrt n$.
Le programme demande d'expérimenter ce critère : il ne donne pas une proportion
universelle exacte, et cette proportion n'est pas une probabilité de certitude.
\end{definition}
\begin{exemple}[Monte-Carlo]
Pour estimer une aire dans le carré unité, tirer des points $(x;y)$ uniformes
dans ce carré et compter la proportion de ceux qui appartiennent à la zone.
Pour le quart de disque $x^2+y^2\leq1$, multiplier cette proportion par 4
donne une estimation de $\pi$.
\end{exemple}
\begin{attention}
Un résultat de simulation dépend du tirage et de la taille de l'échantillon.
Une observation numérique, même répétée, ne démontre pas un théorème général.
\end{attention}

\exercices

\rubrique{Découverte}
\begin{blocLivre}{ex-11-1}
\begin{exercice}[1]{D’un tableau aux indicateurs}
Les valeurs 0, 1, 2 et 3 ont pour effectifs 2, 3, 4 et 1.\begin{questions}
\item Calculer l’effectif total et les fréquences.
\item Calculer la moyenne, la médiane et l’étendue.
\end{questions}
\end{exercice}
\end{blocLivre}
\begin{blocLivre}{ex-11-2}
\begin{exercice}[1]{Lire un diagramme en barres}
\begin{center}
\begin{tikzpicture}\begin{axis}[otLivreGraphe,width=96mm,height=40mm,xmin=8,xmax=18,ymin=0,ymax=5,xtick={10,12,14,16},ytick={1,2,3,4},xlabel={valeur},ylabel={effectif},axis lines=left,every axis x label/.style={at={(axis description cs:1,0)},anchor=north east,yshift=-4mm,font=\small,text=otMuted},every axis y label/.style={at={(axis description cs:0,1)},anchor=south west,yshift=1mm,font=\small,text=otMuted}]
\addplot[ycomb,very thick,otNavy,mark=*]coordinates{(10,2)(12,4)(14,3)(16,1)};
\end{axis}\end{tikzpicture}
\end{center}
\begin{questions}
\item Reconstituer les effectifs et calculer la moyenne.
\item Donner la médiane et les quartiles selon la convention du cours.
\end{questions}
\end{exercice}
\end{blocLivre}
\begin{blocLivre}{ex-11-3}
\begin{exercice}[1]{Tracer une boîte}
On donne la série ordonnée $1,2,2,3,4,4,5,7$.\begin{questions}
\item Déterminer la médiane et les deux quartiles.
\item Tracer une boîte et calculer l’écart interquartile.
\end{questions}
\end{exercice}
\end{blocLivre}
\begin{blocLivre}{ex-11-4}
\begin{exercice}[1]{Une fréquence observée}
Une variable de Bernoulli de paramètre $0{,}3$ a fourni l’échantillon $[0,1,0,0,1,1,0,0,0,1]$.\begin{questions}
\item Calculer sa moyenne observée.
\item Calculer l’écart absolu à l’espérance. Cet écart invalide-t-il le modèle ?
\end{questions}
\end{exercice}
\end{blocLivre}
\begin{blocLivre}{ex-11-5}
\begin{exercice}[1]{Comprendre un tirage simulé}
Une fonction Python renvoie 1 si \code{random() < 0.25}, sinon 0.\begin{questions}
\item Quelle loi cherche-t-elle à simuler ?
\item Quelles valeurs peut prendre la moyenne de quatre appels ?
\item La moyenne de quatre appels vaut-elle nécessairement $0{,}25$ ?
\end{questions}
\end{exercice}
\end{blocLivre}
\clearpage
\rubrique{Entraînement}
\begin{blocLivre}{ex-11-6}
\begin{exercice}[2]{Un histogramme de largeurs inégales}
Les classes $[0;10[$, $[10;30[$ et $[30;60]$ ont pour fréquences $0{,}2$, $0{,}5$ et $0{,}3$.\begin{questions}
\item Calculer les hauteurs d’un histogramme dont les aires représentent les fréquences.
\item Comparer les hauteurs et les fréquences. Pourquoi des hauteurs directement égales aux fréquences seraient-elles trompeuses ?
\end{questions}
\end{exercice}
\end{blocLivre}
\begin{blocLivre}{ex-11-7}
\begin{exercice}[2]{Transformer une série}
Une série $X$ a pour moyenne 12 et pour écart type 3. On définit $Y=2X+5$ et $Z=X-12$.\begin{questions}
\item Calculer les moyennes de $Y$ et de $Z$.
\item Calculer leurs variances et écarts types.
\item Interpréter le centrage $Z=X-12$.
\end{questions}
\end{exercice}
\end{blocLivre}
\begin{blocLivre}{ex-11-8}
\begin{exercice}[2]{Une aire sous une parabole}
Dans 400 tirages de points uniformes dans le carré $[0;1]^2$, 137 vérifient $y\leq x^2$. L’aire théorique, fournie ici, vaut $1/3$.\begin{questions}
\item Donner l’estimation de Monte-Carlo.
\item Calculer l’erreur absolue exacte.
\item Écrire le test Python qui reconnaît un point dans la zone.
\end{questions}
\label{figlivre-r11-08-zone-ex}
\figureLivre{r11-08-zone}

\end{exercice}
\end{blocLivre}
\begin{blocLivre}{ex-11-9}
\begin{exercice}[2]{Estimer le nombre pi}
Parmi 1000 points tirés uniformément dans le carré unité, 786 sont dans le quart de disque $x^2+y^2\leq1$.\begin{questions}
\item Donner l’estimation de $\pi$ et son erreur absolue à $10^{-5}$ près.
\item Expliquer le facteur 4.
\item Un second essai de 1000 points doit-il donner exactement le même résultat ?
\end{questions}
\label{figlivre-r11-09-zone-ex}
\figureLivre{r11-09-zone}

\end{exercice}
\end{blocLivre}
\begin{blocLivre}{ex-11-10}
\begin{exercice}[2]{Comparer des moyennes d'échantillons}
$X$ vaut 0 ou 2 avec probabilités égales. Quatre échantillons de taille 5 donnent les listes $[0,0,2,0,2]$, $[2,2,2,0,2]$, $[0,0,0,0,0]$, $[2,0,2,2,0]$.\begin{questions}
\item Calculer $\mu$ et $\sigma$ pour le modèle.
\item Calculer les quatre moyennes.
\item Quelle proportion satisfait $|m-\mu|\leq2\sigma/\sqrt5$ ?
\end{questions}
\end{exercice}
\end{blocLivre}
\begin{blocLivre}{ex-11-11}
\begin{exercice}[2]{Effet de la taille d'échantillon}
On étudie un modèle d’espérance $\mu=1$ et d’écart type $\sigma=2$.\begin{questions}
\item Calculer le seuil $2\sigma/\sqrt n$ pour $n=16$ puis $n=64$.
\item Le critère est-il satisfait pour $m=2{,}2$, $n=16$ ? Et pour $m=1{,}4$, $n=64$ ?
\item Par quel facteur faut-il multiplier $n$ pour diviser le seuil par 3 ?
\end{questions}
\end{exercice}
\end{blocLivre}
\clearpage
\rubrique{Problèmes}
\begin{blocLivre}{ex-11-12}
\begin{exercice}[3]{Simuler une loi à trois valeurs}
$X$ prend les valeurs $-1$, 2 et 5 avec probabilités $0{,}5$, $0{,}3$ et $0{,}2$.\begin{questions}
\item Calculer $\mu$ et $\sigma$.
\item Écrire une fonction de simulation avec un seul appel à \code{random()}.
\item Écrire une fonction qui renvoie la moyenne d’un échantillon de taille $n\geq1$.
\item Expliquer comment comparer les moyennes pour $n=10$, 100 et 1000 sans confondre résultat observé et preuve.
\end{questions}
\end{exercice}
\end{blocLivre}
\begin{blocLivre}{ex-11-13}
\begin{exercice}[3]{Répéter cinq cents échantillons}
Pour une variable de Bernoulli de paramètre $0{,}4$, on simule 500 échantillons de taille 25.\begin{questions}
\item Calculer l’espérance, l’écart type et le seuil $2\sigma/\sqrt{25}$.
\item Proposer un programme qui compte les moyennes satisfaisant le critère et renvoie leur proportion.
\item Le résultat doit-il obligatoirement être 95\,\% ?
\end{questions}
\end{exercice}
\end{blocLivre}
\begin{blocLivre}{ex-11-14}
\begin{exercice}[3]{Comparer deux groupes et leur réunion}
Le groupe A donne les valeurs $4,5,5,6$, le groupe B les valeurs $1,5,5,9$.\begin{questions}
\item Comparer moyennes et variances descriptives.
\item Calculer la moyenne et la variance de la série réunissant les huit valeurs.
\item Donner la médiane et les quartiles de la série réunie.
\item Expliquer ce que la moyenne seule masque ici.
\end{questions}
\end{exercice}
\end{blocLivre}
\begin{blocLivre}{ex-11-15}
\begin{exercice}[3]{Fréquences des lettres d'un texte}
On utilise les textes sans espaces ni accents \code{BANANE} et \code{ANANAS}.\begin{questions}
\item Calculer les fréquences des lettres présentes dans chaque texte.
\item Proposer une liste des fréquences pour les 26 lettres, avec \code{texte.count(c)} et \code{len(texte)}.
\item On choisit uniformément une position de \code{BANANE} et on note $V$ l’indicatrice d’une voyelle. Donner sa loi, son espérance et sa variance.
\end{questions}
\end{exercice}
\end{blocLivre}

% !TEX root = ../main.tex
\begin{ouverture}[listes]{12}{Algorithmique et programmation}{Partie 5 · Méthodes transversales}
{Variables\enspace·\enspace Conditions\enspace·\enspace Boucles\enspace·\enspace Fonctions et listes}
\objectif{Lire et écrire des programmes Python courts.}
\objectif{Créer, parcourir et modifier une liste.}
\objectif{Programmer des termes, des sommes et des seuils.}
\objectif{Décomposer une simulation en fonctions et vérifier ses résultats.}
\prerequis{Calcul numérique et suites.}
\prerequis{Les exercices de simulation utilisent les chapitres 9 à 11.}
\end{ouverture}
\essentiel
\essSection{1}{Variables, tests et fonctions}
\begin{definition}
Une affectation comme \code{u = 3} associe une valeur à un nom. En langage
naturel, on peut la noter $u\leftarrow3$. Le signe \code{==} teste l'égalité.
Les types courants sont les entiers, les nombres décimaux et les booléens.
Une fonction définie avec \code{def} reçoit des arguments et renvoie une
valeur avec \code{return}. L'indentation organise les blocs.
\end{definition}
\begin{propriete}[Choisir une boucle]
La boucle \code{for k in range(n)} effectue $n$ passages, avec $k$ de 0 à $n-1$.
La boucle \code{while condition} continue tant que la condition est vraie.
Il faut actualiser les variables de la condition et étudier l'arrêt.
Un \code{if} exécute un bloc selon un test ; \code{else} traite l'autre cas.
\end{propriete}
\begin{savoirfaire}[Suivre l'état d'un programme]
Construire un tableau avec une ligne par passage et une colonne par variable.
Distinguer la valeur avant le passage et la valeur après l'affectation.
Pour un compteur, préciser ce qu'il compte : termes, itérations ou succès.
\end{savoirfaire}
\begin{center}
\begin{tikzpicture}[node distance=11mm]
\node[otBloc](a){Initialiser $u$ et $n$};
\node[otTest,below=of a](b){Seuil atteint ?};
\node[otBloc,right=24mm of b](c){Actualiser $u$ et $n$};
\node[otBorne,below=of b](d){Renvoyer $n$};
\draw[otFleche](a)--(b);\draw[otFleche](b)--node[right,font=\scriptsize]{oui}(d);
\draw[otFleche](b)--node[above,font=\scriptsize]{non}(c);
\coordinate (reprise) at ([yshift=4mm]b.north);
\draw[otFleche](c.north)|-(reprise);
\fill[otMuted](reprise)circle(.6pt);
\end{tikzpicture}
\par\legende{Une boucle de seuil demande une initialisation et une actualisation cohérentes.}
\end{center}

\clearpage
\essSection{2}{Créer et manipuler des listes}
\begin{definition}
Une liste regroupe des valeurs ordonnées : \code{L = [3, 1, 4]}.
Les indices commencent à 0 ; \code{L[0]} vaut 3 et \code{L[-1]} vaut 4.
\code{len(L)} donne la longueur et \code{sum(L)} la somme.
\end{definition}
\begin{propriete}
\code{L.append(x)} ajoute $x$ à la fin ; \code{L[i] = x} modifie une valeur ;
\code{L.pop(i)} supprime et renvoie la valeur d'indice $i$.
On peut parcourir directement les valeurs avec \code{for x in L}.
On peut aussi parcourir les indices avec \code{for i in range(len(L))}.
\end{propriete}
\begin{savoirfaire}[Compréhension et filtrage]
\code{[k*k for k in range(6)]} crée les carrés de 0 à 5.
\code{[x for x in L if x > 0]} sélectionne les éléments positifs.
Une liste peut servir à stocker les termes d'une suite, des valeurs d'une fonction
ou les résultats d'une simulation.
\end{savoirfaire}
\essSection{3}{Tester et organiser un programme}
\begin{propriete}
Décomposer une tâche en fonctions rend chaque étape contrôlable.
Tester les cas usuels et les cas limites autorisés : $n=0$, liste à un élément,
seuil déjà atteint. Vérifier les hypothèses : dénominateur non nul, liste non vide,
entier positif. Une simulation ne démontre pas une propriété générale.
\end{propriete}
\begin{attention}
Une boucle sur \code{range(n)} ne va pas jusqu'à $n$ inclus. Pour calculer
$u_0,\ldots,u_n$, il y a $n$ transitions mais $n+1$ valeurs à conserver.
Le signe \code{**} désigne une puissance ; \code{\textasciicircum} ne la remplace pas en Python.
\end{attention}

\exercices

\rubrique{Découverte}
\begin{blocLivre}{ex-12-1}
\begin{exercice}[1]{Suivre des affectations}
Exécuter mentalement le code, puis donner les valeurs finales de \code{a} et \code{b}.
\begin{python}
a = 3
b = 5
a = a + b
b = a - b
a = a - b
\end{python}
Quel effet global a ce programme ?
\end{exercice}
\end{blocLivre}
\begin{blocLivre}{ex-12-2}
\begin{exercice}[1]{Tester une condition}
Que renvoie la fonction pour 7, 8 et 0 ? Expliquer le test.
\begin{python}
def traitement(n):
    if n % 2 == 0:
        return n//2
    return 3*n + 1
\end{python}
\end{exercice}
\end{blocLivre}
\begin{blocLivre}{ex-12-3}
\begin{exercice}[1]{Définir une fonction}
Écrire une fonction Python pour $f(x)=x^2-3x+2$, puis calculer \code{f(0)}, \code{f(2)} et \code{f(5)}. Faut-il modifier le nom de la variable dans la formule mathématique ?
\end{exercice}
\end{blocLivre}
\begin{blocLivre}{ex-12-4}
\begin{exercice}[1]{Modifier une liste}
Donner l'état de la liste après chaque ligne et les valeurs finales de \code{len(L)} et \code{sum(L)}.
\begin{python}
L = [3, 1, 4]
L.append(2)
L[1] = 5
L.pop(0)
\end{python}
\end{exercice}
\end{blocLivre}
\begin{blocLivre}{ex-12-5}
\begin{exercice}[1]{Comprendre une liste en compréhension}
Donner les deux listes créées, puis leurs longueurs.
\begin{python}
L = [k*k for k in range(6)]
P = [x for x in L if x % 2 == 0]
\end{python}
\end{exercice}
\end{blocLivre}
\clearpage
\rubrique{Entraînement}
\begin{blocLivre}{ex-12-6}
\begin{exercice}[2]{Calculer une factorielle}
Pour un entier $n\geq0$, écrire une fonction renvoyant $n!$, avec $0!=1$.\begin{questions}
\item Quel produit accumuler à chaque passage ?
\item Calculer les résultats pour 0, 1 et 5.
\item Pourquoi une initialisation à 0 serait-elle fausse ?
\end{questions}
\end{exercice}
\end{blocLivre}
\begin{blocLivre}{ex-12-7}
\begin{exercice}[2]{Accumuler une somme}
Écrire une fonction \code{somme(n)} qui calcule $1+2+\cdots+n$ pour $n\geq0$, par une boucle. Vérifier pour $n=0$ et $n=10$, puis comparer à $n(n+1)/2$.
\end{exercice}
\end{blocLivre}
\begin{blocLivre}{ex-12-8}
\begin{exercice}[2]{Trouver un premier seuil}
Une suite vérifie $u_0=100$ et $u_{n+1}=u_n+30$.\begin{questions}
\item Écrire une boucle qui renvoie le premier rang tel que $u_n\geq500$.
\item Vérifier le résultat avec la formule explicite.
\item Que renverrait le même principe si le seuil était déjà atteint au rang 0 ?
\end{questions}
\end{exercice}
\end{blocLivre}
\begin{blocLivre}{ex-12-9}
\begin{exercice}[2]{Additionner les termes d'une suite}
On veut calculer les $n$ premiers termes de la suite $u_k=2\times3^k$, de $k=0$ à $k=n-1$.\begin{questions}
\item Écrire une fonction de somme par actualisations successives.
\item Calculer le résultat pour $n=4$.
\item Donner une formule et traiter le cas $n=0$.
\end{questions}
\end{exercice}
\end{blocLivre}
\begin{blocLivre}{ex-12-10}
\begin{exercice}[2]{Compter les lettres sans autre collection}
Pour un texte non vide, écrire une fonction qui renvoie la fréquence d’une lettre, en parcourant le texte avec une boucle. Tester sur \code{BANANE} pour A et B.
\end{exercice}
\end{blocLivre}
\begin{blocLivre}{ex-12-11}
\begin{exercice}[2]{Calculer les moments d'une loi}
Les listes \code{valeurs} et \code{probas} ont la même longueur et décrivent une loi de probabilité valide. Écrire une fonction renvoyant espérance et variance, puis la tester sur $[-1,0,2]$ et $[0{,}2,0{,}5,0{,}3]$.
\end{exercice}
\end{blocLivre}
\clearpage
\rubrique{Problèmes}
\begin{blocLivre}{ex-12-12}
\begin{exercice}[3]{Construire la suite de Fibonacci}
On définit $F_0=0$, $F_1=1$ et $F_{n+2}=F_{n+1}+F_n$.\begin{questions}
\item Écrire une fonction qui renvoie la liste de $F_0$ à $F_n$, pour $n\geq0$.
\item Donner cette liste pour $n=7$.
\item Chercher le premier rang où le terme dépasse 1000 et vérifier le terme précédent.
\end{questions}
\end{exercice}
\end{blocLivre}
\begin{blocLivre}{ex-12-13}
\begin{exercice}[3]{Expérimenter avec Syracuse}
À partir d’un entier positif, on le divise par 2 s’il est pair ; sinon on le remplace par $3n+1$.\begin{questions}
\item Donner les termes à partir de 6 jusqu’à 1.
\item Écrire une fonction qui s’arrête à 1 ou après un nombre maximal d’étapes.
\item Pourquoi la borne sur le nombre d’étapes garantit-elle la fin du programme, indépendamment d’une propriété générale de la suite ?
\end{questions}
\end{exercice}
\end{blocLivre}
\begin{blocLivre}{ex-12-14}
\begin{exercice}[3]{Une approximation par Newton}
On cherche une valeur approchée de $\sqrt2$, racine positive de $f(x)=x^2-2$.\begin{questions}
\item Pour $a>0$, montrer que la tangente à la courbe de $f$ au point d’abscisse $a$ coupe l’axe des abscisses au point d’abscisse $(a+2/a)/2$. La méthode de Newton pose donc $x_0=1$ et $x_{k+1}=(x_k+2/x_k)/2$.
\item Écrire une fonction calculant $x_n$.
\item Calculer $x_1$ et $x_2$ exactement.
\item Expérimenter avec $n=5$ et comparer le carré obtenu à 2.
\item Vérifier que les termes calculés restent positifs, ce qui autorise la division.
\end{questions}
\end{exercice}
\end{blocLivre}
\begin{blocLivre}{ex-12-15}
\begin{exercice}[3]{Un mini-projet Monte-Carlo}
On veut estimer simultanément l’aire sous $y=x^2$ dans le carré unité et le nombre $\pi$.\begin{questions}
\item Écrire une fonction qui tire $n\geq1$ points, avec une graine fixée, et renvoie les deux estimations.
\item Expliquer pourquoi on peut réutiliser les mêmes points pour les deux comptes.
\item Proposer des contrôles automatiques simples des valeurs renvoyées.
\end{questions}
\label{figlivre-r12-15-zones-ex}
\figureLivre{r12-15-zones}

\end{exercice}
\end{blocLivre}

% !TEX root = ../main.tex
\begin{ouverture}[ensembles]{13}{Ensembles et raisonnement}{Partie 5 · Méthodes transversales}
{Ensembles\enspace·\enspace Logique\enspace·\enspace Quantificateurs\enspace·\enspace Démonstrations}
\objectif{Utiliser les opérations sur les ensembles et les intervalles.}
\objectif{Distinguer implication, réciproque et contraposée.}
\objectif{Nier une proposition et identifier les rôles des lettres.}
\objectif{Rédiger des preuves par cas, par contraposée ou par l’absurde.}
\prerequis{Ensembles de nombres, calcul littéral et intervalles.}
\prerequis{Parité des entiers et fonctions du second degré.}
\end{ouverture}
\essentiel
\essSection{1}{Ensembles, intervalles et couples}
\begin{definition}
$x\in A$ signifie que $x$ est un élément de $A$ ; $A\subseteq B$ que tout
élément de $A$ appartient à $B$. L'ensemble vide est noté $\varnothing$.
$A\cap B$ contient les éléments communs ; $A\cup B$ ceux de l'un ou de l'autre.
Dans un univers $E$, le complémentaire est $E\setminus A$, aussi noté $\overline A$.
Le cardinal $\operatorname{Card}(A)$ compte les éléments d'un ensemble fini.
\end{definition}
\begin{propriete}
Un couple est ordonné : $(a;b)$ et $(b;a)$ ne sont pas en général identiques.
Le produit cartésien $A\times B$ est l'ensemble des couples $(a;b)$ avec
$a\in A$ et $b\in B$. Si les deux ensembles sont finis,
$\operatorname{Card}(A\times B)=\operatorname{Card}(A)\operatorname{Card}(B)$.
\end{propriete}
\begin{center}
\begin{tikzpicture}[scale=.8]
\draw[otMuted,rounded corners](0,0)rectangle(7,3);
\fill[otNavy,opacity=.12](2.5,1.5)circle(1.2);
\fill[otGreen,opacity=.2](4.2,1.5)circle(1.2);
\draw[otNavy](2.5,1.5)circle(1.2);\draw[otGreenDark](4.2,1.5)circle(1.2);
\node at(2,1.5){$A$};\node at(4.8,1.5){$B$};\node[above]at(3.35,1.5){$A\cap B$};
\node at(.4,2.6){$E$};
\end{tikzpicture}
\par\legende{La réunion inclut les deux disques ; l’intersection est leur zone commune.}
\end{center}

\begin{attention}
Ne pas confondre appartenance et inclusion : $2\in\{2;4\}$ mais
$\{2\}\subseteq\{2;4\}$. Une borne exclue d'un intervalle n'en est pas un élément.
\end{attention}

\clearpage
\essSection{2}{Implication, réciproque et conditions}
\begin{definition}
L'implication « si $P$, alors $Q$ » affirme que $Q$ est vraie dès que $P$
est vraie. Sa réciproque est « si $Q$, alors $P$ ».
Sa contraposée est « si non $Q$, alors non $P$ » ; elle lui est équivalente.
Une équivalence réunit une implication et sa réciproque.
\end{definition}
\begin{propriete}
Dans $P\Rightarrow Q$, $P$ est une condition suffisante pour $Q$ ; $Q$ est
une condition nécessaire pour $P$. Exemple : être multiple de 4 suffit
pour être pair ; être pair est nécessaire pour être multiple de 4.
\end{propriete}
\essSection{3}{Quantifier, nier et démontrer}
\begin{propriete}
La négation de « pour tout $x$, $P(x)$ » est « il existe $x$ pour lequel $P(x)$
est fausse ». La négation de « il existe $x$ tel que $P(x)$ » est
« pour tout $x$, $P(x)$ est fausse ». Nier « $P$ et $Q$ » donne « non $P$ ou non $Q$ ».
Nier « $P$ ou $Q$ » donne « non $P$ et non $Q$ ».
\end{propriete}
\begin{savoirfaire}
Une preuve directe enchaîne des propriétés connues. Une preuve par cas
traite toutes les possibilités. Une preuve par contraposée démontre la
contraposée. Une preuve par l'absurde suppose le contraire de la conclusion
et obtient une contradiction. Un contre-exemple suffit à réfuter une affirmation universelle.
\end{savoirfaire}
\begin{attention}
Une identité est vraie pour toutes les valeurs autorisées ; une équation
demande de chercher des valeurs de l'inconnue. Dans une famille $f_m(x)$,
$x$ peut être une variable et $m$ un paramètre. Un test sur quelques nombres
ne démontre pas une affirmation « pour tout ».
\end{attention}

\exercices

\rubrique{Découverte}
\begin{blocLivre}{ex-13-1}
\begin{exercice}[1]{Appartenance et opérations}
On donne $E=\{1;2;3;4;5;6\}$, $A=\{2;4;6\}$ et $B=\{1;2;3\}$.\begin{questions}
\item Déterminer $A\cap B$, $A\cup B$ et $E\setminus A$.
\item Corriger si nécessaire : $2\subseteq A$ et $\{2\}\in A$.
\end{questions}
\end{exercice}
\end{blocLivre}
\begin{blocLivre}{ex-13-2}
\begin{exercice}[1]{Réunir des intervalles}
On donne $A=[-2;3]$ et $B=]1;5[$.\begin{questions}
\item Déterminer $A\cap B$ et $A\cup B$.
\item Donner le complémentaire de $A$ dans $\R$.
\end{questions}
\label{figlivre-r13-02-intervalles-ex}
\figureLivre{r13-02-intervalles}

\end{exercice}
\end{blocLivre}
\begin{blocLivre}{ex-13-3}
\begin{exercice}[1]{Construire des couples}
On pose $A=\{a;b\}$ et $B=\{0;1;2\}$.\begin{questions}
\item Lister $A\times B$ et donner son cardinal.
\item Le couple $(1;a)$ appartient-il à $A\times B$ ? À $B\times A$ ?
\end{questions}
\end{exercice}
\end{blocLivre}
\begin{blocLivre}{ex-13-4}
\begin{exercice}[1]{Trouver un contre-exemple}
Réfuter les deux affirmations suivantes :\begin{questions}
\item « Pour tout réel $x$, $x^2\geq x$. »
\item « Si $x^2=4$, alors $x=2$. »
\end{questions}
\end{exercice}
\end{blocLivre}
\begin{blocLivre}{ex-13-5}
\begin{exercice}[1]{Nier une phrase quantifiée}
Pour chacune des phrases suivantes, écrire sa négation puis donner la valeur de vérité des deux phrases :\begin{questions}
\item « Pour tout réel $x$, $x^2\geq0$. »
\item « Il existe un réel $x$ tel que $x^2=-1$. »
\end{questions}
\end{exercice}
\end{blocLivre}
\clearpage
\rubrique{Entraînement}
\begin{blocLivre}{ex-13-6}
\begin{exercice}[2]{Réciproque et contraposée}
On considère l’implication « si $x>2$, alors $x^2>4$ », pour $x$ réel.\begin{questions}
\item Justifier cette implication.
\item Écrire sa réciproque et décider si elle est vraie.
\item Écrire sa contraposée et sa valeur de vérité.
\end{questions}
\end{exercice}
\end{blocLivre}
\begin{blocLivre}{ex-13-7}
\begin{exercice}[2]{Nécessaire ou suffisante}
Pour un entier $n$, comparer les propriétés P : « $n$ est multiple de 4 » et Q : « $n$ est pair ».\begin{questions}
\item Quelle implication est toujours vraie ?
\item Compléter : P est une condition \ldots\ pour Q ; Q est une condition \ldots\ pour P.
\item Donner un exemple montrant que l’équivalence est fausse.
\end{questions}
\end{exercice}
\end{blocLivre}
\begin{blocLivre}{ex-13-8}
\begin{exercice}[2]{Variable, inconnue et paramètre}
\begin{questions}
\item Quel est le statut de $x$ dans l’identité $x^2-1=(x-1)(x+1)$ ?
\item Dans l’équation $x^2-1=0$, quel est le rôle de $x$ et quelles sont ses valeurs solutions ?
\item Dans $f_m(x)=x^2-mx+1$, distinguer variable et paramètre.
\end{questions}
\end{exercice}
\end{blocLivre}
\begin{blocLivre}{ex-13-9}
\begin{exercice}[2]{Complémentaires et connecteurs}
Dans $E=\{1;\ldots;8\}$, A est l’ensemble des entiers pairs et B celui des entiers inférieurs ou égaux à 4.\begin{questions}
\item Calculer $E\setminus(A\cup B)$.
\item Calculer $(E\setminus A)\cap(E\setminus B)$ et comparer.
\item Expliquer l’égalité obtenue en utilisant « et » et « ou ».
\end{questions}
\end{exercice}
\end{blocLivre}
\begin{blocLivre}{ex-13-10}
\begin{exercice}[2]{Lire une quantification implicite}
On rappelle : « f est croissante sur I » signifie que pour tous $a,b\in I$ tels que $a<b$, $f(a)\leq f(b)$.\begin{questions}
\item Écrire la négation de cette définition.
\item Montrer que $f(x)=x^2$ n’est pas croissante sur $\R$.
\item Montrer directement qu’elle est croissante sur $[0;+\infty[$.
\end{questions}
\end{exercice}
\end{blocLivre}
\begin{blocLivre}{ex-13-11}
\begin{exercice}[2]{Démontrer par disjonction des cas}
\begin{questions}
\item Montrer que $n(n+1)$ est pair pour tout entier naturel $n$, en distinguant $n$ pair et $n$ impair.
\item Montrer que $|x|\geq x$ pour tout réel $x$, en distinguant le signe de $x$.
\end{questions}
\end{exercice}
\end{blocLivre}
\clearpage
\rubrique{Problèmes}
\begin{blocLivre}{ex-13-12}
\begin{exercice}[3]{Une preuve par contraposée}
On veut démontrer : « si le carré d’un entier est impair, alors cet entier est impair ».\begin{questions}
\item Écrire la contraposée.
\item La démontrer en écrivant un entier pair sous la forme $2k$.
\item Démontrer aussi la réciproque et formuler l’équivalence obtenue.
\end{questions}
\end{exercice}
\end{blocLivre}
\begin{blocLivre}{ex-13-13}
\begin{exercice}[3]{Une preuve par l'absurde}
On veut montrer que $\sqrt2$ n’est pas rationnel. On suppose au contraire $\sqrt2=p/q$, avec $p,q$ entiers positifs sans facteur commun.\begin{questions}
\item En déduire $p^2=2q^2$ puis montrer que $p$ est pair.
\item Écrire $p=2k$ et montrer que $q$ est pair.
\item Identifier la contradiction et conclure.
\end{questions}
\end{exercice}
\end{blocLivre}
\begin{blocLivre}{ex-13-14}
\begin{exercice}[3]{Une transformation sous condition}
Résoudre $\dfrac{x^2-1}{x-1}=2$.\begin{questions}
\item Préciser le domaine et simplifier sous cette condition.
\item Résoudre l’équation simplifiée puis conclure.
\item Expliquer l’erreur d’une résolution qui multiplierait sans conserver la valeur interdite.
\end{questions}
\end{exercice}
\end{blocLivre}
\begin{blocLivre}{ex-13-15}
\begin{exercice}[3]{Quantifier une famille de trinômes}
Pour un réel $m$, on pose $P_m(x)=x^2-2mx+1$.\begin{questions}
\item Écrire $P_m$ sous forme canonique.
\item Déterminer les $m$ tels que « pour tout réel $x$, $P_m(x)\geq0$ ».
\item Déterminer les $m$ tels que « pour tout réel $x$, $P_m(x)>0$ ».
\item Écrire la négation de la phrase de la question 2 et trouver les $m$ correspondants.
\item Présenter oralement, en deux minutes, la distinction entre les questions 2 et 3.
\end{questions}
\end{exercice}
\end{blocLivre}

% !TEX root = ../main.tex
\begin{ouverture}[boucles]{14}{Automatismes et synthèses}{Partie 5 · Méthodes transversales}
{Calcul\enspace·\enspace Pourcentages\enspace·\enspace Lectures graphiques\enspace·\enspace Problèmes transversaux}
\objectif{Entretenir les calculs et les lectures de représentations.}
\objectif{Choisir rapidement une méthode adaptée.}
\objectif{Combiner plusieurs chapitres dans une même résolution.}
\objectif{Rédiger et contrôler une réponse complète.}
\prerequis{Les exercices 1 à 8 réactivent les acquis antérieurs.}
\prerequis{Les synthèses finales mobilisent les chapitres du livre.}
\end{ouverture}
\essentiel
\essSection{1}{Des rituels à répartir sur l'année}
\begin{savoirfaire}
Les exercices courts de ce chapitre sont une banque de réactivation.
Travailler régulièrement quelques questions de thèmes différents, avant
de passer aux problèmes. Les automatismes se construisent tout au long de
l'année ; ils ne constituent pas une leçon à apprendre en une seule séance.
\end{savoirfaire}
\begin{propriete}[Évolutions en pourcentage]
Une évolution de taux $t$ multiplie par $1+t$. Des évolutions successives
multiplient leurs coefficients : le taux global est le produit moins 1.
Pour revenir à la valeur initiale après un coefficient $c>0$, utiliser $1/c$.
Un taux réciproque n'est donc pas en général l'opposé du taux initial.
\end{propriete}
\begin{exemple}
Une hausse de 25\,\% suivie d'une baisse de 20\,\% donne
$1{,}25\times0{,}8=1$ : retour à la valeur initiale.
Après une baisse de 20\,\%, il faut une hausse de $1/0{,}8-1=25\,\%$ pour revenir.
\end{exemple}
\essSection{2}{Calculer efficacement}
\begin{propriete}
Identités : $(a+b)^2=a^2+2ab+b^2$, $(a-b)^2=a^2-2ab+b^2$,
$a^2-b^2=(a-b)(a+b)$. Un produit est nul si au moins un facteur est nul.
Avant de développer, identifier l'objectif : racines, signe, dérivée ou valeur numérique.
\end{propriete}
\begin{attention}
Une fraction ne se simplifie que par des facteurs communs non nuls.
L'ensemble des solutions doit respecter les valeurs interdites et le domaine du modèle.
\end{attention}

\clearpage
\essSection{3}{Lire et produire des représentations}
\begin{propriete}
La droite passant par $A(x_A;y_A)$ et $B(x_B;y_B)$, avec $x_A\neq x_B$,
a pour pente $m=(y_B-y_A)/(x_B-x_A)$. Son équation est $y=mx+p$,
avec $p=y_A-mx_A$. Dans un diagramme circulaire, une fréquence $f$ correspond
à un angle de $360f$ degrés.
\end{propriete}
\begin{savoirfaire}
Sur une courbe, résoudre $f(x)=k$ revient à lire les intersections avec la
droite horizontale $y=k$ ; pour $f(x)\geq k$, lire les portions au-dessus.
Contrôler les unités, l'origine et l'échelle de chaque axe. Dans un histogramme,
les fréquences sont portées par les aires ; dans une boîte, les repères sont les quartiles.
\end{savoirfaire}
\essSection{4}{Rédiger un problème de synthèse}
\begin{savoirfaire}
Définir les objets et le domaine. Écrire la méthode avant le calcul.
Présenter les étapes et vérifier les bornes. Donner une réponse dans les unités
de la question et distinguer valeur exacte, approximation et conjecture.
Une étude de signe, un tableau de variations ou un arbre peut soutenir l'explication.
\end{savoirfaire}
\begin{attention}
Les quatre derniers exercices sont des sujets originaux de synthèse, pas
des reproductions de sujets officiels. Ils mobilisent plusieurs chapitres.
Chercher sans le corrigé, puis reprendre les notions signalées par les erreurs.
\end{attention}

\exercices

\rubrique{Découverte}
\begin{blocLivre}{ex-14-1}
\begin{exercice}[1]{Évolutions et retour au départ}
Un prix de 80 euros augmente de 25\,\%, puis diminue de 20\,\%.\begin{questions}
\item Calculer le prix après chaque évolution et le taux global.
\item Quel taux de hausse compense exactement une baisse de 20\,\% ?
\end{questions}
\end{exercice}
\end{blocLivre}
\begin{blocLivre}{ex-14-2}
\begin{exercice}[1]{Retrouver une valeur et un taux}
\begin{questions}
\item Un article soldé de 30\,\% coûte 84 euros. Quel était son prix initial ?
\item Une quantité passe de 45 à 54. Calculer son taux d’évolution.
\item Calculer le taux global de deux hausses successives de 10\,\% puis 20\,\%.
\end{questions}
\end{exercice}
\end{blocLivre}
\begin{blocLivre}{ex-14-3}
\begin{exercice}[1]{Développer, factoriser, résoudre}
\begin{questions}
\item Développer $(2x-3)^2$.
\item Factoriser $9x^2-25$.
\item Résoudre $(x-2)(x+1)=0$.
\end{questions}
\end{exercice}
\end{blocLivre}
\begin{blocLivre}{ex-14-4}
\begin{exercice}[1]{Signe et intervalles}
\begin{questions}
\item Résoudre $3-2x\geq0$.
\item Résoudre $(x+2)(x-3)<0$.
\item Donner un exemple de solution commune aux deux inéquations.
\end{questions}
\end{exercice}
\end{blocLivre}
\begin{blocLivre}{ex-14-5}
\begin{exercice}[1]{Équations de droites}
\begin{questions}
\item Déterminer la droite passant par $A(1;3)$ et $B(3;7)$.
\item Pour $d:y=-x/2+4$, donner pente, ordonnée à l’origine et intersection avec l’axe des abscisses.
\end{questions}
\end{exercice}
\end{blocLivre}
\clearpage
\rubrique{Entraînement}
\begin{blocLivre}{ex-14-6}
\begin{exercice}[2]{Lire une courbe par morceaux}
La fonction $f$ est représentée par les segments suivants ; les sommets ont pour coordonnées exactes $(-2;1)$, $(0;-1)$, $(2;3)$, $(4;1)$.\begin{center}
\begin{tikzpicture}\begin{axis}[otLivreGraphe,width=101mm,height=44mm,xmin=-2.5,xmax=4.5,ymin=-1.5,ymax=3.5,xtick={-2,-1,0,1,2,3,4},ytick={-1,1,2,3}]
\addplot[otNavy,thick,mark=*]coordinates{(-2,1)(0,-1)(2,3)(4,1)};
\end{axis}\end{tikzpicture}
\end{center}
\begin{questions}
\item Résoudre $f(x)=1$ et $f(x)\geq1$.
\item Donner les variations et les extremums.
\end{questions}
\end{exercice}
\end{blocLivre}
\begin{blocLivre}{ex-14-7}
\begin{exercice}[2]{Des fréquences conditionnelles}
Dans un groupe de 100 personnes, 50 ont moins de 25 ans. Parmi ces 50, 20 utilisent le bus ; parmi les 50 autres, 30 l’utilisent.\begin{questions}
\item Calculer la fréquence globale d’utilisation du bus.
\item Calculer la fréquence d’utilisation parmi les moins de 25 ans.
\item Dans le modèle de tirage uniforme d’une personne, les deux évènements sont-ils indépendants ?
\end{questions}
\end{exercice}
\end{blocLivre}
\begin{blocLivre}{ex-14-8}
\begin{exercice}[2]{Réactiver la trigonométrie}
\begin{questions}
\item Donner les coordonnées du point-image de $5\pi/3$.
\item Comparer ses coordonnées à celles du point-image de $11\pi/3$.
\item Calculer un arc de rayon 3 et d’angle $2\pi/3$.
\end{questions}
\end{exercice}
\end{blocLivre}
\begin{blocLivre}{ex-14-9}
\begin{exercice}[2]{Relier dérivée et extremums}
$f(x)=x^3-6x^2+9x$ sur $[0;4]$.\begin{questions}
\item Factoriser $f\prime(x)$ et étudier son signe.
\item Donner le tableau de variations avec les valeurs aux bornes.
\item Donner toutes les positions des extremums globaux.
\end{questions}
\end{exercice}
\end{blocLivre}
\begin{blocLivre}{ex-14-10}
\begin{exercice}[2]{Un diagramme circulaire}
Lors d’un vote de 80 personnes, les options A, B et C recueillent respectivement 40, 24 et 16 voix.\begin{questions}
\item Calculer les fréquences et les angles des secteurs.
\item Quelle est la différence entre le nombre de voix pour A et pour B ?
\item Peut-on calculer une moyenne numérique des noms A, B et C ?
\end{questions}
\label{figlivre-r14-10-secteurs-ex}
\figureLivre{r14-10-secteurs}

\end{exercice}
\end{blocLivre}
\begin{blocLivre}{ex-14-11}
\begin{exercice}[2]{Une exponentielle dans une étude de fonction}
\begin{questions}
\item Résoudre $e^{2x}=e^{x+3}$.
\item Pour $g(x)=(2x-1)e^{-x}$ sur $[0;3]$, factoriser $g\prime(x)$.
\item Donner le maximum et le minimum de $g$ sur le segment.
\end{questions}
\end{exercice}
\end{blocLivre}
\clearpage
\rubrique{Problèmes}
\begin{blocLivre}{ex-14-12}
\begin{exercice}[3]{Synthèse : une boîte à dimensions liées}
Une boîte a une base carrée de côté $x$ et une hauteur $6-x$, en décimètres, avec $0<x<6$.\begin{questions}
\item Exprimer son volume $V(x)$ et calculer $V\prime(x)$.
\item Déterminer ses dimensions et son volume maximal.
\item Écrire la tangente à la courbe de $V$ en $x=2$.
\item Vérifier $V(x)-16=-(x-2)(x^2-4x-8)$, puis déterminer quand le volume est au moins 16 litres.
\end{questions}
\label{figlivre-r14-12-boite-ex}
\figureLivre{r14-12-boite}

\end{exercice}
\end{blocLivre}
\begin{blocLivre}{ex-14-13}
\begin{exercice}[3]{Synthèse : triangle, cercle et tangente}
On donne $A(0;0)$, $B(4;0)$ et $C(2;2)$.\begin{questions}
\item Montrer que le triangle est rectangle en $C$ et calculer l’angle $\widehat{BAC}$.
\item Déterminer le cercle circonscrit.
\item Écrire sa tangente en $C$ et projeter $M(5;3)$ sur cette tangente.
\item Calculer la longueur du petit arc $AC$ du cercle.
\end{questions}
\label{figlivre-r14-13-geometrie-ex}
\figureLivre{r14-13-geometrie}

\end{exercice}
\end{blocLivre}
\begin{blocLivre}{ex-14-14}
\begin{exercice}[3]{Synthèse : défauts et gain net}
Deux fournisseurs A et B livrent 60\,\% et 40\,\% des articles, avec des taux de défaut de 1\,\% et 4\,\%. Un article est vendu 20 euros, coûte 12 euros et entraîne 30 euros de dépense supplémentaire s’il est défectueux.\begin{questions}
\item Calculer la probabilité $P(D)$ d’un défaut.
\item Sachant un défaut, calculer la probabilité que l’article provienne de B.
\item Donner la loi du gain net G, puis son espérance et sa variance.
\item Pour quatre articles de défauts indépendants et identiquement distribués, calculer la probabilité d’au moins un défaut.
\end{questions}
\label{figlivre-r14-14-ex}
\figureLivre{r14-14}

\end{exercice}
\end{blocLivre}
\begin{blocLivre}{ex-14-15}
\begin{exercice}[3]{Synthèse : deux modèles d'évolution}
Un stock est modélisé par $u_0=1000$ et $u_{n+1}=0{,}9u_n+50$, au bout de $n$ semaines. On propose aussi le modèle continu $W(t)=500+500e^{-0{,}1t}$.\begin{questions}
\item Poser $v_n=u_n-500$ et en déduire une expression de $u_n$.
\item Justifier la décroissance de $u_n$ et montrer que $u_n>500$.
\item Par essais ou Python, trouver le premier rang tel que $u_n\leq600$.
\item Étudier les variations de W et trouver la première semaine entière où $W(n)\leq600$.
\item En utilisant $e^x\geq1+x$, avec égalité seulement en 0, comparer $u_n$ et $W(n)$ pour $n\geq1$. Expliquer la différence des seuils.
\end{questions}
\end{exercice}
\end{blocLivre}

\clearpage
\gdef\otChapNum{0}\pagestyle{otFront}
\phantomsection\label{defis-debut}\pdfbookmark[0]{Défis d'approfondissement}{bm-defis}
\otPageTitre{Défis d'approfondissement}{Chercher plus loin avec les outils de Première}
\fontsize{11}{16}\selectfont
Les 210 exercices des chapitres installent les bases et les méthodes. Les
14 défis suivants portent le total à \textbf{224 exercices corrigés}.
Ils prolongent chaque thème par une question qui demande de prendre du recul,
de choisir une représentation ou de relier plusieurs notions.\par\vspace{4mm}
\textbf{Avancé} désigne une modélisation ou une synthèse soutenue.
\textbf{Très exigeant} ajoute une discussion, une preuve ou plusieurs changements de point de vue.
Ces indications sont des repères pédagogiques, pas un classement officiel.\par\vspace{4mm}
Les défis sont facultatifs pour un premier passage. Ils utilisent les notions
de Première ; les étapes qui dépasseraient les outils du cours sont explicitement admises.
Les durées proposées sont indicatives.\par\vspace{5mm}
\editionEtape{Chercher avant de lire la solution}
Faire un dessin ou un tableau, tester un cas simple, puis écrire ce qui manque
pour conclure. En cas de blocage, lire seulement le premier indice, puis reprendre
la recherche. Les indications sont séparées des énoncés et des corrigés.
\begin{center}\renewcommand{\arraystretch}{1.25}
\begin{tabular}{@{}lll@{}}\toprule
Défis & Thèmes & Préparation conseillée\\\midrule
D01 à D02 & Algèbre et suites & Chapitres 1, 2, 12, 13\\
D03 à D06 & Analyse et trigonométrie & Chapitres 1 à 6, 8, 12\\
D07 à D08 & Géométrie & Chapitres 1, 7, 8\\
D09 à D11 & Probabilités et statistiques & Chapitres 1, 9 à 11\\
D12 à D14 & Algorithmes et synthèses & Plusieurs chapitres, selon le défi\\\bottomrule
\end{tabular}\end{center}
\editionRetenir{Une solution incomplète mais justifiée est un progrès.
Après la correction, refaire la question clé sans regarder, puis expliquer pourquoi la méthode fonctionne.}

\editionPage{defi-ex-1}{D01 · Deux racines dans un intervalle}{Recherche et rédaction}{Avancé · 35 à 50 min · Chapitres 1 et 13\hfill\hyperref[defi-cor-1-1]{Corrigé p.~\pageref*{defi-cor-1-1}}}
\begin{blocEdition}{defi-ex-1}
Pour $m\in\R$, on pose $P_m(x)=x^2-2mx+m$. On cherche les paramètres pour
lesquels $P_m$ possède \textbf{deux racines distinctes, toutes deux dans $]0;2[$}.
\begin{questions}
\item Calculer le discriminant. Pour quels $m$ existe-t-il deux racines distinctes ?
\item Montrer que deux racines dans $]0;2[$ imposent $P_m(0)>0$, $P_m(2)>0$
et que l'abscisse du sommet appartient à $]0;2[$.
\item Réciproquement, expliquer pourquoi ces trois conditions, ajoutées à un
discriminant strictement positif, suffisent. On peut s'aider de la parabole.
\item Déterminer tous les paramètres recherchés. Étudier les deux valeurs frontières.
\item Les seules conditions $P_m(0)>0$ et $P_m(2)>0$ suffisent-elles ? Donner un contre-exemple.
\end{questions}
La figure fournit des exemples pour conjecturer ; la réponse doit être démontrée.
\label{figedition-defi-ex-1-1}\figureGeo{fr-d01}
\end{blocEdition}

\editionPage{defi-ex-2}{D02 · Le seuil d’une moyenne}{Recherche et rédaction}{Très exigeant · 45 à 65 min · Chapitres 2 et 12\hfill\hyperref[defi-cor-2-1]{Corrigé p.~\pageref*{defi-cor-2-1}}}
\begin{blocEdition}{defi-ex-2}
On définit $u_0=0$ et $u_{n+1}=\frac45u_n+60$ pour $n\geq0$.
Pour $n\geq1$, on note $S_n=\sum_{k=0}^{n-1}u_k$ et $A_n=S_n/n$.
\begin{questions}
\item À l'aide d'une suite auxiliaire géométrique, montrer que
$u_n=300\bigl(1-(4/5)^n\bigr)$. Étudier ses variations.
\item Calculer $S_n$, puis montrer que
$A_n=300-\frac{1500}{n}\bigl(1-(4/5)^n\bigr)$.
\item Expliquer pourquoi $A_n<u_n$, puis établir
$A_{n+1}-A_n=(u_n-A_n)/(n+1)$. En déduire les variations de $A_n$.
\item Écrire une fonction Python qui renvoie le premier $n\geq1$ tel que $A_n\geq240$.
Préciser ce que représentent les variables à chaque passage de la boucle.
\item Le programme conduit à $n=25$. Vérifier le seuil en calculant $A_{24}$ et $A_{25}$.
Pourquoi ces deux seuls contrôles suffisent-ils après la question 3 ?
\end{questions}
La moyenne porte sur les termes $u_0$ à $u_{n-1}$ : elle ne comprend pas $u_n$.
\end{blocEdition}

\editionPage{defi-ex-3}{D03 · Trois tangentes par un point}{Recherche et rédaction}{Très exigeant · 55 à 75 min · Chapitres 3, 4 et 1\hfill\hyperref[defi-cor-3-1]{Corrigé p.~\pageref*{defi-cor-3-1}}}
\begin{blocEdition}{defi-ex-3}
On considère $f(x)=x^3-3x$. Pour $t\in[-1;2]$, on note $T_t$ la tangente
à sa courbe au point d'abscisse $t$. Pour $b\in\R$, on pose $P_b(1;b)$.
\begin{questions}
\item Montrer que $T_t$ a pour équation $y=(3t^2-3)x-2t^3$.
\item Montrer que $T_t$ passe par $P_b$ si et seulement si $g(t)=b$,
où $g(t)=-2t^3+3t^2-3$.
\item Étudier $g$ sur $[-1;2]$ et calculer ses valeurs aux bornes et aux points critiques.
\item Selon $b$, déterminer le nombre de tangentes $T_t$ qui passent par $P_b$.
Soigner les cas où $b$ est une valeur extrême ou une valeur frontière.
\item Pour $b=-5/2$, factoriser $4t^3-6t^2+1$ et déterminer les trois
abscisses de contact, puis les équations des trois tangentes.
\end{questions}
Seules les tangentes dont l'abscisse de contact est dans $[-1;2]$ sont comptées.
\label{figedition-defi-ex-3-1}\figureGeo{fr-d03}
\end{blocEdition}

\editionPage{defi-ex-4}{D04 · La fenêtre la plus lumineuse}{Recherche et rédaction}{Avancé · 35 à 50 min · Chapitres 1 et 4\hfill\hyperref[defi-cor-4-1]{Corrigé p.~\pageref*{defi-cor-4-1}}}
\begin{blocEdition}{defi-ex-4}
Une fenêtre est formée d'un rectangle de largeur $2r$ et de hauteur $h$, surmonté
d'un demi-disque de rayon $r$. La longueur totale de son \textbf{contour extérieur} est de 20 m.
On suppose $r>0$ et $h>0$ et on veut maximiser l'aire vitrée.
\begin{questions}
\item Exprimer le périmètre. Pourquoi le diamètre commun au rectangle et au demi-disque
ne doit-il pas être ajouté au contour extérieur ?
\item Exprimer $h$ en fonction de $r$ et préciser exactement le domaine admissible.
\item Montrer que l'aire est $A(r)=20r-(2+\pi/2)r^2$.
\item Déterminer les dimensions optimales et l'aire maximale, sous forme exacte.
Quelle relation simple lie $h$ et $r$ à l'optimum ?
\item Expliquer pourquoi le maximum calculé correspond à une vraie fenêtre et non à une forme dégénérée.
\end{questions}
\label{figedition-defi-ex-4-1}\figureGeo{fr-d04}
\end{blocEdition}

\editionPage{defi-ex-5}{D05 · Encadrer sans logarithme}{Recherche et rédaction}{Très exigeant · 45 à 60 min · Chapitres 5, 2 et 12\hfill\hyperref[defi-cor-5-1]{Corrigé p.~\pageref*{defi-cor-5-1}}}
\begin{blocEdition}{defi-ex-5}
On veut encadrer $e^{1/2}$ avec des puissances de nombres rationnels.
\begin{questions}
\item Étudier $F(x)=e^x-1-x$ sur $\R$ et montrer que $e^x\geq1+x$,
avec égalité seulement pour $x=0$.
\item Pour $0<t<1$, appliquer ce résultat à $-t$, puis justifier
$1+t<e^t<1/(1-t)$.
\item Pour tout entier $n\geq1$, en déduire
\[L_n=\left(1+\frac1{2n}\right)^n<e^{1/2}<
U_n=\left(1-\frac1{2n}\right)^{-n}.\]
\item Programmer le calcul de $L_n$ et $U_n$, puis donner des encadrements
pour $n=10$, $100$ et $1000$. Arrondir la borne inférieure vers le bas et la borne supérieure vers le haut.
\item Que prouve l'encadrement théorique ? Que ne suffit pas à prouver la seule observation des décimales calculées ?
\end{questions}
On utilise uniquement les propriétés de l'exponentielle ; le logarithme n'est pas nécessaire.
\end{blocEdition}

\editionPage{defi-ex-6}{D06 · Additionner sinus et cosinus}{Recherche et rédaction}{Très exigeant · 45 à 65 min · Chapitres 6, 1 et 8\hfill\hyperref[defi-cor-6-1]{Corrigé p.~\pageref*{defi-cor-6-1}}}
\begin{blocEdition}{defi-ex-6}
Pour $t\in[0;2\pi[$, on pose $x=\cos t$, $y=\sin t$. On cherche à résoudre
$\cos t+\sin t=k$, où $k$ est un paramètre réel.
\begin{questions}
\item Traduire le problème comme une intersection entre un cercle et une droite.
\item En utilisant $(x-y)^2\geq0$, montrer que $-\sqrt2\leq x+y\leq\sqrt2$.
\item Lorsque $x+y=k$, montrer que $xy=(k^2-1)/2$, puis que $x$ et $y$
sont les racines de $X^2-kX+(k^2-1)/2=0$.
\item Classer le nombre de solutions $t$ selon $k$. Donner les coordonnées des points correspondants.
\item Résoudre explicitement pour $k=0$, $k=1$ et $k=\sqrt2$.
\end{questions}
L'intervalle est ouvert en $2\pi$ pour ne pas compter deux fois le même point du cercle.
\label{figedition-defi-ex-6-1}\figureGeo{fr-d06}
\end{blocEdition}

\editionPage{defi-ex-7}{D07 · Un triangle sous deux contraintes}{Recherche et rédaction}{Très exigeant · 40 à 60 min · Chapitres 7 et 8\hfill\hyperref[defi-cor-7-1]{Corrigé p.~\pageref*{defi-cor-7-1}}}
\begin{blocEdition}{defi-ex-7}
Dans un repère orthonormé, $A(-2;0)$, $B(2;0)$ et $M(x;y)$ vérifient
\[\vect{MA}\cdot\vect{MB}=0,\qquad MA^2+2MB^2=k.\]
Le paramètre $k$ est réel.
\begin{questions}
\item Montrer que la première condition équivaut à $x^2+y^2=4$.
\item Sous cette condition, établir $MA^2=8+4x$ et $MB^2=8-4x$.
\item En déduire les coordonnées possibles de $M$ en fonction de $k$ et discuter leur nombre.
\item Pour quels $k$ obtient-on un \textbf{vrai triangle} $AMB$ rectangle en $M$ ?
Pourquoi les valeurs frontières doivent-elles être exclues ?
\item Pour quel $k$ le triangle est-il aussi isocèle en $M$ ? Donner ses deux sommets possibles.
\end{questions}
\label{figedition-defi-ex-7-1}\figureGeo{fr-d07}
\end{blocEdition}

\editionPage{defi-ex-8}{D08 · Une famille de cercles tangents}{Recherche et rédaction}{Très exigeant · 55 à 75 min · Chapitres 8, 7 et 1\hfill\hyperref[defi-cor-8-1]{Corrigé p.~\pageref*{defi-cor-8-1}}}
\begin{blocEdition}{defi-ex-8}
Pour $m\in\R$, on considère $\mathcal C_m$ d'équation
$x^2+y^2-4x-2my=0$ et la droite $d:x+y=6$.
\begin{questions}
\item Déterminer le centre et le rayon de $\mathcal C_m$. Montrer que tous
ces cercles passent par $A(0;0)$ et $B(4;0)$.
\item Substituer $y=6-x$ dans l'équation du cercle. Montrer que le discriminant
de l'équation obtenue est $4(m^2+8m-8)$.
\item Classer, selon $m$, le nombre de points communs à $d$ et $\mathcal C_m$.
Pour quels paramètres y a-t-il tangence ?
\item On note $\Omega_m(2;m)$. Calculer son projeté orthogonal $H_m$ sur $d$.
Vérifier, pour chacun des paramètres de tangence, que $H_m$ appartient au cercle.
\item Expliquer pourquoi la droite $d$ est alors perpendiculaire au rayon au point de contact.
\end{questions}
La figure de l'énoncé représente seulement les points fixes et la droite ; les cercles restent à déterminer.
\label{figedition-defi-ex-8-1}\figureGeo{fr-d08-ex}
\end{blocEdition}

\editionPage{defi-ex-9}{D09 · Choisir une règle de contrôle}{Recherche et rédaction}{Très exigeant · 55 à 75 min · Chapitres 9 et 10\hfill\hyperref[defi-cor-9-1]{Corrigé p.~\pageref*{defi-cor-9-1}}}
\begin{blocEdition}{defi-ex-9}
Dans un modèle fictif, $10\,\%$ des articles présentent un défaut $D$.
On applique deux tests $T_1$ et $T_2$, de résultat positif $+$ ou négatif $-$.
Sur les articles défectueux, chaque test est positif avec probabilité $0{,}9$ ;
sur les autres, avec probabilité $0{,}2$.
\textbf{Dans chacun de ces deux groupes}, le résultat du second test a les mêmes
probabilités quel que soit le résultat du premier.
\begin{questions}
\item Compléter l'arbre à trois niveaux. Les deux tests sont-ils pour autant
indépendants dans la population entière ? Comparer $P(T_1^+\cap T_2^+)$ à $P(T_1^+)P(T_2^+)$.
\item Calculer la probabilité du défaut sachant qu'un test est positif,
puis sachant que les deux sont positifs.
\item La règle U écarte un article dès qu'au moins un test est positif.
La règle B ne l'écarte que si les deux sont positifs. Pour un article choisi dans toute la production,
calculer, pour chaque règle, $P(\overline D\text{ et rejet})$ et $P(D\text{ et acceptation})$.
\item Le premier type d'erreur coûte 1 unité, le second $\lambda$ unités, avec $\lambda\geq0$.
Comparer les coûts moyens des deux règles selon $\lambda$.
\end{questions}
\label{figedition-defi-ex-9-1}\figureGeo{fr-d09}
\end{blocEdition}

\editionPage{defi-ex-10}{D10 · Même moyenne, objectifs différents}{Recherche et rédaction}{Très exigeant · 40 à 60 min · Chapitres 10 et 1\hfill\hyperref[defi-cor-10-1]{Corrigé p.~\pageref*{defi-cor-10-1}}}
\begin{blocEdition}{defi-ex-10}
Un jeu présente trois issues $I_1,I_2,I_3$ de probabilités $1/4,1/2,1/4$.
Pour un paramètre $a\in[0;2]$, le lot brut $L_a$ vaut respectivement
$2a$, $2-a$ et $6$ unités. Le prix d'entrée est $5/2$ ; le gain net est $G_a=L_a-5/2$.
\begin{questions}
\item Donner la loi issue par issue de $L_a$ et calculer son espérance.
À quelle valeur de $a$ deux lots coïncident-ils ? Comment regrouper alors les probabilités ?
\item Calculer $E(G_a)$ et montrer que
$V(G_a)=\frac32a^2-2a+\frac{19}{4}$.
\item Déterminer le paramètre qui minimise cette variance et la valeur minimale.
\item Calculer la probabilité $P(G_a>0)$ selon $a$, en traitant le seuil d'égalité.
Le choix de variance minimale maximise-t-il la probabilité de gagner strictement ?
\item Pour une prévision réelle $r$, développer $E((L_a-r)^2)$.
Déterminer simultanément les valeurs de $a$ et $r$ qui minimisent cette erreur quadratique moyenne.
\end{questions}
\label{figedition-defi-ex-10-1}\figureGeo{fr-d10}
\end{blocEdition}

\editionPage{defi-ex-11}{D11 · La dispersion d’un groupe réuni}{Recherche et rédaction}{Très exigeant · 45 à 65 min · Chapitres 11 et 1\hfill\hyperref[defi-cor-11-1]{Corrigé p.~\pageref*{defi-cor-11-1}}}
\begin{blocEdition}{defi-ex-11}
Le groupe A contient les valeurs 0 et 2 ; le groupe B contient 8 et 10.
La variance d'une série est la moyenne des carrés des écarts à sa moyenne.
\begin{questions}
\item Calculer les moyennes et variances des deux groupes, puis de leur réunion.
Pourquoi la variance de la réunion n'est-elle pas la moyenne des deux variances ?
\item Plus généralement, deux groupes ont des tailles $n_A,n_B>0$, des moyennes
$\mu_A,\mu_B$ et des variances $V_A,V_B$. Poser $p=n_A/(n_A+n_B)$ et exprimer
la moyenne $\mu$ de leur réunion.
\item Dans le groupe A, développer
$(x-\mu)^2=((x-\mu_A)+(\mu_A-\mu))^2$ puis moyenner.
Que devient le terme croisé ? Faire de même dans B.
\item En déduire la formule
\[V=pV_A+(1-p)V_B+p(1-p)(\mu_A-\mu_B)^2.\]
\item Deux groupes ont chacun une variance égale à 1. La variance de leur réunion
peut-elle être aussi grande que l'on veut ? Construire un exemple avec deux groupes de même taille.
\end{questions}
\label{figedition-defi-ex-11-1}\figureGeo{fr-d11}
\end{blocEdition}

\editionPage{defi-ex-12}{D12 · Un algorithme qui encadre une racine}{Recherche et rédaction}{Très exigeant · 45 à 60 min · Chapitres 12, 3 et 4\hfill\hyperref[defi-cor-12-1]{Corrigé p.~\pageref*{defi-cor-12-1}}}
\begin{blocEdition}{defi-ex-12}
On pose $f(x)=x^3+x-1$. \textbf{On admet l'existence} d'une solution $\alpha$ de
$f(x)=0$ dans $[0;1]$ ; il ne s'agit pas ici de démontrer un théorème de continuité.
\begin{questions}
\item Justifier que cette solution est unique.
\item On part de $a=0$, $b=1$, puis on calcule $m=(a+b)/2$.
Si $f(m)<0$, on remplace $a$ par $m$ ; sinon, on remplace $b$ par $m$.
Expliquer pourquoi $a\leq\alpha\leq b$ reste vrai après cette étape.
\item Écrire une fonction qui réalise exactement $N$ étapes et renvoie les deux bornes.
Préciser les entrées autorisées.
\item Montrer que la largeur finale vaut $2^{-N}$.
Quel est le plus petit $N$ qui garantit une largeur au plus $10^{-3}$ ?
\item Pour $N=10$, vérifier l'encadrement exact
$698/1024\leq\alpha\leq699/1024$ en étudiant les signes aux deux bornes.
Quelle erreur maximale garantit le milieu de cet intervalle ?
\end{questions}
\end{blocEdition}

\editionPage{defi-ex-13}{D13 · Une garantie pour tous les x}{Recherche et rédaction}{Très exigeant · 45 à 65 min · Chapitres 13, 1 et 4\hfill\hyperref[defi-cor-13-1]{Corrigé p.~\pageref*{defi-cor-13-1}}}
\begin{blocEdition}{defi-ex-13}
Pour $m\in\R$, on définit $P_m(x)=x^2-2mx+m$, et l'on s'intéresse à $x\in[0;1]$.
\begin{questions}
\item Formuler la négation de : « Pour tout $x\in[0;1]$, $P_m(x)\geq0$. »
\item Montrer que cette propriété impose $0\leq m\leq1$ en regardant les bornes.
Prouver la réciproque avec la forme canonique.
\item On note $q(m)$ le minimum de $P_m$ sur $[0;1]$. Déterminer $q(m)$
pour $m<0$, $0\leq m\leq1$ et $m>1$.
\item Quel paramètre fournit la plus grande garantie $q(m)$ ?
Donner une seconde preuve en utilisant la valeur commune $P_m(1/2)$.
\item Comparer les deux affirmations : « Il existe un $m$ tel que, pour tout $x$,
$P_m(x)\geq1/4$ » et « Pour tout $m$, il existe un $x$ tel que $P_m(x)\geq1/4$ ».
Préciser les domaines et justifier chacune des réponses.
\end{questions}
\end{blocEdition}

\editionPage{defi-ex-14}{D14 · Qualité, coût et risque}{Recherche et rédaction}{Très exigeant · 60 à 80 min · Chapitres 5, 9, 10 et 12\hfill\hyperref[defi-cor-14-1]{Corrigé p.~\pageref*{defi-cor-14-1}}}
\begin{blocEdition}{defi-ex-14}
Dans un atelier fictif, une intensité de contrôle $x\in[0;4]$ entraîne un coût
de $2x$ unités par article et une probabilité de défaut $p(x)=\frac15e^{-x}$.
Avant ce coût, un article rapporte 12 unités s'il est sans défaut et $-8$ sinon.
Le gain net après contrôle est noté $G_x$.
\begin{questions}
\item Donner les deux valeurs de $G_x$, leurs probabilités, puis montrer que
$E(G_x)=g(x)=12-2x-4e^{-x}$.
\item Montrer que le maximum est atteint pour l'unique $\alpha\in]0;1[$ tel que $e^\alpha=2$.
On admet l'existence de ce réel. Aucun logarithme n'est demandé.
\item Écrire un algorithme de dichotomie sur $[0;1]$ pour encadrer $\alpha$ à $10^{-3}$ près.
En déduire une approximation du gain moyen maximal.
\item Montrer que $V(G_x)=80e^{-x}-16e^{-2x}$ et étudier ses variations sur $[0;4]$.
Au point optimal pour l'espérance, calculer la probabilité de défaut et la variance exacte.
\item Deux articles sont contrôlés avec la même intensité $\alpha$ ; leurs défauts sont indépendants.
Dessiner l'arbre, puis calculer la probabilité d'au moins un défaut.
La plus grande intensité maximise-t-elle aussi le gain moyen ?
\end{questions}
\end{blocEdition}


\editionPage{indices-defis}{Indications progressives}{Lire une aide, puis reprendre la recherche}{Les indices orientent la démarche ; les solutions restent dans la partie Corrigés.}
\begin{blocEdition}{indices-defis}
\editionEtape{D01 · Deux racines dans un intervalle}
\textbf{Indice 1.} Repérer le sommet et calculer la valeur du polynôme aux deux bornes.\par\smallskip
\textbf{Indice 2.} Une parabole tournée vers le haut, positive aux bornes et négative à son sommet intérieur, coupe l'axe deux fois entre les bornes.\par\smallskip
\editionEtape{D02 · Le seuil d’une moyenne}
\textbf{Indice 1.} Chercher le point fixe de la relation affine, puis soustraire ce nombre à chaque terme.\par\smallskip
\textbf{Indice 2.} Pour la moyenne, utiliser $S_{n+1}=S_n+u_n$. Au test de boucle, conserver simultanément le rang, la somme $S_n$ et le terme $u_n$.\par\smallskip
\editionEtape{D03 · Trois tangentes par un point}
\textbf{Indice 1.} Écrire la tangente au point d'abscisse t avant de remplacer x par 1.\par\smallskip
\textbf{Indice 2.} Pour compter les solutions, utiliser les trois intervalles de monotonie de g et traiter séparément les hauteurs -3 et -2.\par\smallskip
\editionEtape{D04 · La fenêtre la plus lumineuse}
\textbf{Indice 1.} Tracer au crayon uniquement le contour extérieur ; sa longueur est la somme de quatre morceaux.\par\smallskip
\textbf{Indice 2.} Éliminer h avec le périmètre, puis optimiser le trinôme en r.\par\smallskip
\editionEtape{D05 · Encadrer sans logarithme}
\textbf{Indice 1.} Étudier le signe de la dérivée de $e^x-1-x$.\par\smallskip
\textbf{Indice 2.} Pour la majoration, remplacer $x$ par $-t$ et inverser deux nombres positifs. Puis prendre $t=1/(2n)$.\par\smallskip
\end{blocEdition}

\editionPage{indices-defis-2}{Indications progressives}{Lire une aide, puis reprendre la recherche}{Les indices orientent la démarche ; les solutions restent dans la partie Corrigés.}
\begin{blocEdition}{indices-defis-2}
\editionEtape{D06 · Additionner sinus et cosinus}
\textbf{Indice 1.} Dans le repère, écrire $x^2+y^2=1$ et $x+y=k$.\par\smallskip
\textbf{Indice 2.} Calculer le produit $xy$ grâce au carré de la somme, puis utiliser un trinôme de somme $k$ et de produit $(k^2-1)/2$.\par\smallskip
\editionEtape{D07 · Un triangle sous deux contraintes}
\textbf{Indice 1.} Développer le produit scalaire en coordonnées avant de reconnaître un cercle.\par\smallskip
\textbf{Indice 2.} Sur ce cercle, remplacer $x^2+y^2$ par 4 dans les carrés des distances. Le second lieu devient une droite verticale.\par\smallskip
\editionEtape{D08 · Une famille de cercles tangents}
\textbf{Indice 1.} Compléter deux carrés, puis remplacer $y$ par $6-x$ pour obtenir un trinôme en $x$.\par\smallskip
\textbf{Indice 2.} Pour projeter le centre, se déplacer dans la direction du vecteur normal $(1;1)$ jusqu'à satisfaire $x+y=6$.\par\smallskip
\editionEtape{D09 · Choisir une règle de contrôle}
\textbf{Indice 1.} Commencer par séparer les groupes D et son complément ; multiplier ensuite les poids des deux tests à l'intérieur de chaque groupe.\par\smallskip
\textbf{Indice 2.} Pour U, le seul chemin accepté réunit deux tests négatifs. Pour comparer les règles, calculer une différence de deux fonctions affines du coût d'erreur.\par\smallskip
\editionEtape{D10 · Même moyenne, objectifs différents}
\textbf{Indice 1.} Calculer d'abord l'espérance et le deuxième moment avec les probabilités des issues, même si deux lots coïncident.\par\smallskip
\textbf{Indice 2.} Mettre la variance sous forme canonique ; pour le gain strictement positif, traiter les trois issues une à une.\par\smallskip
\end{blocEdition}

\editionPage{indices-defis-3}{Indications progressives}{Lire une aide, puis reprendre la recherche}{Les indices orientent la démarche ; les solutions restent dans la partie Corrigés.}
\begin{blocEdition}{indices-defis-3}
\editionEtape{D11 · La dispersion d’un groupe réuni}
\textbf{Indice 1.} Dans le calcul de la variance réunie, tous les écarts doivent être mesurés à la même moyenne totale.\par\smallskip
\textbf{Indice 2.} La moyenne des écarts à la moyenne d'un groupe est nulle ; c'est ce qui annule le terme croisé du développement.\par\smallskip
\editionEtape{D12 · Un algorithme qui encadre une racine}
\textbf{Indice 1.} Utiliser la stricte croissance pour relier le signe de $f(m)$ à la position de $m$ par rapport à la racine.\par\smallskip
\textbf{Indice 2.} La largeur est une suite géométrique de raison $1/2$ ; distinguer la largeur de l'intervalle et l'erreur du milieu.\par\smallskip
\editionEtape{D13 · Une garantie pour tous les x}
\textbf{Indice 1.} Le minimum peut être au sommet ou à une borne : commencer par situer le sommet selon m.\par\smallskip
\textbf{Indice 2.} Comparer le minimum à la valeur $P_m(1/2)$, puis chercher un paramètre qui atteint cette borne pour tous les $x$.\par\smallskip
\editionEtape{D14 · Qualité, coût et risque}
\textbf{Indice 1.} Calculer l'espérance avant de dériver ; la différence entre les deux gains est constante.\par\smallskip
\textbf{Indice 2.} Le signe de $g'$ se lit avec la croissance de $e^x$. À l'optimum, remplacer $e^{-\alpha}$ par $1/2$ simplifie les probabilités et la variance.\par\smallskip
\end{blocEdition}


\clearpage
\gdef\otChapNum{0}\pagestyle{otFront}\phantomsection\label{epreuve-guide}
\pdfbookmark[0]{Préparer l'épreuve anticipée}{bm-epreuve}
\otPageTitre{Préparer l'épreuve}{Première générale · Spécialité mathématiques · 2026-2027}
\fontsize{11}{15}\selectfont
\editionEtape{Le cadre officiel}
La note de service du 11 septembre 2026 définit l'épreuve pour l'année 2026-2027,
au titre du baccalauréat session 2028 ; elle remplace celle du 10 juin 2025.
L'épreuve écrite dure \textbf{2 heures}, a pour coefficient 2 et se passe
\textbf{sans calculatrice}. Le candidat de spécialité compose sur le sujet propre à cet enseignement.
\par\smallskip
{\renewcommand{\arraystretch}{1.12}\begin{tabularx}{\linewidth}{@{}p{43mm}Y@{}}\toprule
Première partie · 6 points & Automatismes : liste de questions à \textbf{réponses courtes}.\\
Deuxième partie · 14 points & Deux ou trois exercices indépendants.\\
Maîtrise de la langue & 2 points sur 20 : orthographe, syntaxe, lexique, organisation du propos.\\
Oral de contrôle & 20 minutes de préparation, 20 minutes d'entretien ; au moins deux exercices
sur des parties différentes du programme, sans calculatrice.\\\bottomrule
\end{tabularx}}\par\vspace{2mm}
\href{https://www.education.gouv.fr/bo/2026/Special4/MENE2622640N}{\textcolor{otNavy}{Définition officielle : note de service du 11 septembre 2026, BO spécial n° 4.}}
\editionEtape{Les automatismes évalués}
Ils relèvent de la rubrique « Automatismes » du programme de Première : évolutions et
variations, calcul numérique et algébrique, fonctions et représentations, statistiques,
probabilités ; s'y ajoutent les automatismes de Seconde. Le chapitre 14 et les sujets A et B
les entraînent. Les sujets zéro publiés en 2025 proposaient un QCM : pour la session 2028,
le candidat écrit lui-même chaque réponse courte.
\editionEtape{Deux entraînements originaux}
Les sujets A et B, créés pour ce livre, ne sont pas des sujets officiels. Chacun comporte
12 automatismes, puis deux exercices indépendants. Leur barème, choix pédagogique du livre, donne 0,5 point
par automatisme et 7 points par exercice, dont 1 point de maîtrise de la langue : orthographe,
phrases construites, vocabulaire précis, étapes enchaînées jusqu'à la conclusion.
\editionEtape{Une gestion du temps à adapter}
Environ 20 minutes pour les automatismes, 40 minutes par exercice et 20 minutes pour relire.
Après la séance, classer chaque erreur (cours, calcul, lecture, rédaction) et refaire plus tard
l'exercice de chapitre correspondant.
\editionRetenir{Une valeur exacte, un signe justifié, un domaine précisé et une phrase
qui répond à la question rendent le raisonnement vérifiable.}
\clearpage


\editionPage{sujet-A}{Sujet A · Automatismes}{Première partie · 6 points}{Corrigé p.~\pageref*{sujet-A-cor-auto}}
\begin{blocEdition}{sujet-A}
\textbf{Sujet original · Durée : 2 heures · L’usage de la calculatrice n’est pas autorisé.}
\par\smallskip\textbf{Première partie (6 points) : automatismes.} Pour chacune des 12 questions,
reporter son numéro sur la copie et donner une réponse courte, avec les unités utiles ;
aucune justification n’est demandée. Chaque question vaut 0,5 point.
\par\vspace{2mm}\begin{questions}\item Un prix de 60 euros augmente de $15\,\%$. Quel est le nouveau prix ?
\item Après une hausse de $25\,\%$, un article coûte 75 euros. Quel était son prix avant la hausse ?
\item Un prix baisse de $20\,\%$. Quel pourcentage de hausse permet de revenir au prix initial ?
\item Calculer $(2-\sqrt3)(2+\sqrt3)$.
\item Résoudre dans $\R$ l’équation $(2x-1)(x+3)=0$.
\item Résoudre dans $\R$ l’inéquation $-3x+6>0$.
\item Pour quels réels $x$ a-t-on $(x-1)(x+4)<0$ ?
\item Déterminer le coefficient directeur de la droite passant par $A(-1;3)$ et $B(3;-5)$.
\item On donne $P(A)=0{,}4$ et $P_A(B)=0{,}25$. Calculer $P(A\cap B)$.
\item Donner la médiane de la série $1\,;\,2\,;\,2\,;\,5\,;\,10$.
\item La figure 1 représente une fonction $f$ définie sur $[-2;4]$. Résoudre graphiquement $f(x)\leq1$.
\item Lire graphiquement l’équation réduite de la droite $d$ de la figure 2.\end{questions}
\label{figedition-sujet-A-1}\figureGeo{fr-sa-auto}
\end{blocEdition}

\editionPage{sujet-A-ex-1}{Sujet A · Exercice 1}{Deuxième partie · Une courbe et sa tangente}{7 points, dont 1 point de maîtrise de la langue\hfill Corrigé p.~\pageref*{sujet-A-cor-1}}
\begin{blocEdition}{sujet-A-ex-1}
\textbf{Deuxième partie (14 points) : deux exercices indépendants, notés chacun sur 7 points.}
\par\vspace{2mm}
On définit $f(x)=x^3-6x^2+9x+2$ pour $x\in[0;4]$.
\begin{questions}
\item \mbox{\textit{(1 point)}}\enspace Calculer $f'(x)$ et le factoriser.
\item \mbox{\textit{(2 points)}}\enspace Dresser le tableau de variations, avec les valeurs aux bornes et aux points critiques.
\item \mbox{\textit{(1 point)}}\enspace Donner le minimum et le maximum sur le segment, ainsi que toutes les abscisses où ils sont atteints.
\item \mbox{\textit{(1 point)}}\enspace Écrire la tangente $T$ au point d'abscisse 2.
\item \mbox{\textit{(1 point)}}\enspace Vérifier que $f(x)-(-3x+10)=(x-2)^3$, puis déterminer la position de la courbe par rapport à $T$.
\end{questions}
On attend un tableau lisible et une phrase de conclusion pour les extremums et la position relative.

\end{blocEdition}

\editionPage{sujet-A-ex-2}{Sujet A · Exercice 2}{Deuxième partie · Ateliers, défauts et gain}{7 points, dont 1 point de maîtrise de la langue\hfill Corrigé p.~\pageref*{sujet-A-cor-2}}
\begin{blocEdition}{sujet-A-ex-2}

L'atelier A fournit $60\,\%$ des articles et l'atelier B les autres.
Les taux de défaut sont respectivement $5\,\%$ et $10\,\%$.
On choisit un article au hasard. $D$ désigne l'évènement « l'article présente un défaut ».
\begin{questions}
\item \mbox{\textit{(1 point)}}\enspace Compléter l'arbre et calculer $P(D)$.
\item \mbox{\textit{(1 point)}}\enspace Sachant que l'article présente un défaut, calculer la probabilité qu'il provienne de B.
\item \mbox{\textit{(1 point)}}\enspace Le gain net est 12 unités sans défaut et $-8$ avec défaut. Donner sa loi.
\item \mbox{\textit{(1 point)}}\enspace Calculer l'espérance du gain.
\item \mbox{\textit{(1 point)}}\enspace Calculer sa variance.
\item \mbox{\textit{(1 point)}}\enspace Pour deux articles de défauts indépendants et de même loi, calculer la probabilité d'au moins un défaut.
\end{questions}

\label{figedition-sujet-A-ex-2-1}\figureGeo{fr-sa-arbre}
\end{blocEdition}

\editionPage{sujet-B}{Sujet B · Automatismes}{Première partie · 6 points}{Corrigé p.~\pageref*{sujet-B-cor-auto}}
\begin{blocEdition}{sujet-B}
\textbf{Sujet original · Durée : 2 heures · L’usage de la calculatrice n’est pas autorisé.}
\par\smallskip\textbf{Première partie (6 points) : automatismes.} Pour chacune des 12 questions,
reporter son numéro sur la copie et donner une réponse courte, avec les unités utiles ;
aucune justification n’est demandée. Chaque question vaut 0,5 point.
\par\vspace{2mm}\begin{questions}\item Calculer $\dfrac34+\dfrac16$.
\item Calculer $\dfrac{(2^3)^2}{2^4}$.
\item Factoriser $9x^2-1$.
\item Une quantité passe de 80 à 92. Calculer son taux d’évolution, en pourcentage.
\item Deux hausses successives d’un même taux donnent une hausse globale de $21\,\%$. Quel est ce taux ?
\item Exprimer $h$ en fonction de $V$ et de $r$, sachant que $V=\pi r^2h$ avec $r>0$.
\item La droite $d$ passe par $A(0;3)$ et a pour coefficient directeur $-2$. Donner son équation réduite.
\item Calculer la moyenne des notes 8, 12 et 15, de coefficients respectifs 1, 2 et 1.
\item \begin{minipage}[t]{86mm}Le tableau ci-contre répartit les 200 élèves d’un lycée. On choisit un élève au hasard ; $F$ : « être une fille » et $I$ : « être interne ». Calculer $P_F(I)$.\end{minipage}\hfill
\begin{minipage}[t]{62mm}\vspace{-3.2mm}\small\begin{tabular}[t]{@{}lccc@{}}\toprule
 & Interne & Externe & Total\\\midrule
Fille & 30 & 90 & 120\\
Garçon & 30 & 50 & 80\\
Total & 60 & 140 & 200\\\bottomrule
\end{tabular}\end{minipage}
\item Avec le même tableau, calculer $P(F\cap I)$, puis $P_I(F)$.
\item La figure 1 représente une fonction $g$ définie sur $[-2;4]$. Donner le signe de $g(x)$ selon $x$.
\item La figure 2 résume une série de notes. Donner sa médiane et son écart interquartile.\end{questions}
\label{figedition-sujet-B-1}\figureGeo{fr-sb-auto}
\end{blocEdition}

\editionPage{sujet-B-ex-1}{Sujet B · Exercice 1}{Deuxième partie · Une évolution avec apport}{7 points, dont 1 point de maîtrise de la langue\hfill Corrigé p.~\pageref*{sujet-B-cor-1}}
\begin{blocEdition}{sujet-B-ex-1}
\textbf{Deuxième partie (14 points) : deux exercices indépendants, notés chacun sur 7 points.}
\par\vspace{2mm}
On pose $u_0=80$ et $u_{n+1}=\frac34u_n+10$, pour $n\geq0$.
\begin{questions}
\item \mbox{\textit{(1,5 point)}}\enspace On pose $v_n=u_n-40$. Montrer que $v$ est géométrique et préciser ses paramètres.
\item \mbox{\textit{(1 point)}}\enspace En déduire l'expression de $u_n$.
\item \mbox{\textit{(1 point)}}\enspace Justifier que $u_n>40$ pour tout $n$ et que la suite décroît strictement.
\item \mbox{\textit{(1,5 point)}}\enspace On donne $3^4=81$, $4^4=256$, $3^5=243$ et $4^5=1024$.
Trouver le premier rang tel que $u_n\leq50$.
\item \mbox{\textit{(1 point)}}\enspace Écrire une boucle Python qui renvoie ce premier rang.
\end{questions}
La recherche du premier rang doit comporter une justification, pas seulement un essai réussi.

\end{blocEdition}

\editionPage{sujet-B-ex-2}{Sujet B · Exercice 2}{Deuxième partie · Cercle, tangente et projection}{7 points, dont 1 point de maîtrise de la langue\hfill Corrigé p.~\pageref*{sujet-B-cor-2}}
\begin{blocEdition}{sujet-B-ex-2}

Dans un repère orthonormé, on donne $A(0;0)$, $B(4;0)$, $C(1;3)$ et $M(5;4)$.
\begin{questions}
\item \mbox{\textit{(1,5 point)}}\enspace Déterminer les médiatrices de $[AB]$ et $[AC]$, puis leur point d'intersection $\Omega$.
\item \mbox{\textit{(1 point)}}\enspace En déduire l'équation du cercle circonscrit au triangle $ABC$.
\item \mbox{\textit{(0,5 point)}}\enspace Vérifier que $C$ appartient à ce cercle.
\item \mbox{\textit{(1,5 point)}}\enspace Déterminer la tangente au cercle en $C$.
\item \mbox{\textit{(1,5 point)}}\enspace Projeter orthogonalement $M$ sur cette tangente et calculer la distance correspondante.
\end{questions}

\label{figedition-sujet-B-ex-2-1}\figureGeo{fr-sb-geometrie}
\end{blocEdition}


\clearpage
\gdef\otChapNum{0}
\pagestyle{otFront}
\phantomsection\label{demonstrations-debut}
\pdfbookmark[0]{Démonstrations de référence}{bm-demonstrations}
\otPageTitre{Démonstrations}{Des preuves à comprendre et à savoir expliquer}
\phantomsection\label{dem-second-degre}
\essSection{1}{Résoudre une équation du second degré}
\begin{demonstration}[Faire apparaître le discriminant]
Soit $f(x)=ax^2+bx+c$, avec $a\neq0$, et $\Delta=b^2-4ac$.
En développant, on vérifie l'identité
\[4af(x)=(2ax+b)^2-\Delta.\]
Ainsi $f(x)=0$ équivaut à $(2ax+b)^2=\Delta$.
Si $\Delta<0$, aucun carré réel ne convient. Si $\Delta=0$, on obtient
la seule solution $x=-b/(2a)$. Si $\Delta>0$, les deux possibilités
$2ax+b=\pm\sqrt\Delta$ donnent
\[x=\frac{-b-\sqrt\Delta}{2a}\quad\text{ou}\quad x=\frac{-b+\sqrt\Delta}{2a}.\]
\end{demonstration}
\begin{propriete}[Forme canonique et sommet]
L'identité précédente donne aussi
\[f(x)=a\left(x+\frac b{2a}\right)^2-\frac\Delta{4a}.
\]
Le sommet a donc pour abscisse $\alpha=-b/(2a)$ et pour ordonnée $f(\alpha)$.
Pour une détermination effective dans un cas simple, on complète le carré :
$x^2-6x+5=(x-3)^2-4$.
\end{propriete}
\begin{demonstration}[Somme et produit des racines]
Si les racines sont $r_1$ et $r_2$, la forme factorisée est
$f(x)=a(x-r_1)(x-r_2)$, y compris en cas de racine double.
En développant, $f(x)=ax^2-a(r_1+r_2)x+ar_1r_2$.
L'identification avec $ax^2+bx+c$ donne
$r_1+r_2=-b/a$ et $r_1r_2=c/a$.
\end{demonstration}
\begin{savoirfaire}[Expliquer une preuve]
Repérer l'hypothèse $a\neq0$, l'identité de départ et les trois cas possibles
pour le discriminant. La formule des solutions est une conséquence du calcul
sur un carré, pas une règle isolée à appliquer sans conditions.
\end{savoirfaire}
\clearpage
\otPageTitre{Démonstrations}{Suites arithmétiques, géométriques et sommes}
\phantomsection\label{dem-suites}
\essSection{2}{Termes généraux et sommes}
\begin{demonstration}[Suite arithmétique]
Passer de $u_0$ à $u_n$ nécessite $n$ ajouts de la raison $r$ :
$u_n=u_0+nr$. Entre deux indices $p$ et $n$, il y a $n-p$ ajouts,
donc $u_n=u_p+(n-p)r$.
Pour additionner $N$ termes consécutifs d'une suite arithmétique, on écrit
la somme dans les deux sens. Chaque paire « premier plus dernier »,
« deuxième plus avant-dernier », etc. a la même somme.
On obtient $2S=N(\text{premier}+\text{dernier})$.
\end{demonstration}
\begin{demonstration}[Les entiers de 1 à $n$]
Pour $n\geq1$, écrivons $S=1+2+\cdots+n$, puis $S=n+(n-1)+\cdots+1$.
En additionnant terme à terme, les $n$ colonnes valent toutes $n+1$ :
$2S=n(n+1)$, donc $S=n(n+1)/2$.
Pour $n=0$, la somme vide vaut 0 et la même formule reste valable.
\end{demonstration}
\begin{demonstration}[Suite géométrique]
Passer de $v_0$ à $v_n$ revient à multiplier $n$ fois par $q$, donc
$v_n=v_0q^n$. Pour $q=0$, cette formule donne les termes de rang $n\geq1$,
tous nuls ; le terme de rang 0 reste le terme initial donné.
Plus généralement $v_n=v_pq^{n-p}$ pour $n>p$ ; si l'on veut remonter
vers un indice inférieur, il faut $q\neq0$.
\end{demonstration}
\begin{demonstration}[Une somme géométrique]
Soit $G=1+q+\cdots+q^n$. Alors $qG=q+q^2+\cdots+q^{n+1}$.
En soustrayant, $(1-q)G=1-q^{n+1}$.
Si $q\neq1$, $G=(1-q^{n+1})/(1-q)$. Si $q=1$, il y a $n+1$ termes
égaux à 1, donc $G=n+1$. La séparation de ces deux cas évite de diviser par 0.
\end{demonstration}
\clearpage
\otPageTitre{Démonstrations}{Nombre dérivé et tangente}
\phantomsection\label{dem-derivation}
\essSection{3}{Dériver à partir d'un taux de variation}
\begin{demonstration}[La fonction carré]
Pour $f(x)=x^2$, $h\neq0$ et $a\in\R$,
\[\frac{f(a+h)-f(a)}h=\frac{2ah+h^2}h=2a+h.\]
Quand $h$ se rapproche de 0, ce taux se rapproche de $2a$.
Ainsi $f'(a)=2a$, pour tout réel $a$.
\end{demonstration}
\begin{demonstration}[La fonction inverse]
Pour $a\neq0$, $h\neq0$ et $a+h\neq0$,
\[\frac{1/(a+h)-1/a}{h}=-\frac1{a(a+h)}.\]
Quand $h$ se rapproche de 0, ce taux se rapproche de $-1/a^2$.
Ainsi la dérivée de $x\mapsto1/x$ est $-1/x^2$ sur $\R\setminus\{0\}$.
\end{demonstration}
\begin{demonstration}[Équation de la tangente]
La tangente en $A(a;f(a))$ a pour pente $f'(a)$.
Une droite de pente $m$ passant par $(a;b)$ s'écrit $y-b=m(x-a)$.
En remplaçant $m$ par $f'(a)$ et $b$ par $f(a)$, on obtient
$y=f(a)+f'(a)(x-a)$.
\end{demonstration}
\begin{exemple}
La tangente à $y=x^2$ en $a=2$ a pour équation $y=4+4(x-2)=4x-4$.
L'approximation affine associée donne $(2+h)^2\approx4+4h$ pour un petit $h$.
Le terme exact oublié est $h^2$, ce qui permet ici de mesurer l'erreur.
\end{exemple}
\clearpage
\otPageTitre{Démonstrations}{Produit et absence de dérivabilité}
\essSection{4}{Ce que signifie un nombre dérivé}
\begin{demonstration}[La fonction racine carrée en 0]
La fonction est définie pour $x\geq0$ et vaut 0 en 0. Pour $h>0$,
le taux de variation entre 0 et $h$ vaut
\[\frac{\sqrt h-0}{h}=\frac1{\sqrt h}.\]
Ces taux grandissent sans borne lorsque $h$ positif devient très petit.
Ils ne se rapprochent pas d'un nombre réel fini : la fonction racine carrée
n'est pas dérivable en 0.
\end{demonstration}
\begin{demonstration}[Dérivée d'un produit]
Soient $u$ et $v$ dérivables en $a$. Pour $h\neq0$, ajouter puis retrancher
$u(a+h)v(a)$ dans le numérateur donne
\begin{align*}
&\frac{u(a+h)v(a+h)-u(a)v(a)}h\\
&=u(a+h)\frac{v(a+h)-v(a)}h
+v(a)\frac{u(a+h)-u(a)}h.
\end{align*}
Quand $h$ se rapproche de 0, les deux taux se rapprochent de $v'(a)$ et
$u'(a)$. La dérivabilité de $u$ implique aussi que $u(a+h)$ se rapproche
de $u(a)$ : en effet $u(a+h)-u(a)$ est $h$ multiplié par un taux qui reste
proche du nombre fini $u'(a)$.
On obtient $(uv)'(a)=u(a)v'(a)+u'(a)v(a)$.
\end{demonstration}
\begin{attention}
Le raisonnement utilise la dérivabilité des deux fonctions. Il ne suffit
pas qu'elles soient définies. L'écriture de la dérivée d'un produit doit
conserver les deux contributions : l'une provenant de la variation de $u$,
l'autre de celle de $v$.
\end{attention}
\clearpage
\otPageTitre{Démonstrations}{Valeurs trigonométriques remarquables}
\phantomsection\label{dem-trigonometrie}
\essSection{5}{Construire les valeurs à partir de triangles}
\begin{demonstration}[L'angle $\pi/4$]
Dans un triangle rectangle isocèle dont les deux petits côtés mesurent 1,
Pythagore donne une hypoténuse $\sqrt2$. Les angles aigus mesurent $\pi/4$.
Le rapport « adjacent sur hypoténuse » et le rapport « opposé sur hypoténuse »
valent tous deux $1/\sqrt2=\sqrt2/2$.
Donc $\cos(\pi/4)=\sin(\pi/4)=\sqrt2/2$.
\end{demonstration}
\figureGeo[corrige]{preuve-pi4}
\begin{demonstration}[Les angles $\pi/3$ et $\pi/6$]
Partons d'un triangle équilatéral de côté 2. Une hauteur le partage en deux
triangles rectangles. La demi-base mesure 1, l'hypoténuse 2 et la hauteur
$\sqrt{2^2-1^2}=\sqrt3$.
Pour l'angle $\pi/3$, le côté adjacent mesure 1 et le côté opposé $\sqrt3$,
donc $\cos(\pi/3)=1/2$ et $\sin(\pi/3)=\sqrt3/2$.
Pour l'autre angle, $\pi/6$, les deux rôles s'échangent :
$\cos(\pi/6)=\sqrt3/2$ et $\sin(\pi/6)=1/2$.
\end{demonstration}
\begin{center}\begin{tikzpicture}[x=18mm,y=18mm]
\draw[otNavy,thick](0,0)--(2,0)--(1,1.73205)--cycle;
\draw[otGreenDark,dashed](1,1.73205)--(1,0);
\node[below]at(.5,0){$1$};\node[below]at(1.5,0){$1$};
\node[left]at(.4,.85){$2$};\node[right]at(1.6,.85){$2$};
\node[right,font=\small]at(1,.8){$\sqrt3$};
\draw[otMuted](1,.15)--(1.15,.15)--(1.15,0);
\end{tikzpicture}\end{center}
\begin{propriete}[Relier le triangle au cercle unité]
Diviser toutes les longueurs par l'hypoténuse place le point correspondant
sur le cercle unité. Ses coordonnées sont les rapports définissant cosinus
et sinus. Les symétries du cercle donnent ensuite les angles associés.
\end{propriete}
\clearpage
\otPageTitre{Démonstrations}{Produit scalaire et géométrie}
\phantomsection\label{dem-geometrie}
\essSection{6}{Deux preuves à partir des vecteurs}
\begin{demonstration}[La formule d'Al-Kashi]
Par Chasles, $\vect{BC}=\vect{AC}-\vect{AB}$.
Développer le carré de la norme donne
\[BC^2=AC^2+AB^2-2\vect{AB}\cdot\vect{AC}.
\]
Comme $\vect{AB}\cdot\vect{AC}=AB\,AC\cos\widehat{BAC}$,
on obtient $BC^2=AB^2+AC^2-2AB\,AC\cos\widehat{BAC}$.
Pour un angle droit en $A$, le cosinus est nul : on retrouve Pythagore.
\end{demonstration}
\begin{demonstration}[Le cercle de diamètre $[AB]$]
Soit $I$ le milieu de $[AB]$, avec $A\neq B$.
On a $\vect{MA}=\vect{MI}+\vect{IA}$ et
$\vect{MB}=\vect{MI}-\vect{IA}$.
Par bilinéarité et symétrie,
\[\vect{MA}\cdot\vect{MB}=MI^2-IA^2=MI^2-\frac{AB^2}{4}.
\]
Le produit est nul si et seulement si $MI=AB/2$, puisque les distances sont
positives ou nulles. C'est précisément le cercle de centre $I$ et de rayon
$AB/2$, donc de diamètre $[AB]$.
Les points $A$ et $B$ sont inclus : l'un des vecteurs du produit y est nul.
\end{demonstration}
\begin{center}\legende{Les deux configurations sont illustrées dans les repères visuels du chapitre 7.}\end{center}
\begin{savoirfaire}[Choisir un point intermédiaire]
Le milieu rend deux vecteurs opposés, ce qui supprime les termes croisés.
Le choix d'un point intermédiaire dans Chasles constitue souvent la première
étape utile d'une preuve géométrique sans coordonnées.
\end{savoirfaire}
\clearpage
\otPageTitre{Démonstrations}{Probabilités et moments d'une variable}
\phantomsection\label{dem-probabilites}
\essSection{7}{De sommes finies aux formules usuelles}
\begin{demonstration}[Probabilités totales]
Pour une partition $A_1,\ldots,A_k$, les évènements $A_i\cap B$ sont deux à deux
incompatibles et leur réunion est $B$. Ainsi $P(B)=\sum_i P(A_i\cap B)$.
Lorsque les $P(A_i)$ sont positifs, chaque terme s'écrit $P(A_i)P_{A_i}(B)$,
ce qui donne la formule des probabilités totales.
\end{demonstration}
\begin{demonstration}[Linéarité de l'espérance]
Sur un univers fini, on peut calculer $E(X)$ comme
$\sum_{\omega\in\Omega}X(\omega)P(\{\omega\})$, en regroupant ou non les issues
qui donnent la même valeur. Distribuer une somme finie donne alors
\[E(aX+bY)=\sum_{\omega\in\Omega}(aX(\omega)+bY(\omega))P(\{\omega\})
=aE(X)+bE(Y).\]
Aucune hypothèse d'indépendance n'intervient dans ce calcul.
\end{demonstration}
\begin{demonstration}[König-Huygens et transformation affine]
Avec $\mu=E(X)$, développer le carré donne
$V(X)=E(X^2-2\mu X+\mu^2)=E(X^2)-2\mu^2+\mu^2=E(X^2)-\mu^2$.
Pour $Y=aX+c$, $Y-E(Y)=a(X-E(X))$, donc
$V(Y)=a^2V(X)$ et $\sigma(Y)=|a|\sigma(X)$.
\end{demonstration}
\begin{exemple}[Une variable de Bernoulli]
Si $I$ vaut 0 ou 1 et $P(I=1)=p$, on a $I^2=I$.
Donc $E(I)=p$, $E(I^2)=p$ et $V(I)=p-p^2=p(1-p)$.
\end{exemple}
\clearpage
\pagestyle{otMain}

\partieCorriges
% !TEX root = ../main.tex
\corrigesChapitre{1}{Polynômes du second degré}
\corrigesSouples

\begin{blocLivre}{cor-1-1}
\begin{corrige}{1}{Reconnaître un trinôme}
\methode{développer pour identifier les coefficients}
Les fonctions du second degré sont $3x^2-2x+5$ et $(x-2)^2=x^2-4x+4$.
Leurs coefficients sont respectivement $(3;-2;5)$ et $(1;-4;4)$.
La fonction $4x-1$ est de degré 1 ; $x^3+x$ est de degré 3.
\end{corrige}
\end{blocLivre}
\begin{blocLivre}{cor-1-2}
\begin{corrige}{2}{Développer ou factoriser}
\demarche{Choisir entre double distributivité, identité remarquable et facteur commun.}
\begin{questions}
\item Par double distributivité, $(x-3)(x+2)=x^2+2x-3x-6=\res{x^2-x-6}$.
Puis $(x-1)^2=x^2-2x+1$, donc $2(x-1)^2-8=2x^2-4x+2-8=\res{2x^2-4x-6}$.
\item Une différence de carrés donne $x^2-25=x^2-5^2=\res{(x-5)(x+5)}$.
Le facteur commun $3x$ donne $3x^2-12x=\res{3x(x-4)}$. Développer ces produits permet de les contrôler.
\end{questions}

\end{corrige}
\end{blocLivre}
\begin{blocLivre}{cor-1-3}
\begin{corrige}{3}{Lire la forme canonique}
\demarche{Lire le sommet dans le carré, puis isoler ce carré pour les racines.}
\begin{questions}
\item Le carré est minimal en $x=3$ : sommet $\res{S(3;-8)}$, axe $x=3$, minimum $-8$.
\item $f(0)=18-8=10$. $f(x)=0\iff(x-3)^2=4\iff x-3=\pm2$,
donc $\res{x=1\text{ ou }x=5}$.
\end{questions}
\end{corrige}
\end{blocLivre}
\begin{blocLivre}{cor-1-4}
\begin{corrige}{4}{Calculer un discriminant}
\demarche{Identifier les trois coefficients avant de calculer le discriminant.}
\begin{questions}
\item Ici $(a;b;c)=(1;-5;6)$ et $\Delta=(-5)^2-4\times1\times6=1>0$.
Les solutions sont $(5-1)/2=2$ et $(5+1)/2=3$, donc $\res{S=\{2;3\}}$.
\item Ici $\Delta=4^2-4\times1\times4=0$. La racine double vaut $-4/(2\times1)=-2$ : $\res{S=\{-2\}}$.
On retrouve $(x+2)^2=0$.
\item Ici $(a;b;c)=(2;1;2)$ et $\Delta=1^2-4\times2\times2=-15<0$.
Le trinôme ne s’annule pas sur $\R$ : $\res{S=\varnothing}$.
\end{questions}

\end{corrige}
\end{blocLivre}
\begin{blocLivre}{cor-1-5}
\begin{corrige}{5}{Signe d'une forme factorisée}
\demarche{Étudier les deux facteurs et tenir compte du coefficient négatif.}
Les facteurs s’annulent en $-1$ et 4. Chacun est négatif à gauche de sa racine et positif à droite.
On multiplie leurs signes, puis on inverse le signe obtenu à cause du coefficient $-2$.
Le tableau ci-dessous donne un signe positif entre les racines et négatif à l’extérieur.
Ainsi $\res{P(x)\geq0\iff x\in[-1;4]}$ et
$\res{P(x)<0\iff x\in\mathopen{]}-\infty;-1[\cup]4;+\infty[}$.
Les racines sont incluses pour l’inégalité large, et exclues pour l’inégalité stricte.

\label{figlivre-r01-05-signes-cor}
\figureLivre[corrige]{r01-05-signes}

\end{corrige}
\end{blocLivre}
\begin{blocLivre}{cor-1-6}
\begin{corrige}{6}{Choisir la bonne forme}
\demarche{Choisir la forme qui rend chaque équation ou inéquation la plus courte.}
\begin{questions}
\item $(x-3)^2-4=x^2-6x+5=(x-1)(x-5)$ en développant.
\item $f(x)=5\iff x(x-6)=0$ : $\res{x=0\text{ ou }6}$.
Le trinôme est négatif ou nul entre ses racines : $\res{x\in[1;5]}$ pour $f(x)\leq0$.
\item La forme canonique donne $(x-3)^2\leq1$, donc $-1\leq x-3\leq1$ :
$\res{x\in[2;4]}$.
\end{questions}
\end{corrige}
\end{blocLivre}
\begin{blocLivre}{cor-1-7}
\begin{corrige}{7}{Deux paraboles}
\demarche{Comparer les courbes en étudiant la différence des deux fonctions.}
\begin{questions}
\item Le dessin suggère des intersections pour $x=0$ et $x=3$.
\item $f(x)-g(x)=2x^2-6x=2x(x-3)$.
Les intersections sont donc $\res{(0;3)\text{ et }(3;0)}$.
La différence est négative ou nulle entre 0 et 3, donc $\res{f\leq g\iff x\in[0;3]}$.
\end{questions}
\methode{le dessin suggère ; la factorisation démontre}
\end{corrige}
\end{blocLivre}
\begin{blocLivre}{cor-1-8}
\begin{corrige}{8}{Un quotient et une valeur interdite}
\demarche{Exclure le zéro du dénominateur avant toute étude de signe.}
Le domaine est $\R\setminus\{1\}$ et $x^2-5x+6=(x-2)(x-3)$.
\begin{center}\begin{tableau}
\otTab[2.1]{$x$/0.8,$x-2$/0.7,$x-3$/0.7,$x-1$/0.7,$Q(x)$/0.8}{$-\infty$,$1$,$2$,$3$,$+\infty$}
\tkzTabLine{,-,t,-,z,+,t,+,}
\tkzTabLine{,-,t,-,t,-,z,+,}
\tkzTabLine{,-,z,+,t,+,t,+,}
\tkzTabLine{,-,d,+,z,-,z,+,}
\end{tableau}\end{center}
Ainsi $\res{S=]1;2]\cup[3;+\infty[}$. Le point 1 est exclu ; 2 et 3 sont inclus.
\end{corrige}
\end{blocLivre}
\begin{blocLivre}{cor-1-9}
\begin{corrige}{9}{Retrouver une fonction}
\demarche{Utiliser les racines pour factoriser, puis le point donné pour fixer le coefficient.}
\begin{questions}
\item Les deux racines donnent $f(x)=a(x+2)(x-4)$. Avec $f(0)=-8$,
$-8a=-8$, donc $\res{a=1}$.
\item $f(x)=x^2-2x-8=(x-1)^2-9$. Le sommet est $\res{S(1;-9)}$.
\item $f(x)=-5\iff(x-1)^2=4$, donc $\res{x=-1\text{ ou }3}$.
\end{questions}
\label{figlivre-r01-09-parabole-cor}
\figureLivre[corrige]{r01-09-parabole}

\end{corrige}
\end{blocLivre}
\begin{blocLivre}{cor-1-10}
\begin{corrige}{10}{Somme et produit des racines}
\demarche{Développer le produit associé aux deux racines pour relier somme et produit.}
Si $u$ et $v$ sont les deux nombres cherchés, $(x-u)(x-v)=x^2-(u+v)x+uv$.
\begin{questions}
\item La somme 7 et le produit 10 conduisent à $x^2-7x+10=0$.
Comme $x^2-7x+10=(x-2)(x-5)$, les nombres sont $\res{2\text{ et }5}$.
\item La somme 4 et le produit 5 conduisent à $x^2-4x+5=0$.
Son discriminant vaut $16-20=-4<0$ : il n’existe aucun couple de réels qui convienne.
\item Pour $x^2+x-12$, la somme des racines doit être $-1$ et leur produit $-12$.
Les nombres $-4$ et 3 conviennent, donc $\res{x^2+x-12=(x+4)(x-3)}$.
\end{questions}

\end{corrige}
\end{blocLivre}
\begin{blocLivre}{cor-1-11}
\begin{corrige}{11}{Discuter selon un paramètre}
\demarche{Comparer le minimum du trinôme à zéro, en gardant le paramètre fixé.}
\begin{questions}
\item $x^2-4x+m=\res{(x-2)^2+m-4}$.
\item Résoudre revient à $(x-2)^2=4-m$. Si $m<4$, deux solutions
$2\pm\sqrt{4-m}$ ; si $m=4$, une solution double 2 ; si $m>4$, aucune solution réelle.
\item Pour $m=3$ : $\{1;3\}$ ; pour $m=4$ : $\{2\}$ ; pour $m=5$ : $\varnothing$.
\item Le minimum, atteint en 2, vaut $m-4$. Il est positif ou nul si et seulement si $\res{m\geq4}$.
\end{questions}
\end{corrige}
\end{blocLivre}
\begin{blocLivre}{cor-1-12}
\begin{corrige}{12}{Un jardin contre un mur}
\demarche{Traduire la longueur de clôture avant de chercher le maximum de l’aire.}
\begin{questions}
\item La longueur restante est $40-2x>0$ ; avec $x>0$, on obtient $0<x<20$.
L'aire vaut $A(x)=x(40-2x)=\res{-2x^2+40x}$.
\item $A(x)=-2(x-10)^2+200$. Le carré étant positif ou nul, le maximum est
$\res{200\text{ m}^2}$ pour $x=10$ m ; la longueur vaut 20 m.
\item $A(x)\geq150\iff(x-10)^2\leq25$,
donc $\res{x\in[5;15]}$ mètres, valeurs toutes admissibles.
\end{questions}
\end{corrige}
\end{blocLivre}
\begin{blocLivre}{cor-1-13}
\begin{corrige}{13}{Une vente rentable}
\demarche{Calculer recette moins coût, puis limiter la conclusion au domaine commercial.}
\begin{questions}
\item La recette est le prix multiplié par la quantité :
$R(q)=q(50-0{,}2q)=-0{,}2q^2+50q$.
Le bénéfice est $B=R-C$, donc $\res{B(q)=-0{,}2q^2+40q-500}$.
\item En développant $-0{,}2(q-100)^2+1500$, on retrouve $B$.
Le maximum sur $[0;100]$ vaut $\res{1500\text{ euros}}$, atteint pour $q=100$.
\item $B(q)\geq0\iff(q-100)^2\leq7500$.
Les racines sont $100\pm50\sqrt3$, soit environ $13{,}40$ et $186{,}60$.
En tenant compte du domaine, $q\in[100-50\sqrt3;100]$.
Les quantités entières rentables sont donc $\res{14\leq q\leq100}$.
\end{questions}
\begin{attention}
Le seuil supérieur calculé dépasse le domaine du modèle. On ne prolonge pas
la conclusion économique au-delà de 100 articles.
\end{attention}
\end{corrige}
\end{blocLivre}
\begin{blocLivre}{cor-1-14}
\begin{corrige}{14}{Un ballon au-dessus d'un filet}
\demarche{Étudier la hauteur sur les temps physiquement admissibles.}
\begin{questions}
\item $h(0)=1$ m et $h(t)=-5(t-1)^2+6$ : maximum $\res{6\text{ m}}$ à $t=1$ s.
\item Les solutions sont $1\pm\sqrt{30}/5$. Seule la positive est admissible :
$\res{t_{\text{sol}}=1+\sqrt{30}/5\approx2{,}10\text{ s}}$.
\item $h(t)\geq4\iff(t-1)^2\leq2/5$, donc
$\res{t\in[1-\sqrt{10}/5;1+\sqrt{10}/5]}$, environ $[0{,}37;1{,}63]$ s.
\item $h(1{,}5)=-11{,}25+15+1=4{,}75$ m.
Le ballon franchit le filet avec une marge verticale de $\res{0{,}75\text{ m}}$.
\end{questions}
\label{figlivre-r01-14-ballon-cor}
\figureLivre[corrige]{r01-14-ballon}

\end{corrige}
\end{blocLivre}
\begin{blocLivre}{cor-1-15}
\begin{corrige}{15}{Une famille de paraboles}
\demarche{Étudier le signe du discriminant comme une expression du paramètre.}
\begin{questions}
\item $\Delta=(-2m)^2-4(m+2)=\res{4(m-2)(m+1)}$.
\item Il y a deux racines distinctes si $m<-1$ ou $m>2$ ; une double si
$m=-1$ ou $m=2$ ; aucune si $-1<m<2$.
\item $P_m(1/2)=1/4-m+m+2=9/4$ : le point commun est $\res{A(1/2;9/4)}$.
\item $P_3(x)=x^2-6x+5=(x-1)(x-5)$. Ainsi $P_3(x)\leq0$ pour
$\res{x\in[1;5]}$ et son sommet est $\res{S(3;-4)}$.
\item Le coefficient dominant vaut 1. Le trinôme est strictement positif
partout exactement lorsqu'il n'a aucune racine réelle : $\res{-1<m<2}$.
\end{questions}
\label{figlivre-r01-15-parametre-cor}
\figureLivre[corrige]{r01-15-parametre}

\end{corrige}
\end{blocLivre}

\clearpage
% !TEX root = ../main.tex
\corrigesChapitre{2}{Suites numériques}
\corrigesSouples

\begin{blocLivre}{cor-2-1}
\begin{corrige}{1}{Substituer un indice}
\demarche{Remplacer chaque occurrence de l’indice par la même valeur.}
\begin{questions}
\item On remplace l’indice dans les trois termes : $u_0=0-0+2=\res{2}$,
$u_1=1-3+2=\res{0}$, $u_2=4-6+2=\res{0}$ et $u_5=25-15+2=\res{12}$.
\item $u_{n+1}=(n+1)^2-3(n+1)+2=n^2+2n+1-3n-3+2=n^2-n$.
On soustrait ensuite l’expression entière de $u_n$ :
\[u_{n+1}-u_n=(n^2-n)-(n^2-3n+2)=\res{2n-2}.\]
Les parenthèses évitent de perdre le signe moins devant $u_n$.
\end{questions}

\end{corrige}
\end{blocLivre}
\begin{blocLivre}{cor-2-2}
\begin{corrige}{2}{Calculer de proche en proche}
\demarche{Calculer dans l’ordre : chaque terme réutilise le terme précédent.}
\begin{questions}\item $v_1=3\times2-1=5$, puis $\res{v_2=14,\ v_3=41,\ v_4=122}$.
\item Les différences sont $3$, $9$, $27$ ; les rapports sont $5/2$, $14/5$, $41/14$.
Ils ne sont pas constants : la suite n'est ni arithmétique ni géométrique.
\end{questions}
\end{corrige}
\end{blocLivre}
\begin{blocLivre}{cor-2-3}
\begin{corrige}{3}{Une suite arithmétique}
\demarche{Employer la formule arithmétique, puis résoudre une inégalité sur les indices entiers.}
\begin{questions}\item $\res{u_n=7-3n}$ et $u_8=-17$. Comme la raison est négative, la suite est strictement décroissante.
\item $7-3n\leq-20\iff n\geq9$ : les indices cherchés sont tous les entiers $\res{n\geq9}$.
\end{questions}
\end{corrige}
\end{blocLivre}
\begin{blocLivre}{cor-2-4}
\begin{corrige}{4}{Une suite géométrique}
\demarche{Employer la formule géométrique et compter les termes de la somme.}
\begin{questions}\item $\res{v_n=5\times2^n}$, $v_4=80$ et $v_6=320$.
\item Directement : $5+10+20+40+80=155$.
Avec la formule : $5(1-2^5)/(1-2)=\res{155}$. Il y a cinq termes.
\end{questions}
\end{corrige}
\end{blocLivre}
\begin{blocLivre}{cor-2-5}
\begin{corrige}{5}{Du pourcentage à la formule}
\demarche{Transformer la baisse annuelle en coefficient multiplicateur.}
\begin{questions}\item Le coefficient multiplicateur est $1-0{,}15=0{,}85$ :
$V_{n+1}=0{,}85V_n$, $V_0=800$, donc $\res{V_n=800\times0{,}85^n}$.
\item $V_2=578$ euros. La baisse vaut $1-0{,}85^2=0{,}2775$, soit
$\res{27{,}75\,\%}$. Les pourcentages successifs se multiplient via leurs coefficients.
\end{questions}
\end{corrige}
\end{blocLivre}
\begin{blocLivre}{cor-2-6}
\begin{corrige}{6}{Retrouver les paramètres}
\demarche{Compter les transitions entre les deux indices connus.}
\begin{questions}
\item \methode{utiliser l'écart entre les indices}
$u_8-u_3=5r=20$, donc $r=4$ ; $u_0=u_3-3r=-1$.
Ainsi $\res{u_n=4n-1}$.
\item $v_5=v_2q^3$ donne $q^3=8$, donc $q=2$.
$v_0=12/2^2=3$, d'où $\res{v_n=3\times2^n}$.
\end{questions}
\end{corrige}
\end{blocLivre}
\begin{blocLivre}{cor-2-7}
\begin{corrige}{7}{Compter les termes d'une somme}
\demarche{Identifier premier terme, dernier terme et nombre de termes avant de sommer.}
\begin{questions}
\item $61=4+19\times3$ : vingt termes, de rang 0 à 19.
La somme vaut $20(4+61)/2=\res{650}$.
\item $384=3\times2^7$ : huit termes. La somme vaut
$3(2^8-1)=\res{765}$.
\item De 3 à 10 inclus, huit termes ; le premier vaut 7, le dernier 21.
La somme est $8(7+21)/2=\res{112}$.
\end{questions}
\end{corrige}
\end{blocLivre}
\begin{blocLivre}{cor-2-8}
\begin{corrige}{8}{Comparer des variations}
\demarche{Calculer la différence de deux termes consécutifs pour étudier les variations.}
\begin{questions}
\item $u_{n+1}-u_n=2n-3$, négatif pour $n=0,1$, positif pour $n\geq2$.
La suite décroît jusqu'au rang 2, puis croît ; elle n'est pas monotone sur tous les indices.
\item Au même dénominateur,
$v_{n+1}-v_n=\dfrac{(n+2)^2-(n+1)(n+3)}{(n+2)(n+3)}
=\dfrac1{(n+2)(n+3)}>0$. La suite est strictement croissante.
\item $v_n=(n+2-1)/(n+2)=1-1/(n+2)$.
Les termes restent inférieurs à 1 et s'en rapprochent quand $n$ devient grand.
Il s'agit ici d'une lecture intuitive du comportement à long terme.
\end{questions}
\end{corrige}
\end{blocLivre}
\begin{blocLivre}{cor-2-9}
\begin{corrige}{9}{Lire puis justifier un nuage de points}
\demarche{Séparer la conjecture graphique de la justification donnée par la récurrence.}
\begin{questions}
\item Les termes semblent divisés par 2 à chaque étape : $u_n=16(1/2)^n$.
\item La récurrence donnée est précisément celle d'une suite géométrique
de raison $1/2$, de premier terme 16, ce qui justifie la formule.
Les termes sont positifs et la raison est dans $]0;1[$ : décroissance stricte.
\item $u_7=0{,}125$ et $u_8=0{,}0625$. Par décroissance, le premier rang
tel que $u_n<0{,}1$ est $\res{n=8}$.
\end{questions}
\end{corrige}
\end{blocLivre}
\begin{blocLivre}{cor-2-10}
\begin{corrige}{10}{Comprendre une boucle Python}
\demarche{Distinguer le nombre de transitions et le nombre de valeurs stockées.}
\begin{questions}
\item Les valeurs successives sont 3, 7, 15, 31, 63. La liste finale est
\code{[3, 7, 15, 31, 63]}.
\item $u_0=3$, $u_{n+1}=2u_n+1$. Après quatre passages, on a calculé $u_4$.
\item Remplacer \code{range(4)} par \code{range(10)}.
La liste comprend alors onze termes, de $u_0$ à $u_{10}$.
\end{questions}
\end{corrige}
\end{blocLivre}
\begin{blocLivre}{cor-2-11}
\begin{corrige}{11}{Un seuil d'épargne}
\demarche{Justifier la croissance, puis vérifier le seuil et son prédécesseur.}
\begin{questions}
\item La raison vaut $1{,}05>1$ et le premier terme est positif : croissance stricte.
\item On initialise $n=0$ et $C=1000$, puis on actualise ces deux variables ensemble.
\end{questions}
\begin{python}
n = 0
C = 1000
while C < 1500:
    C = 1.05*C
    n = n + 1
\end{python}
\otLabel{3.}\enspace $C_8\approx1477{,}46$ et $C_9\approx1551{,}33$ euros.
La première année cherchée est donc $\res{n=9}$.

À chaque test, les variables représentent bien $C_n$ et son rang $n$. L’arrêt à 9 donne un candidat ; la croissance et la vérification de $C_8$ prouvent que ce candidat est le premier.

\end{corrige}
\end{blocLivre}
\begin{blocLivre}{cor-2-12}
\begin{corrige}{12}{Des rangées de sièges}
\demarche{Respecter le premier indice 1 et additionner des rangées consécutives.}
\begin{questions}
\item \methode{ici le premier indice est 1}
$u_n=18+2(n-1)=\res{2n+16}$, donc $u_{20}=56$.
\item La capacité est $20(18+56)/2=\res{740\text{ places}}$.
\item Les dix premières rangées contiennent $10(18+36)/2=270$ sièges ;
les dix dernières, $10(38+56)/2=470$.
La recette vaut $9\times270+12\times470=\res{8070\text{ euros}}$.
\item Avec $N$ rangées, la capacité est $S_N=N(18+2N+16)/2=N(N+17)$.
Elle croît avec $N\geq1$. Or $S_{24}=984<1000$ et $S_{25}=1050\geq1000$.
Il faut donc au moins $\res{25\text{ rangées}}$.
\end{questions}
\end{corrige}
\end{blocLivre}
\begin{blocLivre}{cor-2-13}
\begin{corrige}{13}{Un traitement répété}
\demarche{Soustraire la valeur d’équilibre pour obtenir une suite géométrique.}
\begin{questions}
\item On conserve $0{,}8u_n$, puis on ajoute 6 : $u_{n+1}=0{,}8u_n+6$.
$u_1=14$ mg et $u_2=17{,}2$ mg.
\item $v_{n+1}=u_{n+1}-30=0{,}8u_n-24=0{,}8(u_n-30)=0{,}8v_n$.
\item $v_0=-20$, donc $v_n=-20\times0{,}8^n$ et
$\res{u_n=30-20\times0{,}8^n}$.
Puis $u_{n+1}-u_n=4\times0{,}8^n>0$ et $30-u_n=20\times0{,}8^n>0$.
\item $u_{10}\approx27{,}853<28$ et $u_{11}\approx28{,}282\geq28$.
Par croissance, le premier indice est $\res{n=11}$.
\end{questions}
\label{figlivre-r02-13-cycle-cor}
\figureLivre[corrige]{r02-13-cycle}

\end{corrige}
\end{blocLivre}
\begin{blocLivre}{cor-2-14}
\begin{corrige}{14}{Deux contrats de rémunération}
\demarche{Distinguer indice de la suite et numéro de l’année.}
\begin{questions}
\item $A_n=24000+1000n$ et $B_n=24000\times1{,}04^n$.
La sixième année correspond à $n=5$ : $A_5=29000$ euros et
$B_5\approx29199{,}67$ euros ; B est supérieur.
\item Les dix premières années vont de $n=0$ à $n=9$.
\[\sum_{n=0}^9 A_n=10\frac{24000+33000}2=\res{285000},\]
\[\sum_{n=0}^9 B_n=24000\frac{1{,}04^{10}-1}{0{,}04}\approx\res{288146{,}57}.\]
B rapporte environ 3146,57 euros de plus sur dix ans.
\item Pour $n=1,2,3$, les salaires B valent respectivement
24960 ; 25958,40 ; 26996,736 euros, inférieurs à 25000 ; 26000 ; 27000.
Pour $n=4$, $B_4=28076{,}60544>28000=A_4$.
Le premier indice est donc $\res{n=4}$, soit la cinquième année.
\end{questions}
\label{figlivre-r02-14-salaires-cor}
\figureLivre[corrige]{r02-14-salaires}

\end{corrige}
\end{blocLivre}
\begin{blocLivre}{cor-2-15}
\begin{corrige}{15}{Des carrés construits avec des allumettes}
\demarche{Compter les allumettes nouvelles, puis additionner les constructions séparées.}
\begin{questions}
\item Le premier carré utilise 4 allumettes ; chaque carré supplémentaire en ajoute 3.
Ainsi $a_1=4$, $a_{n+1}=a_n+3$ et $\res{a_n=3n+1}$.
\item $3n+1\leq100\iff n\leq33$ : on peut aligner $\res{33\text{ carrés}}$,
avec exactement 100 allumettes.
\item Les constructions sont séparées : $T_N=\sum_{n=1}^N(3n+1)$.
Donc $T_N=3N(N+1)/2+N=\res{(3N^2+5N)/2}$.
\item $T_{25}=1000$ et $T_{26}=1079$. Comme $T_N$ est croissant,
le maximum possible est $\res{N=25}$.
\end{questions}
\end{corrige}
\end{blocLivre}

\clearpage
% !TEX root = ../main.tex
\corrigesChapitre{3}{Dérivation}
\corrigesSouples

\begin{blocLivre}{cor-3-1}
\begin{corrige}{1}{Calculer une pente moyenne}
\demarche{Utiliser la différence des images divisée par la différence des abscisses.}
\begin{questions}
\item $f(1)=-1$, $f(3)=3$, $f(2)=0$. Les taux valent
$(3+1)/(3-1)=\res{2}$ et $(0+1)/(2-1)=\res{1}$.
\item $f(1+h)=h^2-1$, donc le taux est $h^2/h=h$ pour $h\neq0$.
Il se rapproche de 0 avec $h$ : $\res{f'(1)=0}$.
\end{questions}
\end{corrige}
\end{blocLivre}
\begin{blocLivre}{cor-3-2}
\begin{corrige}{2}{Lire la dérivée sur une tangente}
\demarche{La pente de la tangente fournit le nombre dérivé au point de contact.}
\begin{questions}\item $A$ est sur la courbe : $f(2)=3$.
La pente de $(AB)$ vaut $(7-3)/(4-2)=2$, donc $\res{f'(2)=2}$.
\item $y=2(x-2)+3$, soit $\res{y=2x-1}$.
\end{questions}
\label{figlivre-r03-02-tangente-cor}
\figureLivre[corrige]{r03-02-tangente}

\end{corrige}
\end{blocLivre}
\begin{blocLivre}{cor-3-3}
\begin{corrige}{3}{Dérivées immédiates}
\methode{dériver chaque terme sur le domaine autorisé}\par
$\res{f'(x)=20x^3-6x+7}$ sur $\R$.\par
$\res{g'(x)=-3/x^2}$ sur $\R\setminus\{0\}$.\par
$\res{h'(x)=1/(2\sqrt x)+2}$ sur $]0;+\infty[$.
La fonction $h$ est définie en 0, mais la formule de dérivée n'y est pas valable.
Pour la puissance négative, $\res{k'(x)=-3x^{-4}=-3/x^4}$ sur $\R\setminus\{0\}$.
\end{corrige}
\end{blocLivre}
\begin{blocLivre}{cor-3-4}
\begin{corrige}{4}{Une tangente à un polynôme}
\demarche{Calculer séparément l’image et la dérivée avant de former la tangente.}
\begin{questions}\item $f'(x)=2x+3$, $f(1)=3$ et $f'(1)=5$.
\item $y=5(x-1)+3=\res{5x-2}$. La tangente passe bien par $(1;3)$.
\end{questions}
\label{figlivre-r03-04-tangente-cor}
\figureLivre[corrige]{r03-04-tangente}

\end{corrige}
\end{blocLivre}
\begin{blocLivre}{cor-3-5}
\begin{corrige}{5}{Reconnaître l'opération}
\methode{produit puis quotient}
$u'(x)=2(x^2-3)+(2x+1)2x=\res{6x^2+2x-6}$ sur $\R$.
Pour $x\neq2$, $v'(x)=\dfrac{x-2-(x+1)}{(x-2)^2}=\res{-3/(x-2)^2}$.
\end{corrige}
\end{blocLivre}
\begin{blocLivre}{cor-3-6}
\begin{corrige}{6}{Revenir à la définition}
\demarche{Développer le numérateur, puis simplifier seulement lorsque le dénominateur est non nul.}
\begin{questions}
\item $(a+h)^2-a^2=2ah+h^2$. En divisant par $h\neq0$, on obtient $2a+h$.
\item Lorsque $h$ se rapproche de 0, le taux se rapproche de $2a$ :
$\res{f'(a)=2a}$. La tangente est $y=2a(x-a)+a^2=\res{2ax-a^2}$.
\item Pour $a=2$, le taux vaut $4+h$ : on obtient
$\res{5;\ 4{,}1;\ 3{,}9;\ 4{,}01}$, à comparer au nombre dérivé 4.
\item La fonction accumule les pentes dans une liste, pour un pas donné :
\end{questions}
\begin{python}
def carre(x):
    return x**2
def secantes(h, n):
    pentes = []
    for k in range(n):
        x = k*h
        pentes.append((carre(x + h) - carre(x))/h)
    return pentes
\end{python}
D'après la question 1, la pente entre $kh$ et $(k+1)h$ vaut $2kh+h$.
Ainsi \code{secantes(0.5, 4)} renvoie $\res{[0{,}5;\ 1{,}5;\ 2{,}5;\ 3{,}5]}$, soit
\code{[0.5, 1.5, 2.5, 3.5]}, tandis que $f'(kh)=2kh$ donne $0$, $1$, $2$ et $3$.
Chaque pente dépasse le nombre dérivé au point de départ exactement de $h=0{,}5$ :
plus le pas est petit, plus les sécantes se rapprochent des tangentes.
\end{corrige}
\end{blocLivre}
\begin{blocLivre}{cor-3-7}
\begin{corrige}{7}{Dériver trois produits}
\demarche{Dans un produit, dériver successivement chacun des deux facteurs.}
\begin{questions}
\item Avec $u=x^2+1$, $u'=2x$, $v=3x-2$ et $v'=3$, la règle du produit donne
\[f'(x)=2x(3x-2)+3(x^2+1)=6x^2-4x+3x^2+3=\res{9x^2-4x+3}.\]
\item Pour $x>0$, $(x^2)'=2x$ et $(\sqrt x)'=1/(2\sqrt x)$.
Ainsi $g'(x)=2x\sqrt x+x^2/(2\sqrt x)$. Comme $x^2/\sqrt x=x\sqrt x$,
on obtient $\res{g'(x)=\frac52x\sqrt x}$.
\item En développant, $(x-1)^2(x+2)=x^3-3x+2$, donc $\res{h'(x)=3x^2-3}$.
Par produit, on trouve aussi $2(x-1)(x+2)+(x-1)^2=(x-1)(3x+3)$, le même résultat.
\end{questions}

\end{corrige}
\end{blocLivre}
\begin{blocLivre}{cor-3-8}
\begin{corrige}{8}{Dériver trois quotients}
\demarche{Exclure les zéros du dénominateur et simplifier les expressions quand cela raccourcit le calcul.}
\begin{questions}
\item Le dénominateur s’annule en $-3$. Pour $x\neq-3$,
\[f'(x)=\frac{2(x+3)-(2x-1)}{(x+3)^2}=\res{\frac7{(x+3)^2}}.\]
\item Pour $x\neq0$, $g(x)=x+1/x$, donc $\res{g'(x)=1-1/x^2}$.
Cette simplification réduit le calcul ; la valeur interdite 0 reste exclue.
\item Comme $x^2+1>0$, $h$ est dérivable sur $\R$.
Avec $u=x$, $u'=1$, $v=x^2+1$ et $v'=2x$,
\[h'(x)=\frac{(x^2+1)-x(2x)}{(x^2+1)^2}=\res{\frac{1-x^2}{(x^2+1)^2}}.\]
\end{questions}

\end{corrige}
\end{blocLivre}
\begin{blocLivre}{cor-3-9}
\begin{corrige}{9}{Des tangentes parallèles}
\demarche{Des tangentes parallèles ont le même coefficient directeur.}
\begin{questions}
\item $f'(x)=3x^2-3$.
\item La droite proposée a pour pente 9. On résout $f'(a)=9$, soit
$3a^2-3=9\iff a^2=4$. Les abscisses sont $\res{a=-2\text{ et }a=2}$.
\item $f(2)=3$, donc $T_2:y=9(x-2)+3=\res{9x-15}$.
$f(-2)=-1$, donc $T_{-2}:y=9(x+2)-1=\res{9x+17}$.
Les ordonnées à l'origine sont différentes : les droites sont distinctes.
\end{questions}
\label{figlivre-r03-09-paralleles-cor}
\figureLivre[corrige]{r03-09-paralleles}

\end{corrige}
\end{blocLivre}
\begin{blocLivre}{cor-3-10}
\begin{corrige}{10}{Approcher une racine carrée}
\demarche{Utiliser l’approximation affine au voisinage du nombre dont la racine est connue.}
\begin{questions}
\item $f(4)=2$ et $f'(4)=1/4$, donc $T:y=2+(x-4)/4=\res{x/4+1}$.
\item Avec $h=0{,}08$, $f(4+h)\approx2+h/4=\res{2{,}02}$.
\item $\sqrt{4{,}08}\approx2{,}019900988$.
L'erreur absolue est environ $0{,}000099012$, soit $\res{0{,}000099}$ à $10^{-6}$ près.
L'approximation est légèrement supérieure à la valeur exacte.
\end{questions}
\end{corrige}
\end{blocLivre}
\begin{blocLivre}{cor-3-11}
\begin{corrige}{11}{Deux exemples de non-dérivabilité}
\demarche{Comparer les taux de variation près de zéro, et pas seulement les valeurs des fonctions.}
\begin{questions}
\item Le taux vaut $|h|/h$, donc 1 pour $h>0$ et $-1$ pour $h<0$.
Les deux côtés ne donnent pas le même nombre : $|x|$ n'est pas dérivable en 0.
\item $\sqrt h/h=1/\sqrt h$. Ces valeurs grandissent sans borne quand $h>0$
se rapproche de 0 : elles ne se rapprochent pas d'un nombre réel fini.
La fonction racine carrée n'est donc pas dérivable en 0.
\item Les deux fonctions sont définies en 0 et y valent 0.
\textbf{Être définie n'implique pas être dérivable.}
\end{questions}
\label{figlivre-r03-11-nonderivable-cor}
\figureLivre[corrige]{r03-11-nonderivable}

\end{corrige}
\end{blocLivre}
\begin{blocLivre}{cor-3-12}
\begin{corrige}{12}{Tangentes à une parabole}
\demarche{Écrire une tangente générale puis imposer le passage par le point donné.}
\begin{questions}
\item Pour $f(x)=x^2$, $f'(a)=2a$ : $T_a:y=2ax-a^2$.
\item Le point $P(0;-1)$ est sur cette droite si $-1=-a^2$, soit $a^2=1$.
\item On trouve $a=1$ ou $a=-1$. Les contacts sont $A(1;1)$ et $B(-1;1)$.
Les tangentes sont $\res{y=2x-1}$ et $\res{y=-2x-1}$.
Elles passent toutes deux par $P$ et sont symétriques par rapport à l'axe des ordonnées.
\end{questions}
\begin{center}\begin{tikzpicture}[x=14mm,y=9mm]
\draw[otMuted,->](-2.1,0)--(2.4,0)node[right]{$x$};
\draw[otMuted,->](0,-1.6)--(0,3.9)node[above]{$y$};
\draw[otNavy,thick]plot[domain=-1.8:1.8,samples=60](\x,{\x*\x});
\draw[otGreen](-.1,-1.2)--(2,3);\draw[otGreen](-2,3)--(.1,-1.2);
\fill[otNavy](1,1)circle(1.5pt)node[right]{$A$};
\fill[otNavy](-1,1)circle(1.5pt)node[left]{$B$};
\fill[otGreenDark](0,-1)circle(1.5pt)node[right]{$P$};
\end{tikzpicture}\end{center}
\end{corrige}
\end{blocLivre}
\begin{blocLivre}{cor-3-13}
\begin{corrige}{13}{Vitesse moyenne et vitesse instantanée}
\demarche{Distinguer une vitesse moyenne sur un intervalle d’une vitesse à un instant.}
\begin{questions}
\item $s(1)=7$ m et $s(3)=9$ m : la vitesse moyenne vaut $(9-7)/2=\res{1\text{ m/s}}$.
\item $v(t)=3t^2-12t+12=\res{3(t-2)^2}$.
\item $v(t)=0$ uniquement à $\res{t=2\text{ s}}$ ; alors $s(2)=\res{8\text{ m}}$.
\item Un carré est positif ou nul, donc $v(t)\geq0$ sur $[0;4]$.
La vitesse s'annule un instant mais ne devient jamais négative.
\end{questions}
\end{corrige}
\end{blocLivre}
\begin{blocLivre}{cor-3-14}
\begin{corrige}{14}{Retrouver un polynôme grâce à une tangente}
\demarche{Traduire chaque donnée par une équation sur les coefficients.}
\begin{questions}
\item $c=1$, $4a+2b+c=3$ et, puisque $f'(x)=2ax+b$, $4a+b=4$.
\item On obtient $4a+2b=2$. Avec $b=4-4a$, cela donne $a=3/2$ et $b=-2$.
Ainsi $\res{f(x)=\frac32x^2-2x+1}$.
Vérifications : $f(0)=1$, $f(2)=6-4+1=3$, $f'(2)=6-2=4$.
\item $T:y=4(x-2)+3=\res{4x-5}$.
\item $f(x)-(4x-5)=\frac32x^2-6x+6=\res{\frac32(x-2)^2}\geq0$.
La courbe est au-dessus de la tangente et ne la rencontre qu'en $x=2$.
\end{questions}
\label{figlivre-r03-14-position-cor}
\figureLivre[corrige]{r03-14-position}

\end{corrige}
\end{blocLivre}
\begin{blocLivre}{cor-3-15}
\begin{corrige}{15}{Coût marginal et coût supplémentaire}
\demarche{Comparer un accroissement fini à la variation prédite par la dérivée.}
\begin{questions}
\item En dérivant terme à terme, $C'(q)=0{,}06q^2-1{,}2q+10$.
Ainsi $C'(20)=24-24+10=\res{10}$ euros par objet.
\item $C(20)=220$ euros et $C(21)=230{,}62$ euros. L’accroissement exact vaut
$\res{C(21)-C(20)=10{,}62\text{ euros}}$, soit 0,62 euro de plus que la prévision marginale.
\item L’approximation affine est $C(20+h)-C(20)\approx hC'(20)$.
Pour $h=0{,}1$, elle donne $\res{1\text{ euro}}$.
\item On développe $(20+h)^2=400+40h+h^2$ et $(20+h)^3=8000+1200h+60h^2+h^3$.
Après soustraction des termes constants,
\begin{align*}
C(20+h)-C(20)&=0{,}02(1200h+60h^2+h^3)-0{,}6(40h+h^2)+10h\\
&=\res{10h+0{,}6h^2+0{,}02h^3}.
\end{align*}
\item Pour $h=0{,}1$, l’accroissement exact vaut $1{,}00602$ euro.
L’erreur absolue de la prévision 1 euro vaut donc $\res{0{,}00602\text{ euro}}$.
La dérivée mesure une variation instantanée ; sur un intervalle fini, les termes en $h^2$ et $h^3$ interviennent aussi.
\end{questions}

\end{corrige}
\end{blocLivre}

\clearpage
% !TEX root = ../main.tex
\corrigesChapitre{4}{Variations et optimisation}
\corrigesSouples

\begin{blocLivre}{cor-4-1}
\begin{corrige}{1}{Du signe aux variations}
\demarche{Traduire les signes de la dérivée en flèches et comparer aussi les valeurs aux bornes.}
\begin{questions}\item La fonction croît sur $[-3;1]$, puis décroît sur $[1;4]$.
Son maximum est atteint en 1 et vaut $\res{7}$.
\item On compare les deux bornes : $-1<2$, donc le minimum global vaut
$\res{-1}$, atteint en 4.
\end{questions}
\label{figlivre-r04-01-var-cor}
\figureLivre[corrige]{r04-01-var}

\end{corrige}
\end{blocLivre}
\begin{blocLivre}{cor-4-2}
\begin{corrige}{2}{Un trinôme sur un segment}
\demarche{Chercher le zéro de la dérivée à l’intérieur du segment.}
$f'(x)=2x-4$, négative pour $x<2$, nulle en 2, positive pour $x>2$.
Donc $f$ décroît sur $[-1;2]$, de $f(-1)=6$ à $f(2)=-3$, puis croît
sur $[2;5]$ jusqu'à $f(5)=6$.
Le minimum est $\res{-3}$ en 2 ; le maximum est $\res{6}$, aux deux bornes $-1$ et 5.
\label{figlivre-r04-02-var-cor}
\figureLivre[corrige]{r04-02-var}

\end{corrige}
\end{blocLivre}
\begin{blocLivre}{cor-4-3}
\begin{corrige}{3}{Une dérivée ne change pas de signe}
\demarche{Une dérivée strictement positive impose une croissance stricte.}
$g'(x)=3x^2+3>0$ pour tout réel $x$, donc $g$ est strictement croissante sur $\R$.
Sur $[-1;2]$, le minimum vaut $g(-1)=\res{-4}$ et le maximum $g(2)=\res{14}$.
\label{figlivre-r04-03-var-cor}
\figureLivre[corrige]{r04-03-var}

\end{corrige}
\end{blocLivre}
\begin{blocLivre}{cor-4-4}
\begin{corrige}{4}{Un quotient simple}
\demarche{Le dénominateur de la dérivée est positif sur tout le domaine étudié.}
$h'(x)=\dfrac{x+2-(x+1)}{(x+2)^2}=\dfrac1{(x+2)^2}>0$ sur $[0;4]$.
La fonction croît de $\res{h(0)=1/2}$ à $\res{h(4)=5/6}$ : ce sont son
minimum et son maximum.
\label{figlivre-r04-04-var-cor}
\figureLivre[corrige]{r04-04-var}

\end{corrige}
\end{blocLivre}
\begin{blocLivre}{cor-4-5}
\begin{corrige}{5}{Dérivée nulle et extremum}
\demarche{Un zéro isolé de la dérivée n’entraîne un extremum que si le sens de variation change.}
$f'(x)=3x^2\geq0$ et $f'(0)=0$. La dérivée est positive sauf en 0 :
$f$ est strictement croissante sur $\R$. Elle ne change pas de sens de variation
en 0 : $\res{\text{aucun extremum en }0}$.
On peut aussi comparer $f(-\varepsilon)<f(0)<f(\varepsilon)$ pour tout $\varepsilon>0$.
\label{figlivre-r04-05-var-cor}
\figureLivre[corrige]{r04-05-var}

\end{corrige}
\end{blocLivre}
\begin{blocLivre}{cor-4-6}
\begin{corrige}{6}{Un polynôme du troisième degré}
\demarche{Factoriser la dérivée et distinguer les valeurs locales des valeurs globales.}
\begin{questions}
\item $f'(x)=3x^2-6x-9=3(x^2-2x-3)=\res{3(x+1)(x-3)}$.
\item Le produit est positif à l'extérieur des racines et négatif entre elles.
\end{questions}
\begin{center}\begin{tableau}
\otTab[2.5]{$x$/0.8,$f'(x)$/0.8,$f$/1.8}{$-3$,$-1$,$3$,$5$}
\tkzTabLine{,+,z,-,z,+,}
\tkzTabVar{-/$-22$,+/$10$,-/$-22$,+/$10$}
\end{tableau}\end{center}
\otLabel{3.}\enspace À l'intérieur du segment, maximum local 10 en $-1$,
minimum local $-22$ en 3. Le maximum global 10 est atteint en $-1$ \emph{et} 5 ;
le minimum global $-22$ est atteint en $-3$ \emph{et} 3.
\end{corrige}
\end{blocLivre}
\begin{blocLivre}{cor-4-7}
\begin{corrige}{7}{Un minimum avec un quotient}
\demarche{Sur les réels positifs, isoler le seul facteur qui peut changer de signe.}
\begin{questions}
\item $f'(x)=1-4/x^2=\dfrac{(x-2)(x+2)}{x^2}$.
Sur $]0;+\infty[$, son signe est celui de $x-2$.
\item La fonction décroît jusqu'à 2, puis croît. Son minimum est
$\res{f(2)=4}$ ; donc $x+4/x\geq4$ pour tout $x>0$.
\item $f(x)-4=(x^2-4x+4)/x=\res{(x-2)^2/x}\geq0$ car $x>0$.
L'égalité a lieu exactement pour $x=2$.
\end{questions}
\label{figlivre-r04-07-var-cor}
\figureLivre[corrige]{r04-07-var}

\end{corrige}
\end{blocLivre}
\begin{blocLivre}{cor-4-8}
\begin{corrige}{8}{Étudier une fonction paire}
\demarche{Étudier une moitié du domaine puis utiliser la symétrie.}
\begin{questions}
\item $f(-x)=(-x)^4-4(-x)^2=f(x)$ : $f$ est paire.
\item $f'(x)=4x(x^2-2)$. Sur $[0;3]$, elle est négative entre 0 et $\sqrt2$,
puis positive après $\sqrt2$ : $f$ décroît de 0 à $-4$, puis croît jusqu'à 45.
\item Sur $[-3;0]$, par symétrie : décroissance de 45 à $-4$ entre $-3$ et $-\sqrt2$,
puis croissance de $-4$ à 0. Minimum global $\res{-4}$ en $\pm\sqrt2$ ; maximum
global $\res{45}$ en $\pm3$.
\item $g(-x)=-x^3+3x=-g(x)$ : $g$ est impaire et sa courbe est symétrique par rapport à l'origine.
\end{questions}
\label{figlivre-r04-08-var-cor}
\figureLivre[corrige]{r04-08-var}

\end{corrige}
\end{blocLivre}
\begin{blocLivre}{cor-4-9}
\begin{corrige}{9}{Position par rapport à une tangente}
\demarche{La différence entre la fonction et sa tangente donne leur position relative.}
\begin{questions}
\item $f(1)=1$, $f'(1)=3$ : $\res{T:y=3x-2}$.
\item $(x-1)^2(x+2)=(x^2-2x+1)(x+2)=x^3-3x+2$.
\item Sur $[0;3]$, $(x-1)^2\geq0$ et $x+2>0$, donc
$f(x)-(3x-2)\geq0$. La courbe est au-dessus de $T$ et la rencontre uniquement
au point $\res{(1;1)}$ sur cet intervalle.
\end{questions}
\label{figlivre-r04-09-position-cor}
\figureLivre[corrige]{r04-09-position}

\end{corrige}
\end{blocLivre}
\begin{blocLivre}{cor-4-10}
\begin{corrige}{10}{Une inégalité grâce aux variations}
\demarche{Mettre la dérivée au même dénominateur pour lire son signe.}
\begin{questions}
\item $f'(x)=1/(2\sqrt x)-1/2=\dfrac{1-\sqrt x}{2\sqrt x}$.
Pour $x>1$, le dénominateur est positif et le numérateur négatif.
Donc $f$ est strictement décroissante sur $[1;9]$.
\item $f(1)=1/2$ et $f(9)=3-9/2=-3/2$. Maximum $1/2$ en 1 ; minimum $-3/2$ en 9.
\item Toute valeur de $f$ est comprise entre ces extremums :
$\res{-\frac32\leq\sqrt x-\frac x2\leq\frac12}$.
\end{questions}
\label{figlivre-r04-10-var-cor}
\figureLivre[corrige]{r04-10-var}

\end{corrige}
\end{blocLivre}
\begin{blocLivre}{cor-4-11}
\begin{corrige}{11}{Optimiser une quantité entière}
\demarche{Chercher l’optimum réel, puis comparer les entiers qui l’encadrent.}
\begin{questions}
\item $B'(x)=-4x+27$ s'annule en $27/4=6{,}75$.
Le bénéfice croît jusqu'à 6,75 puis décroît.
\item Les deux entiers voisins sont $\res{6\text{ et }7}$.
Tous les autres entiers sont moins favorables d'après les variations.
\item $B(6)=50$ et $B(7)=51$. La quantité optimale est $\res{7\text{ articles}}$,
pour un bénéfice de $\res{51\text{ euros}}$.
\end{questions}
\label{figlivre-r04-11-var-cor}
\figureLivre[corrige]{r04-11-var}

\end{corrige}
\end{blocLivre}
\begin{blocLivre}{cor-4-12}
\begin{corrige}{12}{Une boîte sans couvercle}
\demarche{Exprimer le volume avec une seule variable et contrôler les dimensions positives.}
\begin{questions}
\item La hauteur est $x>0$ et le fond a pour côté $24-2x>0$ : $0<x<12$.
Le volume vaut $\res{V(x)=x(24-2x)^2}$ cm$^3$.
\item En développant, $V(x)=4x^3-96x^2+576x$.
Donc $V'(x)=12x^2-192x+576=\res{12(x-4)(x-12)}$.
La dérivée est positive sur $]0;4[$ et négative sur $]4;12[$.
Le volume croît jusqu'à 4 puis décroît.
\item Le maximum est atteint pour $\res{x=4\text{ cm}}$.
Le fond mesure $16\times16$ cm et la hauteur 4 cm.
Le volume est $4\times16^2=1024$ cm$^3$, soit $\res{1{,}024\text{ L}}$.
\end{questions}
\label{figlivre-r04-12-var-cor}
\figureLivre[corrige]{r04-12-var}

\label{figlivre-r04-12-boite-cor}
\figureLivre[corrige]{r04-12-boite}

\end{corrige}
\end{blocLivre}
\begin{blocLivre}{cor-4-13}
\begin{corrige}{13}{Un rectangle sous une parabole}
\demarche{Relier les coordonnées des sommets à la largeur et à la hauteur.}
\begin{questions}
\item La largeur vaut $2x$ et la hauteur $9-x^2>0$, donc
$\res{A(x)=2x(9-x^2)=18x-2x^3}$.
\item $A'(x)=18-6x^2=6(\sqrt3-x)(\sqrt3+x)$.
Sur $]0;3[$, le signe est positif avant $\sqrt3$, négatif après.
Le maximum est donc atteint en $\res{x=\sqrt3}$.
\item La largeur est $\res{2\sqrt3}$, la hauteur $\res{6}$ et l'aire
$\res{12\sqrt3\approx20{,}78}$ unités d'aire.
\end{questions}
\begin{center}\begin{tikzpicture}\begin{axis}[otLivreGraphe,width=91mm,height=48mm,xmin=-3.3,xmax=3.3,ymin=-0.5,ymax=10,xtick={-3,0,3},ytick={3,6,9}]
\addplot[otNavy,thick,domain=-3:3,samples=70]{9-x^2};
\draw[otGreenDark,fill=otGreen,fill opacity=.13](axis cs:-1.73205,0)rectangle(axis cs:1.73205,6);
\end{axis}\end{tikzpicture}\end{center}
\label{figlivre-r04-13-var-cor}
\figureLivre[corrige]{r04-13-var}

\end{corrige}
\end{blocLivre}
\begin{blocLivre}{cor-4-14}
\begin{corrige}{14}{Un emballage de périmètre imposé}
\demarche{Utiliser le périmètre pour éliminer la hauteur avant de dériver le volume.}
\begin{questions}
\item $2r+2h=40$, donc $h=20-r$. Les deux dimensions étant strictement positives,
$\res{0<r<20}$.
\item $V(r)=\pi r^2(20-r)$, donc $\res{V'(r)=\pi r(40-3r)}$.
\item Sur le domaine, le signe dépend de $40-3r$ : $V$ croît jusqu'à $40/3$,
puis décroît. Le rayon optimal est $\res{40/3\text{ cm}}$ et la hauteur $\res{20/3\text{ cm}}$.
\item Le volume maximal vaut $\res{32000\pi/27\text{ cm}^3}$, environ 3723,37 cm$^3$.
Aux cas limites, le rayon ou la hauteur est nul : le volume est nul et l'objet dégénéré.
\end{questions}
\label{figlivre-r04-14-var-cor}
\figureLivre[corrige]{r04-14-var}

\label{figlivre-r04-14-cylindre-cor}
\figureLivre[corrige]{r04-14-cylindre}

\end{corrige}
\end{blocLivre}
\begin{blocLivre}{cor-4-15}
\begin{corrige}{15}{Distance minimale à un point}
\demarche{Minimiser d’abord la distance au carré, plus simple à dériver.}
\begin{questions}
\item $AM^2=(x-0)^2+(x^2-2)^2=x^2+x^4-4x^2+4
=\res{x^4-3x^2+4}$.
\item La racine carrée est strictement croissante : les valeurs de $AM=\sqrt{D(x)}$
sont minimales exactement lorsque $D$ l'est.
\item $D'(x)=4x^3-6x=2x(2x^2-3)$.
En posant $a=\sqrt{3/2}$, les signes successifs sont $-$, $+$, $-$, $+$
sur les intervalles séparés par $-a$, 0 et $a$.
$D$ décroît de 8 à $7/4$, croît jusqu'à 4, décroît à $7/4$, puis croît jusqu'à 8.
\item Les points optimaux sont $\res{M_1(-\sqrt{3/2};3/2)}$ et
$\res{M_2(\sqrt{3/2};3/2)}$.
\item La distance minimale vaut $\res{\sqrt7/2\approx1{,}323}$,
strictement inférieure à $AO=2$.
\end{questions}
\label{figlivre-r04-15-var-cor}
\figureLivre[corrige]{r04-15-var}

\label{figlivre-r04-15-distance-cor}
\figureLivre[corrige]{r04-15-distance}

\end{corrige}
\end{blocLivre}

\clearpage
% !TEX root = ../main.tex
\corrigesChapitre{5}{Fonction exponentielle}
\corrigesSouples

\begin{blocLivre}{cor-5-1}
\begin{corrige}{1}{Simplifier une expression}
\demarche{Dans un produit, additionner les exposants ; dans un quotient, les soustraire.}
On applique les règles de calcul avant de simplifier les exposants :
\begin{align*}
e^{2x}e^{3-x}&=e^{2x+3-x}=\res{e^{x+3}},\\
\frac{e^{5x}}{e^{2x-1}}&=e^{5x-(2x-1)}=\res{e^{3x+1}},\\
(e^{-x})^3e^{2x}&=e^{-3x}e^{2x}=\res{e^{-x}}.
\end{align*}
Dans le quotient, les parenthèses conservent le signe du terme $-1$.

\end{corrige}
\end{blocLivre}
\begin{blocLivre}{cor-5-2}
\begin{corrige}{2}{Des équations immédiates}
\demarche{Utiliser l’injectivité de l’exponentielle et sa stricte positivité.}
L'exponentielle est strictement croissante, donc injective :
$e^{2x-1}=e^3\iff2x-1=3\iff\res{x=2}$.
Puis $e^{x+2}=1=e^0\iff\res{x=-2}$.
Enfin $e^x>0$ pour tout $x$, donc $e^x=-2$ n'a $\res{\text{aucune solution réelle}}$.
\end{corrige}
\end{blocLivre}
\begin{blocLivre}{cor-5-3}
\begin{corrige}{3}{Comparer des exponentielles}
\demarche{Comparer les exposants grâce à la croissance stricte.}
Par croissance stricte, $e^{3x+1}\leq e^{x+5}\iff3x+1\leq x+5\iff\res{x\leq2}$.
$e^{-2x}>1=e^0\iff-2x>0\iff\res{x<0}$.
Comme $-1<0<1/2$, on a $\res{e^{-1}<1<e^{1/2}}$.
\end{corrige}
\end{blocLivre}
\begin{blocLivre}{cor-5-4}
\begin{corrige}{4}{Dériver et donner le signe}
\demarche{Le signe de la constante devant l’exponentielle détermine celui de la dérivée.}
$f'(x)=2e^{2x}>0$ : $f$ est strictement croissante.
$g'(x)=-3e^{-x}<0$ : $g$ est strictement décroissante.
$h'(x)=e^{x+1}>0$ : $h$ est strictement croissante.
Ces résultats valent sur tout $\R$.
\end{corrige}
\end{blocLivre}
\begin{blocLivre}{cor-5-5}
\begin{corrige}{5}{La tangente en zéro}
\demarche{Une tangente fournit une approximation locale, avec un signe d’approximation.}
$f(0)=1$ et $f'(0)=1$, donc $\res{T:y=x+1}$.
Pour un petit $h$, $e^h\approx1+h$ : $\res{e^{0{,}03}\approx1{,}03}$.
Il s'agit d'une approximation, pas d'une égalité.
\end{corrige}
\end{blocLivre}
\begin{blocLivre}{cor-5-6}
\begin{corrige}{6}{Factoriser avant de résoudre}
\demarche{Factoriser l’exponentielle, toujours strictement positive.}
\begin{questions}
\item $e^{2x}-e^x=e^x(e^x-1)$. Le premier facteur ne s'annule pas :
$e^x=1\iff\res{x=0}$.
\item $(x-2)e^x=0\iff\res{x=2}$.
\item $e^{-x}>0$, donc le signe est celui de
$x^2-3x+2=(x-1)(x-2)$. L'ensemble des solutions est $\res{[1;2]}$.
\end{questions}
\end{corrige}
\end{blocLivre}
\begin{blocLivre}{cor-5-7}
\begin{corrige}{7}{Étudier un produit}
\demarche{Factoriser la dérivée par l’exponentielle et comparer les valeurs aux bornes.}
\begin{questions}
\item Par produit, $f'(x)=e^{-x}-(x+1)e^{-x}=\res{-xe^{-x}}$.
\item Le signe de $f'$ est celui de $-x$.
\end{questions}
\begin{center}\begin{tableau}
\otTab[3.1]{$x$/0.8,$f'(x)$/0.8,$f$/1.8}{$-2$,$0$,$4$}
\tkzTabLine{,+,z,-,}
\tkzTabVar{-/$-e^2$,+/$1$,-/$5e^{-4}$}
\end{tableau}\end{center}
Le maximum est $\res{1}$ en 0. Le minimum sur le segment vaut $\res{-e^2}$ en $-2$,
car l'autre borne donne une valeur positive.
\par\otLabel{3.}\enspace $f(1)=2/e$ et $f'(1)=-1/e$, donc
$\res{T:y=(3-x)/e}$. Le signe de $f$ est celui de $x+1$ : sur le domaine,
$\res{f(x)\geq0\iff x\in[-1;4]}$.
\end{corrige}
\end{blocLivre}
\begin{blocLivre}{cor-5-8}
\begin{corrige}{8}{Une différence d'exponentielles}
\demarche{Faire apparaître soit un produit, soit un carré positif.}
\begin{questions}
\item $g'(x)=2e^{2x}-2e^x=2e^x(e^x-1)$.
Elle est négative pour $x<0$, nulle en 0, positive pour $x>0$.
\item $g$ décroît puis croît : son minimum vaut $\res{g(0)=-1}$.
\item $(e^x-1)^2\geq0$, donc $g(x)\geq-1$ avec égalité exactement lorsque $e^x=1$,
c'est-à-dire $x=0$.
\end{questions}
\label{figlivre-r05-08-var-cor}
\figureLivre[corrige]{r05-08-var}

\end{corrige}
\end{blocLivre}
\begin{blocLivre}{cor-5-9}
\begin{corrige}{9}{Un quotient toujours défini}
\demarche{La positivité du dénominateur autorise les calculs et la comparaison des fractions.}
\begin{questions}
\item Le dénominateur $1+e^x$ est strictement positif sur $\R$.
Les fonctions en jeu sont dérivables, donc le quotient aussi.
\item $f'(x)=\dfrac{e^x(1+e^x)-e^xe^x}{(1+e^x)^2}
=\res{\dfrac{e^x}{(1+e^x)^2}}>0$ : croissance stricte sur $\R$.
\item Le numérateur est positif et inférieur au dénominateur, donc $0<f<1$.
$f(x)=1/2\iff2e^x=1+e^x\iff e^x=1\iff\res{x=0}$.
\end{questions}
\end{corrige}
\end{blocLivre}
\begin{blocLivre}{cor-5-10}
\begin{corrige}{10}{Passer du continu au discret}
\demarche{Échantillonner le modèle continu aux instants entiers.}
\begin{questions}
\item $N(0)=\res{500}$ et $N(5)=\res{500e\approx1359}$ individus.
L'arrondi concerne la valeur calculée ; le modèle continu peut donner un nombre non entier.
\item $u_n=500(e^{0{,}2})^n$ : suite géométrique de premier terme 500 et
de raison $\res{e^{0{,}2}}$.
\item Le taux annuel vaut $e^{0{,}2}-1\approx0{,}2214028$, soit
$\res{22{,}1\,\%}$ à 0,1 point près. Ce taux n'est pas 20\,\%.
\end{questions}
\end{corrige}
\end{blocLivre}
\begin{blocLivre}{cor-5-11}
\begin{corrige}{11}{Construire une approximation d'Euler}
\demarche{Chaque étape d’Euler utilise la pente au début de cette étape.}
\begin{questions}
\item $y_{k+1}=(1+1/N)y_k$, donc la suite est géométrique de raison $1+1/N$.
Après $N$ étapes, $\res{y_N=(1+1/N)^N}$.
\item Pour $N=2$ : $y_2=(3/2)^2=2{,}25$ ; pour $N=10$ :
$y_{10}=1{,}1^{10}\approx2{,}593742$.
Les deux valeurs sont inférieures à $e\approx2{,}718282$ ; celle avec dix étapes est plus proche.
Les termes de la suite $\big((1+1/n)^n\big)$ fournissent ainsi des valeurs approchées de $e$.
\item Une boucle reproduit directement l'actualisation d'Euler :
\end{questions}
\begin{python}
def euler(N):
    h = 1/N
    y = 1.0
    for k in range(N):
        y = y + h*y
    return y
\end{python}

Les $N$ étapes de longueur $1/N$ couvrent l’intervalle de 0 à 1. Dans la mise à jour, la pente utilisée est la valeur courante de $y$, puisque la fonction étudiée vérifie $y^\prime=y$.

\label{figlivre-r05-11-euler-cor}
\figureLivre[corrige]{r05-11-euler}

\end{corrige}
\end{blocLivre}
\begin{blocLivre}{cor-5-12}
\begin{corrige}{12}{Le refroidissement d'une boisson}
\demarche{Justifier la décroissance avant de chercher la première minute admissible.}
\begin{questions}
\item $T(0)=90\,{}^\circ$C.
$T'(t)=-7e^{-0{,}1t}<0$ : la température est strictement décroissante.
\item $70e^{-0{,}1t}>0$, donc $T(t)>20$ pour tout instant fini $t\geq0$.
\item $T(12)\approx41{,}08\,{}^\circ$C et $T(13)\approx39{,}08\,{}^\circ$C.
La première minute entière où $T\leq40$ est donc $\res{13}$.
\item Le programme suivant teste les minutes entières dans l'ordre :
\end{questions}
\begin{python}
from math import exp
t = 0
while 20 + 70*exp(-0.1*t) > 40:
    t = t + 1
\end{python}
À la sortie de la boucle, \code{t} vaut 13.
\end{corrige}
\end{blocLivre}
\begin{blocLivre}{cor-5-13}
\begin{corrige}{13}{Une concentration après injection}
\demarche{La dérivée permet de localiser et de prouver le pic de concentration.}
\begin{questions}
\item $C'(t)=8e^{-0{,}5t}-4te^{-0{,}5t}=\res{4(2-t)e^{-0{,}5t}}$.
La concentration croît sur $[0;2]$, puis décroît sur $[2;8]$.
Les valeurs remarquables sont $C(0)=0$, $C(2)=16/e$, $C(8)=64e^{-4}$.
\item Maximum à $\res{t=2\text{ h}}$, de valeur $\res{16/e\approx5{,}886\text{ mg/L}}$.
\item $C(0)=0$ et $C'(0)=8$ : la tangente est $\res{y=8t}$.
La pente est une vitesse initiale d'augmentation de 8 mg/L par heure.
\item Le taux moyen vaut $(C(3)-C(2))/(3-2)=24e^{-1{,}5}-16/e$
$\approx\res{-0{,}53\text{ mg/L par heure}}$. Son signe négatif traduit une diminution.
\end{questions}
\label{figlivre-r05-13-var-cor}
\figureLivre[corrige]{r05-13-var}

\label{figlivre-r05-13-concentration-cor}
\figureLivre[corrige]{r05-13-concentration}

\end{corrige}
\end{blocLivre}
\begin{blocLivre}{cor-5-14}
\begin{corrige}{14}{Prouver une inégalité sans convexité}
\demarche{Étudier la différence avec la droite pour prouver une inégalité sur tout le domaine.}
\begin{questions}
\item $F'(x)=e^x-1$, négative pour $x<0$, nulle en 0 et positive pour $x>0$.
\item $F$ décroît jusqu'à 0 puis croît : son minimum est $F(0)=0$.
Donc $F(x)\geq0$, soit $\res{e^x\geq1+x}$.
La stricte monotonie de part et d'autre assure que l'égalité a lieu seulement pour $x=0$.
\item Pour $x=0{,}1$, on obtient $e^{0{,}1}\geq1{,}1$ ; pour $x=-0{,}1$,
on obtient $e^{-0{,}1}\geq0{,}9$.
\item Comme les nombres sont positifs, $1/e^{0{,}1}\geq0{,}9$ entraîne
$e^{0{,}1}\leq1/0{,}9=10/9$. Ainsi $\res{1{,}1\leq e^{0{,}1}\leq10/9}$.
\item L'approximation affine en 0 vaut 1,1. Elle est ici une borne inférieure.
L'encadrement garantit une erreur d'au plus $10/9-1{,}1=1/90\approx0{,}0111$.
\end{questions}
\label{figlivre-r05-14-var-cor}
\figureLivre[corrige]{r05-14-var}

\label{figlivre-r05-14-inegalite-cor}
\figureLivre[corrige]{r05-14-inegalite}

\end{corrige}
\end{blocLivre}
\begin{blocLivre}{cor-5-15}
\begin{corrige}{15}{Choisir une durée de campagne}
\demarche{Étudier les variations, puis convertir les centaines d’euros en euros.}
\begin{questions}
\item Par produit, $G'(t)=100e^{-0{,}5t}-50(t+1)e^{-0{,}5t}
=\res{50(1-t)e^{-0{,}5t}}$.
\item La dérivée a le signe de $1-t$ : croissance sur $[0;1]$, décroissance sur $[1;6]$.
Le maximum est atteint après $\res{1\text{ semaine}}$ et vaut
$200e^{-1/2}$ centaines d'euros, soit $\res{20000/\sqrt e\approx12130{,}61\text{ euros}}$.
\item $G(0)=100$ et $G(1)=200/\sqrt e>100$ car $e<4$.
\item $G(2)=300/e\approx110{,}36$ et $G(3)=400e^{-1{,}5}\approx89{,}25$.
Sur $]0;1]$, la croissance stricte depuis $G(0)=100$ donne $G>100$.
Sur $[1;2]$, la décroissance donne $G(t)\geq G(2)=300/e>100$ car $e<3$.
\item Les durées entières favorables sont $\res{1\text{ et }2\text{ semaines}}$.
À partir de 3, le gain est déjà inférieur à 100 et continue de décroître jusqu'à 6.
\end{questions}
\label{figlivre-r05-15-var-cor}
\figureLivre[corrige]{r05-15-var}

\label{figlivre-r05-15-campagne-cor}
\figureLivre[corrige]{r05-15-campagne}

\end{corrige}
\end{blocLivre}

\clearpage
% !TEX root = ../main.tex
\corrigesChapitre{6}{Trigonométrie}
\corrigesSouples


\begin{blocLivre}{cor-6-1}
\begin{corrige}{1}{Convertir des angles}
\demarche{Un demi-tour vaut 180 degrés, soit pi radians.}
\begin{questions}
\item Pour convertir des degrés en radians, on multiplie par $\pi/180$ :
$30\pi/180=\res{\pi/6}$ ; $45\pi/180=\res{\pi/4}$ ;
$120\pi/180=\res{2\pi/3}$ ; $270\pi/180=\res{3\pi/2}$.
\item Dans l’autre sens, on multiplie par $180/\pi$ :
$(\pi/3)(180/\pi)=\res{60^\circ}$ ;
$(5\pi/6)(180/\pi)=\res{150^\circ}$ ;
$(7\pi/4)(180/\pi)=\res{315^\circ}$.
\end{questions}

\end{corrige}
\end{blocLivre}
\begin{blocLivre}{cor-6-2}
\begin{corrige}{2}{Même point-image}
\demarche{Deux mesures donnent le même point lorsque leur différence est un nombre entier de tours.}
\begin{questions}
\item La différence vaut $4\pi=2\times2\pi$ : oui.
\item $3\pi/2-(-\pi/2)=2\pi$ : oui. En revanche $5\pi/4-\pi/4=\pi$ : les points sont opposés, donc distincts.
\end{questions}
\label{figlivre-r06-02-images-cor}
\figureLivre[corrige]{r06-02-images}

\end{corrige}
\end{blocLivre}
\begin{blocLivre}{cor-6-3}
\begin{corrige}{3}{Lire les valeurs remarquables}
\demarche{Sur le cercle unité, l’abscisse est le cosinus et l’ordonnée est le sinus.}
\begin{questions}
\item Les coordonnées sont $\res{(1;0)}$, $\res{(\sqrt3/2;1/2)}$, $\res{(\sqrt2/2;\sqrt2/2)}$ et $\res{(0;1)}$.
\item On lit les points sur les axes : $\res{-1;\ 0;\ 0;\ -1}$ respectivement.
\end{questions}
\label{figlivre-r06-03-valeurs-cor}
\figureLivre[corrige]{r06-03-valeurs}

\end{corrige}
\end{blocLivre}
\begin{blocLivre}{cor-6-4}
\begin{corrige}{4}{Utiliser une symétrie}
\demarche{Les symétries conservent une coordonnée et changent le signe de l’autre.}
\begin{questions}
\item La réflexion par rapport à l’axe horizontal conserve l’abscisse et oppose l’ordonnée.
Ainsi $\cos(-\pi/3)=\cos(\pi/3)=\res{1/2}$ et
$\sin(-\pi/3)=-\sin(\pi/3)=\res{-\sqrt3/2}$.
\item Comme $2\pi/3=\pi-\pi/3$, la réflexion par rapport à l’axe vertical oppose l’abscisse et conserve l’ordonnée.
Donc $\cos(2\pi/3)=\res{-1/2}$ et $\sin(2\pi/3)=\res{\sqrt3/2}$.
Le point du deuxième quadrant a bien un cosinus négatif et un sinus positif.
\end{questions}

\label{figlivre-r06-04-symetries-cor}
\figureLivre[corrige]{r06-04-symetries}

\end{corrige}
\end{blocLivre}
\begin{blocLivre}{cor-6-5}
\begin{corrige}{5}{Une longueur d'arc}
\demarche{Utiliser les radians dans la formule longueur égale rayon multiplié par angle.}
La formule $L=R\theta$ donne $\res{2\pi\text{ cm}}$ et $\res{9\pi/2\text{ cm}}$. Pour un tour entier, $L=6\times2\pi=\res{12\pi\text{ cm}}$.
\label{figlivre-r06-05-arcs-cor}
\figureLivre[corrige]{r06-05-arcs}

\end{corrige}
\end{blocLivre}
\begin{blocLivre}{cor-6-6}
\begin{corrige}{6}{Réduire des grandes mesures}
\demarche{Ajouter ou retirer des tours complets avant de lire une valeur remarquable.}
\begin{questions}
\item On choisit les nombres entiers de tours qui ramènent les mesures dans $\mathopen{]}-\pi;\pi]$ :
\[\frac{17\pi}6-2\pi=\frac{5\pi}6,\qquad
-\frac{11\pi}4+2\pi=-\frac{3\pi}4,\qquad
\frac{29\pi}3-10\pi=-\frac\pi3.\]
\item Les points-images sont conservés. Les symétries et les valeurs remarquables donnent, dans le même ordre,
\[\res{\left(-\frac{\sqrt3}2;\frac12\right)},\quad
\res{\left(-\frac{\sqrt2}2;-\frac{\sqrt2}2\right)},\quad
\res{\left(\frac12;-\frac{\sqrt3}2\right)}.\]
Chaque couple est écrit dans l’ordre $(\cos;\sin)$ ; les signes se contrôlent avec le quadrant.
\end{questions}

\end{corrige}
\end{blocLivre}
\begin{blocLivre}{cor-6-7}
\begin{corrige}{7}{Retrouver une coordonnée}
\demarche{Calculer le carré du sinus, puis choisir son signe avec le quadrant.}
\begin{questions}
\item $\sin^2t=1-9/25=16/25$. Dans le deuxième quadrant, le sinus est positif : $\res{\sin t=4/5}$.
\item Pour $-t$ : $\res{(-3/5;-4/5)}$ ; pour $t+\pi$ : $\res{(3/5;-4/5)}$.
\end{questions}
\label{figlivre-r06-07-coordonnee-cor}
\figureLivre[corrige]{r06-07-coordonnee}

\end{corrige}
\end{blocLivre}
\begin{blocLivre}{cor-6-8}
\begin{corrige}{8}{Deux angles associés}
\demarche{Repérer le quadrant pour contrôler les signes des deux coordonnées.}
\begin{questions}
\item Les deux angles sont dans le troisième quadrant. Les valeurs sont $-\sqrt2/2$, $-\sqrt2/2$, $-\sqrt3/2$ et $-1/2$.
\item $A=-\sqrt2-\sqrt2/2=\res{-3\sqrt2/2}$.
\end{questions}
\end{corrige}
\end{blocLivre}
\begin{blocLivre}{cor-6-9}
\begin{corrige}{9}{Retrouver des angles sur le cercle}
\demarche{Une condition sur le cosinus fixe une abscisse ; une condition sur le sinus fixe une ordonnée.}
\begin{questions}
\item Les abscisses égales à $1/2$ correspondent à $\res{t=\pi/3\text{ ou }5\pi/3}$.
\item Les ordonnées égales à $-\sqrt2/2$ donnent $\res{t=5\pi/4\text{ ou }7\pi/4}$.
\item Les points sur l’axe vertical donnent $\res{t=\pi/2\text{ ou }3\pi/2}$.
\end{questions}

\label{figgeo-06-09-cor}
\figureGeo[corrige]{06-09}
\end{corrige}
\end{blocLivre}
\begin{blocLivre}{cor-6-10}
\begin{corrige}{10}{Une portion de cercle}
\demarche{Calculer la fraction d’un tour avant l’aire et le périmètre du secteur.}
\begin{questions}
\item La fraction est $(2\pi/5)/(2\pi)=\res{1/5}$.
\item L’aire vaut $100\pi/5=\res{20\pi\text{ cm}^2}$. L’arc mesure $10\times2\pi/5=4\pi$ cm ; le périmètre est $\res{20+4\pi\text{ cm}}$.
\end{questions}

\label{figgeo-06-10-cor}
\figureGeo[corrige]{06-10}
\end{corrige}
\end{blocLivre}
\begin{blocLivre}{cor-6-11}
\begin{corrige}{11}{Une identité à démontrer}
\demarche{Développer les deux carrés et regrouper les termes opposés.}
\begin{questions}
\item Le premier carré vaut $\cos^2t+2\cos t\sin t+\sin^2t$ ; le second vaut
$\cos^2t-2\cos t\sin t+\sin^2t$. En les additionnant, les deux termes croisés se compensent.
\item La somme devient $2\cos^2t+2\sin^2t=2(\cos^2t+\sin^2t)=\res{2}$,
car le point du cercle unité vérifie $\cos^2t+\sin^2t=1$.
\item Pour $t=\pi/4$, le sinus et le cosinus valent $\sqrt2/2$.
Le premier carré vaut $(\sqrt2)^2=2$ et le second vaut $0^2=0$, ce qui confirme le résultat.
\end{questions}

\end{corrige}
\end{blocLivre}
\begin{blocLivre}{cor-6-12}
\begin{corrige}{12}{Une grande roue}
\demarche{Ajouter l’altitude du centre à la composante verticale du rayon.}
\begin{questions}
\item On obtient $\res{(12;14)}$, $\res{(0;26)}$, $\res{(-12;14)}$ et $\res{(0;2)}$.
\item Les hauteurs sont $14+12/2=\res{20\text{ m}}$ et $14-12/2=\res{8\text{ m}}$.
\item La durée est $48\times(5\pi/6)/(2\pi)=\res{20\text{ s}}$, sous l’hypothèse d’une rotation uniforme.
\end{questions}

\label{figgeo-06-12-cor}
\figureGeo[corrige]{06-12}
\end{corrige}
\end{blocLivre}
\begin{blocLivre}{cor-6-13}
\begin{corrige}{13}{Un hexagone régulier}
\demarche{Découper l’hexagone en six triangles centraux équilatéraux.}
\begin{questions}
\item Les angles sont $k\pi/3$ pour $k=0,\ldots,5$. Les coordonnées sont $(1;0)$, $(1/2;\sqrt3/2)$, $(-1/2;\sqrt3/2)$, $(-1;0)$, $(-1/2;-\sqrt3/2)$, $(1/2;-\sqrt3/2)$.
\item Le carré de la distance vaut $(1/2)^2+(\sqrt3/2)^2=1$ : chaque côté mesure $\res{1}$.
\item Le périmètre vaut $\res{6}$. Les six triangles équilatéraux ont chacun pour aire $\sqrt3/4$ : aire totale $\res{3\sqrt3/2}$.
\end{questions}

\label{figgeo-06-13-cor}
\figureGeo[corrige]{06-13}

Pour contrôler l’aire, chaque triangle équilatéral de côté 1 a pour hauteur $\sqrt{1-(1/2)^2}=\sqrt3/2$. Son aire est donc $\sqrt3/4$, à multiplier par six.

\end{corrige}
\end{blocLivre}
\begin{blocLivre}{cor-6-14}
\begin{corrige}{14}{Une zone autorisée}
\demarche{Traduire les deux contraintes par des zones du cercle, puis prendre leur intersection.}
\begin{questions}
\item Le point est dans le demi-plan supérieur : $t\in[0;\pi]$.
\item Le point est à gauche de la droite $x=-1/2$ : $t\in[2\pi/3;4\pi/3]$.
\item L’intersection est $\res{[2\pi/3;\pi]}$. L’angle parcouru vaut $\pi/3$ ; sur le cercle unité, la longueur de cet arc est donc $\res{\pi/3}$.
\end{questions}

\label{figgeo-06-14-cor}
\figureGeo[corrige]{06-14}
\end{corrige}
\end{blocLivre}
\begin{blocLivre}{cor-6-15}
\begin{corrige}{15}{Approcher un cercle par des polygones}
\demarche{Partager un triangle central en deux triangles rectangles.}
\begin{questions}
\item Le rayon du cercle unité vaut 1. Dans le demi-triangle rectangle, le demi-côté vaut
$\sin(\pi/n)$ : un côté entier mesure $2\sin(\pi/n)$. Pour $n$ côtés,
$\res{P_n=2n\sin(\pi/n)}$.
\item $P_6=12\sin(\pi/6)=\res{6}$.
Avec la valeur fournie, $P_{12}=24(\sqrt6-\sqrt2)/4=\res{6(\sqrt6-\sqrt2)}\approx6{,}21166$.
\item Le cercle a pour longueur $2\pi\approx6{,}28319$.
Les deux périmètres calculés sont plus petits ; celui à douze côtés est plus proche.
Ces deux comparaisons numériques ne démontrent pas une convergence pour tous les $n$.
\item La relation admise double le nombre de côtés sans utiliser $\pi$ :
\end{questions}
\begin{python}
from math import sqrt
def archimede(k):
    n, c = 6, 1.0
    for i in range(k):
        c = sqrt(2 - sqrt(4 - c*c))
        n = 2*n
    return n*c/2
\end{python}
Pour $k=1$, $c_{12}=\sqrt{2-\sqrt3}=(\sqrt6-\sqrt2)/2$ : on retrouve $P_{12}$.
Pour $k=4$, le polygone a 96 côtés et la fonction renvoie environ $\res{3{,}14103}$.
Chaque côté est plus court que l'arc qu'il sous-tend, donc $P_n<2\pi$ : ce demi-périmètre
est une valeur approchée par défaut de $\pi$ ;
Archimède en déduisait $\pi>3+10/71$. La relation admise se démontre avec Al-Kashi :
$c_{2n}^2=2-2\cos(\pi/n)$ et $\cos(\pi/n)=\sqrt{1-c_n^2/4}$.
\label{figgeo-06-15-cor}
\figureGeo[corrige]{06-15}
\end{corrige}
\end{blocLivre}

\clearpage
% !TEX root = ../main.tex
\corrigesChapitre{7}{Vecteurs et produit scalaire}
\corrigesSouples


\begin{blocLivre}{cor-7-1}
\begin{corrige}{1}{Coordonnées et norme}
\demarche{Soustraire les coordonnées dans l’ordre indiqué par le nom du vecteur.}
\begin{questions}
\item $\res{\vect{AB}=(3;4)}$ et $\res{\vect{BA}=(-3;-4)}$.
\item $AB=\sqrt{3^2+4^2}=\res{5}$. Le milieu est $\res{(5/2;4)}$.
\end{questions}

\label{figgeo-07-01-cor}
\figureGeo[corrige]{07-01}
\end{corrige}
\end{blocLivre}
\begin{blocLivre}{cor-7-2}
\begin{corrige}{2}{Un calcul avec les coordonnées}
\demarche{Le produit scalaire utilise les coordonnées ; la norme utilise la somme de leurs carrés.}
\begin{questions}
\item $u\cdot v=6-4=\res{2}$ ; $u\cdot u=5$ ; $v\cdot v=25$.
\item Les normes sont $\res{\sqrt5}$ et $\res{5}$.
\end{questions}

\label{figgeo-07-02-cor}
\figureGeo[corrige]{07-02}
\end{corrige}
\end{blocLivre}
\begin{blocLivre}{cor-7-3}
\begin{corrige}{3}{Longueurs et angle}
\demarche{Le cosinus de l’angle détermine le signe du produit scalaire.}
On utilise $u\cdot v=30\cos\theta$. Les valeurs sont $\res{15}$, $\res{0}$ et $\res{-15}$. Le signe indique si l’angle est aigu, droit ou obtus.
\label{figgeo-07-03-cor}
\figureGeo[corrige]{07-03}
\end{corrige}
\end{blocLivre}
\begin{blocLivre}{cor-7-4}
\begin{corrige}{4}{Orthogonaux ou colinéaires}
\demarche{Tester séparément orthogonalité et colinéarité, qui correspondent à deux critères différents.}
\begin{questions}
\item $u\cdot v=6-6=0$, donc $u\perp v$. $u\cdot w=-18-8=-26\neq0$, donc $u$ et $w$ ne sont pas orthogonaux.
\item $\res{w=-2u}$ : $u$ et $w$ sont colinéaires de sens opposés.
\end{questions}

\label{figgeo-07-04-cor}
\figureGeo[corrige]{07-04}
\end{corrige}
\end{blocLivre}
\begin{blocLivre}{cor-7-5}
\begin{corrige}{5}{Lire une projection}
\demarche{Identifier le projeté sur la droite horizontale puis vérifier le calcul par coordonnées.}
\begin{questions}
\item L’axe $(AB)$ est horizontal : $\res{H(2;0)}$.
\item Par projection, $AB\times AH=6\times2=\res{12}$. Par coordonnées, $(6;0)\cdot(2;3)=12$.
\end{questions}

\label{figgeo-07-05-cor}
\figureGeo[corrige]{07-05}
\end{corrige}
\end{blocLivre}
\begin{blocLivre}{cor-7-6}
\begin{corrige}{6}{Les normes d'une somme}
\demarche{Développer les carrés des normes avec le produit scalaire.}
\begin{questions}
\item Par bilinéarité,
$\|u+v\|^2=\|u\|^2+2u\cdot v+\|v\|^2=9+12+16=\res{37}$.
Avec une différence, le terme croisé change de signe :
$\|u-v\|^2=9-12+16=\res{13}$.
\item Une norme est positive ou nulle : $\res{\|u+v\|=\sqrt{37}}$ et $\res{\|u-v\|=\sqrt{13}}$.
Les carrés des diagonales ont pour somme $37+13=50$, égale à $2\|u\|^2+2\|v\|^2=2\times9+2\times16$.
\end{questions}
\label{figgeo-07-06-cor}
\figureGeo[corrige]{07-06}
\end{corrige}
\end{blocLivre}
\begin{blocLivre}{cor-7-7}
\begin{corrige}{7}{Calculer un angle de triangle}
\demarche{Isoler le cosinus dans la définition du produit scalaire.}
\begin{questions}
\item Le produit vaut $4$, $AB=4$ et $AC=2$.
\item $\cos\widehat{BAC}=4/(4\times2)=1/2$, donc $\res{\widehat{BAC}=\pi/3}$.
\item $BC=\sqrt{(-3)^2+(\sqrt3)^2}=2\sqrt3$. Comme $AC^2+BC^2=4+12=16=AB^2$, le triangle est rectangle en $C$.
\end{questions}

\label{figgeo-07-07-cor}
\figureGeo[corrige]{07-07}
\end{corrige}
\end{blocLivre}
\begin{blocLivre}{cor-7-8}
\begin{corrige}{8}{Déterminer un paramètre}
\demarche{Traduire chaque propriété par une équation sur le paramètre.}
\begin{questions}
\item L’orthogonalité équivaut à $u_m\cdot v=0$.
Or $u_m\cdot v=3m-2$, donc $3m-2=0\iff\res{m=2/3}$.
\item La colinéarité équivaut à un déterminant nul :
$m(-1)-2\times3=-m-6=0\iff\res{m=-6}$.
Dans ce cas $u_{-6}=(-6;2)=-2v$, ce qui contrôle le résultat.
\item Pour $m=0$, le produit vaut $-2$, et les normes valent 2 et $\sqrt{10}$.
La définition donne $\cos\theta=(-2)/(2\sqrt{10})=\res{-1/\sqrt{10}}$.
Ce cosinus négatif correspond bien à un angle obtus.
\end{questions}
\label{figgeo-07-08-cor}
\figureGeo[corrige]{07-08}
\end{corrige}
\end{blocLivre}
\begin{blocLivre}{cor-7-9}
\begin{corrige}{9}{Les diagonales d'un rectangle}
\demarche{Décomposer les diagonales à l’aide des deux côtés du rectangle.}
\begin{questions}
\item $\vect{AC}=u+v$ et $\vect{BD}=v-u$.
\item Comme $u\cdot v=0$, $(u+v)\cdot(v-u)=\|v\|^2-\|u\|^2=\res{-20}$. Puis $u\cdot(u+v)=\res{36}$.
\item Le produit des diagonales n’est pas nul : elles ne sont pas perpendiculaires.
\end{questions}

\label{figgeo-07-09-cor}
\figureGeo[corrige]{07-09}
\end{corrige}
\end{blocLivre}
\begin{blocLivre}{cor-7-10}
\begin{corrige}{10}{Un parallélogramme non rectangle}
\demarche{Calculer d’abord le produit des côtés, puis les carrés des diagonales.}
\begin{questions}
\item Le produit vaut $3\times4\times1/2=\res{6}$.
\item $AC^2=9+16+12=37$ et $BD^2=9+16-12=13$, donc $\res{AC=\sqrt{37},\ BD=\sqrt{13}}$.
\item $37+13=50=2\times9+2\times16$ : l’identité du parallélogramme est vérifiée.
\end{questions}

\label{figgeo-07-10-cor}
\figureGeo[corrige]{07-10}
\end{corrige}
\end{blocLivre}
\begin{blocLivre}{cor-7-11}
\begin{corrige}{11}{Projeter sans coordonnées}
\demarche{Relier le produit scalaire à la projection avant d’utiliser Pythagore.}
\begin{questions}
\item $8\times5\times1/2=\res{20}$.
\item $8\times AH=20$, donc $AH=5/2$. Dans $AHC$, $CH=\sqrt{25-25/4}=\res{5\sqrt3/2}$.
\item $BC^2=64+25-40=49$, donc $\res{BC=7}$.
\end{questions}

\label{figgeo-07-11-cor}
\figureGeo[corrige]{07-11}
\end{corrige}
\end{blocLivre}
\begin{blocLivre}{cor-7-12}
\begin{corrige}{12}{Un triangle et sa hauteur}
\demarche{Vérifier les angles droits par les produits scalaires, puis calculer les longueurs.}
\begin{questions}
\item $\vect{CA}=(-1;-3)$ et $\vect{CB}=(3;-1)$ ont un produit $-3+3=0$.
\item $\vect{AB}=(4;2)$, $\vect{AC}=(1;3)$ : produit $10$, normes $\sqrt{20}$ et $\sqrt{10}$. Le cosinus vaut $1/\sqrt2$, donc l’angle est $\pi/4$.
\item $H=(3;2)$ et $\vect{CH}=(1;-2)$, donc $\vect{CH}\cdot\vect{AB}=4-4=0$.
\item $CH=\sqrt5$. L’aire vaut $AB\times CH/2=\sqrt{20}\sqrt5/2=\res{5}$.
\end{questions}

\label{figgeo-07-12-cor}
\figureGeo[corrige]{07-12}
\end{corrige}
\end{blocLivre}
\begin{blocLivre}{cor-7-13}
\begin{corrige}{13}{Le cercle de diamètre donné}
\demarche{Comparer la distance au centre avec le rayon pour interpréter le signe du produit.}
\begin{questions}
\item $(-2-x)(2-x)+(-y)(-y)=\res{x^2+y^2-4}$.
\item Le produit est nul pour $x^2+y^2=4$, cercle de centre $O$ et de rayon 2.
\item Il est négatif à l’intérieur strict du disque, positif à l’extérieur strict.
\item Pour $P$, il vaut 0 : $P$ est sur le cercle. Pour $Q$, il vaut $-2$ : $Q$ est à l’intérieur.
\end{questions}

\label{figgeo-07-13-cor}
\figureGeo[corrige]{07-13}
\end{corrige}
\end{blocLivre}
\begin{blocLivre}{cor-7-14}
\begin{corrige}{14}{Travail d'une force}
\demarche{Écrire le déplacement avec ses deux composantes et garder les unités physiques.}
\begin{questions}
\item Pour le déplacement horizontal, $W=(30;40)\cdot(6;0)=30\times6+40\times0=\res{180\text{ J}}$.
La norme de la force est $\sqrt{30^2+40^2}=\res{50\text{ N}}$.
\item Un déplacement de norme 6 et d’angle $\pi/6$ a pour coordonnées
$(6\cos(\pi/6);6\sin(\pi/6))=(3\sqrt3;3)$.
Ainsi $W=30(3\sqrt3)+40\times3=\res{90\sqrt3+120\text{ J}}$.
\item Si $\theta$ est l’angle entre la force et le déplacement,
$W=50\times6\cos\theta\leq300$ puisque $\cos\theta\leq1$.
L’égalité est atteinte dans la direction et le sens de $F$ :
$d=6F/\|F\|=6(30;40)/50=\res{(18/5;24/5)}$ m.
Le travail maximal est $\res{300\text{ J}}$.
\end{questions}
\label{figgeo-07-14-cor}
\figureGeo[corrige]{07-14}
\end{corrige}
\end{blocLivre}
\begin{blocLivre}{cor-7-15}
\begin{corrige}{15}{Une identité avec le milieu}
\demarche{Décomposer les vecteurs en passant par le milieu pour supprimer les termes croisés.}
\begin{questions}
\item Comme $I$ est le milieu, $\vect{IB}=-\vect{IA}$.
Les décompositions de Chasles donnent
\begin{align*}
\vect{MA}\cdot\vect{MB}
&=(\vect{MI}+\vect{IA})\cdot(\vect{MI}-\vect{IA})\\
&=MI^2-\vect{MI}\cdot\vect{IA}+\vect{IA}\cdot\vect{MI}-IA^2\\
&=MI^2-IA^2=\res{MI^2-AB^2/4}.
\end{align*}
La symétrie du produit scalaire fait disparaître les termes croisés.
\item Pour $A(0;0)$ et $B(6;0)$, $I=(3;0)$ et $AB^2/4=9$.
Le produit vaut 5 exactement lorsque $MI^2-9=5$, soit $MI=\sqrt{14}$.
Le lieu est donc le cercle de centre $\res{I(3;0)}$ et de rayon $\res{\sqrt{14}}$.
\item Comme $MI^2\geq0$, le produit est au moins $-9$.
Cette valeur est atteinte si et seulement si $MI=0$, donc uniquement en $\res{M=I}$.
\end{questions}
\label{figgeo-07-15-cor}
\figureGeo[corrige]{07-15}
\end{corrige}
\end{blocLivre}

\clearpage
% !TEX root = ../main.tex
\corrigesChapitre{8}{Géométrie repérée}
\corrigesSouples


\begin{blocLivre}{cor-8-1}
\begin{corrige}{1}{Lire une équation de droite}
\demarche{Lire les coefficients de la droite, puis tester les points par substitution.}
\begin{questions}
\item Un normal est $\res{(3;-2)}$ ; un directeur est $\res{(2;3)}$. Leur produit scalaire vaut 0.
\item Pour $A$, $3-8+5=0$ : oui. Pour $B$, $0-4+5=1\neq0$ : non.
\end{questions}

\label{figgeo-08-01-cor}
\figureGeo[corrige]{08-01}
\end{corrige}
\end{blocLivre}
\begin{blocLivre}{cor-8-2}
\begin{corrige}{2}{Construire une équation}
\demarche{Un normal et un point permettent d’écrire une équation cartésienne.}
$3(x-2)+4(y+1)=0$, donc $\res{3x+4y-2=0}$. Un directeur est $(-4;3)$. En posant $x=0$, on obtient $y=1/2$ : l’ordonnée à l’origine vaut $\res{1/2}$.
\label{figgeo-08-02-cor}
\figureGeo[corrige]{08-02}
\end{corrige}
\end{blocLivre}
\begin{blocLivre}{cor-8-3}
\begin{corrige}{3}{Écrire un cercle}
\demarche{Comparer les carrés des distances au carré du rayon.}
\begin{questions}
\item $\res{(x-2)^2+(y+1)^2=9}$.
\item Les carrés des distances au centre valent 9, 9 et 5. $A$ et $B$ sont sur le cercle ; $C$ est à l’intérieur strict.
\end{questions}

\label{figgeo-08-03-cor}
\figureGeo[corrige]{08-03}
\end{corrige}
\end{blocLivre}
\begin{blocLivre}{cor-8-4}
\begin{corrige}{4}{Reconnaître une équation de cercle}
\demarche{Compléter les deux carrés avant de lire le centre et le rayon.}
On complète chaque carré :
$x^2-4x=(x-2)^2-4$ et $y^2+6y=(y+3)^2-9$.
L’équation devient $(x-2)^2+(y+3)^2-4-9-12=0$, donc
$\res{(x-2)^2+(y+3)^2=25}$.
Le centre est $\res{\Omega(2;-3)}$ et le rayon $\res{5}$.
Sur la droite $x=2$, il reste $(y+3)^2=25$ :
$y+3=5$ ou $y+3=-5$, donc $y=2$ ou $y=-8$.
Les deux intersections sont $\res{P(2;2)\text{ et }Q(2;-8)}$.
\label{figgeo-08-04-cor}
\figureGeo[corrige]{08-04}
\end{corrige}
\end{blocLivre}
\begin{blocLivre}{cor-8-5}
\begin{corrige}{5}{Deux projections simples}
\demarche{Le pied d’une projection appartient à la droite et le segment de projection lui est perpendiculaire.}
\begin{questions}
\item L’axe des abscisses est horizontal : le projeté garde l’abscisse 3 et prend l’ordonnée 0.
Ainsi $\res{H_1(3;0)}$ et $PH_1=|4-0|=\res{4}$.
\item Un point de la droite $y=x$ s’écrit $H_2(t;t)$.
Le vecteur $\vect{AH_2}=(t-3;t-1)$ doit être orthogonal au directeur $(1;1)$ :
$(t-3)+(t-1)=0$, donc $2t-4=0$ et $t=2$.
Ainsi $\res{H_2(2;2)}$ et $AH_2=\sqrt{(-1)^2+1^2}=\res{\sqrt2}$.
\end{questions}
\label{figgeo-08-05-cor}
\figureGeo[corrige]{08-05}
\end{corrige}
\end{blocLivre}
\begin{blocLivre}{cor-8-6}
\begin{corrige}{6}{Projeter et mesurer une distance}
\demarche{Déplacer le point dans la direction d’un normal, puis imposer l’appartenance à la droite.}
\begin{questions}
\item Le normal est $(2;1)$, donc $H=(4-2\lambda;3-\lambda)$.
Imposer $H\in d$ donne
$2(4-2\lambda)+(3-\lambda)-5=0$, soit $6-5\lambda=0$.
On trouve $\lambda=6/5$.
\item En remplaçant $\lambda$, $\res{H(8/5;9/5)}$.
On vérifie l’appartenance : $2(8/5)+9/5-5=0$.
De plus $\vect{AH}=(-12/5;-6/5)=-(6/5)(2;1)$ est normal à $d$.
Le point satisfait donc les deux conditions d’une projection orthogonale.
\item $AH=\sqrt{(12/5)^2+(6/5)^2}=(6/5)\sqrt5=\res{6/\sqrt5}$.
Cette longueur est la distance du point à la droite.
\end{questions}
\label{figgeo-08-06-cor}
\figureGeo[corrige]{08-06}
\end{corrige}
\end{blocLivre}
\begin{blocLivre}{cor-8-7}
\begin{corrige}{7}{Droites parallèles et perpendiculaires}
\demarche{Comparer les vecteurs normaux pour obtenir les directions des droites.}
\begin{questions}
\item Les normaux de $d$ et $e$ sont proportionnels, donc les droites sont parallèles. Doubler l’équation de $d$ donne un terme constant 2, différent de $-7$ : elles sont distinctes.
\item Le produit des normaux vaut $3\times2-2\times3=0$.
\item On prend pour normal $(2;3)$ : $2(x-1)+3(y-2)=0$, soit $\res{2x+3y-8=0}$.
\end{questions}

\label{figgeo-08-07-cor}
\figureGeo[corrige]{08-07}
\end{corrige}
\end{blocLivre}
\begin{blocLivre}{cor-8-8}
\begin{corrige}{8}{Intersections avec les axes décalés}
\demarche{Substituer l’équation de chaque droite dans celle du cercle.}
\begin{questions}
\item Pour $y=1$, $(x-1)^2=16$, donc $\res{(-3;1)\text{ et }(5;1)}$.
\item Pour $x=1$, $(y+2)^2=25$, donc $\res{(1;3)\text{ et }(1;-7)}$.
\item Le centre est $(1;-2)$, de rayon 5. Les deux points de la seconde question sont les extrémités du diamètre vertical ; les deux autres sont symétriques par rapport à $x=1$.
\end{questions}

\label{figgeo-08-08-cor}
\figureGeo[corrige]{08-08}
\end{corrige}
\end{blocLivre}
\begin{blocLivre}{cor-8-9}
\begin{corrige}{9}{Une famille d'ensembles}
\demarche{Le second membre après complétion des carrés fixe l’existence du cercle.}
\begin{questions}
\item $\res{(x-1)^2+(y+2)^2=5-m}$.
\item Si $m<5$, cercle de centre $(1;-2)$ et de rayon $\sqrt{5-m}$. Si $m=5$, seul le point $(1;-2)$. Si $m>5$, aucun point.
\item Pour $m=1$, le rayon est $\res{2}$.
\end{questions}

\label{figgeo-08-09-cor}
\figureGeo[corrige]{08-09}
\end{corrige}
\end{blocLivre}
\begin{blocLivre}{cor-8-10}
\begin{corrige}{10}{Déterminer une tangente}
\demarche{Le rayon au point de contact est normal à la tangente.}
\begin{questions}
\item $\Omega A^2=3^2+4^2=25$ : $A$ est sur le cercle.
\item Le rayon $(3;4)$ est normal à la tangente : $3(x-4)+4(y-6)=0$, soit $\res{3x+4y-36=0}$.
\item $A$ appartient à la tangente et $\vect{\Omega A}$ lui est normal : ces deux conditions caractérisent le projeté.
\end{questions}

\label{figgeo-08-10-cor}
\figureGeo[corrige]{08-10}
\end{corrige}
\end{blocLivre}
\begin{blocLivre}{cor-8-11}
\begin{corrige}{11}{Une médiatrice}
\demarche{La médiatrice passe par le milieu du segment et lui est perpendiculaire.}
\begin{questions}
\item $\res{I(1;2)}$ et $\vect{AB}=(6;2)$.
\item La médiatrice passe par $I$ et a $\vect{AB}$ pour normal : $6(x-1)+2(y-2)=0$, soit $\res{3x+y-5=0}$.
\item $CA^2=2^2+(-4)^2=20$ et $CB^2=4^2+(-2)^2=20$, donc $CA=CB$. On retrouve aussi $3\times0+5-5=0$.
\end{questions}

\label{figgeo-08-11-cor}
\figureGeo[corrige]{08-11}
\end{corrige}
\end{blocLivre}
\begin{blocLivre}{cor-8-12}
\begin{corrige}{12}{Hauteurs et cercle circonscrit}
\demarche{Les hauteurs se coupent à l’orthocentre ; les médiatrices au centre du cercle circonscrit.}
\begin{questions}
\item La droite $(AB)$ est horizontale, donc la hauteur issue de $C$ est $x=2$.
La hauteur issue de $A$ a pour normal $\vect{BC}=(-4;4)$ : son équation est $-4x+4y=0$, soit $y=x$.
Les deux hauteurs se coupent en $\res{H(2;2)}$.
\item Le pied $K$ est sur $(BC):x+y-6=0$ et sur la hauteur $y=x$.
Ainsi $2x=6$, donc $\res{K(3;3)}$.
\item La médiatrice de $[AB]$ est $x=3$.
Le milieu de $[AC]$ est $(1;2)$ et $\vect{AC}=(2;4)$ est normal à sa médiatrice :
$2(x-1)+4(y-2)=0$, soit $x+2y-5=0$.
Avec $x=3$, on obtient $y=1$ : $\res{\Omega(3;1)}$.
\item Le rayon vérifie $R^2=\Omega A^2=(0-3)^2+(0-1)^2=10$.
L’équation est donc $\res{(x-3)^2+(y-1)^2=10}$.
Les points $H$ et $\Omega$ viennent de deux constructions distinctes : hauteurs pour $H$, médiatrices pour $\Omega$.
\end{questions}
\label{figgeo-08-12-cor}
\figureGeo[corrige]{08-12}
\end{corrige}
\end{blocLivre}
\begin{blocLivre}{cor-8-13}
\begin{corrige}{13}{Une droite coupe un cercle}
\demarche{Résoudre le système droite-cercle par substitution, puis vérifier les points.}
\begin{questions}
\item $(x-2)^2+(3-x)^2=5$, donc $2x^2-10x+8=0$, soit $x^2-5x+4=0$.
\item Les racines sont 1 et 4, donc les points sont $\res{(1;3)\text{ et }(4;0)}$.
\item Pour $\Omega(2;1)$ et le normal $(1;1)$, $\lambda=(2+1-4)/2=-1/2$. Le projeté est $(5/2;3/2)$ ; la distance vaut $\res{1/\sqrt2}<\sqrt5$. La droite est bien sécante.
\end{questions}

\label{figgeo-08-13-cor}
\figureGeo[corrige]{08-13}
\end{corrige}
\end{blocLivre}
\begin{blocLivre}{cor-8-14}
\begin{corrige}{14}{Un cercle et des droites horizontales}
\demarche{Le signe d’un carré égal à une expression de k détermine le nombre d’intersections.}
\begin{questions}
\item Dans l’équation $(x-1)^2+(y-3)^2=9$, on remplace $y$ par $k$ :
$\res{(x-1)^2=9-(k-3)^2}$.
\item Un carré ne peut pas être négatif. Si $9-(k-3)^2>0$, il y a deux racines :
$|k-3|<3\iff0<k<6$.
Si le second membre est nul, il y a une seule racine, pour $k=0$ ou $k=6$.
S’il est négatif, donc pour $k<0$ ou $k>6$, il n’y a aucune intersection.
\item Dans les deux cas de tangence, $x=1$. Les points sont $\res{T_0(1;0)}$ et $\res{T_6(1;6)}$.
\item Pour $k=2$, $(x-1)^2=9-1=8$, donc $x=1\pm\sqrt8=1\pm2\sqrt2$.
On obtient $\res{I(1-2\sqrt2;2)}$ et $\res{J(1+2\sqrt2;2)}$.
\end{questions}
\label{figgeo-08-14-cor}
\figureGeo[corrige]{08-14}
\end{corrige}
\end{blocLivre}
\begin{blocLivre}{cor-8-15}
\begin{corrige}{15}{Une parabole comme lieu géométrique}
\demarche{Écrire l’égalité de deux distances non négatives avant de les élever au carré.}
\begin{questions}
\item Le projeté sur l’axe est $(x;0)$, donc la distance est $|y|$. L’égalité cherchée est $x^2+(y-2)^2=y^2$.
\item En développant, $x^2-4y+4=0$, soit $\res{y=x^2/4+1}$. Les deux distances sont positives ou nulles, donc l’égalité des carrés est équivalente à celle des distances.
\item $5=x^2/4+1$ donne $x^2=16$, donc $\res{(-4;5)\text{ et }(4;5)}$.
\item $f\prime(x)=x/2$, donc la pente en 4 vaut 2 et la tangente est $\res{y=2x-3}$. Pour $M(4;5)$, $MF=\sqrt{16+9}=5$, égal à la distance à l’axe.
\end{questions}

\label{figgeo-08-15-cor}
\figureGeo[corrige]{08-15}
\end{corrige}
\end{blocLivre}

\clearpage
% !TEX root = ../main.tex
\corrigesChapitre{9}{Probabilités conditionnelles}
\corrigesSouples

\begin{blocLivre}{cor-9-1}
\begin{corrige}{1}{Appliquer la définition}
\demarche{Identifier la population conditionnante avant de choisir le dénominateur.}
\begin{questions}
\item Comme $P(A)=0{,}4>0$, on peut conditionner par $A$ :
$P_A(B)=P(A\cap B)/P(A)=0{,}1/0{,}4=\res{0{,}25}$.
Dans le même groupe $A$, les évènements $B$ et $\overline B$ sont complémentaires,
donc $P_A(\overline B)=1-0{,}25=\res{0{,}75}$.
\item Les évènements $A\cap B$ et $A\cap\overline B$ se partagent $A$.
Ainsi $P(A\cap\overline B)=0{,}4-0{,}1=\res{0{,}3}$.
Le produit $0{,}4\times0{,}75$ donne le même résultat.
\end{questions}

\end{corrige}
\end{blocLivre}
\begin{blocLivre}{cor-9-2}
\begin{corrige}{2}{Lire un tableau d'effectifs}
\demarche{Compléter les effectifs et garder les totaux de lignes et de colonnes.}
\begin{questions}
\item Il y a $200-120=80$ garçons.
Parmi les filles, $120-50=70$ ne suivent pas l’option ; parmi les garçons, $80-30=50$ ne la suivent pas.
On en déduit les totaux 80 avec option et 120 sans option, indiqués dans le tableau.
\item Le choix parmi les 200 élèves donne $P(F)=120/200=\res{3/5}$.
Sachant $F$, la population de référence comporte 120 filles : $P_F(A)=50/120=\res{5/12}$.
Sachant $A$, elle comporte 80 élèves avec option : $P_A(F)=50/80=\res{5/8}$.
Les numérateurs sont les mêmes, mais les dénominateurs correspondent à deux groupes différents.
\end{questions}

\label{figlivre-r09-02-cor}
\figureLivre[corrige]{r09-02}

\end{corrige}
\end{blocLivre}
\begin{blocLivre}{cor-9-3}
\begin{corrige}{3}{Deux évènements indépendants}
\demarche{Employer l’indépendance pour l’intersection, puis la formule de la réunion.}
\begin{questions}
\item $P(A\cap B)=0{,}3\times0{,}5=\res{0{,}15}$. La réunion vaut $0{,}3+0{,}5-0{,}15=\res{0{,}65}$.
\item $P_A(B)=P(B)=\res{0{,}5}$ et $P(A\cap\overline B)=0{,}3\times0{,}5=\res{0{,}15}$.
\end{questions}
\end{corrige}
\end{blocLivre}
\begin{blocLivre}{cor-9-4}
\begin{corrige}{4}{Incompatibles ne signifie pas indépendants}
\demarche{Comparer la probabilité de l’intersection au produit des probabilités.}
\begin{questions}
\item L’intersection est vide : $P(A\cap B)=\res{0}$, donc $P_A(B)=\res{0}$.
\item $P(A)P(B)=0{,}06\neq0=P(A\cap B)$ : ils ne sont pas indépendants.
\end{questions}
\end{corrige}
\end{blocLivre}
\begin{blocLivre}{cor-9-5}
\begin{corrige}{5}{Exploiter un arbre}
\demarche{Multiplier sur un chemin, additionner entre chemins incompatibles, puis conditionner.}
\begin{questions}
\item Les poids sortant d’un même nœud ont une somme égale à 1.
On complète donc $P(\overline A)=1-0{,}6=0{,}4$,
$P_A(\overline B)=1-0{,}2=0{,}8$ et $P_{\overline A}(\overline B)=1-0{,}5=0{,}5$.
\item Sur le chemin $A$ puis $B$, on multiplie : $P(A\cap B)=0{,}6\times0{,}2=0{,}12$.
Sur l’autre chemin qui arrive à $B$, $P(\overline A\cap B)=0{,}4\times0{,}5=0{,}20$.
Ces chemins sont incompatibles et couvrent $B$, donc $P(B)=0{,}12+0{,}20=\res{0{,}32}$.
Enfin, parmi les issues de $B$, la part issue de $A$ vaut
$P_B(A)=P(A\cap B)/P(B)=0{,}12/0{,}32=\res{3/8}$.
\end{questions}

\label{figlivre-r09-05-cor}
\figureLivre[corrige]{r09-05}

\end{corrige}
\end{blocLivre}
\begin{blocLivre}{cor-9-6}
\begin{corrige}{6}{Deux ateliers de fabrication}
\demarche{Séparer les deux origines des pièces avant de regrouper les défauts.}
\begin{questions}
\item Les ateliers A et B forment une partition. Les deux contributions au défaut sont
$P(A\cap D)=0{,}7\times0{,}02=0{,}014$ et $P(B\cap D)=0{,}3\times0{,}05=0{,}015$.
Par probabilités totales, $P(D)=0{,}014+0{,}015=\res{0{,}029}$, soit 2,9\,\%.
\item Le groupe conditionnant est celui des pièces défectueuses.
On divise la contribution de B par la probabilité totale de défaut :
$P_D(B)=0{,}015/0{,}029=\res{15/29\approx0{,}5172}$.
Environ 51,72\,\% des pièces défectueuses proviennent donc de B dans ce modèle.
\end{questions}

\label{figlivre-r09-06-cor}
\figureLivre[corrige]{r09-06}

\end{corrige}
\end{blocLivre}
\begin{blocLivre}{cor-9-7}
\begin{corrige}{7}{Une partition en trois groupes}
\demarche{Additionner les trois contributions des machines à l’évènement défaut.}
\begin{questions}
\item Chaque contribution est une probabilité de chemin :
$P(A\cap D)=0{,}5\times0{,}01=0{,}005$,
$P(B\cap D)=0{,}3\times0{,}02=0{,}006$,
$P(C\cap D)=0{,}2\times0{,}04=0{,}008$.
Le total vaut $P(D)=0{,}005+0{,}006+0{,}008=\res{0{,}019}$.
\item Parmi les pièces défectueuses, la proportion venant de C est
$P_D(C)=0{,}008/0{,}019=\res{8/19\approx0{,}4211}$.
\item Les trois contributions se comparent sur la même population totale :
$0{,}008>0{,}006>0{,}005$. C contribue le plus aux défauts ; son petit volume de production ne suffit pas à compenser son taux de défaut élevé.
\end{questions}

\label{figlivre-r09-07-cor}
\figureLivre[corrige]{r09-07}

\end{corrige}
\end{blocLivre}
\begin{blocLivre}{cor-9-8}
\begin{corrige}{8}{Inverser un conditionnement}
\demarche{Compter d’abord les positifs dans chaque groupe pour inverser le conditionnement.}
On note $C$ la présence du caractère et $T$ un résultat positif.
\begin{questions}
\item Sur 10000 personnes, le groupe $C$ compte $10000\times0{,}01=100$ personnes,
dont $100\times0{,}95=95$ positives.
Le groupe $\overline C$ compte 9900 personnes, dont $9900\times0{,}04=396$ positives.
\item Le nombre total de positifs est $95+396=491$.
Parmi eux, seuls 95 ont le caractère : $P_T(C)=\res{95/491\approx19{,}35\,\%}$.
Avec l’arbre, on retrouve $P(C\cap T)=0{,}0095$ et $P(T)=0{,}0095+0{,}0396=0{,}0491$.
\item Le taux 95\,\% est $P_C(T)$ : on connaît le caractère et on cherche la fréquence de positifs.
La question demande $P_T(C)$ : on connaît le résultat positif.
La rareté du caractère et les positifs dans l’autre groupe modifient fortement cette proportion.
\end{questions}

\label{figlivre-r09-08-cor}
\figureLivre[corrige]{r09-08}

\end{corrige}
\end{blocLivre}
\begin{blocLivre}{cor-9-9}
\begin{corrige}{9}{Deux épreuves de Bernoulli}
\demarche{Lister les quatre chemins et identifier ceux qui réalisent chaque évènement.}
\begin{questions}
\item À chaque nœud, $P(S)=0{,}4$ et $P(E)=0{,}6$ par indépendance.
Les quatre produits sont $P(SS)=0{,}4^2=\res{0{,}16}$,
$P(SE)=P(ES)=0{,}4\times0{,}6=\res{0{,}24}$ et $P(EE)=0{,}6^2=\res{0{,}36}$.
Leur somme $0{,}16+0{,}24+0{,}24+0{,}36=1$ contrôle l’arbre.
\item Exactement un succès correspond à deux chemins incompatibles, $SE$ et $ES$ :
$P=0{,}24+0{,}24=\res{0{,}48}$.
Au moins un succès est le contraire de $EE$, donc $P=1-0{,}36=\res{0{,}64}$.
\end{questions}

\label{figlivre-r09-09-cor}
\figureLivre[corrige]{r09-09}

\end{corrige}
\end{blocLivre}
\begin{blocLivre}{cor-9-10}
\begin{corrige}{10}{Tirer sans remise}
\demarche{Après un tirage sans remise, actualiser les deux effectifs et le total.}
\begin{questions}
\item Après une rouge, il reste 2 rouges et 2 bleues ; après une bleue, 3 rouges et 1 bleue.
Les probabilités du deuxième niveau ont donc pour dénominateur 4.
\[P(RR)=\frac35\frac24=\frac3{10},\qquad P(RB)=\frac35\frac24=\frac3{10},\]
\[P(BR)=\frac25\frac34=\frac3{10},\qquad P(BB)=\frac25\frac14=\frac1{10}.\]
\item Deux couleurs différentes correspondent à $RB$ ou $BR$ : $P=3/10+3/10=\res{3/5}$.
\item Si $R_1$ et $R_2$ sont les évènements rouge au premier et au deuxième tirage,
$P(R_2)=P(RR)+P(BR)=3/5$, mais $P_{R_1}(R_2)=2/4=1/2$.
Ces probabilités sont différentes : $R_1$ et $R_2$ ne sont pas indépendants.
\end{questions}

\label{figlivre-r09-10-cor}
\figureLivre[corrige]{r09-10}

\end{corrige}
\end{blocLivre}
\begin{blocLivre}{cor-9-11}
\begin{corrige}{11}{Tirer avec remise}
\demarche{La remise rétablit la composition de l’urne avant le second tirage.}
\begin{questions}
\item La remise rétablit les 3 rouges et les 2 bleues. Le modèle donne à nouveau $3/5$ pour rouge et $2/5$ pour bleue, quel que soit le premier résultat : les tirages sont indépendants.
\item Pour deux rouges, $P(RR)=(3/5)^2=\res{9/25}$.
Pour deux couleurs différentes, les deux chemins sont incompatibles et ont le même poids :
$P(RB)+P(BR)=(3/5)(2/5)+(2/5)(3/5)=\res{12/25}$.
\item Avec remise, $9/25=0{,}36>0{,}30$ pour deux rouges.
Pour deux couleurs différentes, $12/25=0{,}48<0{,}60$.
La comparaison doit conserver l’évènement étudié ; la remise modifie les poids du deuxième niveau.
\end{questions}

\label{figlivre-r09-11-cor}
\figureLivre[corrige]{r09-11}

\end{corrige}
\end{blocLivre}
\begin{blocLivre}{cor-9-12}
\begin{corrige}{12}{Trois tentatives indépendantes}
\demarche{Distinguer exactement deux succès, au moins un succès et premier succès tardif.}
\begin{questions}
\item L’arbre a trois niveaux et huit chemins. Exactement deux succès correspond à $SSE$, $SES$ et $ESS$.
Chaque chemin a deux facteurs $0{,}7$ et un facteur $0{,}3$ : son poids est $0{,}7^2\times0{,}3=0{,}147$.
\item Ces trois chemins sont incompatibles : leur probabilité totale est $3\times0{,}147=\res{0{,}441}$.
\item L’évènement contraire d’au moins un succès est aucun succès, soit $EEE$.
Ainsi $P=1-0{,}3^3=1-0{,}027=\res{0{,}973}$.
\item Pour que le premier succès arrive au troisième essai, il faut deux échecs puis un succès :
seul $EES$ convient, de poids $0{,}3^2\times0{,}7=\res{0{,}063}$.
\end{questions}

\label{figlivre-r09-12-cor}
\figureLivre[corrige]{r09-12}

\end{corrige}
\end{blocLivre}
\begin{blocLivre}{cor-9-13}
\begin{corrige}{13}{Quatre essais et six chemins}
\demarche{Compter les chemins au lieu d’utiliser une formule de dénombrement hors de ce chapitre.}
\begin{questions}
\item Les six chemins sont $SSEE$, $SESE$, $SEES$, $ESSE$, $ESES$ et $EESS$.
Ils sont repérés en vert à droite de l’arbre.
\item Chacun comporte deux succès et deux échecs, donc a pour poids
$(1/4)^2(3/4)^2=9/256$.
En les additionnant, on obtient $6\times9/256=\res{27/128}$.
\item Aucun succès correspond au seul chemin $EEEE$, de poids $(3/4)^4=81/256$.
Le contraire donne $P(\text{au moins un succès})=1-81/256=\res{175/256}$.
\item Au plus un succès réunit zéro succès et exactement un succès.
Il y a quatre chemins avec un succès, selon sa position ; chacun a pour poids
$(1/4)(3/4)^3=27/256$.
Le total vaut $81/256+4\times27/256=\res{189/256}$.
\end{questions}

\label{figlivre-r09-13-cor}
\figureLivre[corrige]{r09-13}

\end{corrige}
\end{blocLivre}
\begin{blocLivre}{cor-9-14}
\begin{corrige}{14}{Retrouver une proportion inconnue}
\demarche{Utiliser les probabilités totales pour déterminer la proportion inconnue.}
\begin{questions}
\item Les groupes $A$ et $\overline A$ forment une partition.
Les probabilités totales donnent $P(B)=0{,}1x+0{,}4(1-x)$.
Ainsi $0{,}4-0{,}3x=0{,}22$, donc $-0{,}3x=-0{,}18$ et $\res{x=0{,}6}$,
valeur bien dans $]0;1[$.
\item $P(A\cap B)=P(A)P_A(B)=0{,}6\times0{,}1=0{,}06$.
En conditionnant par $B$, $P_B(A)=0{,}06/0{,}22=\res{3/11}$.
\item Première méthode : $P_A(B)=0{,}1\neq0{,}22=P(B)$.
Deuxième méthode : $P(A\cap B)=0{,}06$, alors que $P(A)P(B)=0{,}6\times0{,}22=0{,}132$.
Les deux critères montrent l’absence d’indépendance.
\end{questions}

\label{figlivre-r09-14-cor}
\figureLivre[corrige]{r09-14}

\end{corrige}
\end{blocLivre}
\begin{blocLivre}{cor-9-15}
\begin{corrige}{15}{Comparer des taux sans se tromper}
\demarche{Comparer les taux dans chaque parcours, puis tenir compte des effectifs pour le total.}
\begin{questions}
\item Dans le parcours facile, A obtient $90/100=90\,\%$ et B $19/20=95\,\%$.
Dans le parcours difficile, A obtient $20/100=20\,\%$ et B $18/60=30\,\%$.
B a le meilleur taux dans chacun des deux parcours.
\item A compte $90+20=110$ réussites sur 200 personnes : $\res{55\,\%}$.
B compte $19+18=37$ réussites sur 80 personnes : $\res{46{,}25\,\%}$.
La comparaison globale favorise donc A.
\item La part du parcours facile est $100/200=\res{50\,\%}$ chez A et $20/80=\res{25\,\%}$ chez B.
\item Les taux globaux sont des moyennes pondérées :
$0{,}5(0{,}90)+0{,}5(0{,}20)=0{,}55$ chez A,
et $0{,}25(0{,}95)+0{,}75(0{,}30)=0{,}4625$ chez B.
B accueille davantage de parcours difficiles en proportion.
Le renversement est expliqué par ces poids ; ces seuls taux ne prouvent pas un effet causal du centre.
\end{questions}

\label{figlivre-r09-15-cor}
\figureLivre[corrige]{r09-15}

\end{corrige}
\end{blocLivre}

\clearpage
% !TEX root = ../main.tex
\corrigesChapitre{10}{Variables aléatoires}
\corrigesSouples

\begin{blocLivre}{cor-10-1}
\begin{corrige}{1}{Une variable indicatrice}
\demarche{Regrouper les issues qui donnent la même valeur de la variable.}
\begin{questions}
\item Trois issues sont paires sur six : $P(X=1)=P(X=0)=\res{1/2}$.
\item $E(X)=0\times1/2+1\times1/2=\res{1/2}$. $P(X\leq0)=P(X=0)=\res{1/2}$.
\end{questions}
\label{figlivre-r10-01-loi-cor}
\figureLivre[corrige]{r10-01-loi}

\end{corrige}
\end{blocLivre}
\begin{blocLivre}{cor-10-2}
\begin{corrige}{2}{Lire et exploiter une loi}
\demarche{Sélectionner les valeurs autorisées par chaque évènement avant d’additionner leurs probabilités.}
\begin{questions}
\item On additionne les masses concernées : $\res{0{,}7}$ et $\res{0{,}3}$.
\item $E(X)=-0{,}2+0{,}6=\res{0{,}4}$. $E(X^2)=0{,}2+4\times0{,}3=\res{1{,}4}$.
\end{questions}
\end{corrige}
\end{blocLivre}
\begin{blocLivre}{cor-10-3}
\begin{corrige}{3}{Compléter une loi}
\demarche{Compléter les probabilités à 1 avant de calculer les moments.}
\begin{questions}
\item La somme vaut 1, donc $p=1-0{,}1-0{,}2-0{,}4=\res{0{,}3}$.
\item $E(X)=0{,}2+0{,}6+1{,}2=\res{2}$. $E(X^2)=0{,}2+1{,}2+3{,}6=5$. Par König-Huygens, $V(X)=5-2^2=\res{1}$.
\end{questions}
\label{figlivre-r10-03-loi-cor}
\figureLivre[corrige]{r10-03-loi}

\end{corrige}
\end{blocLivre}
\begin{blocLivre}{cor-10-4}
\begin{corrige}{4}{Une loi de Bernoulli}
\demarche{Utiliser les deux valeurs possibles dans les définitions de l’espérance et de la variance.}
\begin{questions}
\item À partir des deux valeurs possibles,
$E(X)=0\times0{,}7+1\times0{,}3=\res{0{,}3}$.
La définition de la variance donne
\[V(X)=(0-0{,}3)^2\times0{,}7+(1-0{,}3)^2\times0{,}3
=0{,}063+0{,}147=\res{0{,}21}.\]
On peut contrôler avec $E(X^2)-E(X)^2=0{,}3-0{,}09$.
\item L’écart type est la racine positive de la variance :
$\sigma(X)=\sqrt{0{,}21}=\res{\sqrt{21}/10\approx0{,}458}$.
\end{questions}

\end{corrige}
\end{blocLivre}
\begin{blocLivre}{cor-10-5}
\begin{corrige}{5}{Un jeu équitable}
\demarche{Le gain net inclut la perte éventuelle ; un jeu équitable a une espérance nulle.}
\begin{questions}
\item $P(G=5)=1/6$ et $P(G=-1)=5/6$. $E(G)=5/6-5/6=\res{0}$.
\item $E(G^2)=25/6+5/6=5$, donc $\res{V(G)=5}$. Le gain moyen théorique est nul ; cela ne garantit ni gain ni perte nulle sur une partie.
\end{questions}
\label{figlivre-r10-05-loi-cor}
\figureLivre[corrige]{r10-05-loi}

\end{corrige}
\end{blocLivre}
\begin{blocLivre}{cor-10-6}
\begin{corrige}{6}{Calculer la dispersion de deux façons}
\demarche{Comparer les deux expressions de la variance avec des fractions exactes.}
\begin{questions}
\item $E(X)=(-1)/5+0/2+2(3/10)=2/5$ et $E(X^2)=1/5+4(3/10)=7/5$.
Par König-Huygens, $V(X)=7/5-(2/5)^2=35/25-4/25=\res{31/25}$.
\item Les trois écarts à $2/5$ sont $-7/5$, $-2/5$ et $8/5$.
Leurs carrés pondérés donnent
\[V(X)=\frac{49}{25}\frac15+\frac4{25}\frac12+\frac{64}{25}\frac3{10}
=\frac{49+10+96}{125}=\frac{31}{25}.\]
Les deux méthodes donnent bien la même variance.
\item $\sigma(X)=\sqrt{31/25}=\res{\sqrt{31}/5}$.
La racine porte sur toute la variance.
\end{questions}

\end{corrige}
\end{blocLivre}
\begin{blocLivre}{cor-10-7}
\begin{corrige}{7}{Un changement d'unité}
\demarche{Transformer les valeurs et appliquer les règles de l’espérance et de la variance.}
\begin{questions}
\item La linéarité donne $E(Y)=3E(X)-2=3(2/5)-2=\res{-4/5}$.
La variance est multipliée par le carré de 3 :
$V(Y)=3^2V(X)=9(31/25)=\res{279/25}$.
Donc $\sigma(Y)=\sqrt{279/25}=\res{3\sqrt{31}/5}$.
La translation de $-2$ ne modifie pas la dispersion.
\item Les images de $-1$, 0 et 2 par $x\mapsto3x-2$ sont $-5$, $-2$ et 4.
Ces trois valeurs restent distinctes : chacune garde la probabilité de son antécédent, comme dans le tableau.
\end{questions}

\label{figlivre-r10-07-loi-cor}
\figureLivre[corrige]{r10-07-loi}

\end{corrige}
\end{blocLivre}
\begin{blocLivre}{cor-10-8}
\begin{corrige}{8}{Même moyenne, risques différents}
\demarche{Comparer les montants autour de leur moyenne, puis les probabilités de gain.}
\begin{questions}
\item Les deux espérances valent 0. Pour A, $V(A)=4$ et $\sigma(A)=2$ ; pour B, $V(B)=1$ et $\sigma(B)=1$.
\item Les deux probabilités sont $\res{1/2}$.
\item A présente des écarts plus grands à la moyenne, sans avoir une plus grande probabilité de gain. La variance mesure la dispersion des montants.
\end{questions}
\end{corrige}
\end{blocLivre}
\begin{blocLivre}{cor-10-9}
\begin{corrige}{9}{Linéarité sans indépendance}
\demarche{La linéarité de l’espérance reste valable lorsque les variables dépendent l’une de l’autre.}
\begin{questions}
\item $E(X)=(1+\cdots+6)/6=\res{7/2}$. Par linéarité $E(Y)=7-E(X)=\res{7/2}$.
\item $X+Y=7$ pour toute issue : $E(X+Y)=7$ et $V(X+Y)=0$. On retrouve $E(X)+E(Y)=7$.
\item Les deux évènements coïncident : leur intersection a probabilité $1/6$, différente de $(1/6)^2=1/36$. La linéarité a donc été utilisée sans indépendance.
\end{questions}
\end{corrige}
\end{blocLivre}
\begin{blocLivre}{cor-10-10}
\begin{corrige}{10}{Deux variables sur le même univers}
\demarche{Suivre les deux variables sur les mêmes quatre issues avant de les combiner.}
\begin{questions}
\item Les deux pièces indépendantes donnent quatre issues équiprobables, notées dans l’ordre des lancers.
Pour chaque ligne du tableau, on compte les faces pour $X$, on lit la première pièce pour $Y$, puis on calcule $Z=X-2Y$.
Par exemple, sur $FP$, $X=1$ et $Y=1$, donc $Z=-1$.
\item $E(X)=1$ et $E(Y)=1/2$. La linéarité ne demande pas d’indépendance entre $X$ et $Y$ :
$E(Z)=E(X)-2E(Y)=1-2(1/2)=\res{0}$.
\item La valeur $-1$ apparaît sur $FP$, la valeur 1 sur $PF$, et la valeur 0 sur $PP$ ou $FF$.
Donc $P(Z=-1)=1/4$, $P(Z=0)=1/2$, $P(Z=1)=1/4$.
On vérifie $E(Z)=(-1)/4+1/4=0$ et
$V(Z)=E(Z^2)-0^2=1/4+1/4=\res{1/2}$.
\end{questions}

\label{figlivre-r10-10-issues-cor}
\figureLivre[corrige]{r10-10-issues}

\end{corrige}
\end{blocLivre}
\begin{blocLivre}{cor-10-11}
\begin{corrige}{11}{Choisir une valeur de prévision}
\demarche{Développer le carré de l’erreur et utiliser la linéarité.}
\begin{questions}
\item $E(X)=-1/4+1/2+3/4=\res{1}$. $E(X^2)=1/4+1/2+9/4=\res{3}$.
\item Par développement et linéarité, $E((X-a)^2)=E(X^2)-2aE(X)+a^2=3-2a+a^2=(a-1)^2+2$.
\item Le carré est minimal pour $\res{a=1}$ ; l’erreur minimale vaut $\res{2}$, qui est la variance de $X$.
\end{questions}
\end{corrige}
\end{blocLivre}
\begin{blocLivre}{cor-10-12}
\begin{corrige}{12}{Une loterie et son prix d'entrée}
\demarche{Distinguer le lot brut, le prix d’entrée et le gain net.}
\begin{questions}
\item Il reste $100-3-10=87$ tickets sans lot.
La loi de $L$ est donnée dans le tableau et son espérance vaut
$E(L)=0(87/100)+5(10/100)+20(3/100)=\res{1{,}10\text{ euro}}$.
\item Le prix est payé à chaque partie : $G=L-m$, donc $E(G)=E(L)-m=1{,}1-m$.
Un jeu équitable vérifie $E(G)=0$, ce qui donne $\res{m=1{,}10\text{ euro}}$.
\item Pour $m=2$, $E(G)=\res{-0{,}90\text{ euro}}$.
Soustraire une constante conserve la variance : $V(G)=V(L)$.
Comme $E(L^2)=0+25(0{,}10)+400(0{,}03)=14{,}5$,
$V(G)=14{,}5-1{,}1^2=\res{13{,}29}$ euros carrés.
\item Un gain net positif impose $L>2$ : seuls les lots de 5 et 20 euros conviennent.
La probabilité est donc $10/100+3/100=\res{13/100}$.
\end{questions}

\label{figlivre-r10-12-loi-cor}
\figureLivre[corrige]{r10-12-loi}

\end{corrige}
\end{blocLivre}
\begin{blocLivre}{cor-10-13}
\begin{corrige}{13}{Un jeu sur trois pièces}
\demarche{Compter les huit issues équiprobables et les regrouper par nombre de faces.}
\begin{questions}
\item Les huit issues ont chacune la probabilité $1/8$.
Le tableau les regroupe par nombre de faces : on trouve respectivement 1, 3, 3 et 1 issues,
donc les probabilités $1/8$, $3/8$, $3/8$, $1/8$ pour $X=0,1,2,3$.
\item $E(X)=[0\times1+1\times3+2\times3+3\times1]/8=\res{3/2}$.
De même $E(X^2)=(3+12+9)/8=3$, donc $V(X)=3-(3/2)^2=\res{3/4}$.
\item Le gain net est $G=4X-5$ : la mise est retranchée quel que soit le résultat.
Ainsi $E(G)=4(3/2)-5=\res{1\text{ euro}}$ et $V(G)=4^2(3/4)=\res{12}$ euros carrés.
\item Pour une mise $m$, l’équité demande $4E(X)-m=0$, donc $\res{m=6\text{ euros}}$.
Avec la mise 5, $4X-5>0\iff X>5/4$. Comme $X$ est entier, cela équivaut à $X\geq2$,
de probabilité $(3+1)/8=\res{1/2}$.
\end{questions}

\label{figlivre-r10-13-issues-cor}
\figureLivre[corrige]{r10-13-issues}

\end{corrige}
\end{blocLivre}
\begin{blocLivre}{cor-10-14}
\begin{corrige}{14}{Un modèle de garantie}
\demarche{Retirer le coût de réparation du prix reçu, puis calculer la dispersion.}
\begin{questions}
\item Sans réparation, le résultat net vaut 15 euros.
Avec réparation, il vaut $15-500=-485$ euros.
La loi du tableau donne $E(G)=15(0{,}98)-485(0{,}02)=\res{5\text{ euros}}$.
\item Posons $I=1$ si une réparation est nécessaire, et $I=0$ sinon.
Alors $G=15-500I$, avec $V(I)=0{,}02\times0{,}98$.
Donc $V(G)=500^2\times0{,}02\times0{,}98=\res{4900}$ euros carrés et
$\sigma(G)=\sqrt{4900}=\res{70\text{ euros}}$.
La dépense rare de 500 euros explique cette forte dispersion autour d’une espérance de 5 euros.
\item Si $G_1,\ldots,G_{100}$ sont les résultats des garanties,
$E(G_1+\cdots+G_{100})=100\times5=\res{500\text{ euros}}$.
Cette addition des espérances ne demande pas d’indépendance.
\end{questions}

\label{figlivre-r10-14-loi-cor}
\figureLivre[corrige]{r10-14-loi}

\end{corrige}
\end{blocLivre}
\begin{blocLivre}{cor-10-15}
\begin{corrige}{15}{Une aire aléatoire}
\demarche{Exprimer l’aire en fonction de la hauteur avant d’établir sa loi.}
\begin{questions}
\item Le triangle est rectangle en $A$, avec base 1 et hauteur $H$.
Ainsi $S=1\times H/2=H/2$.
Les trois valeurs possibles sont $1/2$, 1 et 2, avec les probabilités du tableau.
\item En utilisant directement la loi de $S$,
$E(S)=(1/2)(1/2)+1(1/4)+2(1/4)=\res{1}$.
Puis $E(S^2)=(1/4)(1/2)+1(1/4)+4(1/4)=11/8$.
Donc $V(S)=11/8-1^2=\res{3/8}$ et $P(S\geq1)=1/4+1/4=\res{1/2}$.
\item Développer le carré donne
$E((S-a)^2)=E(S^2)-2aE(S)+a^2=11/8-2a+a^2=\res{(a-1)^2+3/8}$.
\item Le carré $(a-1)^2$ est minimal en $a=1$.
La meilleure aire fixe est donc $\res{1}$ unité d’aire, pour une erreur quadratique moyenne de $3/8$.
\end{questions}

\label{figlivre-r10-15-loi-cor}
\figureLivre[corrige]{r10-15-loi}

\label{figlivre-r10-15-triangle-cor}
\figureLivre[corrige]{r10-15-triangle}

\end{corrige}
\end{blocLivre}

\clearpage
% !TEX root = ../main.tex
\corrigesChapitre{11}{Statistiques et simulations}
\corrigesSouples

\begin{blocLivre}{cor-11-1}
\begin{corrige}{1}{D’un tableau aux indicateurs}
\demarche{Calculer les effectifs cumulés pour repérer les rangs de la médiane.}
\begin{questions}
\item $N=10$ ; les fréquences sont $\res{0{,}2;\ 0{,}3;\ 0{,}4;\ 0{,}1}$.
\item $\overline x=(3+8+3)/10=\res{1{,}4}$. Les valeurs aux rangs 5 et 6 sont 1 et 2 : médiane $\res{1{,}5}$. Étendue $\res{3}$.
\end{questions}
\end{corrige}
\end{blocLivre}
\begin{blocLivre}{cor-11-2}
\begin{corrige}{2}{Lire un diagramme en barres}
\demarche{Lire les effectifs sur le graphique, puis ordonner les observations.}
\begin{questions}
\item Les effectifs sont 2, 4, 3 et 1, donc $N=10$. $\overline x=(20+48+42+16)/10=\res{12{,}6}$.
\item Les valeurs centrales sont 12 et 12 : médiane 12. Les rangs des quartiles sont 3 et 8, donc $\res{Q_1=12,\ Q_3=14}$.
\end{questions}
\end{corrige}
\end{blocLivre}
\begin{blocLivre}{cor-11-3}
\begin{corrige}{3}{Tracer une boîte}
\demarche{Repérer les rangs des quartiles avant de placer les cinq repères de la boîte.}
\begin{questions}
\item La série a huit valeurs. Les deux valeurs centrales sont celles des rangs 4 et 5, soit 3 et 4 :
la médiane vaut $(3+4)/2=\res{3{,}5}$.
Les rangs des quartiles sont $\lceil8/4\rceil=2$ et $\lceil3\times8/4\rceil=6$.
On lit donc $\res{Q_1=2}$ et $\res{Q_3=4}$.
\item On place le minimum 1, les quartiles 2 et 4, la médiane 3,5 et le maximum 7.
Le trait de médiane se trouve à l’intérieur de la boîte ; les moustaches relient les bornes de la boîte aux extrêmes.
L’écart interquartile vaut $Q_3-Q_1=4-2=\res{2}$.
\end{questions}

\label{figlivre-r11-03-boite-cor}
\figureLivre[corrige]{r11-03-boite}

\end{corrige}
\end{blocLivre}
\begin{blocLivre}{cor-11-4}
\begin{corrige}{4}{Une fréquence observée}
\demarche{Une moyenne observée peut différer de l’espérance sans contredire le modèle.}
\begin{questions}
\item Quatre succès sur dix : moyenne $\res{0{,}4}$.
\item L’espérance vaut $0{,}3$, donc l’écart est $\res{0{,}1}$. Une fluctuation sur un échantillon ne suffit pas à invalider le modèle.
\end{questions}
\end{corrige}
\end{blocLivre}
\begin{blocLivre}{cor-11-5}
\begin{corrige}{5}{Comprendre un tirage simulé}
\demarche{Distinguer le résultat d’un appel et le nombre de succès sur plusieurs appels.}
\begin{questions}
\item Une loi de Bernoulli de paramètre $\res{1/4}$.
\item Le nombre de succès peut être $0,1,2,3,4$, donc la moyenne est dans $\res{\{0;1/4;1/2;3/4;1\}}$.
\item Non : $0{,}25$ est l’espérance, pas une valeur imposée à tout échantillon.
\end{questions}
\end{corrige}
\end{blocLivre}
\begin{blocLivre}{cor-11-6}
\begin{corrige}{6}{Un histogramme de largeurs inégales}
\demarche{Pour un histogramme, c’est l’aire qui porte la fréquence, pas la hauteur seule.}
\begin{questions}
\item L’aire d’un rectangle est largeur multipliée par hauteur.
Pour représenter une fréquence $f$, on choisit donc hauteur $=f/\text{largeur}$.
Les classes ont pour largeurs $10$, $20$, $30$, d’où les hauteurs
$0{,}2/10=\res{0{,}02}$, $0{,}5/20=\res{0{,}025}$ et $0{,}3/30=\res{0{,}01}$.
On vérifie les aires : $10(0{,}02)=0{,}2$, $20(0{,}025)=0{,}5$, $30(0{,}01)=0{,}3$.
\item La deuxième classe est la plus haute et représente la fréquence la plus importante.
Mais les hauteurs seules ne donnent pas les fréquences lorsque les largeurs diffèrent.
Prendre 0,2, 0,5 et 0,3 comme hauteurs donnerait des aires 2, 10 et 9, dont les proportions ne sont pas celles des données.
\end{questions}

\label{figlivre-r11-06-histogramme-cor}
\figureLivre[corrige]{r11-06-histogramme}

\end{corrige}
\end{blocLivre}
\begin{blocLivre}{cor-11-7}
\begin{corrige}{7}{Transformer une série}
\demarche{Un changement d’échelle multiplie la dispersion ; une translation conserve cette dispersion.}
\begin{questions}
\item $\overline Y=2\times12+5=\res{29}$ ; $\overline Z=\res{0}$.
\item La variance initiale vaut 9. Pour Y : variance 36 et écart type 6. Pour Z : variance 9 et écart type 3.
\item Le centrage retranche la moyenne à chaque valeur ; il déplace la série sans modifier sa dispersion.
\end{questions}
\end{corrige}
\end{blocLivre}
\begin{blocLivre}{cor-11-8}
\begin{corrige}{8}{Une aire sous une parabole}
\demarche{L’aire du carré unité vaut 1, donc la fréquence estime directement l’aire.}
\begin{questions}
\item Un point est uniforme dans un carré d’aire 1 ; la probabilité d’entrer dans la zone est donc l’aire de cette zone.
La fréquence observée fournit l’estimation $\res{137/400=0{,}3425}$.
\item L’erreur absolue par rapport à l’aire fournie est
$|137/400-1/3|=|(411-400)/1200|=\res{11/1200\approx0{,}009167}$.
La fraction reste exacte avant l’arrondi.
\item Après avoir tiré indépendamment \code{x=random()} et \code{y=random()}, le test est \code{y <= x*x}.
Les tirages appartiennent déjà au carré unité : le test supplémentaire porte uniquement sur la zone sous la parabole.
\end{questions}

\label{figlivre-r11-08-zone-cor}
\figureLivre[corrige]{r11-08-zone}

\end{corrige}
\end{blocLivre}
\begin{blocLivre}{cor-11-9}
\begin{corrige}{9}{Estimer le nombre pi}
\demarche{Le carré unité contient un quart du disque, ce qui explique le facteur 4.}
\begin{questions}
\item La fréquence dans le quart de disque est $786/1000=0{,}786$.
On estime donc $\pi$ par $4(0{,}786)=\res{3{,}144}$.
L’erreur absolue $|3{,}144-\pi|\approx0{,}002407346$ s’arrondit à $\res{0{,}00241}$.
\item Le disque de rayon 1 a pour aire $\pi$ ; le quart de disque a pour aire $\pi/4$.
Comme le carré d’échantillonnage a pour aire 1, sa fréquence estime $\pi/4$, d’où le facteur 4.
\item Un nouveau tirage peut donner un autre nombre de points dans la zone.
La taille identique ne supprime pas les fluctuations ; elle ne force pas une répétition exacte du résultat.
\end{questions}

\label{figlivre-r11-09-zone-cor}
\figureLivre[corrige]{r11-09-zone}

\end{corrige}
\end{blocLivre}
\begin{blocLivre}{cor-11-10}
\begin{corrige}{10}{Comparer des moyennes d'échantillons}
\demarche{Comparer chaque écart à la même tolérance avant de compter les moyennes retenues.}
\begin{questions}
\item $E(X)=0(1/2)+2(1/2)=1$ et $E(X^2)=0+4(1/2)=2$.
Ainsi $V(X)=2-1^2=1$, donc $\res{\mu=1,\ \sigma=1}$.
\item Les sommes des quatre échantillons sont 4, 8, 0 et 6.
Après division par 5, les moyennes valent $\res{0{,}8;\ 1{,}6;\ 0;\ 1{,}2}$.
\item Le seuil commun est $2/\sqrt5\approx0{,}8944$.
Les écarts à 1 sont $0{,}2$, $0{,}6$, 1 et $0{,}2$ : le troisième seul dépasse le seuil.
La proportion retenue est $3/4=\res{75\,\%}$.
Ce résultat décrit ces quatre échantillons, sans imposer ce pourcentage à d’autres tirages.
\end{questions}

\label{figlivre-r11-10-critere-cor}
\figureLivre[corrige]{r11-10-critere}

\end{corrige}
\end{blocLivre}
\begin{blocLivre}{cor-11-11}
\begin{corrige}{11}{Effet de la taille d'échantillon}
\demarche{Le seuil diminue comme l’inverse de la racine de la taille.}
\begin{questions}
\item Pour $n=16$, le seuil vaut $4/4=\res{1}$ ; pour $n=64$, $4/8=\res{1/2}$.
\item $|2{,}2-1|=1{,}2>1$ : non. $|1{,}4-1|=0{,}4\leq0{,}5$ : oui.
\item Le seuil varie comme $1/\sqrt n$. Pour le diviser par 3, il faut multiplier $n$ par $\res{9}$.
\end{questions}
\end{corrige}
\end{blocLivre}
\begin{blocLivre}{cor-11-12}
\begin{corrige}{12}{Simuler une loi à trois valeurs}
\demarche{Partager un seul tirage uniforme entre des intervalles de longueurs prescrites.}
\begin{questions}
\item $\mu=-0{,}5+0{,}6+1=\res{1{,}1}$. $E(X^2)=0{,}5+1{,}2+5=6{,}7$, donc $V(X)=\res{5{,}49}$ et $\sigma=\res{\sqrt{549}/10\approx2{,}343}$.
\item On partage $[0;1[$ aux seuils 0,5 et 0,8. Le code ci-dessous utilise le même tirage pour les deux tests.
\item Pour $n\geq1$, on moyenne une liste de $n$ simulations.
\item Répéter les expériences à chaque taille, puis étudier les écarts absolus à 1,1. Une moyenne plus grande ou plus petite que 1,1 reste possible.
\end{questions}
\begin{python}
from random import random
def tirage():
    u = random()
    if u < 0.5:
        return -1
    if u < 0.8:
        return 2
    return 5
def moyenne(n):
    return sum(tirage() for k in range(n))/n
\end{python}

Le deuxième test n’est atteint que si le premier a échoué : la valeur 2 correspond donc à $0{,}5\leq u<0{,}8$, un intervalle de longueur $0{,}3$. Un deuxième tirage aléatoire modifierait le modèle.

\label{figlivre-r11-12-partition-cor}
\figureLivre[corrige]{r11-12-partition}

\end{corrige}
\end{blocLivre}
\begin{blocLivre}{cor-11-13}
\begin{corrige}{13}{Répéter cinq cents échantillons}
\demarche{Simuler de nouveaux échantillons et distinguer le compteur de la proportion finale.}
\begin{questions}
\item $\mu=0{,}4$, $\sigma=\sqrt{0{,}24}$ et le seuil vaut $\res{2\sqrt{0{,}24}/5\approx0{,}19596}$.
\item Chaque moyenne est obtenue à partir de 25 nouveaux tirages. Le compteur est ensuite divisé par 500.
\item Non : il s’agit d’une proportion expérimentale dépendant du modèle et des tirages, pas d’une valeur universelle imposée.
\end{questions}
\begin{python}
from random import Random
from math import sqrt
rng = Random(2026)
compteur = 0
for j in range(500):
    m = sum(rng.random() < 0.4 for k in range(25))/25
    if abs(m - 0.4) <= 2*sqrt(0.24)/5:
        compteur = compteur + 1
proportion = compteur/500
\end{python}

Avec la graine 2026 et ce programme, on obtient 470 moyennes retenues sur 500, soit une proportion $0{,}94$. La graine rend cet exemple reproductible ; elle ne transforme pas 94\,\% en règle générale.

\end{corrige}
\end{blocLivre}
\begin{blocLivre}{cor-11-14}
\begin{corrige}{14}{Comparer deux groupes et leur réunion}
\demarche{Comparer les écarts à un centre commun pour mesurer ce que la moyenne masque.}
\begin{questions}
\item Les deux moyennes valent $\res{5}$. Pour A, $v_A=(1+0+0+1)/4=\res{1/2}$ ; pour B, $v_B=(16+0+0+16)/4=\res{8}$.
\item La moyenne reste 5. La somme des carrés des écarts vaut $2+32=34$, donc la variance est $\res{34/8=17/4}$.
\item La série ordonnée est $1,4,5,5,5,5,6,9$. Médiane 5 ; $Q_1$ au rang 2 vaut 4 ; $Q_3$ au rang 6 vaut 5.
\item Le groupe B est beaucoup plus dispersé, bien que sa moyenne soit identique. Les valeurs extrêmes sont invisibles si on ne rapporte que la moyenne.
\end{questions}
\label{figlivre-r11-14-groupes-cor}
\figureLivre[corrige]{r11-14-groupes}

\end{corrige}
\end{blocLivre}
\begin{blocLivre}{cor-11-15}
\begin{corrige}{15}{Fréquences des lettres d'un texte}
\demarche{Compter chaque lettre puis diviser par la longueur totale du texte.}
\begin{questions}
\item Dans BANANE : B et E ont fréquence $1/6$, A et N fréquence $1/3$. Dans ANANAS : A $1/2$, N $1/3$, S $1/6$. Les autres fréquences sont nulles.
\item Pour un texte non vide, on peut écrire \code{[texte.count(c)/len(texte) for c in alphabet]}, où \code{alphabet} est la chaîne des 26 lettres majuscules. La somme vaut 1 si tous les caractères sont des lettres de cet alphabet.
\item Il y a deux A et un E sur six positions : $P(V=1)=P(V=0)=1/2$, donc $\res{E(V)=1/2}$ et $\res{V(V)=1/4}$.
\end{questions}
\end{corrige}
\end{blocLivre}

\clearpage
% !TEX root = ../main.tex
\corrigesChapitre{12}{Algorithmique et programmation}
\corrigesSouples

\begin{blocLivre}{cor-12-1}
\begin{corrige}{1}{Suivre des affectations}
\demarche{Construire une table de trace et utiliser les valeurs mises à jour à chaque ligne.}
Après les trois actualisations, les couples $(a,b)$ sont $(8,5)$, $(8,3)$ puis $\res{(5,3)}$. Le programme échange les deux valeurs initiales. Il faut utiliser la valeur de \code{a} déjà actualisée dans les deux dernières lignes.
\label{figlivre-r12-01-trace-cor}
\figureLivre[corrige]{r12-01-trace}

\end{corrige}
\end{blocLivre}
\begin{blocLivre}{cor-12-2}
\begin{corrige}{2}{Tester une condition}
\demarche{Tester la parité avant de choisir la branche de la fonction.}
Le reste \code{n \% 2} vaut 0 si $n$ est pair. On obtient $\res{22}$ pour 7, $\res{4}$ pour 8 et $\res{0}$ pour 0. \code{//} désigne la division entière, ici exacte pour les nombres pairs.
\end{corrige}
\end{blocLivre}
\begin{blocLivre}{cor-12-3}
\begin{corrige}{3}{Définir une fonction}
\demarche{Traduire la formule avec les opérations de Python et tester quelques entrées.}
On traduit directement les opérations, avec \code{**} pour la puissance.
\begin{python}
def f(x):
    return x**2 - 3*x + 2
\end{python}
Les résultats sont $\res{2}$, $\res{0}$ et $\res{12}$. Le nom $x$ est un paramètre local ; un autre nom conviendrait si on le remplaçait partout de façon cohérente.
\end{corrige}
\end{blocLivre}
\begin{blocLivre}{cor-12-4}
\begin{corrige}{4}{Modifier une liste}
\demarche{Suivre les indices et les décalages provoqués par la suppression.}
Le tableau donne les quatre états de la liste.
\code{append(2)} ajoute une valeur à la fin, sans changer les anciennes positions.
Puis \code{L[1] = 5} modifie le deuxième élément, puisque les indices commencent à 0.
Enfin \code{L.pop(0)} retire le premier élément ; les autres sont décalés d’une position.
La liste finale est \code{[5, 4, 2]}, d’où $\res{\operatorname{len}(L)=3}$ et $\res{\operatorname{sum}(L)=11}$.

\label{figlivre-r12-04-trace-cor}
\figureLivre[corrige]{r12-04-trace}

\end{corrige}
\end{blocLivre}
\begin{blocLivre}{cor-12-5}
\begin{corrige}{5}{Comprendre une liste en compréhension}
\demarche{Créer la liste initiale avant d’appliquer le filtre à ses valeurs.}
\code{range(6)} donne 0 à 5 : \code{L = [0, 1, 4, 9, 16, 25]}, de longueur 6. Le filtre garde les valeurs paires : \code{P = [0, 4, 16]}, de longueur 3.
\end{corrige}
\end{blocLivre}
\begin{blocLivre}{cor-12-6}
\begin{corrige}{6}{Calculer une factorielle}
\demarche{Initialiser un produit à l’élément neutre 1, y compris pour la boucle vide.}
On initialise le produit à 1, puis on multiplie par les entiers de 1 à $n$.
\begin{python}
def factorielle(n):
    produit = 1
    for k in range(1, n+1):
        produit = produit*k
    return produit
\end{python}
Les résultats sont $\res{1;1;120}$. Pour $n=0$, aucun passage n'a lieu.
Une initialisation à 0 rendrait tous les produits nuls.

Après le passage utilisant $k$, la variable contient le produit de 1 à $k$. Pour $n=0$, la boucle est vide et la valeur initiale 1 reste la réponse correcte.

\end{corrige}
\end{blocLivre}
\begin{blocLivre}{cor-12-7}
\begin{corrige}{7}{Accumuler une somme}
\demarche{Initialiser une somme à zéro et contrôler la borne exclue de range.}
Le total part de 0 ; la borne supérieure de \code{range} est exclue.
\begin{python}
def somme(n):
    total = 0
    for k in range(1, n+1):
        total = total + k
    return total
\end{python}
On obtient $\res{0}$ pour $n=0$ et $\res{55}$ pour $n=10$, en accord avec la formule.
Ces tests contrôlent des exemples ; la démonstration générale est donnée dans les annexes.
\end{corrige}
\end{blocLivre}
\begin{blocLivre}{cor-12-8}
\begin{corrige}{8}{Trouver un premier seuil}
\demarche{Faire évoluer ensemble le rang et le terme jusqu’au premier seuil.}
\begin{python}
def seuil():
    n = 0
    u = 100
    while u < 500:
        u = u + 30
        n = n + 1
    return n, u
\end{python}
$u_n=100+30n$ donne $n\geq40/3$, donc le premier entier est $\res{14}$.
$u_{13}=490<500$ et $u_{14}=520$. Si le seuil est déjà atteint, la boucle ne s'exécute pas et le rang renvoyé est 0.

Après chaque passage, $n$ compte les transitions effectuées et $u$ vaut le terme de ce rang. Le test est effectué avant le passage : un seuil déjà atteint laisse directement $n=0$.

\end{corrige}
\end{blocLivre}
\begin{blocLivre}{cor-12-9}
\begin{corrige}{9}{Additionner les termes d'une suite}
\demarche{Ajouter le terme courant au total avant de préparer le suivant.}
\begin{python}
def somme_geometrique(n):
    u = 2
    total = 0
    for k in range(n):
        total = total + u
        u = 3*u
    return total
\end{python}
Pour $n=4$, le résultat vaut $2+6+18+54=\res{80}$.
La formule est $2(3^n-1)/(3-1)=\res{3^n-1}$, qui donne aussi 0 pour la somme vide $n=0$.

Le tableau montre qu’après quatre passages, $u=162$ est déjà préparé pour le passage suivant, mais n’a pas été ajouté. Le total 80 contient exactement les quatre termes demandés.

\label{figlivre-r12-09-trace-cor}
\figureLivre[corrige]{r12-09-trace}

\end{corrige}
\end{blocLivre}
\begin{blocLivre}{cor-12-10}
\begin{corrige}{10}{Compter les lettres sans autre collection}
\demarche{Compter les correspondances exactes en parcourant le texte une seule fois.}
\begin{python}
def frequence(texte, lettre):
    compteur = 0
    for c in texte:
        if c == lettre:
            compteur = compteur + 1
    return compteur/len(texte)
\end{python}
Les résultats sont $\res{1/3}$ pour A et $\res{1/6}$ pour B.
Les lettres majuscules et minuscules sont distinctes ; les hypothèses du texte doivent rester cohérentes.
\end{corrige}
\end{blocLivre}
\begin{blocLivre}{cor-12-11}
\begin{corrige}{11}{Calculer les moments d'une loi}
\demarche{Accumuler séparément le premier et le deuxième moment, puis soustraire le carré de la moyenne.}
\begin{python}
def moments(valeurs, probas):
    mu = 0
    carre = 0
    for k in range(len(valeurs)):
        mu = mu + valeurs[k]*probas[k]
        carre = carre + valeurs[k]**2*probas[k]
    return mu, carre - mu**2
\end{python}
Le résultat théorique est $\res{(0{,}4;1{,}24)}$.
Les deux listes doivent être non vides et de même longueur : chaque valeur est associée à la probabilité de même indice. On contrôle que les probabilités sont positives ou nulles et de somme 1.
La formule employée est celle de König-Huygens.

Ici les accumulateurs finaux valent $\mu=0{,}4$ et $\texttt{carre}=1{,}4$. Le retour calcule $1{,}4-0{,}4^2=1{,}24$. Une variance nulle est possible pour une loi constante.

\label{figlivre-r12-11-moments-cor}
\figureLivre[corrige]{r12-11-moments}

\end{corrige}
\end{blocLivre}
\begin{blocLivre}{cor-12-12}
\begin{corrige}{12}{Construire la suite de Fibonacci}
\demarche{Traiter les petits indices avant de prolonger la liste de Fibonacci.}
\begin{python}
def fibonacci(n):
    if n == 0:
        return [0]
    L = [0, 1]
    for k in range(2, n+1):
        L.append(L[-1] + L[-2])
    return L
\end{python}
Pour $n=7$, la liste est \code{[0, 1, 1, 2, 3, 5, 8, 13]}.
En poursuivant, $F_{16}=987$ et $F_{17}=1597$ : le premier rang cherché est $\res{17}$.
Les termes sont croissants à partir de $F_1$, car chaque terme est augmenté d'un terme positif ou nul.

La liste est initialisée avec les deux premiers termes. Chaque passage ajoute la somme des deux derniers sans les modifier ; le cas $n=0$ est traité avant cette initialisation.

\end{corrige}
\end{blocLivre}
\begin{blocLivre}{cor-12-13}
\begin{corrige}{13}{Expérimenter avec Syracuse}
\demarche{Borner le nombre de passages garantit un arrêt informatique.}
\begin{python}
def syracuse(n, max_etapes=1000):
    L = [n]
    for k in range(max_etapes):
        if n == 1:
            break
        if n % 2 == 0:
            n = n//2
        else:
            n = 3*n + 1
        L.append(n)
    return L
\end{python}
Pour 6 : \code{[6, 3, 10, 5, 16, 8, 4, 2, 1]}, soit huit transitions.
La boucle \code{for} a un nombre borné de passages ; l'arrêt informatique est donc garanti.
Cet exemple ne démontre pas un comportement pour tous les entiers positifs.

Le mot-clé \code{break} arrête plus tôt si 1 est atteint. Sinon, la boucle s’arrête après \code{max\_etapes} passages : la liste peut alors se terminer par une autre valeur.

\end{corrige}
\end{blocLivre}
\begin{blocLivre}{cor-12-14}
\begin{corrige}{14}{Une approximation par Newton}
\demarche{Remplacer la courbe par sa tangente, puis actualiser la moyenne de deux nombres positifs.}
Pour $a>0$, $f'(a)=2a$ et la tangente a pour équation $y=2a(x-a)+a^2-2$.
Elle coupe l'axe des abscisses lorsque $2a(x-a)=2-a^2$, soit
$x=a+\dfrac{2-a^2}{2a}=\res{\dfrac{a^2+2}{2a}=\dfrac12\left(a+\dfrac2a\right)}$.
La méthode de Newton remplace ainsi $a$ par l'abscisse de ce point.
\begin{python}
def newton(n):
    x = 1.0
    for k in range(n):
        x = (x + 2/x)/2
    return x
\end{python}
$x_1=\res{3/2}$ et $x_2=\res{17/12}$. Pour $n=5$, on obtient environ
$1{,}414213562373095$ ; son carré est très proche de 2.
À chaque étape, si $x$ est positif, $2/x$ l'est aussi et leur moyenne est positive.
Cela justifie que les calculs successifs sont définis ; ce n'est pas à lui seul une preuve de convergence.

Le cas $n=0$ renvoie 1 sans passage. Pour $n=1$, la mise à jour donne $(1+2)/2=3/2$ ; pour $n=2$, elle donne $(3/2+4/3)/2=17/12$.

\end{corrige}
\end{blocLivre}
\begin{blocLivre}{cor-12-15}
\begin{corrige}{15}{Un mini-projet Monte-Carlo}
\demarche{Réutiliser les coordonnées, avec deux tests et deux compteurs distincts.}
\begin{python}
from random import Random
def monte_carlo(n, graine=2026):
    rng = Random(graine)
    parabole = 0
    disque = 0
    for k in range(n):
        x, y = rng.random(), rng.random()
        if y <= x*x:
            parabole = parabole + 1
        if x*x + y*y <= 1:
            disque = disque + 1
    return parabole/n, 4*disque/n
\end{python}
Chaque point est uniforme dans le carré et peut être testé dans les deux zones : les deux estimations
peuvent dépendre l'une de l'autre, ce qui ne les empêche pas d'être calculées.
L'aire estimée doit appartenir à $[0;1]$ et l'estimation de $\pi$ à $[0;4]$.
Même graine et même taille doivent reproduire le même résultat.

Les deux tests sont deux instructions \code{if} distinctes : un même point peut appartenir aux deux zones et doit alors incrémenter les deux compteurs. Une instruction \code{elif} ne conviendrait pas.

\label{figlivre-r12-15-zones-cor}
\figureLivre[corrige]{r12-15-zones}

\end{corrige}
\end{blocLivre}

\clearpage
% !TEX root = ../main.tex
\corrigesChapitre{13}{Ensembles et raisonnement}
\corrigesSouples

\begin{blocLivre}{cor-13-1}
\begin{corrige}{1}{Appartenance et opérations}
\demarche{Séparer les éléments communs, les éléments de la réunion et les éléments extérieurs.}
\begin{questions}
\item $\res{A\cap B=\{2\}}$, $\res{A\cup B=\{1;2;3;4;6\}}$ et $\res{E\setminus A=\{1;3;5\}}$.
\item On écrit $\res{2\in A}$ et $\res{\{2\}\subseteq A}$. Le nombre 2 est un élément, le singleton $\{2\}$ est un sous-ensemble.
\end{questions}
\label{figlivre-r13-01-venn-cor}
\figureLivre[corrige]{r13-01-venn}

\end{corrige}
\end{blocLivre}
\begin{blocLivre}{cor-13-2}
\begin{corrige}{2}{Réunir des intervalles}
\demarche{Comparer les bornes et vérifier, pour chacune, si elle est incluse.}
\begin{questions}
\item Les valeurs communes donnent $\res{]1;3]}$. La réunion couvre sans trou $\res{[-2;5[}$.
\item $\res{\R\setminus A=\mathopen{]}-\infty;-2[\cup]3;+\infty[}$. Les bornes $-2$ et 3 sont exclues du complément car elles appartiennent à A.
\end{questions}
\label{figlivre-r13-02-intervalles-cor}
\figureLivre[corrige]{r13-02-intervalles}

\end{corrige}
\end{blocLivre}
\begin{blocLivre}{cor-13-3}
\begin{corrige}{3}{Construire des couples}
\demarche{Respecter l’ordre des deux coordonnées d’un couple.}
\begin{questions}
\item $A\times B=\{(a;0),(a;1),(a;2),(b;0),(b;1),(b;2)\}$, de cardinal $\res{6}$.
\item $(1;a)$ n’appartient pas à $A\times B$, car sa première coordonnée n’est pas dans A. Il appartient à $B\times A$.
\end{questions}
\label{figlivre-r13-03-couples-cor}
\figureLivre[corrige]{r13-03-couples}

\end{corrige}
\end{blocLivre}
\begin{blocLivre}{cor-13-4}
\begin{corrige}{4}{Trouver un contre-exemple}
\demarche{Un seul exemple qui contredit une affirmation universelle suffit à la réfuter.}
\begin{questions}
\item Pour $x=1/2$, $x^2=1/4<1/2$ : l’affirmation est fausse.
\item $x=-2$ vérifie $x^2=4$, mais pas $x=2$. Les deux racines de $x^2=4$ sont $-2$ et 2.
\end{questions}
\end{corrige}
\end{blocLivre}
\begin{blocLivre}{cor-13-5}
\begin{corrige}{5}{Nier une phrase quantifiée}
\demarche{Échanger les quantificateurs et nier la propriété avec la bonne inégalité.}
\begin{questions}
\item La négation est « il existe un réel $x$ tel que $x^2<0$ ». La phrase initiale est vraie, sa négation fausse.
\item La négation est « pour tout réel $x$, $x^2\neq-1$ ». La phrase initiale est fausse, sa négation vraie.
\end{questions}
\end{corrige}
\end{blocLivre}
\begin{blocLivre}{cor-13-6}
\begin{corrige}{6}{Réciproque et contraposée}
\demarche{Distinguer l’implication, sa réciproque et sa contraposée.}
\begin{questions}
\item Pour $x>2$, $x+2>0$ et $x-2>0$, donc $x^2-4=(x-2)(x+2)>0$.
\item La réciproque est « si $x^2>4$, alors $x>2$ ». Elle est fausse pour $x=-3$.
\item La contraposée est « si $x^2\leq4$, alors $x\leq2$ ». Elle est vraie, car équivalente à l’implication initiale.
\end{questions}
\end{corrige}
\end{blocLivre}
\begin{blocLivre}{cor-13-7}
\begin{corrige}{7}{Nécessaire ou suffisante}
\demarche{Identifier quel énoncé suffit à garantir l’autre.}
\begin{questions}
\item $P\Rightarrow Q$ : si $n=4k$, alors $n=2(2k)$ est pair.
\item P est \textbf{suffisante} pour Q ; Q est \textbf{nécessaire} pour P.
\item $n=6$ est pair sans être multiple de 4. Q n’est donc pas suffisante pour P.
\end{questions}
\end{corrige}
\end{blocLivre}
\begin{blocLivre}{cor-13-8}
\begin{corrige}{8}{Variable, inconnue et paramètre}
\demarche{Le rôle de la lettre dépend de la tâche : décrire, résoudre ou choisir une fonction.}
\begin{questions}
\item $x$ est une variable ; l’identité est vraie pour tout réel $x$.
\item $x$ est une inconnue dont on cherche les valeurs : $\res{x=-1\text{ ou }1}$.
\item Pour une fonction fixée de la famille, $x$ est la variable et $m$ un paramètre constant. Faire varier $m$ change la fonction étudiée.
\end{questions}
\end{corrige}
\end{blocLivre}
\begin{blocLivre}{cor-13-9}
\begin{corrige}{9}{Complémentaires et connecteurs}
\demarche{Passer de la négation d’un ou à la conjonction de deux négations.}
\begin{questions}
\item $A\cup B=\{1;2;3;4;6;8\}$, donc le complémentaire est $\res{\{5;7\}}$.
\item $E\setminus A=\{1;3;5;7\}$ et $E\setminus B=\{5;6;7;8\}$ ; l’intersection est aussi $\res{\{5;7\}}$.
\item Ne pas être « dans A ou dans B » signifie être « hors de A et hors de B ». C’est une loi de De Morgan.
\end{questions}
\label{figlivre-r13-09-venn-cor}
\figureLivre[corrige]{r13-09-venn}

\end{corrige}
\end{blocLivre}
\begin{blocLivre}{cor-13-10}
\begin{corrige}{10}{Lire une quantification implicite}
\demarche{Pour réfuter une croissance sur un domaine, exhiber deux points dans le mauvais ordre.}
\begin{questions}
\item Il existe $a,b\in I$ avec $a<b$ et $f(a)>f(b)$.
\item Prendre $a=-2$ et $b=-1$ : $a<b$ mais $f(a)=4>1=f(b)$.
\item Pour $0\leq a<b$, $b^2-a^2=(b-a)(b+a)>0$. Donc $f(a)<f(b)$ : la fonction y est même strictement croissante.
\end{questions}
\end{corrige}
\end{blocLivre}
\begin{blocLivre}{cor-13-11}
\begin{corrige}{11}{Démontrer par disjonction des cas}
\demarche{Traiter tous les cas possibles, puis formuler une conclusion commune.}
\begin{questions}
\item Si $n$ est pair, $n=2k$ et le produit est multiple de 2. Si $n$ est impair, $n+1$ est pair, donc le produit est encore multiple de 2. Les deux cas couvrent tous les entiers naturels.
\item Si $x\geq0$, $|x|=x$. Si $x<0$, $|x|=-x>x$. Ainsi l’inégalité est vraie dans les deux cas.
\end{questions}
\end{corrige}
\end{blocLivre}
\begin{blocLivre}{cor-13-12}
\begin{corrige}{12}{Une preuve par contraposée}
\demarche{Prouver la contraposée revient à prouver l’implication de départ.}
\begin{questions}
\item La contraposée est : « si un entier est pair, alors son carré est pair ».
\item Si $n=2k$, $n^2=4k^2=2(2k^2)$ est pair. La contraposée étant vraie, l’implication initiale l’est aussi.
\item Si $n=2k+1$, $n^2=4k^2+4k+1=2(2k^2+2k)+1$ est impair. Ainsi $\res{n^2\text{ impair}\iff n\text{ impair}}$.
\end{questions}
\end{corrige}
\end{blocLivre}
\begin{blocLivre}{cor-13-13}
\begin{corrige}{13}{Une preuve par l'absurde}
\demarche{Supposer une fraction irréductible, puis obtenir un facteur commun impossible.}
\begin{questions}
\item En élevant $\sqrt2=p/q$ au carré, on obtient $p^2=2q^2$ : le carré $p^2$ est pair.
Si $p$ était impair, son carré serait impair d’après la propriété précédente.
Donc $p$ est pair et s’écrit $p=2k$ pour un entier positif $k$.
\item En remplaçant $p$, $4k^2=2q^2$, donc $q^2=2k^2$.
Le même argument impose que $q$ soit pair.
\item Alors 2 divise simultanément $p$ et $q$.
Cela contredit le choix initial d’une fraction irréductible.
L’hypothèse $\sqrt2=p/q$ est donc impossible, et $\res{\sqrt2\notin\mathbb Q}$.
\end{questions}

\end{corrige}
\end{blocLivre}
\begin{blocLivre}{cor-13-14}
\begin{corrige}{14}{Une transformation sous condition}
\demarche{Conserver la valeur interdite pendant toutes les transformations de l’équation.}
\begin{questions}
\item Le domaine est $\R\setminus\{1\}$. Sur ce domaine, $(x^2-1)/(x-1)=x+1$.
\item $x+1=2$ donne $x=1$, qui est interdit. L’ensemble des solutions est $\res{\varnothing}$.
\item La multiplication conduit à $(x-1)^2=0$, mais cette équation n’est équivalente à l’originale que pour $x\neq1$. Oublier cette condition ajoute une fausse solution.
\end{questions}
\end{corrige}
\end{blocLivre}
\begin{blocLivre}{cor-13-15}
\begin{corrige}{15}{Quantifier une famille de trinômes}
\demarche{Fixer le paramètre et comparer le minimum du trinôme à zéro.}
\begin{questions}
\item $x^2-2mx+1=(x-m)^2+1-m^2$.
Pour chaque $m$ fixé, le carré est minimal en $x=m$ : le minimum de $P_m$ vaut $1-m^2$.
\item La condition pour tout réel $x$ impose que ce minimum soit positif ou nul :
$1-m^2\geq0\iff m^2\leq1\iff\res{m\in[-1;1]}$.
\item Pour être strictement positif en tout réel, le minimum doit être strictement positif :
$1-m^2>0\iff\res{m\in\mathopen{]}-1;1[}$.
\item La négation est : il existe un réel $x$ tel que $P_m(x)<0$.
Elle se produit pour $\res{m<-1\text{ ou }m>1}$ : le choix explicite $x=m$ donne alors $P_m(m)=1-m^2<0$.
\item Pour $m=1$ ou $m=-1$, le minimum vaut 0 et est réellement atteint en $x=m$.
Ces deux paramètres satisfont la condition large, mais pas la condition stricte.
L’explication orale doit donc distinguer positivité du minimum et possibilité d’une valeur nulle.
\end{questions}

\end{corrige}
\end{blocLivre}

\clearpage
% !TEX root = ../main.tex
\corrigesChapitre{14}{Automatismes et synthèses}
\corrigesSouples

\begin{blocLivre}{cor-14-1}
\begin{corrige}{1}{Évolutions et retour au départ}
\demarche{Composer les coefficients multiplicateurs plutôt que les taux.}
\begin{questions}
\item $80\times1{,}25=100$ euros, puis $100\times0{,}8=\res{80}$ euros. Le coefficient global vaut 1, donc le taux global est $\res{0\,\%}$.
\item Le coefficient réciproque est $1/0{,}8=1{,}25$, soit une hausse de $\res{25\,\%}$.
\end{questions}
\end{corrige}
\end{blocLivre}
\begin{blocLivre}{cor-14-2}
\begin{corrige}{2}{Retrouver une valeur et un taux}
\demarche{Retrouver la valeur initiale par division et calculer les taux relativement à cette valeur.}
\begin{questions}
\item $0{,}7P=84$, donc $\res{P=120\text{ euros}}$.
\item $(54-45)/45=9/45=\res{20\,\%}$.
\item $1{,}1\times1{,}2=1{,}32$, donc le taux global est $\res{32\,\%}$.
\end{questions}
\end{corrige}
\end{blocLivre}
\begin{blocLivre}{cor-14-3}
\begin{corrige}{3}{Développer, factoriser, résoudre}
\demarche{Reconnaître l’identité remarquable ou le produit nul adapté.}
\begin{questions}
\item $\res{4x^2-12x+9}$.
\item Différence de carrés : $\res{(3x-5)(3x+5)}$.
\item Un facteur doit être nul : $\res{x=2\text{ ou }x=-1}$.
\end{questions}
\end{corrige}
\end{blocLivre}
\begin{blocLivre}{cor-14-4}
\begin{corrige}{4}{Signe et intervalles}
\demarche{Distinguer le signe d’un facteur et celui d’un produit.}
\begin{questions}
\item On obtient $\res{x\leq3/2}$ en divisant par le coefficient négatif $-2$ et en inversant le sens de l’inégalité.
\item Le produit est négatif strictement entre les racines : $\res{x\in\mathopen{]}-2;3[}$.
\item Par exemple $\res{x=0}$ satisfait les deux conditions. L’ensemble commun est $\mathopen{]}-2;3/2]$.
\end{questions}
\end{corrige}
\end{blocLivre}
\begin{blocLivre}{cor-14-5}
\begin{corrige}{5}{Équations de droites}
\demarche{Calculer la pente avec deux points, puis utiliser l’un d’eux pour l’ordonnée à l’origine.}
\begin{questions}
\item La pente vaut $(7-3)/(3-1)=2$. Avec A, $3=2+p$, donc $\res{y=2x+1}$.
\item La pente est $-1/2$, l’ordonnée à l’origine 4. $y=0$ donne $x=8$, donc l’intersection est $\res{(8;0)}$.
\end{questions}
\label{figlivre-r14-05-droites-cor}
\figureLivre[corrige]{r14-05-droites}

\end{corrige}
\end{blocLivre}
\begin{blocLivre}{cor-14-6}
\begin{corrige}{6}{Lire une courbe par morceaux}
\demarche{Repérer les intersections avec la droite horizontale et les portions au-dessus.}
\begin{questions}
\item Les intersections avec $y=1$ sont $\res{x=-2;1;4}$. L’inéquation donne $\res{\{-2\}\cup[1;4]}$. La valeur 1 sur le segment croissant se vérifie avec $y=2x-1$.
\item Décroissance sur $[-2;0]$, croissance sur $[0;2]$, décroissance sur $[2;4]$. Minimum $-1$ en 0, maximum 3 en 2.
\end{questions}
\label{figlivre-r14-06-lecture-cor}
\figureLivre[corrige]{r14-06-lecture}

\end{corrige}
\end{blocLivre}
\begin{blocLivre}{cor-14-7}
\begin{corrige}{7}{Des fréquences conditionnelles}
\demarche{La population de référence change entre fréquence globale et fréquence conditionnelle.}
\begin{questions}
\item Il y a $20+30=50$ utilisateurs : fréquence $\res{1/2}$.
\item Parmi les 50 plus jeunes, la fréquence est $20/50=\res{2/5}$.
\item La probabilité conditionnelle $2/5$ diffère de la probabilité globale $1/2$ : pas d’indépendance dans ce modèle.
\end{questions}
\label{figlivre-r14-07-cor}
\figureLivre[corrige]{r14-07}

\end{corrige}
\end{blocLivre}
\begin{blocLivre}{cor-14-8}
\begin{corrige}{8}{Réactiver la trigonométrie}
\demarche{Retirer un tour pour comparer les points-images, et utiliser les radians pour les arcs.}
\begin{questions}
\item $5\pi/3=2\pi-\pi/3$ : $\res{(1/2;-\sqrt3/2)}$.
\item La différence entre les deux angles vaut $2\pi$ : les coordonnées sont les mêmes.
\item $L=R\theta=3\times2\pi/3=\res{2\pi}$.
\end{questions}
\label{figlivre-r14-08-trigo-cor}
\figureLivre[corrige]{r14-08-trigo}

\end{corrige}
\end{blocLivre}
\begin{blocLivre}{cor-14-9}
\begin{corrige}{9}{Relier dérivée et extremums}
\demarche{Factoriser la dérivée et comparer tous les candidats, y compris les bornes.}
\begin{questions}
\item $f\prime(x)=3x^2-12x+9=\res{3(x-1)(x-3)}$, de signes $+$, $-$, $+$ successivement.
\item $f$ croît de $f(0)=0$ à $f(1)=4$, décroît à $f(3)=0$, puis croît à $f(4)=4$.
\item Minimum $\res{0}$ en 0 et 3 ; maximum $\res{4}$ en 1 et 4. Il faut comparer aussi les bornes.
\end{questions}
\label{figlivre-r14-09-var-cor}
\figureLivre[corrige]{r14-09-var}

\end{corrige}
\end{blocLivre}
\begin{blocLivre}{cor-14-10}
\begin{corrige}{10}{Un diagramme circulaire}
\demarche{Convertir les fréquences en angles tout en conservant le total de 360 degrés.}
\begin{questions}
\item Les fréquences sont $\res{1/2;3/10;1/5}$. Les angles sont $\res{180^\circ;108^\circ;72^\circ}$ et totalisent $360^\circ$.
\item $40-24=\res{16\text{ voix}}$.
\item Non : les noms sont des catégories qualitatives, sans valeur numérique à moyenner. On peut comparer leurs effectifs ou leurs fréquences.
\end{questions}
\label{figlivre-r14-10-secteurs-cor}
\figureLivre[corrige]{r14-10-secteurs}

\end{corrige}
\end{blocLivre}
\begin{blocLivre}{cor-14-11}
\begin{corrige}{11}{Une exponentielle dans une étude de fonction}
\demarche{Factoriser la dérivée par une exponentielle positive.}
\begin{questions}
\item Par injectivité, $2x=x+3$, donc $\res{x=3}$.
\item $g\prime(x)=2e^{-x}-(2x-1)e^{-x}=\res{(3-2x)e^{-x}}$.
\item La fonction croît jusqu’à $3/2$, puis décroît. Le maximum vaut $\res{2e^{-3/2}}$. Aux bornes, $g(0)=-1$ et $g(3)=5e^{-3}>0$ : minimum $\res{-1}$ en 0.
\end{questions}
\label{figlivre-r14-11-var-cor}
\figureLivre[corrige]{r14-11-var}

\end{corrige}
\end{blocLivre}
\begin{blocLivre}{cor-14-12}
\begin{corrige}{12}{Synthèse : une boîte à dimensions liées}
\demarche{Relier le schéma du solide à son volume, puis étudier un signe sur le domaine géométrique.}
\begin{questions}
\item La base carrée a pour aire $x^2$ et la hauteur vaut $6-x$.
Donc $V(x)=x^2(6-x)=-x^3+6x^2$ et $V'(x)=-3x^2+12x=\res{3x(4-x)}$.
\item Sur $]0;6[$, le facteur $3x$ est positif. La dérivée a le signe de $4-x$ :
le maximum est atteint en $x=4$, avec $V(4)=16\times2=\res{32\text{ L}}$.
La base mesure $4\times4$ dm et la hauteur 2 dm ; $1$ dm$^3=1$ litre.
\item $V(2)=16$ et $V'(2)=12$. L’équation de la tangente est
$y=12(x-2)+16=\res{12x-8}$.
\item En développant, $-(x-2)(x^2-4x-8)=-x^3+6x^2-16=V(x)-16$.
Le second facteur a pour racines $2-2\sqrt3<0$ et $2+2\sqrt3\in]0;6[$.
Sur le domaine géométrique, le tableau de signes ci-dessous donne
$V(x)-16\geq0\iff\res{x\in[2;2+2\sqrt3]}$.
Les deux bornes sont incluses, car elles donnent exactement 16 litres.
\end{questions}

\label{figlivre-r14-12-var-cor}
\figureLivre[corrige]{r14-12-var}

\label{figlivre-r14-12-signes-cor}
\figureLivre[corrige]{r14-12-signes}

\label{figlivre-r14-12-boite-cor}
\figureLivre[corrige]{r14-12-boite}

\end{corrige}
\end{blocLivre}
\begin{blocLivre}{cor-14-13}
\begin{corrige}{13}{Synthèse : triangle, cercle et tangente}
\demarche{Relier produit scalaire, cercle de diamètre donné, tangente et longueur d’arc.}
\begin{questions}
\item $\vect{CA}=(-2;-2)$ et $\vect{CB}=(2;-2)$ ont pour produit $-4+4=0$ : l’angle en $C$ est droit.
De plus $AC^2=BC^2=8$, donc le triangle est rectangle isocèle.
Ses deux autres angles sont égaux et se partagent $\pi/2$ : $\res{\widehat{BAC}=\pi/4}$.
\item Dans un triangle rectangle, le milieu de l’hypoténuse est le centre du cercle circonscrit.
Ici le diamètre est $[AB]$, le centre $\Omega(2;0)$ et le rayon 2 :
$\res{(x-2)^2+y^2=4}$.
\item Le rayon $\Omega C$ est vertical, donc la tangente en $C$ est horizontale : $\res{y=2}$.
La projection verticale de $M(5;3)$ conserve son abscisse et prend l’ordonnée de la tangente : $\res{H(5;2)}$.
\item Les rayons $\vect{\Omega A}=(-2;0)$ et $\vect{\Omega C}=(0;2)$ sont perpendiculaires.
L’angle central du petit arc vaut donc $\pi/2$ radians.
Sa longueur est $R\theta=2(\pi/2)=\res{\pi}$ unités de longueur.
\end{questions}

\label{figlivre-r14-13-geometrie-cor}
\figureLivre[corrige]{r14-13-geometrie}

\end{corrige}
\end{blocLivre}
\begin{blocLivre}{cor-14-14}
\begin{corrige}{14}{Synthèse : défauts et gain net}
\demarche{Lier l’arbre des fournisseurs à la loi du gain net.}
\begin{questions}
\item Les deux fournisseurs forment une partition.
$P(A\cap D)=0{,}6\times0{,}01=0{,}006$ et $P(B\cap D)=0{,}4\times0{,}04=0{,}016$.
Donc $P(D)=0{,}006+0{,}016=\res{0{,}022}$.
\item Parmi les articles défectueux, la part provenant de B vaut
$P_D(B)=P(B\cap D)/P(D)=0{,}016/0{,}022=\res{8/11}$.
\item Un article sans défaut donne $20-12=8$ euros ; avec défaut, il donne $20-12-30=-22$ euros.
Le tableau récapitule cette loi.
Ainsi $E(G)=8(0{,}978)-22(0{,}022)=\res{7{,}34\text{ euros}}$.
Le deuxième moment vaut $E(G^2)=64(0{,}978)+484(0{,}022)=73{,}24$,
d’où $V(G)=73{,}24-7{,}34^2=\res{19{,}3644}$ euros carrés.
\item L’évènement contraire est aucun défaut sur les quatre articles.
Par indépendance, il a pour probabilité $0{,}978^4$.
On obtient $P(\text{au moins un défaut})=\res{1-0{,}978^4\approx0{,}08514}$.
\end{questions}

\label{figlivre-r14-14-cor}
\figureLivre[corrige]{r14-14}

\label{figlivre-r14-14-loi-cor}
\figureLivre[corrige]{r14-14-loi}

\end{corrige}
\end{blocLivre}
\begin{blocLivre}{cor-14-15}
\begin{corrige}{15}{Synthèse : deux modèles d'évolution}
\demarche{Comparer les deux modèles aux instants entiers en vérifiant les seuils et leurs prédécesseurs.}
\begin{questions}
\item $v_{n+1}=0{,}9v_n$ et $v_0=500$, donc $\res{u_n=500+500(0{,}9)^n}$.
\item $u_{n+1}-u_n=-50(0{,}9)^n<0$ et $u_n-500=500(0{,}9)^n>0$.
\item $u_{15}\approx602{,}946>600$ et $u_{16}\approx592{,}651\leq600$ : premier rang $\res{16}$.
\item $W\prime(t)=-50e^{-0{,}1t}<0$. $W(16)\approx600{,}948$ et $W(17)\approx591{,}342$ : première semaine entière $\res{17}$.
\item Pour $x=-0{,}1\neq0$, $e^{-0{,}1}>0{,}9$. En élevant à la puissance entière positive n, $(e^{-0{,}1})^n>(0{,}9)^n$, donc $\res{W(n)>u_n}$. Le modèle continu décroît moins vite aux instants entiers, d’où le seuil plus tardif.
\end{questions}

Pour justifier que les seuils trouvés sont les premiers, on utilise la décroissance de chaque modèle, puis les valeurs de part et d’autre : $u_{15}>600\geq u_{16}$ et $W(16)>600\geq W(17)$.

\label{figlivre-r14-15-modeles-cor}
\figureLivre[corrige]{r14-15-modeles}

\end{corrige}
\end{blocLivre}

\editionPage[\pdfbookmark[0]{Corrigés des défis}{bm-corriges-defis}]{defi-cor-1-1}{D01 · Corrigé}{Deux racines dans un intervalle}{\hyperref[defi-ex-1]{Énoncé p.~\pageref*{defi-ex-1}} · Démarche et vérifications}
\begin{blocEdition}{defi-cor-1-1}
\editionEtape{1. Traduire la position des racines}
On calcule $\Delta=4m(m-1)$. Il y a deux racines distinctes si et seulement si
$m<0$ ou $m>1$. Le coefficient de $x^2$ est positif : la parabole est tournée vers le haut.
Si les deux racines $r_1<r_2$ sont dans $]0;2[$, les bornes 0 et 2 sont à l'extérieur
de l'intervalle entre les racines. On a donc
\[P_m(0)=m>0,\qquad P_m(2)=4-3m>0.\]
Le sommet a pour abscisse $(r_1+r_2)/2=m$, donc $0<m<2$.
\editionEtape{2. Justifier la réciproque}
Supposons $\Delta>0$, $P_m(0)>0$, $P_m(2)>0$ et $0<m<2$.
L'ordonnée du sommet vaut $m-m^2=-\Delta/4<0$. Sur $[0;m]$, la fonction
décroît d'une valeur positive à une valeur négative ; la première racine se trouve
donc entre 0 et $m$. De même, la seconde se trouve entre $m$ et 2.
On peut aussi utiliser les racines explicites $m\pm\sqrt{m(m-1)}$.
\editionEtape{3. Réunir les conditions}
Les conditions de signe donnent $0<m<4/3$, et $\Delta>0$ impose alors $m>1$.
La condition sur le sommet est automatiquement satisfaite. Ainsi
\[\res{m\in\left]1;\frac43\right[}.\]
Pour $m=1$, la racine 1 est double. Pour $m=4/3$, les racines sont $2/3$ et 2 :
la seconde est exclue du domaine demandé. Aucun des deux paramètres frontières ne convient.
\editionEtape{4. Repérer une condition insuffisante}
Pour $m=1/2$, $P_m(0)=1/2$ et $P_m(2)=5/2$ sont positifs, mais $\Delta=-1<0$ :
il n'existe aucune racine réelle.
\editionRetenir{Une liste de conditions nécessaires n'est pas encore une équivalence.
La réciproque doit être justifiée, puis les bornes du domaine doivent être contrôlées.}
\end{blocEdition}

\editionPage{defi-cor-2-1}{D02 · Corrigé}{Le seuil d’une moyenne}{\hyperref[defi-ex-2]{Énoncé p.~\pageref*{defi-ex-2}} · Démarche et vérifications}
\begin{blocEdition}{defi-cor-2-1}
\editionEtape{1. Déplacer le point fixe}
Le nombre 300 vérifie $300=(4/5)300+60$. En posant $v_n=u_n-300$,
on obtient $v_{n+1}=(4/5)v_n$ et $v_0=-300$.
Ainsi $u_n=300-300(4/5)^n$. La puissance $(4/5)^n$ décroît strictement :
$u_n$ croît strictement, reste inférieur à 300 et commence à 0.
\editionEtape{2. Additionner le bon nombre de termes}
Il y a $n$ termes, des indices 0 à $n-1$. La somme géométrique donne
\[S_n=300n-300\frac{1-(4/5)^n}{1-4/5}
=300n-1500\bigl(1-(4/5)^n\bigr).\]
En divisant par $n>0$, on obtient l'expression demandée de $A_n$.
\editionEtape{3. Comparer une moyenne au terme suivant}
Chacun des $n$ termes $u_0,\ldots,u_{n-1}$ est strictement inférieur à $u_n$.
Leur moyenne vérifie donc $A_n<u_n$.
Comme $S_{n+1}=S_n+u_n=nA_n+u_n$,
\[A_{n+1}-A_n=\frac{nA_n+u_n}{n+1}-A_n
=\frac{u_n-A_n}{n+1}>0.\]
La suite des moyennes est strictement croissante.
\editionRetenir{Pour étudier une moyenne de termes successifs, ajouter le nouveau terme
à la somme est souvent plus simple que soustraire deux formules compliquées.}
\end{blocEdition}

\editionPage{defi-cor-2-2}{D02 · Corrigé (suite)}{Le seuil d’une moyenne}{\hyperref[defi-ex-2]{Énoncé p.~\pageref*{defi-ex-2}} · Démarche et vérifications}
\begin{blocEdition}{defi-cor-2-2}
\editionEtape{4. Programmer sans décaler les indices}
\begin{python}
def seuil_moyenne():
    n = 1
    u = 60.0
    total = 0.0
    while total / n < 240:
        total = total + u
        u = 0.8 * u + 60
        n = n + 1
    return n
\end{python}
Au test de boucle, \code{total} est $S_n$ et \code{u} est $u_n$.
Au départ, $S_1=u_0=0$ et $u_1=60$. On ajoute $u_n$ à la somme,
puis on prépare $u_{n+1}$ et on augmente $n$ : les trois variables restent cohérentes.
La fonction renvoie \res{25}.
\editionEtape{5. Valider le premier franchissement}
Avec la formule exacte de la première page,
\[A_{24}\approx237{,}795\quad\text{et}\quad A_{25}\approx240{,}227.\]
Ainsi $A_{24}<240\leq A_{25}$. La croissance stricte prouve que tous les rangs
antérieurs à 25 échouent : il ne suffit pas d'obtenir une valeur qui dépasse le seuil,
il faut aussi justifier qu'elle est la première.
\editionRetenir{Le programme repère le rang candidat. La formule et la monotonie
expliquent pourquoi ce candidat répond exactement à la question.}
\end{blocEdition}

\editionPage{defi-cor-3-1}{D03 · Corrigé}{Trois tangentes par un point}{\hyperref[defi-ex-3]{Énoncé p.~\pageref*{defi-ex-3}} · Démarche et vérifications}
\begin{blocEdition}{defi-cor-3-1}
\editionEtape{1. Remplacer une question géométrique par une équation}
Comme $f'(t)=3t^2-3$, la formule de la tangente donne
\[y=f'(t)(x-t)+f(t)=(3t^2-3)x-2t^3.\]
En prenant $x=1$, son ordonnée vaut $g(t)=-2t^3+3t^2-3$.
Chercher les tangentes passant par $P_b$ revient donc à compter les solutions de $g(t)=b$.
\editionEtape{2. Compter les intersections avec une horizontale}
On a $g'(t)=6t(1-t)$. Ainsi $g$ décroît sur $[-1;0]$, croît sur $[0;1]$
et décroît sur $[1;2]$, avec
\[g(-1)=2,\quad g(0)=-3,\quad g(1)=-2,\quad g(2)=-7.\]
\begin{center}\begin{tabular}{@{}lc@{}}\toprule
Valeur de $b$ & Nombre de tangentes\\\midrule
$b<-7$ ou $b>2$ & $0$\\
$-7\leq b<-3$ ou $-2<b\leq2$ & $1$\\
$b=-3$ ou $b=-2$ & $2$\\
$-3<b<-2$ & $3$\\\bottomrule\end{tabular}\end{center}
À une valeur critique, deux branches du tableau aboutissent au \emph{même} point :
ce point ne doit être compté qu'une fois.
Des paramètres distincts donnent ici des droites distinctes : si deux tangentes avaient
le même terme constant $-2t^3$, leurs paramètres seraient égaux.
\end{blocEdition}

\editionPage{defi-cor-3-2}{D03 · Corrigé (suite)}{Trois tangentes par un point}{\hyperref[defi-ex-3]{Énoncé p.~\pageref*{defi-ex-3}} · Démarche et vérifications}
\begin{blocEdition}{defi-cor-3-2}
\editionEtape{3. Trouver exactement les trois tangentes}
Pour $b=-5/2$, l'équation devient $4t^3-6t^2+1=0$.
Le nombre $1/2$ est racine et
\[4t^3-6t^2+1=(2t-1)(2t^2-2t-1).\]
Les trois abscisses de contact sont
\[\res{\frac{1-\sqrt3}{2},\quad\frac12,\quad\frac{1+\sqrt3}{2}}.\]
Elles appartiennent bien à $[-1;2]$. Pour $t=1/2$, la pente vaut $-9/4$.
Pour les deux autres valeurs, $t^2=t+1/2$, donc la pente vaut $3t-3/2=\pm3\sqrt3/2$.
Toutes les droites passent par $(1;-5/2)$ ; leurs équations sont donc
\[\res{y=-\frac94x-\frac14},\qquad
y=\frac{3\sqrt3}{2}(x-1)-\frac52,\qquad
y=-\frac{3\sqrt3}{2}(x-1)-\frac52.\]
\editionRetenir{Le paramètre de la tangente est une abscisse de contact.
Une droite horizontale dans le graphique de $g$ permet de compter les tangentes dans le graphique de $f$.}
\label{figedition-defi-cor-3-2-1}\figureGeo[corrige]{fr-d03}
\end{blocEdition}

\editionPage{defi-cor-4-1}{D04 · Corrigé}{La fenêtre la plus lumineuse}{\hyperref[defi-ex-4]{Énoncé p.~\pageref*{defi-ex-4}} · Démarche et vérifications}
\begin{blocEdition}{defi-cor-4-1}
\editionEtape{1. Modéliser le contour avant de calculer}
Le contour comprend la base $2r$, deux côtés $h$ et le demi-cercle $\pi r$.
Le segment séparant les deux parties est intérieur. Ainsi
\[2h+(2+\pi)r=20,\qquad h=10-\frac{2+\pi}{2}r.\]
Les contraintes donnent $0<r<20/(2+\pi)$.
\editionEtape{2. Réduire à une seule variable}
L'aire est celle du rectangle plus celle du demi-disque :
\[A(r)=2rh+\frac{\pi r^2}{2}=20r-\frac{4+\pi}{2}r^2.\]
C'est un trinôme de coefficient dominant négatif. Sa dérivée vaut
$A'(r)=20-(4+\pi)r$ : elle est positive avant $r_*=20/(4+\pi)$ et négative après.
\editionEtape{3. Contrôler et interpréter l'optimum}
Le rayon $r_*$ appartient au domaine, puisque $4+\pi>2+\pi$.
En remplaçant dans la formule de $h$,
\[\res{r_*=h_*=\frac{20}{4+\pi}\ \mathrm{m}},\qquad
\res{A_{\max}=\frac{200}{4+\pi}\ \mathrm{m}^2}.\]
La largeur vaut $40/(4+\pi)$ m. Les deux dimensions sont strictement positives :
le maximum est atteint à l'intérieur du domaine géométrique.
\editionRetenir{Un bon modèle distingue le contour extérieur des segments intérieurs.
Après l'optimisation, revenir aux dimensions et aux unités donne son sens au résultat.}
\label{figedition-defi-cor-4-1-1}\figureGeo[corrige]{fr-d04}
\end{blocEdition}

\editionPage{defi-cor-5-1}{D05 · Corrigé}{Encadrer sans logarithme}{\hyperref[defi-ex-5]{Énoncé p.~\pageref*{defi-ex-5}} · Démarche et vérifications}
\begin{blocEdition}{defi-cor-5-1}
\editionEtape{1. Établir une inégalité valable sur tout le domaine}
La dérivée $F'(x)=e^x-1$ est négative pour $x<0$, nulle en 0 et positive pour $x>0$.
Ainsi $F$ admet son minimum $F(0)=0$. Il vient $e^x\geq1+x$, et l'égalité
n'est possible qu'en 0.
\editionEtape{2. Inverser sans changer le sens à tort}
Pour $0<t<1$, l'inégalité stricte donne $e^t>1+t$ et $e^{-t}>1-t>0$.
Sur les nombres strictement positifs, prendre les inverses renverse l'ordre :
$e^t<1/(1-t)$. La positivité de $1-t$ est indispensable.
\editionEtape{3. Passer aux puissances entières}
On prend $t=1/(2n)$, qui est dans $]0;1[$. Les trois quantités sont positives ;
la puissance entière $n$ conserve donc les inégalités. Comme $(e^{1/(2n)})^n=e^{1/2}$,
on obtient $L_n<e^{1/2}<U_n$.
\editionRetenir{Pour produire un encadrement, combiner une inégalité générale
avec des opérations dont le sens et le domaine sont explicitement contrôlés.}
\end{blocEdition}

\editionPage{defi-cor-5-2}{D05 · Corrigé (suite)}{Encadrer sans logarithme}{\hyperref[defi-ex-5]{Énoncé p.~\pageref*{defi-ex-5}} · Démarche et vérifications}
\begin{blocEdition}{defi-cor-5-2}
\editionEtape{4. Calculer et arrondir dans le bon sens}
\begin{python}
def bornes_exponentielle(n):
    bas = (1 + 1 / (2 * n)) ** n
    haut = (1 - 1 / (2 * n)) ** (-n)
    return bas, haut
\end{python}
Des arrondis extérieurs à cinq décimales donnent :
\begin{center}\begin{tabular}{@{}ccc@{}}\toprule
$n$ & Borne inférieure & Borne supérieure\\\midrule
10 & $1{,}62889$ & $1{,}67019$\\
100 & $1{,}64666$ & $1{,}65080$\\
1000 & $1{,}64851$ & $1{,}64893$\\\bottomrule\end{tabular}\end{center}
Chaque ligne enferme $e^{1/2}$. Un arrondi usuel au plus proche pourrait
déplacer une borne vers l'intérieur ; c'est pourquoi on arrondit ici volontairement vers l'extérieur.
\editionEtape{5. Distinguer preuve et observation}
La preuve garantit l'encadrement \emph{pour chaque entier} $n\geq1$.
Les calculs suggèrent qu'il se resserre. À eux seuls, quelques essais ne démontrent
ni une propriété pour tous les rangs, ni une convergence.
La fonction Python travaille en nombres approchés : pour certifier une précision arbitraire,
on utiliserait des calculs rationnels ou des bornes d'erreur adaptées.
\editionRetenir{Le calcul numérique illustre une preuve ; il ne remplace pas
le contrôle du sens des inégalités et des arrondis.}
\end{blocEdition}

\editionPage{defi-cor-6-1}{D06 · Corrigé}{Additionner sinus et cosinus}{\hyperref[defi-ex-6]{Énoncé p.~\pageref*{defi-ex-6}} · Démarche et vérifications}
\begin{blocEdition}{defi-cor-6-1}
\editionEtape{1. Changer de représentation}
Les coordonnées vérifient $x^2+y^2=1$ ; l'autre condition est la droite $x+y=k$.
De $(x-y)^2\geq0$, on déduit $2xy\leq1$, puis
$(x+y)^2=1+2xy\leq2$. Il faut donc $|k|\leq\sqrt2$.
\editionEtape{2. Reconstruire les coordonnées}
En développant $k^2=(x+y)^2$, on trouve $xy=(k^2-1)/2$.
Le trinôme dont les racines sont $x$ et $y$ a pour discriminant $2-k^2$.
Pour $|k|\leq\sqrt2$, les points cherchés ont pour coordonnées
\[\left(\frac{k+\sqrt{2-k^2}}2;\frac{k-\sqrt{2-k^2}}2\right)
\quad\text{et}\quad
\left(\frac{k-\sqrt{2-k^2}}2;\frac{k+\sqrt{2-k^2}}2\right).\]
Ces couples vérifient bien les deux équations ; la condition est donc suffisante.
Si $|k|<\sqrt2$, les deux points sont distincts : il existe deux solutions $t$.
Si $k=\pm\sqrt2$, les deux couples coïncident : une solution.
Si $|k|>\sqrt2$, aucune solution.
\editionRetenir{Une condition nécessaire devient une caractérisation seulement
après avoir vérifié que les coordonnées construites satisfont les équations de départ.}
\end{blocEdition}

\editionPage{defi-cor-6-2}{D06 · Corrigé (suite)}{Additionner sinus et cosinus}{\hyperref[defi-ex-6]{Énoncé p.~\pageref*{defi-ex-6}} · Démarche et vérifications}
\begin{blocEdition}{defi-cor-6-2}
\editionEtape{3. Revenir aux angles remarquables}
Pour $k=0$, on a $y=-x$ et les points
$(-\sqrt2/2;\sqrt2/2)$ et $(\sqrt2/2;-\sqrt2/2)$ :
\[\res{t\in\{3\pi/4;7\pi/4\}}.\]
Pour $k=1$, les coordonnées sont $(1;0)$ et $(0;1)$, donc
\[\res{t\in\{0;\pi/2\}}.\]
Pour $k=\sqrt2$, il n'y a qu'un point, $(\sqrt2/2;\sqrt2/2)$, donc
\[\res{t=\pi/4}.\]
Pour une valeur générale de $k$, les coordonnées exactes suffisent à décrire les points :
on ne suppose pas que leurs angles soient des angles remarquables.
\editionRetenir{Les points du cercle et les angles ne sont pas des réponses interchangeables :
préciser le domaine de l'angle permet de compter correctement les solutions.}
\label{figedition-defi-cor-6-2-1}\figureGeo[corrige]{fr-d06}
\end{blocEdition}

\editionPage{defi-cor-7-1}{D07 · Corrigé}{Un triangle sous deux contraintes}{\hyperref[defi-ex-7]{Énoncé p.~\pageref*{defi-ex-7}} · Démarche et vérifications}
\begin{blocEdition}{defi-cor-7-1}
\editionEtape{1. Reconnaître le lieu d'orthogonalité}
Les vecteurs sont $(-2-x;-y)$ et $(2-x;-y)$ ; leur produit scalaire vaut
$x^2+y^2-4$. Le lieu est le cercle de centre $O$ et de rayon 2, de diamètre $[AB]$.
Sur ce cercle, $MA^2=(x+2)^2+y^2=8+4x$ et $MB^2=8-4x$.
\editionEtape{2. Réduire le second lieu à une droite}
La seconde condition devient $24-4x=k$, donc $x=6-k/4$.
Il reste
\[y^2=4-\left(6-\frac k4\right)^2.\]
Cette quantité est positive si $16<k<32$, nulle si $k=16$ ou $32$, négative sinon.
Il existe donc deux points pour $16<k<32$ :
\[\res{M\left(6-\frac k4;\ \pm\sqrt{4-(6-k/4)^2}\right)}.\]
Aux bornes, un seul point subsiste : $B$ pour $k=16$, $A$ pour $k=32$.
Hors de $[16;32]$, aucun point ne convient.

\end{blocEdition}

\editionPage{defi-cor-7-2}{D07 · Corrigé (suite)}{Un triangle sous deux contraintes}{\hyperref[defi-ex-7]{Énoncé p.~\pageref*{defi-ex-7}} · Démarche et vérifications}
\begin{blocEdition}{defi-cor-7-2}
\editionEtape{3. Séparer orthogonalité et angle d'un triangle}
En $A$ ou en $B$, l'un des vecteurs est nul. Son produit scalaire avec l'autre est nul,
mais le triangle est dégénéré : l'angle demandé n'existe pas.
Les vrais triangles correspondent exactement à \res{16<k<32}.
Enfin $MA=MB$ équivaut à $x=0$, d'où \res{k=24} et \res{M(0;2)\text{ ou }M(0;-2)}.
\editionRetenir{Un produit scalaire nul autorise un vecteur nul. Pour conclure sur
un angle droit, vérifier que les trois sommets forment effectivement un triangle.}
\label{figedition-defi-cor-7-2-1}\figureGeo[corrige]{fr-d07}
\end{blocEdition}

\editionPage{defi-cor-8-1}{D08 · Corrigé}{Une famille de cercles tangents}{\hyperref[defi-ex-8]{Énoncé p.~\pageref*{defi-ex-8}} · Démarche et vérifications}
\begin{blocEdition}{defi-cor-8-1}
\editionEtape{1. Identifier les cercles sans oublier le rayon}
On complète les carrés :
\[(x-2)^2+(y-m)^2=4+m^2.\]
Le centre est $\Omega_m(2;m)$ et le rayon $\sqrt{4+m^2}$, strictement positif pour tout $m$.
Les points $(0;0)$ et $(4;0)$ vérifient l'équation initiale, quel que soit $m$.
\editionEtape{2. Compter les intersections par le discriminant}
Après substitution,
\[2x^2+(2m-16)x+36-12m=0.\]
Son discriminant vaut
\[(2m-16)^2-8(36-12m)=4\bigl((m+4)^2-24\bigr).\]
Posons $m_-=-4-2\sqrt6$ et $m_+=-4+2\sqrt6$.
Il existe deux intersections si $m<m_-$ ou $m>m_+$, aucune si $m_-<m<m_+$,
et une seule si $m=m_-$ ou $m=m_+$.
Les deux paramètres de tangence sont donc \res{m=-4\pm2\sqrt6}.
\editionEtape{3. Trouver le point de contact par projection}
Un vecteur normal à $d$ est $(1;1)$. Le projeté est de la forme
$H_m=(2-\lambda;m-\lambda)$. La condition $x_H+y_H=6$ donne
$\lambda=(m-4)/2$, donc
\[\res{H_m\left(4-\frac m2;2+\frac m2\right)}.\]
\end{blocEdition}

\editionPage{defi-cor-8-2}{D08 · Corrigé (suite)}{Une famille de cercles tangents}{\hyperref[defi-ex-8]{Énoncé p.~\pageref*{defi-ex-8}} · Démarche et vérifications}
\begin{blocEdition}{defi-cor-8-2}
\editionEtape{4. Vérifier les deux constructions}
La distance au carré entre le centre et son projeté est
\[\Omega_mH_m^2=\frac{(m-4)^2}{2}.\]
Pour $m=m_-$ ou $m=m_+$, l'identité $m^2+8m-8=0$ entraîne
$(m-4)^2/2=m^2+4$. La distance au centre est donc égale au rayon : le projeté
appartient bien au cercle. Ses coordonnées explicites sont
\[H_-\bigl(6+\sqrt6;-\sqrt6\bigr),\qquad
H_+\bigl(6-\sqrt6;\sqrt6\bigr).\]
Le vecteur $\vect{\Omega_mH_m}$ est un multiple de $(1;1)$ : il est normal à $d$.
C'est donc le rayon perpendiculaire à la tangente au point de contact.
\editionRetenir{Le discriminant compte les intersections ; la projection construit
le point de contact. Faire concorder les deux démarches fournit une vérification indépendante.}
\label{figedition-defi-cor-8-2-1}\figureGeo[corrige]{fr-d08-cor}
\end{blocEdition}

\editionPage{defi-cor-9-1}{D09 · Corrigé}{Choisir une règle de contrôle}{\hyperref[defi-ex-9]{Énoncé p.~\pageref*{defi-ex-9}} · Démarche et vérifications}
\begin{blocEdition}{defi-cor-9-1}
\editionEtape{1. Respecter la population conditionnante}
Dans le groupe $D$, le chemin $++$ a pour poids $0{,}9^2=0{,}81$ ; dans le groupe
$\overline D$, son poids est $0{,}2^2=0{,}04$.
En ajoutant les deux origines, $P(T_1^+\cap T_2^+)=0{,}1(0{,}81)+0{,}9(0{,}04)=0{,}117$.
Chaque test a une probabilité marginale $0{,}1(0{,}9)+0{,}9(0{,}2)=0{,}27$ d'être positif.
Or $0{,}27^2=0{,}0729\ne0{,}117$ : les tests ne sont pas indépendants dans la population entière.
\editionEtape{2. Inverser le conditionnement}
Pour un seul test positif, la part de défauts est
\[\frac{P(D\cap T_1^+)}{P(T_1^+)}=\frac{0{,}09}{0{,}27}=\res{\frac13}.\]
Pour deux tests positifs, elle devient
\[\frac{P(D\cap T_1^+\cap T_2^+)}{P(T_1^+\cap T_2^+)}
=\frac{0{,}081}{0{,}117}=\res{\frac9{13}}.\]
\editionRetenir{Des poids identiques au second test dans chaque groupe ne donnent
pas automatiquement l'indépendance dans le mélange des groupes.}
\label{figedition-defi-cor-9-1-1}\figureGeo[corrige]{fr-d09}
\end{blocEdition}

\editionPage{defi-cor-9-2}{D09 · Corrigé (suite)}{Choisir une règle de contrôle}{\hyperref[defi-ex-9]{Énoncé p.~\pageref*{defi-ex-9}} · Démarche et vérifications}
\begin{blocEdition}{defi-cor-9-2}
\editionEtape{3. Compter les deux erreurs séparément}
Pour U, on rejette un article sans défaut sauf si les deux tests sont négatifs :
$P(\overline D\text{ et rejet})=0{,}9(1-0{,}8^2)=0{,}324$.
Un défaut passe seulement par le chemin $D--$, de poids $0{,}1(0{,}1^2)=0{,}001$.
Pour B, les erreurs ont les probabilités $0{,}9(0{,}2^2)=0{,}036$ et
$0{,}1(1-0{,}9^2)=0{,}019$.
\editionEtape{4. Comparer des espérances, pas seulement des pourcentages}
Les coûts moyens valent
\[C_U=0{,}324+0{,}001\lambda,\qquad C_B=0{,}036+0{,}019\lambda.\]
Leur différence est $C_U-C_B=0{,}288-0{,}018\lambda$.
Pour $0\leq\lambda<16$, B coûte moins ; pour $\lambda>16$, U coûte moins ;
pour $\lambda=16$, les coûts sont égaux.
Par exemple, pour $\lambda=5$, on obtient $C_U=0{,}329$ et $C_B=0{,}131$.
\editionRetenir{La règle qui détecte le plus de défauts n'est pas nécessairement
celle qui minimise un coût donné. Le choix dépend du modèle et des conséquences des deux erreurs.}
\end{blocEdition}

\editionPage{defi-cor-10-1}{D10 · Corrigé}{Même moyenne, objectifs différents}{\hyperref[defi-ex-10]{Énoncé p.~\pageref*{defi-ex-10}} · Démarche et vérifications}
\begin{blocEdition}{defi-cor-10-1}
\editionEtape{1. Séparer les issues et les valeurs}
Les trois issues ont pour lots $2a$, $2-a$ et 6. Les deux premiers coïncident
pour $a=2/3$ ; on regroupe alors leurs probabilités en $3/4$.
Dans tous les cas,
\[E(L_a)=\frac14(2a)+\frac12(2-a)+\frac14(6)=\frac52,\qquad \res{E(G_a)=0}.\]
\editionEtape{2. Optimiser la dispersion}
On calcule $E(L_a^2)=a^2+(2-a)^2/2+9=3a^2/2-2a+11$.
Soustraire le carré $25/4$ de la moyenne donne
\[V(G_a)=V(L_a)=\frac32\left(a-\frac23\right)^2+\frac{49}{12}.\]
Le minimum est donc \res{49/12}, atteint uniquement en \res{a=2/3}.

\end{blocEdition}

\editionPage{defi-cor-10-2}{D10 · Corrigé (suite)}{Même moyenne, objectifs différents}{\hyperref[defi-ex-10]{Énoncé p.~\pageref*{defi-ex-10}} · Démarche et vérifications}
\begin{blocEdition}{defi-cor-10-2}
\editionEtape{3. Changer d'objectif peut changer la réponse}
Le deuxième gain net vaut $-a-1/2<0$ ; le troisième vaut $7/2>0$.
Le premier, $2a-5/2$, est strictement positif seulement si $a>5/4$.
Ainsi $P(G_a>0)=1/4$ sur $[0;5/4]$ et $1/2$ sur $]5/4;2]$.
Le choix $a=2/3$ ne maximise donc pas cette probabilité.
\editionEtape{4. Choisir aussi la prévision}
En développant et en utilisant $E(L_a)=5/2$,
\[E((L_a-r)^2)=V(L_a)+(r-5/2)^2.
\]
Les deux carrés sont minimaux pour \res{a=2/3\text{ et }r=5/2}.
\editionRetenir{Espérance, variance, chance de gain et erreur de prévision répondent
à des questions différentes. Une moyenne identique ne suffit pas à départager deux choix.}
\label{figedition-defi-cor-10-2-1}\figureGeo[corrige]{fr-d10}
\end{blocEdition}

\editionPage{defi-cor-11-1}{D11 · Corrigé}{La dispersion d’un groupe réuni}{\hyperref[defi-ex-11]{Énoncé p.~\pageref*{defi-ex-11}} · Démarche et vérifications}
\begin{blocEdition}{defi-cor-11-1}
\editionEtape{1. Calculer un exemple qui oblige à changer d'idée}
Dans A, $\mu_A=1$ et $V_A=1$ ; dans B, $\mu_B=9$ et $V_B=1$.
La réunion est $0,2,8,10$, de moyenne 5. Sa variance vaut
\[V=\frac{25+9+9+25}{4}=\res{17}.\]
La moyenne des variances internes vaut seulement 1. Elle ignore l'écart entre les centres des deux groupes.
\editionEtape{2. Décomposer les écarts à la moyenne totale}
La moyenne totale est $\mu=p\mu_A+(1-p)\mu_B$.
Dans A, la moyenne de $x-\mu_A$ est nulle. Le terme croisé du développement
s'annule donc après la moyenne :
\[\text{moyenne dans A de }(x-\mu)^2=V_A+(\mu_A-\mu)^2.\]
Dans B, on obtient $V_B+(\mu_B-\mu)^2$.
Les écarts des moyennes au centre total sont
\[\mu_A-\mu=(1-p)(\mu_A-\mu_B),\qquad
\mu_B-\mu=-p(\mu_A-\mu_B).\]
\end{blocEdition}

\editionPage{defi-cor-11-2}{D11 · Corrigé (suite)}{La dispersion d’un groupe réuni}{\hyperref[defi-ex-11]{Énoncé p.~\pageref*{defi-ex-11}} · Démarche et vérifications}
\begin{blocEdition}{defi-cor-11-2}
\editionEtape{3. Réunir en pondérant par les effectifs}
Une proportion $p$ des observations vient de A ; l'autre vient de B. Ainsi
\begin{align*}
V&=pV_A+(1-p)V_B+p(\mu_A-\mu)^2+(1-p)(\mu_B-\mu)^2\\
 &=pV_A+(1-p)V_B+p(1-p)\bigl((1-p)+p\bigr)(\mu_A-\mu_B)^2.
\end{align*}
On retrouve la formule demandée. Le dernier terme est positif ou nul et s'annule
quand les deux moyennes sont égales.
Pour l'exemple, $p=1/2$ et l'écart des moyennes vaut 8 : $V=1+8^2/4=17$.
\editionEtape{4. Construire une famille de contre-exemples}
Prenons A $=\{-1;1\}$ et B $=\{M-1;M+1\}$, avec $M>0$.
Chaque groupe a une variance égale à 1 ; leurs moyennes sont 0 et $M$.
La réunion a donc pour variance
\[\res{V=1+\frac{M^2}{4}}.\]
En choisissant $M$ suffisamment grand, cette variance dépasse n'importe quelle borne fixée.
\editionRetenir{Une dispersion totale associe la dispersion à l'intérieur des groupes
et la séparation de leurs moyennes. Les effectifs déterminent les pondérations.}
\label{figedition-defi-cor-11-2-1}\figureGeo[corrige]{fr-d11}
\end{blocEdition}

\editionPage{defi-cor-12-1}{D12 · Corrigé}{Un algorithme qui encadre une racine}{\hyperref[defi-ex-12]{Énoncé p.~\pageref*{defi-ex-12}} · Démarche et vérifications}
\begin{blocEdition}{defi-cor-12-1}
\editionEtape{1. Dire précisément ce qui est admis et ce qui est démontré}
On a $f'(x)=3x^2+1>0$ : $f$ est strictement croissante, donc elle s'annule au plus une fois.
L'existence est donnée dans l'énoncé ; on en déduit l'unicité.
\editionEtape{2. Conserver un invariant}
Si $f(m)<0=f(\alpha)$, la croissance donne $m<\alpha$ : on peut avancer la borne gauche.
Si $f(m)\geq0$, on a $m\geq\alpha$ : on peut reculer la borne droite.
Dans les deux cas, le nouvel intervalle contient toujours $\alpha$.
\begin{python}
def encadrer_racine(N):
    a, b = 0.0, 1.0
    for k in range(N):
        m = (a + b) / 2
        if m ** 3 + m - 1 < 0:
            a = m
        else:
            b = m
    return a, b
\end{python}
L'entrée $N$ est un entier naturel, y compris 0. À chaque passage,
$a\leq\alpha\leq b$ est conservé et la largeur est divisée par 2.
\end{blocEdition}

\editionPage{defi-cor-12-2}{D12 · Corrigé (suite)}{Un algorithme qui encadre une racine}{\hyperref[defi-ex-12]{Énoncé p.~\pageref*{defi-ex-12}} · Démarche et vérifications}
\begin{blocEdition}{defi-cor-12-2}
\editionEtape{3. Garantir une précision avant de lancer le calcul}
La largeur initiale vaut 1. Après $N$ divisions par 2, elle vaut $2^{-N}$.
Comme $2^9=512<1000$ et $2^{10}=1024\geq1000$, le premier nombre d'étapes
garantissant une largeur au plus $10^{-3}$ est \res{N=10}.
\editionEtape{4. Vérifier en arithmétique exacte}
On trouve
\[f\left(\frac{698}{1024}\right)=\frac{-1767384}{1073741824}<0,
\qquad f\left(\frac{699}{1024}\right)=\frac{744899}{1073741824}>0.\]
La croissance confirme que $\alpha$ est dans cet intervalle.
Son milieu est $1397/2048$. La distance d'un point de l'intervalle au milieu
ne dépasse pas la demi-largeur, d'où
\[\res{\left|\alpha-\frac{1397}{2048}\right|\leq\frac1{2048}<5\times10^{-4}}.\]
\editionRetenir{Un algorithme d'encadrement associe un invariant, une règle de réduction
et une précision garantie. L'affichage de nombreuses décimales ne suffit pas à garantir leur justesse.}
\label{figedition-defi-cor-12-2-1}\figureGeo[corrige]{fr-d12}
\end{blocEdition}

\editionPage{defi-cor-13-1}{D13 · Corrigé}{Une garantie pour tous les x}{\hyperref[defi-ex-13]{Énoncé p.~\pageref*{defi-ex-13}} · Démarche et vérifications}
\begin{blocEdition}{defi-cor-13-1}
\editionEtape{1. Respecter la portée des quantificateurs}
La négation est : « Il existe $x\in[0;1]$ tel que $P_m(x)<0$. »
La positivité partout impose $P_m(0)=m\geq0$ et $P_m(1)=1-m\geq0$.
Réciproquement, pour $m\in[0;1]$,
$P_m(x)=(x-m)^2+m(1-m)\geq0$. Les paramètres sont donc exactement \res{m\in[0;1]}.
\editionEtape{2. Trouver le minimum avec la position du sommet}
Pour $m<0$, $P'_m(x)=2(x-m)>0$ sur $[0;1]$, donc $q(m)=P_m(0)=m$.
Pour $0\leq m\leq1$, le sommet est dans le segment : $q(m)=m(1-m)$.
Pour $m>1$, la fonction décroît sur le segment : $q(m)=P_m(1)=1-m$.
Sur $[0;1]$, $m(1-m)=1/4-(m-1/2)^2$ ; à l'extérieur, $q(m)<0$.
La meilleure garantie est donc \res{q(1/2)=1/4}, atteinte uniquement pour \res{m=1/2}.
\editionEtape{3. Vérifier la garantie et distinguer deux énoncés}
Pour chaque $m\in\R$, $P_m(1/2)=1/4$, donc le minimum ne peut dépasser $1/4$.
Le choix $m=1/2$ donne $P_{1/2}(x)=(x-1/2)^2+1/4\geq1/4$ pour tout $x\in[0;1]$ :
la première affirmation est vraie.
La seconde est vraie aussi : pour \emph{tout} $m\in\R$, choisir $x=1/2$ suffit.
Mais elle autorise une seule valeur de $x$, et ne garantit aucune borne sur les autres valeurs.
Par exemple, pour $m=2$, $P_2(1)=-1<1/4$.
\editionRetenir{« Une borne valable partout » et « un point où la borne est dépassée »
expriment deux exigences différentes, même lorsque les deux affirmations sont vraies.}
\end{blocEdition}

\editionPage{defi-cor-14-1}{D14 · Corrigé}{Qualité, coût et risque}{\hyperref[defi-ex-14]{Énoncé p.~\pageref*{defi-ex-14}} · Démarche et vérifications}
\begin{blocEdition}{defi-cor-14-1}
\editionEtape{1. Passer de la loi au gain moyen}
Les deux gains sont $12-2x$, de probabilité $1-p$, et $-8-2x$, de probabilité $p$.
Ainsi $E(G_x)=(12-2x)(1-p)+(-8-2x)p=12-2x-20p=12-2x-4e^{-x}$.
\editionEtape{2. Étudier le signe sans logarithme}
On a $g'(x)=-2+4e^{-x}$ ; son signe est positif si $e^x<2$ et négatif si $e^x>2$.
La stricte croissance de l'exponentielle donne l'unicité de $\alpha$.
Comme $e^0=1<2<e^1$, on a $0<\alpha<1$, et $g$ croît jusqu'à $\alpha$, puis décroît.
\[\res{g(\alpha)=10-2\alpha}\qquad\text{car }e^{-\alpha}=1/2.\]
\editionEtape{3. Encadrer le paramètre de réglage}
\begin{python}
from math import exp
def encadrer_reglage(N=10):
    a, b = 0.0, 1.0
    for k in range(N):
        m = (a + b) / 2
        if exp(m) < 2:
            a = m
        else:
            b = m
    return a, b
\end{python}
Pour $N=10$, la largeur est $1/1024<10^{-3}$ ; on obtient
$709/1024<\alpha<710/1024$. En poursuivant les calculs, $\alpha\approx0{,}693147$
et le gain moyen maximal vaut environ $8{,}614$ unités.
\end{blocEdition}

\editionPage{defi-cor-14-2}{D14 · Corrigé (suite)}{Qualité, coût et risque}{\hyperref[defi-ex-14]{Énoncé p.~\pageref*{defi-ex-14}} · Démarche et vérifications}
\begin{blocEdition}{defi-cor-14-2}
\editionEtape{4. Distinguer gain moyen et dispersion}
Les deux gains diffèrent de 20. En calculant les deux écarts à la moyenne,
ou en développant le deuxième moment,
\[V(G_x)=400p(1-p)=80e^{-x}-16e^{-2x}.\]
La dépense $2x$, constante pour un réglage fixé, ne change pas la variance.
La dérivée vaut
\[V'(x)=16e^{-x}(2e^{-x}-5)<0\quad\text{sur }[0;4],\]
car $e^{-x}\leq1$. La dispersion décroît donc sur tout le segment.
Pour $x=\alpha$, on a \res{p(\alpha)=1/10} et \res{V(G_\alpha)=36}.
\editionEtape{5. Répéter deux contrôles indépendants}
L'arbre donne la probabilité de deux articles sans défaut : $(9/10)^2=81/100$.
Par complément,
\[\res{P(\text{au moins un défaut})=1-81/100=19/100}.\]
L'intensité 4 minimise ici la variance et le défaut, mais elle ne maximise pas l'espérance,
puisque $4>\alpha$ et $g$ décroît strictement sur $[\alpha;4]$.
\editionRetenir{Le meilleur réglage dépend du critère choisi. Un modèle doit
préciser les coûts, les probabilités et l'indépendance avant que l'on optimise quoi que ce soit.}
\label{figedition-defi-cor-14-2-1}\figureGeo[corrige]{fr-d14}
\end{blocEdition}


\editionPage[\pdfbookmark[0]{Corrigés des sujets blancs}{bm-corriges-sujets}]{sujet-A-cor-auto}{Sujet A · Corrigé}{Automatismes : réponses et méthode}{\hyperref[sujet-A]{Énoncé p.~\pageref*{sujet-A}}}
\begin{blocEdition}{sujet-A-cor-auto}
\begin{questions}\item $60\times1{,}15=\res{69\text{ euros}}$.
\item On divise par le coefficient multiplicateur : $75/1{,}25=\res{60\text{ euros}}$.
\item Le coefficient réciproque est $1/0{,}8=1{,}25$ : \res{25\,\%}.
\item Par différence de deux carrés : $4-3=\res{1}$.
\item Un produit est nul si l’un de ses facteurs est nul : $\res{S=\{-3\,;\,1/2\}}$.
\item On divise par $-3$ en changeant le sens : \res{x<2}, soit $\mathopen{]}-\infty;2[$.
\item Le coefficient de $x^2$ vaut 1 : le produit est strictement négatif entre ses racines $-4$ et 1, donc $\res{x\in\mathopen{]}-4;1[}$.
\item $m=\dfrac{-5-3}{3-(-1)}=\dfrac{-8}{4}=\res{-2}$.
\item $P(A\cap B)=P(A)\times P_A(B)=0{,}4\times0{,}25=\res{0{,}1}$.
\item La série ordonnée a 5 valeurs ; la troisième est la médiane : \res{2}.
\item La courbe est sur ou sous la droite d’équation $y=1$ pour les abscisses de $-1$ à 3 : $\res{S=[-1;3]}$.
\item L’ordonnée à l’origine vaut 3 ; de $(0;3)$ à $(2;2)$, $y$ diminue de 1 quand $x$ augmente de 2 : $\res{y=-\tfrac12x+3}$.\end{questions}
\label{figedition-sujet-A-cor-auto-1}\figureGeo[corrige]{fr-sa-auto}
\end{blocEdition}

\editionPage{sujet-A-cor-1}{Sujet A · Corrigé 1}{Une courbe et sa tangente}{\hyperref[sujet-A-ex-1]{Énoncé p.~\pageref*{sujet-A-ex-1}} · Raisonnement et rédaction}
\begin{blocEdition}{sujet-A-cor-1}

\editionEtape{1. Dérivée, signes et variations}
On obtient $f'(x)=3x^2-12x+9=3(x-1)(x-3)$.
Sur $[0;4]$, les signes sont successivement $+$, $-$, $+$, avec des zéros en 1 et 3.
\begin{center}\begin{tableau}\otTab[2.1]{$x$/1,$f'(x)$/1,$f$/2}{0,1,3,4}
\tkzTabLine{,+,z,-,z,+,}\tkzTabVar{-/2,+/6,-/2,+/6}\end{tableau}\end{center}
Le minimum est \res{2}, atteint en 0 et 3 ; le maximum est \res{6}, atteint en 1 et 4.
Les bornes du segment sont des candidats aux extremums globaux et doivent être comparées aux points critiques.
\editionEtape{2. Tangente et position relative}
On a $f(2)=4$ et $f'(2)=-3$. Ainsi $T:y=-3(x-2)+4=-3x+10$.
Le développement de $(x-2)^3$ donne $x^3-6x^2+12x-8$, égal à $f(x)-(-3x+10)$.
Ce cube est négatif pour $x<2$, nul en 2, positif pour $x>2$.
La courbe est donc sous la tangente sur $[0;2[$ et au-dessus sur $]2;4]$.
Elle la traverse au point de contact : une tangente ne reste pas forcément du même côté de la courbe.

\label{figedition-sujet-A-cor-1-1}\figureGeo[corrige]{fr-sa-fonction}
\end{blocEdition}

\editionPage{sujet-A-cor-2}{Sujet A · Corrigé 2}{Ateliers, défauts et gain}{\hyperref[sujet-A-ex-2]{Énoncé p.~\pageref*{sujet-A-ex-2}} · Raisonnement et rédaction}
\begin{blocEdition}{sujet-A-cor-2}

\editionEtape{1. Reconstituer les probabilités}
Les poids manquants sont $19/20$ et $9/10$. Par probabilités totales,
\[P(D)=\frac35\frac1{20}+\frac25\frac1{10}=\frac3{100}+\frac4{100}=\res{\frac7{100}}.\]
Les articles défectueux provenant de B représentent $4/100$ de tous les articles.
En se limitant aux défauts, on obtient
$P_D(B)=(4/100)/(7/100)=\res{4/7}$.
\editionEtape{2. Passer de la loi aux deux moments}
Le gain $G$ prend $-8$ avec probabilité $7/100$ et 12 avec probabilité $93/100$.
Ainsi $E(G)=(-8)(7/100)+12(93/100)=\res{53/5=10{,}6}$.
Le deuxième moment vaut $64(7/100)+144(93/100)=692/5$.
Il vient
\[V(G)=\frac{692}{5}-\left(\frac{53}{5}\right)^2=\res{\frac{651}{25}=26{,}04}.\]
\editionEtape{3. Utiliser l'évènement contraire}
Deux articles sans défaut ont pour probabilité $(93/100)^2$ grâce à l'indépendance.
La probabilité demandée est donc $1-(93/100)^2=\res{1351/10000=0{,}1351}$.
L'évènement « au moins un » comprend un ou deux défauts ; compter seulement un défaut serait incomplet.

\label{figedition-sujet-A-cor-2-1}\figureGeo[corrige]{fr-sa-arbre}
\end{blocEdition}

\editionPage{sujet-B-cor-auto}{Sujet B · Corrigé}{Automatismes : réponses et méthode}{\hyperref[sujet-B]{Énoncé p.~\pageref*{sujet-B}}}
\begin{blocEdition}{sujet-B-cor-auto}
\begin{questions}\item Au dénominateur 12 : $\dfrac9{12}+\dfrac2{12}=\res{\dfrac{11}{12}}$.
\item $(2^3)^2=2^6$, donc le quotient vaut $2^{6-4}=\res{4}$.
\item Différence de deux carrés : $(3x)^2-1^2=\res{(3x-1)(3x+1)}$.
\item $\dfrac{92-80}{80}=\dfrac{12}{80}=0{,}15$ : hausse de \res{15\,\%}.
\item $(1+t)^2=1{,}21=1{,}1^2$ et $1+t>0$, donc \res{t=10\,\%}.
\item On divise par $\pi r^2\neq0$ : $\res{h=\dfrac{V}{\pi r^2}}$.
\item L’ordonnée à l’origine est l’ordonnée de $A$, soit 3, et la pente vaut $-2$ : $\res{y=-2x+3}$.
\item $\dfrac{8+2\times12+15}{1+2+1}=\dfrac{47}4=\res{11{,}75}$.
\item Parmi les 120 filles, 30 sont internes : $P_F(I)=\dfrac{30}{120}=\res{\dfrac14}$.
\item $P(F\cap I)=\dfrac{30}{200}=\res{\dfrac3{20}}$ ; parmi les 60 internes, 30 sont des filles : $P_I(F)=\res{\dfrac12}$.
\item $g$ s’annule en $-1$ et en 3 ; \res{g(x)>0} sur $\mathopen{]}-1;3[$ et \res{g(x)<0} sur $[-2;-1[$ et sur $]3;4]$.
\item La médiane vaut \res{12} ; l’écart interquartile vaut $Q_3-Q_1=15-9=\res{6}$.\end{questions}
\label{figedition-sujet-B-cor-auto-1}\figureGeo[corrige]{fr-sb-auto}
\end{blocEdition}

\editionPage{sujet-B-cor-1}{Sujet B · Corrigé 1}{Une évolution avec apport}{\hyperref[sujet-B-ex-1]{Énoncé p.~\pageref*{sujet-B-ex-1}} · Raisonnement et rédaction}
\begin{blocEdition}{sujet-B-cor-1}

\editionEtape{1. Déplacer la valeur d'équilibre}
On calcule $v_{n+1}=u_{n+1}-40=(3/4)(u_n-40)=(3/4)v_n$ et $v_0=40$.
Ainsi $v_n=40(3/4)^n$ et \res{u_n=40+40(3/4)^n}.
La puissance est positive et strictement décroissante ; on a $u_n>40$ et la suite décroît strictement.
\editionEtape{2. Prouver le premier seuil}
L'inégalité $u_n\leq50$ équivaut à $(3/4)^n\leq1/4$.
Au rang 4, $81/256>64/256=1/4$ ; au rang 5, $243/1024<256/1024=1/4$.
La décroissance établie précédemment prouve donc que le premier rang est \res{5}.
\begin{python}
def premier_rang():
    n, u = 0, 80.0
    while u > 50:
        u = 0.75 * u + 10
        n = n + 1
    return n
\end{python}
Le test \code{u > 50} continue tant que le seuil n'est pas atteint.
À la sortie, $u=u_n\leq50$ ; le compteur correspond au nombre d'actualisations effectuées.

\end{blocEdition}

\editionPage{sujet-B-cor-2}{Sujet B · Corrigé 2}{Cercle, tangente et projection}{\hyperref[sujet-B-ex-2]{Énoncé p.~\pageref*{sujet-B-ex-2}} · Raisonnement et rédaction}
\begin{blocEdition}{sujet-B-cor-2}

\editionEtape{1. Construire le centre du cercle}
La médiatrice de $[AB]$ est $x=2$.
Pour un point $(x;y)$, l'égalité des distances à $A$ et à $C$ s'écrit
$x^2+y^2=(x-1)^2+(y-3)^2$, soit $x+3y=5$.
L'intersection des deux médiatrices est \res{\Omega(2;1)}.
Le rayon au carré est $\Omega A^2=5$ : le cercle est
\[\res{(x-2)^2+(y-1)^2=5}.\]
Pour $C(1;3)$, le membre de gauche vaut $1+4=5$ : l'appartenance est vérifiée.
\editionEtape{2. Utiliser la direction normale du rayon}
Le rayon $\vect{\Omega C}=(-1;2)$ est normal à la tangente en $C$.
Son équation est $-(x-1)+2(y-3)=0$, donc \res{-x+2y-5=0}.
\editionEtape{3. Projeter et contrôler}
Le projeté s'écrit $H=M-\lambda(-1;2)=(5+\lambda;4-2\lambda)$.
L'équation de la tangente impose $-(5+\lambda)+2(4-2\lambda)-5=0$,
soit $-2-5\lambda=0$. Ainsi \res{H(23/5;24/5)}.
On vérifie $-23/5+48/5-5=0$. De plus $\vect{MH}=(-2/5;4/5)$ est parallèle au normal.
Les deux conditions du projeté sont donc satisfaites et
\[MH=\sqrt{4/25+16/25}=\res{\frac{2\sqrt5}{5}}.\]

\label{figedition-sujet-B-cor-2-1}\figureGeo[corrige]{fr-sb-geometrie}
\end{blocEdition}


\clearpage
\clearpage
\gdef\otChapNum{0}
\pagestyle{otFront}
\phantomsection\label{couverture-programme}
\pdfbookmark[0]{Repères du programme}{bm-couverture-programme}
\otPageTitre{Repères du programme}{Retrouver les contenus, les méthodes et les preuves}
{\fontsize{10}{14}\selectfont
Ce tableau relie les contenus et capacités du programme de Première générale
spécialité aux cours et à des exercices du livre. Il sert à organiser les révisions.
Les approfondissements facultatifs du programme sont sélectionnés ; ils ne
sont pas tous nécessaires à l'apprentissage des attendus. Les références du
texte officiel figurent à la fin du livre.}\par\vspace{3mm}
\begingroup\fontsize{9}{12}\selectfont
\setlength{\tabcolsep}{3pt}
\begin{longtable}{@{}p{101mm}p{58mm}@{}}
\toprule\textbf{Contenu ou capacité travaillée}&\textbf{Où travailler dans le livre}\\\midrule
\endfirsthead
\toprule\textbf{Contenu ou capacité travaillée}&\textbf{Où travailler dans le livre}\\\midrule
\endhead
\multicolumn{2}{@{}l}{\otSemi\color{otNavy}Algèbre}\\*[2mm]
Modes de génération, indices et passage entre formule, récurrence, graphique et motifs & \hyperref[ess-2]{Cours p.~\pageref*{ess-2}}\newline \hyperref[exo-2-1]{Ex.~2.1}\newline \hyperref[exo-2-2]{Ex.~2.2}\\[2.3mm]
Suites arithmétiques : terme général, variations et sommes de termes consécutifs & \hyperref[ess-2]{Cours p.~\pageref*{ess-2}}\newline \hyperref[exo-2-3]{Ex.~2.3}\newline \hyperref[exo-2-6]{Ex.~2.6}\newline \hyperref[dem-suites]{Preuve p.~\pageref*{dem-suites}}\\[2.3mm]
Suites géométriques : terme général, variations et sommes de termes consécutifs & \hyperref[ess-2]{Cours p.~\pageref*{ess-2}}\newline \hyperref[exo-2-4]{Ex.~2.4}\newline \hyperref[exo-2-5]{Ex.~2.5}\newline \hyperref[dem-suites]{Preuve p.~\pageref*{dem-suites}}\\[2.3mm]
Modèles discrets à accroissement ou à taux constant ; calcul de seuils & \hyperref[ess-2]{Cours p.~\pageref*{ess-2}}\newline \hyperref[exo-2-5]{Ex.~2.5}\newline \hyperref[exo-2-11]{Ex.~2.11}\\[2.3mm]
Comportements à long terme : valeur approchée, croissance sans borne et oscillation & \hyperref[ess-2]{Cours p.~\pageref*{ess-2}}\newline \hyperref[exo-2-8]{Ex.~2.8}\newline \hyperref[exo-2-9]{Ex.~2.9}\\[2.3mm]
Sommes des entiers et des puissances successives : justifications & \hyperref[ess-2]{Cours p.~\pageref*{ess-2}}\newline \hyperref[exo-2-7]{Ex.~2.7}\newline \hyperref[exo-2-12]{Ex.~2.12}\newline \hyperref[dem-suites]{Preuve p.~\pageref*{dem-suites}}\\[2.3mm]
Formes développée, canonique et factorisée ; somme et produit des racines & \hyperref[ess-1]{Cours p.~\pageref*{ess-1}}\newline \hyperref[exo-1-2]{Ex.~1.2}\newline \hyperref[exo-1-3]{Ex.~1.3}\newline \hyperref[dem-second-degre]{Preuve p.~\pageref*{dem-second-degre}}\\[2.3mm]
Trinômes s'annulant en deux réels donnés ; courbe de $x\mapsto f(x-m)$ par translation & \hyperref[ess-1]{Cours p.~\pageref*{ess-1}}\newline \hyperref[exo-1-9]{Ex.~1.9}\newline \hyperref[exo-1-15]{Ex.~1.15}\\[2.3mm]
Discriminant, équations, signe du trinôme et inéquations & \hyperref[ess-1]{Cours p.~\pageref*{ess-1}}\newline \hyperref[exo-1-4]{Ex.~1.4}\newline \hyperref[exo-1-5]{Ex.~1.5}\newline \hyperref[dem-second-degre]{Preuve p.~\pageref*{dem-second-degre}}\\[2.3mm]
Choix de forme adapté à une comparaison, une aire ou un problème de maximum & \hyperref[ess-1]{Cours p.~\pageref*{ess-1}}\newline \hyperref[exo-1-6]{Ex.~1.6}\newline \hyperref[exo-1-7]{Ex.~1.7}\\[2.3mm]
\multicolumn{2}{@{}l}{\otSemi\color{otNavy}Analyse}\\*[2mm]
Taux de variation, sécantes, nombre dérivé et interprétation en contexte & \hyperref[ess-3]{Cours p.~\pageref*{ess-3}}\newline \hyperref[exo-3-1]{Ex.~3.1}\newline \hyperref[exo-3-2]{Ex.~3.2}\newline \hyperref[dem-derivation]{Preuve p.~\pageref*{dem-derivation}}\\[2.3mm]
Algorithme : liste des coefficients directeurs des sécantes pour un pas donné & \hyperref[ess-3]{Cours p.~\pageref*{ess-3}}\newline \hyperref[exo-3-6]{Ex.~3.6}\\[2.3mm]
Équation et construction de tangente ; approximation affine & \hyperref[ess-3]{Cours p.~\pageref*{ess-3}}\newline \hyperref[exo-3-4]{Ex.~3.4}\newline \hyperref[exo-3-9]{Ex.~3.9}\newline \hyperref[dem-derivation]{Preuve p.~\pageref*{dem-derivation}}\\[2.3mm]
Dérivées de référence, puissances entières et opérations sur les fonctions & \hyperref[ess-3]{Cours p.~\pageref*{ess-3}}\newline \hyperref[exo-3-3]{Ex.~3.3}\newline \hyperref[exo-3-5]{Ex.~3.5}\newline \hyperref[dem-derivation]{Preuve p.~\pageref*{dem-derivation}}\\[2.3mm]
Non-dérivabilité de la valeur absolue et de la racine carrée en zéro & \hyperref[ess-3]{Cours p.~\pageref*{ess-3}}\newline \hyperref[exo-3-11]{Ex.~3.11}\newline \hyperref[dem-derivation]{Preuve p.~\pageref*{dem-derivation}}\\[2.3mm]
Signe de la dérivée, variations, extremums et optimisation & \hyperref[ess-4]{Cours p.~\pageref*{ess-4}}\newline \hyperref[exo-4-1]{Ex.~4.1}\newline \hyperref[exo-4-2]{Ex.~4.2}\\[2.3mm]
Nombre dérivé nul en un extremum intérieur ; trinôme étudié par sa dérivée & \hyperref[ess-4]{Cours p.~\pageref*{ess-4}}\newline \hyperref[exo-4-2]{Ex.~4.2}\newline \hyperref[exo-4-5]{Ex.~4.5}\\[2.3mm]
Fonctions paires et impaires ; symétries de leurs courbes & \hyperref[ess-4]{Cours p.~\pageref*{ess-4}}\newline \hyperref[exo-4-8]{Ex.~4.8}\\[2.3mm]
Comparaison de courbes, tangentes et inégalités par les variations & \hyperref[ess-4]{Cours p.~\pageref*{ess-4}}\newline \hyperref[exo-4-7]{Ex.~4.7}\newline \hyperref[exo-4-9]{Ex.~4.9}\\[2.3mm]
Exponentielle : définition, propriétés algébriques, positivité et croissance & \hyperref[ess-5]{Cours p.~\pageref*{ess-5}}\newline \hyperref[exo-5-1]{Ex.~5.1}\newline \hyperref[exo-5-2]{Ex.~5.2}\\[2.3mm]
Dérivée exponentielle affine, étude de fonctions et modèles continus & \hyperref[ess-5]{Cours p.~\pageref*{ess-5}}\newline \hyperref[exo-5-6]{Ex.~5.6}\newline \hyperref[exo-5-7]{Ex.~5.7}\\[2.3mm]
Lien avec les suites géométriques ; méthode d'Euler et valeurs approchées de $e$ & \hyperref[ess-5]{Cours p.~\pageref*{ess-5}}\newline \hyperref[exo-5-10]{Ex.~5.10}\newline \hyperref[exo-5-11]{Ex.~5.11}\\[2.3mm]
Cercle trigonométrique, radians, points-images et longueurs d'arcs & \hyperref[ess-6]{Cours p.~\pageref*{ess-6}}\newline \hyperref[exo-6-1]{Ex.~6.1}\newline \hyperref[exo-6-2]{Ex.~6.2}\\[2.3mm]
Sinus, cosinus, valeurs remarquables et symétries des angles & \hyperref[ess-6]{Cours p.~\pageref*{ess-6}}\newline \hyperref[exo-6-3]{Ex.~6.3}\newline \hyperref[exo-6-4]{Ex.~6.4}\newline \hyperref[dem-trigonometrie]{Preuve p.~\pageref*{dem-trigonometrie}}\\[2.3mm]
Algorithme : approximation de $\pi$ par la méthode d'Archimède & \hyperref[ess-6]{Cours p.~\pageref*{ess-6}}\newline \hyperref[exo-6-15]{Ex.~6.15}\\[2.3mm]
\multicolumn{2}{@{}l}{\otSemi\color{otNavy}Géométrie}\\*[2mm]
Calcul vectoriel dans le plan, coordonnées, normes et bases orthonormées & \hyperref[ess-7]{Cours p.~\pageref*{ess-7}}\newline \hyperref[exo-7-1]{Ex.~7.1}\newline \hyperref[exo-7-2]{Ex.~7.2}\\[2.3mm]
Produit scalaire par les coordonnées, les normes, un angle ou une projection & \hyperref[ess-7]{Cours p.~\pageref*{ess-7}}\newline \hyperref[exo-7-2]{Ex.~7.2}\newline \hyperref[exo-7-3]{Ex.~7.3}\\[2.3mm]
Bilinéarité, symétrie, carrés de normes et calculs sans coordonnées & \hyperref[ess-7]{Cours p.~\pageref*{ess-7}}\newline \hyperref[exo-7-6]{Ex.~7.6}\newline \hyperref[exo-7-9]{Ex.~7.9}\\[2.3mm]
Orthogonalité, longueurs et angles ; formule d'Al-Kashi & \hyperref[ess-7]{Cours p.~\pageref*{ess-7}}\newline \hyperref[exo-7-4]{Ex.~7.4}\newline \hyperref[exo-7-7]{Ex.~7.7}\newline \hyperref[dem-geometrie]{Preuve p.~\pageref*{dem-geometrie}}\\[2.3mm]
Expression avec un milieu et cercle de diamètre donné & \hyperref[ess-7]{Cours p.~\pageref*{ess-7}}\newline \hyperref[exo-7-13]{Ex.~7.13}\newline \hyperref[exo-7-15]{Ex.~7.15}\newline \hyperref[dem-geometrie]{Preuve p.~\pageref*{dem-geometrie}}\\[2.3mm]
Vecteurs normaux et équations cartésiennes de droites & \hyperref[ess-8]{Cours p.~\pageref*{ess-8}}\newline \hyperref[exo-8-1]{Ex.~8.1}\newline \hyperref[exo-8-2]{Ex.~8.2}\\[2.3mm]
Coordonnées d'un projeté orthogonal et configuration repérée & \hyperref[ess-8]{Cours p.~\pageref*{ess-8}}\newline \hyperref[exo-8-5]{Ex.~8.5}\newline \hyperref[exo-8-6]{Ex.~8.6}\\[2.3mm]
Équations de cercles : reconnaissance, centre, rayon et utilisation & \hyperref[ess-8]{Cours p.~\pageref*{ess-8}}\newline \hyperref[exo-8-3]{Ex.~8.3}\newline \hyperref[exo-8-4]{Ex.~8.4}\\[2.3mm]
\multicolumn{2}{@{}l}{\otSemi\color{otNavy}Probabilités et statistiques}\\*[2mm]
Probabilités conditionnelles, arbres et tableaux ; inversion du conditionnement & \hyperref[ess-9]{Cours p.~\pageref*{ess-9}}\newline \hyperref[exo-9-1]{Ex.~9.1}\newline \hyperref[exo-9-2]{Ex.~9.2}\\[2.3mm]
Partitions, probabilités totales et modèles à plusieurs groupes & \hyperref[ess-9]{Cours p.~\pageref*{ess-9}}\newline \hyperref[exo-9-5]{Ex.~9.5}\newline \hyperref[exo-9-6]{Ex.~9.6}\newline \hyperref[dem-probabilites]{Preuve p.~\pageref*{dem-probabilites}}\\[2.3mm]
Indépendance, incompatibilité et évènements contraires & \hyperref[ess-9]{Cours p.~\pageref*{ess-9}}\newline \hyperref[exo-9-3]{Ex.~9.3}\newline \hyperref[exo-9-4]{Ex.~9.4}\\[2.3mm]
Succession de deux épreuves et répétition de Bernoulli jusqu'à quatre épreuves & \hyperref[ess-9]{Cours p.~\pageref*{ess-9}}\newline \hyperref[exo-9-9]{Ex.~9.9}\newline \hyperref[exo-9-10]{Ex.~9.10}\\[2.3mm]
Variable aléatoire comme fonction d'un univers fini ; évènements et lois & \hyperref[ess-10]{Cours p.~\pageref*{ess-10}}\newline \hyperref[exo-10-1]{Ex.~10.1}\newline \hyperref[exo-10-2]{Ex.~10.2}\\[2.3mm]
Espérance, variance, écart type et formule de König-Huygens & \hyperref[ess-10]{Cours p.~\pageref*{ess-10}}\newline \hyperref[exo-10-2]{Ex.~10.2}\newline \hyperref[exo-10-3]{Ex.~10.3}\newline \hyperref[dem-probabilites]{Preuve p.~\pageref*{dem-probabilites}}\\[2.3mm]
Linéarité de l'espérance, transformations et problèmes de jeux équitables & \hyperref[ess-10]{Cours p.~\pageref*{ess-10}}\newline \hyperref[exo-10-5]{Ex.~10.5}\newline \hyperref[exo-10-7]{Ex.~10.7}\newline \hyperref[dem-probabilites]{Preuve p.~\pageref*{dem-probabilites}}\\[2.3mm]
Simulation de lois, échantillons et moyennes observées & \hyperref[ess-11]{Cours p.~\pageref*{ess-11}}\newline \hyperref[exo-11-4]{Ex.~11.4}\newline \hyperref[exo-11-5]{Ex.~11.5}\\[2.3mm]
Répétition d'échantillons et proportion des écarts à la moyenne sous un seuil donné & \hyperref[ess-11]{Cours p.~\pageref*{ess-11}}\newline \hyperref[exo-11-10]{Ex.~11.10}\newline \hyperref[exo-11-11]{Ex.~11.11}\\[2.3mm]
Estimation d'aires et de $\pi$ par Monte-Carlo & \hyperref[ess-11]{Cours p.~\pageref*{ess-11}}\newline \hyperref[exo-11-8]{Ex.~11.8}\newline \hyperref[exo-11-9]{Ex.~11.9}\\[2.3mm]
\multicolumn{2}{@{}l}{\otSemi\color{otNavy}Algorithmique et programmation}\\*[2mm]
Variables, affectations, tests, fonctions et programmation modulaire & \hyperref[ess-12]{Cours p.~\pageref*{ess-12}}\newline \hyperref[exo-12-1]{Ex.~12.1}\newline \hyperref[exo-12-2]{Ex.~12.2}\\[2.3mm]
Listes en extension, par ajouts et en compréhension ; modification et parcours & \hyperref[ess-12]{Cours p.~\pageref*{ess-12}}\newline \hyperref[exo-12-4]{Ex.~12.4}\newline \hyperref[exo-12-5]{Ex.~12.5}\\[2.3mm]
Termes, sommes, seuils et factorielle ; Fibonacci et Syracuse & \hyperref[ess-12]{Cours p.~\pageref*{ess-12}}\newline \hyperref[exo-12-6]{Ex.~12.6}\newline \hyperref[exo-12-7]{Ex.~12.7}\\[2.3mm]
Méthode de Newton : tangente et approximation de $\sqrt2$ & \hyperref[ess-12]{Cours p.~\pageref*{ess-12}}\newline \hyperref[exo-12-14]{Ex.~12.14}\\[2.3mm]
Moments d'une loi, fréquences de lettres et simulations d'aires & \hyperref[ess-12]{Cours p.~\pageref*{ess-12}}\newline \hyperref[exo-12-10]{Ex.~12.10}\newline \hyperref[exo-12-11]{Ex.~12.11}\\[2.3mm]
\multicolumn{2}{@{}l}{\otSemi\color{otNavy}Vocabulaire et logique}\\*[2mm]
Appartenance, inclusion, ensembles de nombres, opérations et intervalles & \hyperref[ess-13]{Cours p.~\pageref*{ess-13}}\newline \hyperref[exo-13-1]{Ex.~13.1}\newline \hyperref[exo-13-2]{Ex.~13.2}\\[2.3mm]
Couples, produit cartésien et cardinal d'un ensemble fini & \hyperref[ess-13]{Cours p.~\pageref*{ess-13}}\newline \hyperref[exo-13-3]{Ex.~13.3}\\[2.3mm]
Implication, équivalence, réciproque, contraposée et conditions nécessaires ou suffisantes & \hyperref[ess-13]{Cours p.~\pageref*{ess-13}}\newline \hyperref[exo-13-6]{Ex.~13.6}\newline \hyperref[exo-13-7]{Ex.~13.7}\\[2.3mm]
Quantification, négation, contre-exemples et statut des égalités ou des lettres & \hyperref[ess-13]{Cours p.~\pageref*{ess-13}}\newline \hyperref[exo-13-4]{Ex.~13.4}\newline \hyperref[exo-13-5]{Ex.~13.5}\\[2.3mm]
Preuves directes, par cas, par contraposée et par l'absurde & \hyperref[ess-13]{Cours p.~\pageref*{ess-13}}\newline \hyperref[exo-13-10]{Ex.~13.10}\newline \hyperref[exo-13-11]{Ex.~13.11}\\[2.3mm]
\multicolumn{2}{@{}l}{\otSemi\color{otNavy}Automatismes}\\*[2mm]
Valeurs initiales ou finales, taux simples, composés et réciproques & \hyperref[ess-14]{Cours p.~\pageref*{ess-14}}\newline \hyperref[exo-14-1]{Ex.~14.1}\newline \hyperref[exo-14-2]{Ex.~14.2}\\[2.3mm]
Calcul littéral, produit nul et signes d'expressions & \hyperref[ess-14]{Cours p.~\pageref*{ess-14}}\newline \hyperref[exo-14-3]{Ex.~14.3}\newline \hyperref[exo-14-4]{Ex.~14.4}\\[2.3mm]
Lectures graphiques, pentes et équations de droites & \hyperref[ess-14]{Cours p.~\pageref*{ess-14}}\newline \hyperref[exo-14-5]{Ex.~14.5}\newline \hyperref[exo-14-6]{Ex.~14.6}\newline \hyperref[sujet-A]{Sujet A p.~\pageref*{sujet-A}}\newline \hyperref[sujet-B]{Sujet B p.~\pageref*{sujet-B}}\\[2.3mm]
Automatismes de Seconde : fractions, puissances, proportions, variable isolée & \hyperref[ess-14]{Cours p.~\pageref*{ess-14}}\newline \hyperref[exo-14-2]{Ex.~14.2}\newline \hyperref[exo-14-3]{Ex.~14.3}\newline \hyperref[sujet-B]{Sujet B p.~\pageref*{sujet-B}}\\[2.3mm]
Tableaux croisés, probabilités conditionnelles et représentations statistiques & \hyperref[ess-14]{Cours p.~\pageref*{ess-14}}\newline \hyperref[exo-14-7]{Ex.~14.7}\newline \hyperref[exo-14-10]{Ex.~14.10}\newline \hyperref[sujet-A]{Sujet A p.~\pageref*{sujet-A}}\newline \hyperref[sujet-B]{Sujet B p.~\pageref*{sujet-B}}\\[2.3mm]
Histogrammes, diagrammes en barres ou en boîte et indicateurs descriptifs & \hyperref[ess-11]{Cours p.~\pageref*{ess-11}}\newline \hyperref[exo-11-1]{Ex.~11.1}\newline \hyperref[exo-11-2]{Ex.~11.2}\\[2.3mm]
\multicolumn{2}{@{}l}{\otSemi\color{otNavy}Compétences transversales}\\*[2mm]
Chercher, modéliser, représenter, raisonner, calculer et communiquer dans des synthèses & \hyperref[ess-14]{Cours p.~\pageref*{ess-14}}\newline \hyperref[exo-14-12]{Ex.~14.12}\newline \hyperref[exo-14-13]{Ex.~14.13}\\[2.3mm]
\bottomrule\end{longtable}\endgroup\clearpage

\thispagestyle{otFront}
\phantomsection\pdfbookmark[0]{Références et repères éditoriaux}{bm-references}
\otPageTitre{Références}{France · Première générale · Spécialité mathématiques}
\fontsize{10.5}{14.5}\selectfont
\otIntertitre{Le programme de référence}
Le livre s'appuie sur le programme publié au \textbf{BO n° 14 du 2 avril 2026}
(arrêté du 26 février 2026, NOR MENE2602917A), applicable à la rentrée 2026-2027.
Les repères de l'annexe relient ses contenus aux cours et exercices du livre.
Les approfondissements facultatifs sont sélectionnés, sans prétendre les reproduire tous.\par\vspace{3mm}
\href{https://www.education.gouv.fr/bo/2026/Hebdo14/MENE2602917A}{\textcolor{otNavy}{Programme officiel et annexe sur education.gouv.fr}}\par\vspace{5mm}
\otIntertitre{L'épreuve anticipée}
La préparation suit la \textbf{note de service du 11 septembre 2026}, NOR MENE2622640N,
publiée au BO spécial n° 4 du 17 septembre 2026. Elle s'applique en 2026-2027 aux
épreuves présentées au titre du baccalauréat session 2028 et remplace la note de service
du 10 juin 2025 (NOR MENE2515469N). Les automatismes y prennent la forme de réponses
courtes. Le programme de Première précise que sa liste d'automatismes s'ajoute à celle
de Seconde, à entretenir en Première. Le guide du livre distingue ce cadre officiel du barème
choisi pour ses sujets d'entraînement.\par\vspace{3mm}
\href{https://www.education.gouv.fr/bo/2026/Special4/MENE2622640N}{\textcolor{otNavy}{Définition officielle de l'épreuve anticipée}}\par\vspace{2mm}
\href{https://eduscol.education.gouv.fr/5688/epreuve-anticipee-de-mathematiques-aux-baccalaureats-general-et-technologique}{\textcolor{otNavy}{Présentation et ressources de préparation sur Éduscol}}\par\vspace{5mm}
\otIntertitre{Les choix pédagogiques}
Les 210 exercices de chapitre, 14 défis et deux sujets blancs sont des créations pour ce livre.
Les sujets blancs ne reproduisent pas un sujet officiel. Les défis très exigeants sont
facultatifs et explicitent leurs prérequis ; ils ne nécessitent pas de cours de Terminale.
L'ordre du livre peut être adapté par le professeur.\par\vspace{3mm}
Les figures, courbes, arbres et tableaux sont vectoriels. Les schémas sans échelle
imposée sont signalés. Les programmes Python peuvent être exécutés séparément.
Les résultats numériques illustrent les calculs théoriques ; les arrondis et les
hypothèses d'indépendance sont précisés lorsqu'ils interviennent.\par\vspace{5mm}
{\fontsize{9}{12}\selectfont\color{otMuted}Sources officielles consultées le 23 octobre 2026.
La charte reprend le modèle OptiTAB de Terminale : couleurs, typographie et repères de difficulté.}
\clearpage

%% =====================================================================
%%  4e de couverture
%% =====================================================================
\clearpage
\thispagestyle{empty}
\begin{tikzpicture}[remember picture,overlay]
\begin{scope}[shift={(current page.south west)}]
  % Fond
  \fill[otNavy] (0,0) rectangle (210mm,297mm);
  \shade[bottom color=otNavyDark,top color=otNavy] (0,0) rectangle (210mm,160mm);
  \draw[white,opacity=0.06,line width=0.3pt,step=7mm] (0,0) grid (210mm,297mm);

  % Accroche et texte
  \node[anchor=north west,inner sep=0pt,text width=150mm,align=left] at (24mm,270mm)
    {{\otSemi\fontsize{9}{11}\selectfont\color{otGreen}%
       \otCaps[18]{\otTitreLivre\ · \otNiveau}}\par\vspace{7mm}
     {\otSemi\fontsize{25}{31}\selectfont\color{white}%
       Réussir sa Première,\\exercice après exercice.}\par\vspace{9mm}
     {\fontsize{11.5}{16}\selectfont\color{white!88!otNavy}%
       Ce livre rassemble 224 exercices de spécialité mathématiques :
       210 exercices dans 14 chapitres et 14 défis d'approfondissement. Chaque
       exercice est accompagné d'un corrigé détaillé qui explique la méthode
       à chaque étape : de quoi travailler en autonomie, réviser efficacement
       et préparer l'épreuve anticipée avec deux sujets blancs originaux.\par}\vspace{9mm}
     {\otRegular\fontsize{10.5}{14}\selectfont\color{white}%
       \otCheck\hspace{3mm}Des bases accessibles aux défis très exigeants\par\vspace{3.4mm}
       \otCheck\hspace{3mm}Des corrigés rédigés, avec la méthode à retenir\par\vspace{3.4mm}
       \otCheck\hspace{3mm}Des tableaux de variations et des courbes soignés\par\vspace{3.4mm}
       \otCheck\hspace{3mm}Un tableau de bord pour suivre sa progression\par}};

  % Bloc « en ligne »
  \fill[white,rounded corners=3mm] (24mm,62mm) rectangle (186mm,98mm);
  \node[anchor=west,inner sep=0pt] at (30mm,80mm)
    {\color{otNavy}\qrcode[height=26mm,nolink]{\otURL}};
  \node[anchor=west,inner sep=0pt,text width=110mm] at (64mm,80mm)
    {{\otSemi\fontsize{12}{15}\selectfont\color{otNavy}Prolonge ton travail en ligne}\par\vspace{2mm}
     {\fontsize{10.5}{14}\selectfont\color{otInk}Cours, fiches de synthèse et exercices
       sur \textbf{\otSite}. Scanne le code pour y accéder.}};

  % Logo, prix et code-barres
  \node[anchor=south west,inner sep=0pt] at (24mm,20mm)
    {\includegraphics[width=46mm]{logo-blanc}};
  \fill[white] (138mm,16mm) rectangle (186mm,44mm);
  \node[text=otMuted,align=center,font=\otLight\fontsize{7}{9}\selectfont]
    at (162mm,30mm) {ANNÉE COMPLÈTE\\14 chapitres\\Édition 2026};
\end{scope}
\end{tikzpicture}

\end{document}
